\documentclass[10pt,a4paper,twoside,bg=print,twocolumn,openany,nodeprecatedcode]{dndbook}
\usepackage[utf8]{inputenc}
\usepackage{tabulary}
\usepackage[normalem]{ulem}
\usepackage{color,soul}
\usepackage{float}
\usepackage{caption}
\usepackage{wrapfig}
\usepackage{tocloft}
\usepackage[toc]{multitoc}
\usepackage[hidelinks]{hyperref}
\raggedbottom
\extrafloats{2000}
\cftsetindents{section}{0em}{0em}
\cftsetindents{subsection}{0em}{0em}
\cftsetindents{subsubsection}{0.2em}{0em}
\setcounter{tocdepth}{1}
\renewcommand*{\multicolumntoc}{3}
\setlist[description]{nolistsep,listparindent=3pt,topsep=6pt}
\date{}
\begin{document}
\frontmatter
\tableofcontents
\mainmatter
\addcontentsline{toc}{chapter}{Introduction}
\markboth{INTRODUCTION}{INTRODUCTION}
\chapter*{Introduction}
\begin{DndSidebar}[float=!b]{}
Pitter patter pitter patter



Run run run child



Dive beneath the shadows now



Or seek shelter in the trees.



It will avail you nothing



The shadow fey are coming



With their hounds and their arrows and bright silver spears



Though you be the canniest fox or the swiftest hawk



They will find you and bring you to their dark queen



Woven deep in her shadow court



Run child run



Find your shelter and hold still as death



The shadow fey are coming!



\end{DndSidebar}

Shadow is everything and nothing. It is gloom and delight. If you try to seize it, it ripples away, laughing.

So yes, the Plane of Shadows is an uncertain place—which makes it a very difficult place to define and capture in maps and descriptions, in known quantities or certainties. Everyone knows that the demonic creatures of the hells are vile and secure in their malice. Likewise, the shining angels of the celestial halls are unyielding in their eternal pilgrimage for righteousness and law. These realms are well known and understood, though one is bright and the other stygian.

\DndQuote%
{}
{''The creatures of shadow are far less defined and far more flexible—and no less dangerous for all that. Even their language, Umbral, is a fluid thing, full of double meanings and wordplay. They resist definitions, but their malleable, quicksilver nature defines them. You should always trust the word of an angel, rarely given but given with certainty. The word of a shadow creature, however, is given quickly and easily but broken just as readily.''}
{''Creatures of the Shadow Realm are like unto shadows themselves, flickering and shifting with the sun and moon, changing and moving. And their realm is also a place of frequent change, including the violent overthrow of courts and dynasties and the slow erosion of places by the shifting ebon tides, drifting like sandbars or vanishing like a hillside after the rains.''}
{''This book explores the Shadow Realm and attempts to define it, as understood by its inhabitants. Its people and their gods, its magic and the distinctive way it works in the world, and even the monsters, nations, and great cities of the Shadow Realm are all defined here. Perhaps defined a bit more than the inhabitants themselves would like, for they prefer to slip always from one definition to another, to a slightly more favorable one—as the mood strikes them.''}
{''Welcome to the Shadow Realm. Its high-flown powers and its lowborn scoundrels alike have a proposal for you: swear your sword and wands to their cause and shift the ebon tides in their favor. Join us in the shadows!''}
{Wolfgang Baur Kirkland, WA November 2021}
\chapter{1. Overview and History}
For such a vast realm, surprisingly little is known about the Plane of Shadow, often also called the Shadow Realm. Herewith, a short précis and overview for those who have not previously drunk deep of its twilight power. The whole plane is often described as a distortion and a shadow of the mortal realm, and certainly there are resemblances, as any shadow resembles its source and origin—and dark reflections and mirrored echoes are frequent enough. No shadow exists without a light and a source, after all.

At the same time, parts of the Shadow Realm are distinct unto the plane itself, especially those portions worked by shadow fey, goblins, bearfolk, humans, and darakhul as their realms and kingdoms. These places have been shaped by magic, as creatures manipulated the raw material of shadow into lands and settlements that please their makers. Great mages and clerics bend the land and the structures upon it into new forms, and the Shadow Realm itself often alters its own shape and direction in unpredictable ways.

\section{Elements of Shadow}
Here are the elements of the Shadow Realm covered in this volume.


\begin{description}
\item[Overview of Shadow Realm's Nature and Travel.]
A short history of what the Shadow Realm is, how it became an inhabited plane, and various questions of its lore.
\item[Umbral People and Heroes.]
The peoples of the Shadow Realm are surprisingly many for such a strange plane. Shadow fey and darakhul find it congenial, as do sable elves, wyrd gnomes, bearfolk, shadow goblins, umbral humans, and even the small erina and alseid more commonly found in deep forests. With relatively few cities and wide stretches of wilderness, the Shadow Realm is not a place of great kingdoms but rather small settlements and fiefdoms loosely bound to charismatic kings and queens. The various peoples and heroes of the Shadow Realm can be found in \textit{Chapter 2} and \textit{Chapter 3}.
\item[Shadow Magic.]
A discussion of the nature of umbral magic, plus new spells of illumination, trickery, and shadow. \textit{Chapter 4} also includes all-new fey magic and bits of weather magic and an expansion of the school of dooms.
\item[Land, Water, and Shadow Walking.]
The Shadow Realm itself is best known to mortals through the shadow roads, the shortcuts the fey take from one place to another through shadow. These roads have their own rituals, guardians, and mysteries, as described in \textit{Chapter 5}. This chapter also includes a discussion of the Black River Styx and other hazards.
\item[Shadow Realm and the Courts of the Fey.]
The Shadow Realm is ruled by undead, fey, and a scattered group of their allies and friends. They create regions where the shadow stuff of the plane is relatively stable, where crops grow and magic flows normally and where roads are stable. The realms of the fey are described in \textit{Chapter 6}, and the realms of the bearfolk, darakhul, shades, and humans in \textit{Chapter 7}.
\item[Umbral Pantheon.]
Some gods of the Shadow Realm are familiar to mortals, though many have slightly different features, and they rely more on shadow goblins and shadow fey than the human pantheons do. Bengta the Bear Maiden is a power here and so is Ailuros—clearly a form of Bastet—and Hecate and Charun. Additionally, minor godlings rise and fall with some frequency, including the many Santerr, once-mortal figures raised to an immortal sainthood through the devotion of their followers and twilight magic. Of course, it must be noted that demon lords and dark gods are common here, often gathering great power for a time. Vardesain and Alquam both have large followings, as does a distinct form of Loki. All these can be found in detail in \textit{Chapter 8}.
\item[Shadow Creatures.]
The monsters of the Shadow Realm—including quicksteps and fey sorcerers as well as undead, exquisite rose golems, birch sirens, and other oddities from the shadow forests and quicksilver marshes—are provided in \textit{Chapter 9}. Encounter tables can be found in the \textit{Appendix}.
\item[Umbral Magic Items and Tricks.]
The best-known umbral magic items often focus on shapeshifting, radiance and illusion, shadow magic, and shadow warping. You will also find related fey magic, such as memory philters, illusion seeds, and road magic. Spells related to these are described in \textit{Chapter 4}. Magic items and the gear and equipment most commonly found in the umbral lands are found in \textit{Chapter 10}.
\end{description}

\section{A Brief History of Shadow}
The Shadow Realm has existed for as long as the mortal world has. It is a place of shadows where twilight rules, and its various peoples hold close their nefarious plans and draw power from the dark while avoiding its corruption and the identity-draining waters of the afterlife.

As a place of mortals mixed with minor planar powers, it has been congenial to the quicksilver fey, the nihilistic satarre, and ambitious humans willing to leave sunlit and wholesome lands for a more malleable, more dangerous, and more powerful realm.

\subsection{Making of the Shadow Plane}
In the Shadow Realm, the world is what you make it. And the first great power to discover this was not one of the fey but Hecate, the goddess of the moon.

Most natives of the plane and most visitors think of shadow as a land of reasonably well-fixed, or at least solid, plains and forests. And yet it was not always so, for Hecate found it as an ocean of light and darkness, a morass of radiance beams, trickling shadow mires, and undulating chaotic hills and waves.

As is the way of a goddess of moon and magic, she first conjured a place for herself to stand by making a stone with radiant runes (and thus inventing the use of foundation stones). Then she brought the tides to still the waters, to make them flow and echo rather than crash and storm. The hymns praising Hecate in every one of her temples in the Shadow Realm refer to her as She Who Steadied the Tides, and she is indeed the maker and founder of inhabitation of the plane. To make room for her followers among the fey and humans, she conjured white birch groves, tall red pines on the hills, and moon stones. The waters she diverted into gentle ponds and deep lakes to reflect her glory.

When it was to her liking, she shared the place with her most dutiful followers among the goblins, shadow fey, some exiles from the elf lands, and then the bearfolk and a few night-loving humans. Over time, more and more of the land became solid. All around its shores, even now, the shadow has no end, a roiling ocean of raw chaos and deepest night, leavened with moonlight.

The Shadow Realm is still very much a place of arcane dreams and shapings, where forms and thoughts flow like water, where nothing lasts forever, where dark tides carry away towns and mountains—where skill and kindness win over brute strength or cruelty. Except, of course, for those days and nights when everything changes and dissolves back into mystery, and new lights shine to please new shapers.

\subsection{Early Visitors}
The power of the Shadow Realm has always drawn the ambitious and the striving—sometimes just for a season, long enough to master a few useful secrets—but others have come and stayed. Hecate is widely acknowledged as she who invited the first visitors, though some claim a more ancient race once made their home in the realm. She brought the shadow fey here, and then the sable elves came, and both of these fey groups brought servitor races. In the case of the shadow fey, these were shadow goblins, humans, and gnomes. In the case of the sable elves, they were courtfolk halflings, erina, and occasionally alseids, quicksteps, and various sprites and pixies. The land was soon overrun with fey courts and kingdoms in a distant time that both sides refer to as the First Flowering of Shadow.

\subsubsection{Shadow Fey and Servitors}
The Scáthesidhe, or shadow fey, will always be quick to tell you that they were the first to arrive in the realm. They do not usually mention the reason why: they broke with the elves of the Summer Lands in pursuit of power and found it in the land Hecate showed them. Their understanding and mastery of the realm goes back to a time exceeding recorded history, though their numbers have never been great.

Indeed, they brought with them the elven custom of employing servitor races as their eyes, ears, and hands in remaking and ruling the Shadow Realm. Where the elves retained halflings, humans, bearfolk, alseid, and gnomes, the shadow fey servitors were somewhat different. Gnomes—more precisely, those most distinguished of arcanists, the wyrd gnomes—certainly retain influence among the shadow fey. Also goblins, which came with the shadow fey and became attuned to moonlight magic and fey trickery. And the humans who serve are relatively numerous, especially in the cities and the companies of the umbral human settlements.

\subsubsection{Humans}
Many humans have come to the Shadow Realm as servants or dupes of the shadow fey, and some have fallen into the realm by misadventure, straying from a shadow road. Most, however, descend from those who sought the place out to acquire power, study magic, or escape troubles of their own. This is the foundation of the cities and baronies ruled and inhabited by humans. Umbral humans, as they are called, are often considered ambitious, greedy, or slightly treacherous compared to more reliable servants like goblins and gnomes.

\subsubsection{Darakhul and Other Undead}
Not surprisingly, the darakhul, with their hatred of daylight in the mortal world, found the Shadow Realm almost pleasant, and they established a new empire. Their emperor here rules a greater and far more populated region of undead with ties to Evermaw and to the mortal Empire of Ghouls. Undead dragons, their unhatched wyrmlings, and lich-kings are part of the \textit{Twilight Empire}, and the darakhul are among the greatest rivals to the shadow fey and their courts. Their realms in shadow also host a much higher number of shades than those found elsewhere.

\subsubsection{Bearfolk}
The bearfolk came to shadow not because they were drawn to its power, but because others were. The shadow chewers are said to have first pursued shadow fey into darkness and found their planar bolt-hole. They soon realized the war against the undead and the shadow fey would never succeed unless they could at least understand the Shadow Realm. In time, the first druids, rangers, and other bearfolk visitors to the Shadow Realm built settlements full of light, built warding stones, and grew entire circles of moonlit oaks, gathering the strength of their intimate understanding of darkness to prevent evil from ruling the shadows without rivals. The moon wells, the groves, and the mere presence of the bearfolk have all given the Shadow Realm places of refuge.

\begin{DndSidebar}[float=!b]{Peoples of Midgard}
The Shadow Realm has many useful connections to the Midgard campaign setting. The people of the Midgard Crossroads, Morgau, and Nuria have all found their way to the Shadow Realm, as did a surprising number of the people of the Magocracies after the mage wars, the bearfolk as a whole, and a splinter faction of the Ghoul Empire. Likewise, the trollkin share an ancient bond of blood with creatures who dwell in shadow, and the elves of Arbonesse walk secret paths that lead beyond the material world.



However, most of the people of Midgard know little or nothing of the Shadow Realm other than it's a "place of elves and hostile fey" or similar. They have enough worries with the everyday dangers of Midgard to avoid seeking out additional magical realms.



\end{DndSidebar}

\subsubsection{Minor Races}
Trollkin, ravenfolk, alseid, erina, quicksteps, and others arrived in the Shadow Realm by many paths and in much smaller numbers than those races listed above. Shades are relatively well known, though few in numbers, and the verminous satarre are often found seeking to unravel the plane by various acts of sabotage. Ratatosks seem to enjoy the company of elves, bearfolk, and gnomes and are sometimes visitors in their cities and settlements. Hobgoblins arrived often as mercenaries or leaders among the shadow goblins.

Least commonly, occasionally one can find gnolls in the realm, generally servants of the gods of death, such as Anu-Akma, and subek, tosculi, orcs, dragonborn—even catfolk, often on some pilgrimage for Hecate or Ailuros. Aasimar and tieflings visit as planar travelers but are considered strangers, never natives of the plane.

\subsection{Connections to the Mortal World}
The Shadow Realm is connected to the mortal world primarily by the shadow roads, well known to geomancers, elves, and the shadow fey as arcane connections and fey portals. While most travelers using such roads desire nothing more than to move swiftly through shadow to their destination, they sometimes grow lost or are tempted to leave the known path in pursuit of some vision, treasure, or refuge just off the road. Once they step off the shadow road, they are in the Shadow Realm itself, and the less skilled or less magically adept find it difficult to return to the path... and so they remain in the Shadow Realm.

Connections to other spheres are a bit less common. They fall into three rough categories: rituals and summoning, the roots of Yggdrasil, and cities fallen into shadow.

\subparagraph{Rituals and Summoning}
Rituals can summon creatures from the Shadow Realm—and powerful casters in the Shadow Realm can summon mortals there. For samples of this magic, see \textit{Chapter 4} and the \textit{shadow portal scroll} in \textit{Chapter 10}. In a few cases, a caster deliberately warps a shadow walk spell or speaks a glamour to travelers on a shadow road to bring them deeper into shadow. Most creatures brought into shadow in this way arrive against their will, so getting back to the mortal world becomes their primary motivation.

\subparagraph{Shadow Roots of Yggdrasil}
The World Trees that climb up to the celestial realms and to distant spheres also sometimes have hollow sections of trunk that lead down to Yggdrasil's roots, passing through the Shadow Realm on the way. While this superficially seems similar to taking a shadow road into the realm, there are important differences in the two approaches. Entering the roots of Yggdrasil takes one out of the mortal realm and into a twisting set of interplanar passageways that wind up in various hells, planes of the undead, ghostly afterlives of brooding northern pantheons, and other dark realms at the root of the World Tree.

If one is not wise in the ways of the planes, it is possible to enter planes entirely unrelated to the Shadow Realm, which tends to cling close to the mortal plane. Exactly how one steps aside into shadow is a crucial skill known to the trollkin, bearfolk, ravenfolk, and gnomes, among others, who often use the World Tree for this purpose. Other races, including humans, goblins, darakhul, and the shadow fey, avoid this method of moving between planes.

\subparagraph{Cities in Shadow}
Multiple cities have fallen into the Shadow Realm over the years, starting with Corremel in the days of the Gods War, and later the city of \textit{Oshragora} and parts of the \textit{Twilight Empire}. Cities linked to the human god-kings and the lich-kings of the undead princes seem especially prone to this, and most cities fall into shadow because a mortal caster draws so much power from the Shadow Realm that it ultimately consumes them and everything around them, literally pulling the place out of the mortal world through the sheer hold that shadow corruption has on the place. In the case of Corremel, the process has left a distinct set of links between the city of Corremel-in-Light (as those of the Shadow Realm refer to its rebuilt counterpart in the mortal world) and the city of Corremel-in-Shadow (ruled by shadow fey on a dark river and described in \textit{Chapter 6}).

\begin{DndSidebar}[float=!b]{The Stones of the Ancients / from the Book of Ebon Tides}
Some ancient sagas and fragments of early script found on strange foundation stones mention the Ylgaar, a race of river giants who controlled the plane before the shadow fey but whose attempts to tame the tides ultimately failed. Some driftstone islands of theirs remain, which sometimes wash up on the Saltwood strand, or less often, a goblin explorer or fisherman will find their ruins at sea.



From time to time, mad prophets and keen arcane divinations known only to the shadow fey reveal a whole city or region of empty roads, crumbled keeps, or echoing hills inhabited by a few gaunt giants who seem to know nothing of any earlier age of glory. The greatest of these lost lands was Olderain, said to be anchored around a grove of golden apples and a bright white tower-tomb of the last of the giants' wizards. Many goblins will happily sell you a map to it and a boat to get there.



\end{DndSidebar}

\section{A Land of Ebon Tides}
When learning about the Shadow Realm, the first element that must be addressed is not the people of the realm but the mutable land itself. Most people of the mortal world think of the ground, the forests, and the mountains as permanent, and we think of shadows and light as transitory. And they are, usually, except when some bright eternal fire wells up from the deep earth or some cavern remains shrouded in darkness for aeons.

In the Shadow Realm though, shadows are eternal while land and rivers are—if not entirely transitory—notably less certain than they are in the mortal lands. In the lands of the shadow fey and bearfolk, of sable elves and wyrd gnomes, the forests walk, and the mountains wander from time to time. It happens slowly enough that few notice it before their eyes, but over days and weeks, when the tides shift, the land itself shoves rivers and even cities some distance from their usual and well-accustomed sites.

These ripples in the land itself are called the ebon tides, the stretches and compressions in reality that shadow folk know and plan for. Any survey map is useful for a season or sometimes for a lifetime—at least if that lifetime is not overlong. Hills can crumble quickly, exposed to eroding moonlight and washed out to sea. Harbors silt up, and forests quiver and rearrange themselves around riverbeds that dance slowly, ever so slowly, like liquid clocks shifting first one way and then another.

The Shadow Realm drives cartographers a bit mad, and it guarantees their work is constant. The road wardens work hard to ensure that the shadow roads stay open, connecting the castles, towns, and settlements in a familiar way, even as details change year after year. Change is relentless among the shadows, but certain elements of this plane remain fixed, true, and certain.

These fixed elements, the landmarks, major cities, headlands, and elder rivers, hold true to one another, bound together first by the shadow roads and later by all three tools of oak, road, and stone. Some ancient roads shift just one grain of sand in a century, for the roads themselves have binding properties. Some towers keep a forest in their orbit. Some windswept plains retain their slow rise and chosen, distant hills. The rangers and wanderers, the wise women and the unerring elves, know their way from place to place. For others, for visitors, or for those foolish enough to anger a geomancer, the plane and its movements are a constant struggle.

\subsection{Drifting Towers and Cities}
While some things are clearly anchored, other cities are unmoored by roads or forests. The most famous of these unmoored locales is Oshragora, the city of psychic vampires and dark deeds, said to have once been a kingdom of wizards or a region so overrun with magic that it became a magical wasteland. In either case, it was pulled entirely into the Shadow Realm by its corruption and by improvident magic.

The city has never had a road to elsewhere in the Shadow Realm (for the vampires who inhabit it seem unable to build one), and they have never placed foundation stones for their buildings or planted earthbound trees to hold it down. Thus, the city drifts slowly across the lands: sometimes near the forests of the shadow fey, other times near the town of Wormwood, and occasionally near the darakhul pits. Some claim this serves the vampires' will, for it brings them always to new places where fresh victims might live. (See \textit{Chapter 7} for details on the City Fallen into Shadow.)

\subsection{Drifting Courts}
Other famous drifting locales include the Court of Wine and Roses (built before the use of foundation stones was understood), and the Lost Court of the Triple Moon, said to have been built with a luminous tower that outshone the true moons and stars with a silver, arcane light. The Court of Wine and Roses seems entangled in the southern forests and is not difficult to find if one knows roughly where to seek. The Court of the Triple Moon has not been seen in centuries and was likely unmade by the Lord of the Hunt or one of the darker fey lords. Maps purporting to show a path to lost courts are a common grift by the more unscrupulous shadow goblins and shadow fey with visitors to the realm.

\subsection{Driftwood Villages}
Sometimes a single black oak tree or a cherry orchard or a stand of birches drifts away from the rest of the realm, taking nearby homes with it and becoming a sort of wandering village. When this happens to famous cities like Oshragora, everyone notices, but smaller settlements sometimes fade into the plains or even part from the seashore and are not seen again for many years. For most of them, the village needs a new road to bind them to neighbors and re-establish trade; others seem to thrive without the burden of gifts to the fey or tolls to the darakhul.

Several of these freeholds are well-known and seen every few years in the plains, as a river island, or elsewhere. The most famous include the fishing village of Willow Song, the woodworkers' hamlet of Rowan's Heart, the cheesemakers' hilltop of Perreville, the druids' retreat of Owlbridge, and the shepherds' village of Cloverfield. They are never marked on maps, for they never connect to a road.

\section{Oak, Road, and Stone}
The peoples of the realm have harnessed three enchanted tools to keep lands stable and to build their realms to withstand shifts in the ebon tides. These are the trinity of oak, road, and stone—the three elements that can bind shadow landscapes into a fixed form.

The rich forests first raised by Hecate and later planted by the shadow fey and bearfolk (and more recently by the darakhul) hold land together, allowing for farming, building, and hunting in regions that surround the forested lands. Their roots and branches seem to always steady the land around them.

The roads are used to bind places to one another so that they do not drift apart and become separated lands with a sea between them. Building new roads of fine flagstone or marking roads with milestones and guideposts or simply cutting through the deep brush of a white birch forest all bind together the communities and peoples of the land. Road building and road warding are both quite honorable professions.

Further, glyph-marked foundation stones are also carved and enchanted with a particular radiance to make the basis of any permanent structure. These bind a structure, such as a wall, tower, or bridge, to a place, keeping huts, castles, and palaces from moving if the shadow landscape should be drawn in one direction or another.

Natives of the Shadow Realm use this as an invocation and exclamation—"By oak, road, and stone!"—because they know the importance of maintaining them to keep their realm whole. Destroying any of them is considered an act of chaos at the very least as it unravels the habitable regions and makes navigation even more difficult than normal. Cutting roads, chopping down more timber than the druids permit, or defacing foundation stones are crimes among all the civilized races. Though they are also considered a form of sport or entertainment by the more chaotic and feral fey of the realm.

\subsection{Shadow Roads}
The shadow roads were the first tool used by natives of the Shadow Realm to keep settlements from drifting apart. One might say that their creation made the settlement of the plane by the fey possible, for without them, the courts could never find one another easily or consistently, and each new court would have remained isolated. This also explains why the shadow fey, shadow goblins, bearfolk, humans, and even the darakhul expend considerable effort to patrol and maintain the shadow roads. They also go out each spring to "walk the roads," repairing any fallen milestones, checking on bridges, restoring weak enchantments, and so forth.

However, tricksters and evildoers often work against the roads, hoping to either carve a road of their own or to deny an enemy the use of one. If a shadow road is cut by terrain-shifting magic or by the simple, muscle-powered expedient of digging a ditch or tearing up stones, the two locations that road connects on the plane become unmoored. While the ends of the cut road may bend and curl and reconnect by happenstance (or with arcane help), if the cut road is not addressed, it fades back into the mire and muck of unformed shadow or sinks into the earth, leaving no sign of its presence. And this means the places it once connected drift apart.

\begin{DndSidebar}[float=!b]{The Making of the First Roads / from the Book of Ebon Tides}
Like all those who seek to build a new road, I know that the construction of the first shadow roads is attributed to Hargen the Builder, a gnomish engineer and surveyor. This bearded arcanist learned to fix shadow matter with rune-marked milestones, carefully enchanted paving stones, and celestially inspired geometry. The task has since become rote: utilizing theodolite, lead weights, and measuring chains, I send a road builder to mark out a course, place milestones at key points (such as curves, hills, and bridges), and seed the whole with a well-seasoned course of moonstones and silver filings. To bind it, I command a priest to invoke the blessing of Hecate and the Road Gods, calling down the deep roots of the moon. In terms of the cost to the lords and ladies, each mile in simple terrain requires at least 100 gp and as much as a month's work to build and almost as much to maintain—and five or even ten times that in marsh, hills, or deep forest. The expense eats at our treasury!



\end{DndSidebar}

\subsection{Creating a New Road}
To pave a new shadow road or even to mark it out with milestones and crossroads is a work of many hands. The roles are generally divided into three parts: the survey, the laying of stones, and the binding of points.

\subparagraph{The Survey}
To begin the work, a surveyor (often a wizard or cleric) takes silver chains, a theodolite, and a notebook to the foundation stone where the road is meant to begin. This is often a crossroads to a new path or an informal forest path, a trail in the tall grass, or a muddy stretch of river or hills that has been useful for a long enough period that the inhabitants will pay for a more permanent stretch of road to be mapped out. The survey involves putting down silver markers, measuring the current position of all curves, turns, and crossroads, and often leaving a long thread or, in the case of more ambitious roads, a bit of thin silver chain to mark the path the road is meant to take.

\subparagraph{Laying of Stones}
Once the course is plotted, time becomes critical. The road's grasses, trees, and swampy patches must be addressed with traditional gravel fill, full flagstones over a roadbed, a rumbling way of wooden slats, making for a jarring route for wagons and carts, or simply the clearing of brush and the erection of markers, waypoint stones, and guide poles to show the way. This work is usually done by shadow goblin or hobgoblin gangs, when the road is a new one for the fey, or by young bearfolk, erina, and even molefolk if the road is meant to connect bearfolk lands.

\subparagraph{Binding of Points}
The last step is the arcane one that binds the foundation stones, the silver filings, and a crumble of prayers and chunks of opal or moonstone into the road, making it impervious to destruction by the slow warping of the ebon tides. The shadow road can still stretch, bend, warp, and turn to accommodate the drift of the two or more points that its binding ritual mentions. The ritual is most often performed by shadow fey clerics of Hecate together with shadow goblin mooncatchers and moon priests of various kinds, though in recent years this binding has also become an important facet of the rise of the bearfolk, whose druids understand its depths and difficulties, and even the darakhul, whose binding involves a bit more bone and blood and a bit less moon and silver. (See also the \textit{deep roots of the moon} spell in \textit{Chapter 4}.)

\subparagraph{Life of a Shadow Road}
Once the binding is complete, the road remains until destroyed or unbound. In some cases, these shadow roads are little more than forest trails, connecting a village or standing stones to a larger settlement. At other times, the road is entirely paved and remains so even when the relative positions of various settlements changes or when new hills, forests, or rivers cut through the land. The road itself is a living thing, and it can (and does) stretch to accommodate changes in the shadow terrain.

Note that various void cults and the satarre have often attempted to bind Oshragora into the larger realm by such a road, but so far, all such attempts have been thwarted by either the fey or the bearfolk. The prospects of binding it to the greater Shadow Realm seem ill-omened, at best.

\subsection{Road Worship}
Some shadow roads are sites of worship, associated with ancient gods that predate the arrival of the fey. These roads are called "pilgrim roads" or "old paths," and they include some of the oldest, best warded, and most revered roads of the realm, including the road from Fandeval to the Court of Night and Shadow, and the small stretch from the Hawthorne Grove to the city of Solvheim, which is built on the ruins of an ancient fey court.

The shrines themselves often combine a small pantheon, with the Old Roads worshipped next to Hecate, Charun, and sometimes Bengta or Ailuros. They may be little more than a small stone- or tile-roofed cabinet at the side of the road or a small cave that provides shelter. Some resemble a shepherd's hut or an especially convenient boulder or cliffside. Offerings include worn-out boots and sandals, candles, and similar objects.

For more information on travel and use of shadow roads over long distances, see \textit{Chapter 5}.

\subsection{Forest Anchors}
Forest anchors are a way to secure a region filled with streams, villages, and other scattered landmarks from going drift. Most trees able to grow in the Shadow Realm also hold the land together, though not in the same way that roads do. Forests are less like a cord binding two places that float on the ebon tides and more like a matt of reeds or kelp, keeping a larger area calmer, disturbed only by great storms, and moving as a single unit. A great forest might be said to float on the ebon tides, but it drifts with all its trees, courts, castles, villages, streams, and everything else with them—an island. For this reason, even the darakhul plant their strange, crackling forests and similar structures.

Spells such \textit{hallucinatory terrain} and \textit{mirage arcane} can be used to generate forest lands that will hold for longer than usual in the Shadow Realm (3 days rather than 24 hours for \textit{hallucinatory terrain} and a month rather than 10 days for \textit{mirage arcane}). Bearfolk druids and priests of Hecate can "lock in" the results of such spells to create new shadow lands with the \textit{deep roots of the moon} spell (see \textit{Chapter 4}).

\subsection{Foundation Stones}
Fortunately, courts, castles, and villages can be built to withstand the ebon tides rather than being shifted or cracked by the earth's slow movement. This requires the use of radiant foundation stones—enchanted stone pilings marked with the elements of a \textit{daylight} or a \textit{darkness} spell and used either as single stones under the floor (for a hut, tavern, or small barn, for instance) or in sets of two or four stones at the corners (for a large inn, manor house, or a great chapterhouse or keep) or even more (for many-towered castles, large monasteries, great courtyard homes or cloisters, and arcane colleges).

Breaking a small foundation stone requires a \textit{shatter} spell or 20 points of bludgeoning damage. Breaking a medium-size foundation stones requires a \textit{dimension door} spell or 50 points of bludgeoning damage. Breaking a large one requires a \textit{disintegrate} spell or 150 points of bludgeoning damage. In all cases, when a foundation stone crumbles, creatures within the structure feel a quivering, shaking motion for a moment, like a small earthquake.

\section{Shadow Realm in Motion}
How does a group of travelers know when the black tides are running and when the forces of shadow are warping the mountains, rivers, and roads? There are many common signs, all well known to natives of the Shadow Realm. They may indicate a small re-alignment of a nearby creek or the vast shift of an entire set of grasslands or the motion of a village or city that has come unmoored from road or forest. The following list is useful for identifying an ebon tide.

\begin{DndTable}[header=Table 1-1: Ebon Tide Manifestation]{cX}

d10 & Result\\
1 & Strange chatter in the wind and voices just out of hearing.\\
2 & Currents under stone and the sound of water far from any river.\\
3 & Echoes of the past: ghostly figures, forests, even cities appear and disappear.\\
4 & Black wind, a bit like a very thick fog, reducing vision to 30 feet.\\
5 & Rolling earth, like a mild earthquake or the presence of an earth elemental, and terrain spins or tilts.\\
6 & Shadow trembling: all light and shadow seem to quiver as if someone is shaking the moon.\\
7 & Sound echoes through the sky, feeling like half a crashing wave, half the banging of a massive gong.\\
8 & Shadow thickening: all shadows in the region grow more dense and make \textit{terrain difficult} for an hour.\\
9 & A new radiant well opens nearby, casting a beam of light into the sky.\\
10 & A blaze of light passes through the party like a wave, blinded everyone for 1 round.\\
\end{DndTable}

\subparagraph{Reading the Shift}
A successful DC 14 Wisdom (Survival) check by a member of the party or their guide will confirm how far and in what direction the ground of the Shadow Realm is shifting. On a failure, the party becomes lost for at least a full day of travel. A new check may be made the following day as new landmarks come into view. This ability check is made with advantage if the party is traveling on a shadow road.

\begin{DndSidebar}[float=!b]{Where is the Shadow Realm?}
While some sidebars in this volume refer to the mortal world as Midgard, the choice of which planes are close by is entirely up to the GM. It could be a published campaign world that you especially enjoy or a setting entirely of your own creation. The gloom and flicker of the Plane of Shadow does, after all, reach any place where light and darkness mingle. As GM, you can make room for shadow fey and shadow roads next to any existing campaign.



\end{DndSidebar}

\section{Five High Holidays of Shadow}
With several major and distinct cultures spread across the Shadow Realm, it's rare that holidays are celebrated in all places in the same way. However, five main holidays are widely acknowledged by most in some way—or at least accorded some level of polite observance. These are described here.

\subparagraph{Twin Moon Festival}
All peoples of the Shadow Realm rely on the twin moons for light, seasons, and harvests, and the occasion of the twin full moon of both Celene and Lindar is always cause for a celebration of life, hope, and radiant feasting. The shadow fey, goblins, and bearfolk strive to outdo each other in celebrations with bonfires and dances. Humans and darakhul are less keen but observe the day's events: usually a torchlight procession through a settlement to the fields or to a shrine of Hecate, followed by a short blessing and then a celebratory feast featuring sweet cakes.

\subparagraph{Double Dark Blessing}
When both moons are new and raw and darkness covers the plane, the dark gods and their cults celebrate as do many evil spirits, haunts, and shades. The festival is celebrated with cold wines in Soriglass and with fiery offerings in the temples of Charun, Alquam, and Anu-Akma. Elsewhere, it is celebrated in much darker fashions—with blood and bone and with animal sacrifices.

\subparagraph{River Soul Harvest}
When the River Styx crests above its usual height, the priests of Charun often declare a River Soul Harvest Day, when moonlight nets and silver hooks are used to fish various items from the waters and some daring followers of the boatman ask for a sprinkling of river water, all to provide good fortune for months to come. Fish, oysters, clams, and seaweed are common feast foods and offerings. The celebration is especially well-marked by the shadow goblins of Fandeval, though others partake as well, if only for a bit of witchlight or enchanted flotsam to brighten one's home.

\subparagraph{Marsh Mustering (Witchlight Day)}
Popular among children near any marsh or body of water, the Marsh Mustering is said to have begun as a rite of the King of Toads, an old fey lord otherwise largely neglected. Held in spring when the marshes are active with water birds, the druids encourage the pursuit of marsh creatures by the young. A goodly catch of active, healthy frogs, dragonflies, witchlights, and other creatures is an omen of a good harvest to come. Cakes and candies made in the shape of stars, frogs, and ducks are often part of the celebrations along with sugared eggs and sweets made with marsh roots.

\subparagraph{Feast of All Souls}
Celebrated primarily by humans, goblins, and bearfolk, the Feast of All Souls is a remembrance and memorial for the dead, when grave goods are offered at various graveyards and burial grounds and candles burn in vast numbers in memory of the departed. For some shadow fey, this is an occasion to show one's status and wealth by burning candles embellished with gold leaf or silver pearls. The darakhul offer a very different version of this holiday, seeing it more as a celebration of rebirth, of shedding one's mortal days for a second life.

\begin{DndSidebar}[float=!b]{Famous Shadow Roads and Portals of Midgard}
The most famous shadow roads are the original Black Road to Zobeck, the Lantern Road between the two sides of \textit{Corremel}, and the Cat's Road near \textit{Fandeval}. But many other such roads exist as well as portals to Siwal, Morgau, Triolo, and the Scarlet Citadel.



For those shadow roads known by more than a rarified few, the best-known ones are these four:




\begin{itemize}
\item Mage's Road from the \textit{Court of Midnight Teeth} to the mage city of Bemmea
\item Ailuros Road from that goddess's temple in \textit{Merrymead} to Bourgund
\item Flower Road from the \textit{Court of Pale Roses} to Carnessa
\item Moon Road from \textit{Fandeval} to Kammae
\end{itemize}



Other roads are decidedly more obscure, such as the Trickster's Road to Skaldhome (known primarily among the bardic College of Shadows and some of the elderly bearfolk), the Maze Road from the Echoing Shore to Capleon (known among minotaurs and some shadow goblins), the River Road in the Arbonesse (which continues on to brighter realms), and the road to the gnomish city of Hexen in Niemheim, best known as ending in a mirror in the Court of Midnight Teeth.



Not all roads or portals lead to the mortal lands of Midgard, of course. Connections to the planes are also well represented, including to Silendora and the Summer Lands of the elves, to Yggdrasil through the \textit{Court of the Golden Oak} and various bearfolk roads, and frighteningly direct connections to Evermaw, Niflheim, and the Luminous Hell of Fire and Sulfur via roads in the \textit{Twilight Empire}.



\end{DndSidebar}

\subsection{Tides and Travel}
If you learn how to walk the shadow roads, the Shadow Realm can be your playground. Anger the roads, and the wind will always be at your face, the forest always dark, the path always sorrowful. Dark magic infuses the realm, and that means that temptations corrupt the mortals who visit here (see \textit{Shadow Corruption} in \textit{Chapter 5}).

In the Shadow Realm, you must listen, listen hard, to hear the tides coming. When the ebon tides run, distances change, just as shadows bend and stretch as sources of light move. So travel from one point to another can be quite predictable for weeks, months, or even years and then change suddenly during a time of strong tides. In other cases, the ebb and flow of the shifting land is much less obvious. Travel times between familiar places might change by a few minutes or even an hour but not enough to change patterns of trade or familiar routes for road wardens or pilgrims.

Natives refer to "high tide" travel when conditions stretch the roads, and the distances are longer than usual. They refer to "low tide" travel when conditions make distances collapse and visitors arrive faster than expected. Neither is considered especially unusual.

When traveling along shadow roads, these effects are much less notable than when traveling overland. For wilderness travel in either mode, make a Wisdom (Survival) check on the \textbf{Drift and Tides Travel} table when the ebon tides are strong. (This table can and should be ignored for travel during times when the plane is quiet and shifts are minimal.)

\begin{DndTable}[header=Table 1-2: Drift and Tides Travel]{cXX}

d20 + #$prompt_number:min=0,title=Survival Check$# & Shadow Road Variation & Overland Variation\\
1 & Travel time increased by 100\\% & Travel time increased by 200\\%\\
2–5 & Travel time increased by 50\\% & Travel time increased by 100\\%\\
6–9 & Travel time increased by 10\\% & Travel time increased by 50\\%\\
10–14 & Travel time as usual & Travel time increased by 10\\%\\
15–18 & Travel time reduced by 10\\% & Travel time as usual\\
19 & Travel time reduced by 20\\% & Travel time reduced by 10\\%\\
20 & Travel time reduced by 30\\% & Travel time reduced by 20\\%\\
21+ & Travel time reduced by 50\\% & Travel time reduced by 30\\%\\
\end{DndTable}

\begin{DndSidebar}[float=!b]{Tidal Maps}
The best-loved maps of the Shadow Realm are tidal maps, so called because they show the lands both at their widest-flung extent, called High Tide, as well as how the various landmarks are positioned during years or months of scrunched-together conditions where sites cluster nearer, called Low Tide. These maps are enchanted to show the current state of such things, but the fey consider them fairly mundane objects: useful surely but not the sort of magical item that confers much status. Visitors to the realm consider them a vital necessity to get around during times of high tidal flux.



\end{DndSidebar}

\subsection{Sky and Waters}
Just as the land of the ebon tides is different, so are its sky and waters. Moonlight and twilight are the nearest analogies. There's never a full noon sun or even proper daylight, though muddy light is common enough. Likewise, a coal-black night is almost impossible: stars, comets, luminous streams and clouds, and even a diffuse glow from fog and land all renders some light. Colors wash out quickly, except in places lit by magic or firelight. Colorful items brought to the shadows from the mortal world sometimes fade permanently, remaining touched by the muted tones of shadow even when taken back to mortal lands.

\subparagraph{Shadow Waters}
Some of the Shadow Realm's waterways—streams, lakes, and rivers—are a bit more disquieting and certainly more dangerous than the sky. These streams, lakes, and rivers contain rushing cascades, or ribbons, of pure shadow stuff, especially in the spring when shadow itself seems to flow down the hills with each gray rainfall. Two major waterways twist like lengths of black silk across the plane: the Rivers Styx and Lethe. Their tributaries and minor outlets reach nearly every corner of the realm, providing accessways if one can survive the perils of a waterborne journey. The shores of the Shadow Realm look out on the inky expanse of the Shade Sea, an endless, roiling stretch of wavering darkness to the south, and the chill waters of the Frost Gulf to the north.

A DC 13 Wisdom (Survival) check is enough for a native of the Shadow Realm to distinguish hazardous conditions in a body of water (this check is made with disadvantage for visitors to the plane). A success reveals the danger while a failure means that the water is thought to be normal, potentially leading to shadow corruption (see \textit{Shadow Corruption} in \textit{Chapter 5}) and similar ill effects.

\subsection{Mirrors and Echoes}
If one could look down from a star's-eye view, the Shadow Realm would resemble the mortal world as seen through a smoked-glass mirror. Prominent landmarks have echoes in the Shadow Realm, often playing host to similar or sometimes opposite essences. For instance, a population center in the mortal world might correspond to a city or town where natives of the umbral plane live and scheme. Or it could hold a blighted desolation of tumbled stones and salted fields where a settlement once stood.

\subparagraph{Warped Resemblances}
Despite the overall similarities, a keen understanding of mortal world geography does not provide a key to navigating the Shadow Realm. The plane stretches and twists, distorted in dimension and shape when compared to places in the material world. Distances can deceive an unwary traveler, and directions become confused. Relationships between places in the Shadow Realm even change over time, the way a shadow lengthens and shrinks as the sun travels through the day. A mountain pass might lead to a vast forest one year and to a desert of black sand the next. While the distortion of the landed regions of the Shadow Realm can be vexing enough, the vast emptiness of the Shade Sea is completely unpredictable. Few crews brave or unfortunate enough to venture out on the open Shade Sea have ever returned, and those that did were never the same.

Rare are those visitors that willingly travel to the Shadow Realm, and most have at least some understanding of the place before they cross that threshold. From mighty wizards who weave magic to bear them beyond the light to brash explorers who discover the secrets of a forgotten passage into darkness, these adventurers should know what awaits them. Some cultures nurture a relationship with the Shadow Realm, though these are often misunderstood or rightly feared.

\begin{DndSidebar}[float=!b]{Gods of the Old Roads}
The ancient roads of the Shadow Realm are one of the secrets to understanding how the Shadow Realm's inhabitants thrive, and they have a number of quirks and customs. In particular, the roads themselves have minor gods, road saints akin to the river spirits and fey lords and ladies that rule small domains, landscapes, or regions. Each has a name, and they are often found in pairs. (See \textit{Chapter 8} for more information on their names and powers.)



The road saints are called the "gods of the Old Roads," and most of them seem to be very old indeed. While they have no priesthood other than the Doorlord, they are still revered in small ways by travelers seeking respite. The voices of the gods of the Old Roads can be heard at their shrines, small statues with stone rooves, found most often at crossroads where a shadow road splits into a shadow path and a side path, often called a shadow door or portal path. At these places, travelers can rest in relative safety by making a small offering at the shrine and remaining close to it. Some part of the consecration of the sites still holds strong.



A few of the gods of the Old Roads have old names given them in ages past and marked in runes on some of their shrines: Clauveg, Eiber, Hune, Mantellon, Raggel, Viacorra, Vanta. When their names are invoked, it is often in the context of a \textit{shadow portal} spell (see \textit{Chapter 4}).



\end{DndSidebar}

\section{Etiquette and Status in the Shadow Realm}
The Shadow Realm is home to many unique creatures and species, but it is arguably best known as the home of the shadow fey. The shadow fey resemble other humanoids, such as humans or elves, but their culture is quite alien to most who live in the mortal world, being full of changing and confounding rituals, edicts, and rules. While this is true, that some shadow fey preferences and mores change over time, certain facts do remain constant.

\subparagraph{The Courts of the Fey are the Foundation of the Realm}
Many lords and ladies of the shadow fey, sable elves, gnomes, and others claim to rule the Shadow Realm, but most such claims only apply to a fragment of the whole or to a particular time or season. (For instance, the Moonlit King is in power during the Winter Court's rule, and Sarastra is in power during the Summer Court's rule.) If an individual acts in a manner as prescribed by the fey rulers, its time in the courts and when dealing with members of the courts is often easier and safer.

\subparagraph{Wheels Within Wheels}
When dealing with the fey, the wiser heads realize that what they say is not necessarily what they mean. The fey are immortal creatures of whimsy and caprice. They enjoy watching mortals struggle to comprehend their true intent.

\subparagraph{Status Matters}
The courts are impenetrable labyrinths or empty shells to creatures who simply lack the status and importance required to interact with them. There are ways to increase or decrease one's esteem in the eyes of the courts (see the \textit{Shadow Realm Status Adjustments} table below and the Midgard Worldbook for more on the optional Status rule).

\subparagraph{The Fey Envy Mortals}
While few of the capricious creatures would admit to it, the shadow fey and gnomes are jealous of mortals. They view mortal experiences as being both truer and more primal than the experiences of longer-lived, more sophisticated races.

\subparagraph{Harming a Servant is Dishonorable}
While the shadow fey may or may not like or trust their servants, they do rely upon the creatures to make their lives comfortable. Harming or killing a servant is therefore abhorrent to them.

The following are a dozen examples of shadow fey etiquette and mores that are mutable and might change depending on the week or month or on which court is in power at that time:


\begin{itemize}
\item Do not eat root vegetables, for they are disgusting and offensive to the sensibilities of all civilized creatures.
\item Do not harm any shadow fey outside of an honorable duel.
\item Do not run, for it causes one to sweat like a common laborer.
\item You must have your presence announced in song before entering a room.
\item Civilized folk wear powdered wigs.
\item Clothing made of sackcloth is all the rage, and anyone not wearing a sackcloth dress or breeches is to be considered a fool.
\item True gentlefolk duel with insults, not steel.
\item Hunting is only to be done on the sixth day of the week.
\item Only pretentious twits play cards.
\item All the ladies of the court must use illusions to appear human.
\item Boasting of one's own achievements is never acceptable on the second day of the week.
\item One above a certain station must have exactly thirteen servants.
\end{itemize}

\begin{DndTable}[header=Table 1-3: Shadow Realm Status Adjustments]{cX}

Status & Action\\
−8 & Threaten or attack a fey noble.\\
−3 & Perpetrate arson or vandalism against a member of a court or their holdings.\\
−3 & Get caught brawling or otherwise causing a violent scene in the court.\\
−2 & Be fooled by a shadow fey illusion.\\
−2 & Flee a member of the court who seeks your arrest or imprisonment.\\
−2 & Disobey a direct order or deny a request from a member of the court.\\
−2 & Be seen sneaking over the palace walls.\\
−1 & Be the victim of a shadow fey prank.\\
−1 & Kill more than one shadow goblin or other shadow fey servant.\\
−1 & Lose a fey consort's affections.\\
−1 & Threaten a courtier to gain an audience with the court.\\
+1 & Defeat a fey creature with a challenge rating equal to or greater than your level in single combat. This can be done a maximum of 5 times.\\
+1 & Successfully use Charisma (Deception, Intimidation, or Persuasion) on a fey creature with a challenge rating equal to or greater than your level to gain information or a favor. This can be done no more than once per shadow fey and no more than 5 times total.\\
+1 & Successfully learn a fact about a specific shadow fey or about the workings of the courts with an Intelligence or Charisma check. This can be done no more than 5 times total.\\
+1 & Successfully play a prank on a fey creature with a challenge rating equal to or greater than your level. This can be done no more than once per shadow fey. Continually pranking the same shadow fey likely results in reprisal and the potential loss of status.\\
+1 & Take a shadow fey consort.\\
+1 & Receive a gift from a member of the court.\\
+1 & Receive an invitation to dine or hunt with a member of the court. This benefit can be gained only once from any given member of the court.\\
+2 & Join a shadow fey organization.\\
+2 & Save a member of the court's life or dignity.\\
+5 & Be acclaimed by the \textbf{Queen of Night and Magic} as her champion. This benefit can be gained only once.\\
\end{DndTable}

\chapter{2. Umbral People and Heroes}
The Shadow Realm is home to both strange creatures and ordinary folks, all of them shaped by the land in one way or another. Among them, five groups have risen as the primary powers, each struggling to control the ebon tides and maintain their holdings. These are the \textit{bearfolk}, \textit{darakhul}, \textit{shadow fey}, \textit{shadow goblins}, and \textit{umbral humans}. Each group rules their corners of the realm, governing settlements, constructing roads, and generally existing in large enough numbers to maintain their towns and cities.

Other groups—\textit{quicksteps}, \textit{erina}, and \textit{wyrd gnomes}—had a foothold in the realm from the beginning as servitors of the shadow fey, and while many still fill that role, they have also carved out more modest holdings of their own. The \textit{ratatosk} of course have a small presence everywhere that the World Tree's branches reach, and the \textit{Stygian shades} are simply the Shadow Realm's unique form of the shades that appear anywhere that mortals dwell. And of course, a few peoples have come new to the shadows, only recently attracted to its mystery and power for various reasons, including the \textit{ravenfolk} and the \textit{satarre}.

\begin{DndSidebar}[float=!b]{Peoples of Fantasy}
Every fantasy world is filled with a wide variety of peoples and civilizations. Here you'll find new races and cultures from bearfolk to darakhul to shadow goblins. You'll also find new subraces of existing races, such as the wyrd gnomes and the sable elves.



Each race has statistics that provide a guideline for a typical member of that race. Not every individual exhibits these traits, be it in alignment, training, or disposition. If you want to create a character different from a typical member of their race, go for it! Players, be sure to work with your GM to determine which ability scores, alignment, and traits would best fit that character and the story you and your GM want to tell about the character.



\end{DndSidebar}


\begin{itemize}
\begin{multicols}{3}
\item Bearfolk
\item Bearfolk (Shadowborn)
\item Darakhul
\item Darakhul (Bearfolk Heritage)
\item Darakhul (Shadow Goblin Heritage)
\item Elf (Shadow Fey)
\item Elf (Lunar)
\item Elf (Sable)
\item Shadow Goblin
\item Umbral Human
\item Umbral Human (Changeling)
\item Umbral Human (Gifted)
\end{multicols}

\end{itemize}

\section{Other Peoples of the Shadow Realm}
The Shadow Realm attracts a motley spectrum of the strange, the dispossessed, the curious, and the lost, and some of those races have created small communities. While not great in overall number, these groups are often influential in a particular court, city, or region. They can be brought into the spotlight or ignored by the GM in any particular Shadow Realm adventure or campaign.


\begin{itemize}
\begin{multicols}{3}
\item Quickstep
\item Ratatosk
\item Ratatosk (Ekorre; Clever Tusk)
\item Ratatosk (Tradvakt; Tree Protector)
\item Erina (Spiritfarer)
\item Shade (Stygian)
\item Ravenfolk (Sublime)
\item Satarre (Unbound)
\item Gnome (Wyrd)
\end{multicols}

\end{itemize}

\chapter{3. Heroes from the Shadows}
Heroes native to the Shadow Realm are often quite distinct. Their powers are attuned to drawing from shadow or to resisting it, and many heroes know the allure of darkness all too well. This chapter provides ten new heroic subclasses appropriate to the plane.


\begin{itemize}
\begin{multicols}{3}
\item Druid
\item Bard
\item Cleric
\item Sorcerer
\item Warlock
\item Wizard
\item Cleric
\item Barbarian
\item Rogue
\item Monk
\end{multicols}

\end{itemize}

\chapter{4. Magic in the Shadow Realm}
The Shadow Realm resonates with magic, and the land itself is but shadow matter, a source of arcane shaping for those who understand it. And creatures born into or corrupted by shadow may develop unique magical abilities, possessing great power over illusion, enchantments, and the manipulation of shadow.

With magic woven into its very nature, the plane is fertile ground for some schools and styles of magic. At the same time, spells and other magic used within its borders can produce unexpected effects. Most of these effects appear unsettling but are merely cosmetic. Spells that create light, for instance, glow with a sickly green, blue, or purple illumination. Energy created by spells pulsates with black veins, writhes in dark flames, or crackles with black lightning. Conjured animals are drawn from those typically found in the Shadow Realm, possessing pale or dark coats and coldly gleaming eyes.

\section{The Magic of Shadow and Light}
Apart from leaving its mark on magic brought into the domain, the Shadow Realm also spawns new and enigmatic magic. Spellcasters seeking paths to power and shadow fey living in the depths of shadow have developed two unique disciplines of magic: the studies of illumination and shadow magic. In addition, the darakhul of the Shadow Realm have a number of necromantic spells unique to the realm, and bearfolk and shadow goblins command several lunar spells involving the arcane use of moonlight and the manipulation of wind, rain, and fog. The Shadow Domain is also unique to the plane (see \textit{Chapter 3} for details).

\subsection{Umbral Magic: Empowered Illusions}
The raw matter of shadow has always been a source of arcane power for the schools of illusion and shadow magic. This shadowstuff's role in illusion and phantasm is much greater in the Shadow Realm, where a caster can draw on the material all around them to make their spells stronger, longer-lasting, harder to see through, or simply useful in new and unexpected ways. Here are some optional effects for existing spells, empowering illusions and shadow magic—but only when such spells are cast in the Shadow Realm.

Many spells with the \textbf{shadow} or \textbf{illumination} keywords in their description are empowered when cast in the Shadow Realm. This includes such spells from this volume as well as those of shadow and illumination magic found in \textit{Deep Magic} (which are marked below with an asterisk):


\begin{itemize}
\item \textbf{\textit{Black Ribbons*}.} This spell affects a 30-foot cube instead of a 20-foot cube.
\item \textit{Evard's Black Tentacles}. This spell affects a 30-foot cube instead of a 20-foot cube.
\item \textit{Blur}. This spell is effective against opponents with truesight.
\item \textit{Color Spray}. At higher levels, you add 2d12 rather than 2d10 to its damage.
\item \textit{Creation}. You can create up to a 10-foot cube of new material. The material you create lasts longer, as shown here: vegetable matter, 1 week; stone or crystal, 1 day; precious metals, 4 hours; gems, 1 hour; adamantine or mithral, 2 minutes.
\item \textit{Dancing Lights}. In the Shadow Realm, you can create up to six lights and combine them into two glowing humanoids. Further, a caster of 5th Level or higher can split them into two groups.
\item \textbf{\textit{Dark Path*}.} This spell creates a path or bridge of 100 feet rather than 50 feet.
\item \textbf{\textit{Douse*}.} This spell can affect a campfire in addition to a lantern or torch.
\item \textit{Hallucinatory Terrain}. This spell can be used to generate terrain that will hold for longer than usual in the Shadow Realm (3 days rather than 24 hours). In addition, one built structure in the area of effect can be hidden by fog or shadows by the use of this spell.
\item \textbf{\textit{Legion*}.} Your shadowy soldiers fill two 10-foot cubes rather than one.
\item \textit{Minor Illusion}. This effect can last up to 2 minutes rather than 1. In addition, objects created can create a simple sound, such as creaking, chirping, or moaning (but not language).
\item \textit{Mirage Arcane}. This spell generates terrain that holds for longer than usual in the Shadow Realm (a month rather than 10 days). Any piece of the illusory terrain (such as a rock or stick) that is removed from the spell's area remains solid, made real by the power of shadow.
\item \textit{Phantasmal Killer}. This spell does 4d12 rather than 4d10 damage. In addition, when you cast this spell using a spell slot of 5th Level or higher, the damage increases by 1d12 for each slot level above 4th.
\item \textit{Phantom Steed}. The creature uses the statistics for a \textbf{warhorse}, except it has a speed of 120 feet and can travel 12 miles in an hour or 15 miles at a fast pace.
\item \textit{Programmed Illusion}. This spell can be used to create an illusory program of up to 10 minutes rather than 5, and speaking creatures within the program appear to listen when addressed.
\item \textit{Seeming}. This spell changes creatures' appearance for up to 1 day rather than 8 hours.
\item \textit{Shadow Gateway}\textbf{\textit{*}.} Casting this in the Shadow Realm opens a portal to the mortal world rather than a portal to the Shadow Realm.
\item \textit{Shadow Step}\textbf{\textit{*}.} You can move up to 120 feet per shadow step, and this spell does not require \textit{concentration} in the Shadow Realm.
\item \textbf{\textit{Shield of Starlight*}.} This spell grants protection from both radiant and necrotic damage.
\item \textit{Silent Image}. This spell does not require \textit{concentration} when cast in the Shadow Realm.
\item \textit{Simulacrum}. The simulacrum has 3/4 the creature's hit point maximum and is formed with clothes and one simple tool.
\item \textit{Sunburst}. Casting this spell creates a radiant well (see \textit{Chapter 5}) in the square where it is centered. This lasts for 1d6 days, and it may become permanent if further measures are taken to strengthen it (such as with the \textit{daylight} spell or the like).
\item \textit{True Seeing}. This spell requires \textit{concentration} to maintain in the Shadow Realm.
\end{itemize}

\begin{DndSidebar}[float=!b]{Shadow \& Illumination Magic}
Illumination magic deals in duality, befitting its origin among the shadow fey. Utilizing the extremes of light and darkness, illumination magic draws on portents written in the stars to glimpse the future, channel starfire, and more safely control the powers of shadow.



The dark practice of shadow magic, however, owes its origin to mortal sorcerers, wizards, and warlocks who craved power and thought they understood the risks inherent in obtaining it. Unlike the light and darkness bent through illumination magic, shadow magic taps into the raw stuff of shadow and brings it to bear. This magic is never without a price.



\end{DndSidebar}

\subsection{GM Note: Manipulating Shadow Landscapes}
Spellcasters in the Shadow Realm often describe their magic in terms of gathering shadow or radiant energy and of pulling arcane power from the land or from the skies. In particular, much of shadow magic involves gathering the flow of shadow matter into a caster's hand, pulling it into a vessel, and then bending or shaping that matter to suit one's purpose. You should feel free to describe villain, monster, and NPC magic in similar terms.

Additionally, the landscape itself sometimes holds the power of raw shadow, a thrumming sense of potential, of deception or wild transformation. Spellcasters from other lands claim that magic can be smelled or almost touched in the Shadow Realm, especially illusion, illumination, and necromantic magics. And in truth, spellcasters using \textit{detect magic} can see raw pools of shadow, often in dark water or shadow mire, along luminous paths, or spread across a marsh like a web. Witchlights and will-o'-wisps seem drawn to these forms of raw shadow power. (For more on dark water and other hazards of the Shadow Realm, see \textit{Chapter 5}.)

More settled power can be seen around shadow roads, foundation stones, and similar binding places, including \textit{hallucinatory terrain}, \textit{mirage arcane}, and the like. When such settled shadow is given form, it tends to hold and keep that form unless unbound and dissolved by radiant spells or by detect magic. (For more on the creation of shadow roads, see \textit{Chapter 1}.)

\begin{DndSidebar}[float=!b]{Optional Rule: Radiance Dispelling Shadow}
As an optional rule to reduce the power of shadow mages in the Shadow Realm, you might decide that a \textit{daylight} spell can act as a \textit{dispel magic} against all spells with the \textbf{shadow} keyword in its area of effect. Producing light this way is certainly considered hostile among shadow goblins and shadow fey—but they use this trick among themselves as well.



\end{DndSidebar}

\section{Spells Lists}
The spells presented here were created by the shadow fey, bearfolk, darakhul, shadow goblins, and others within the Shadow Realm. Check with your GM before selecting these spells.

\subsection{Bard Spells}

\begin{description}
\begin{multicols}{2}
{\par \vspace { 9pt plus 3pt minus 3pt } \noindent  \DndFontTableTitle{Cantrips (0 Level)} \nopagebreak }

\begin{description}
\item \textit{Bright Sparks}
\item \textit{Drayfn's Bane of Excellence}
\item \textit{Obfuscate Object}
\end{description}

{\par \vspace { 9pt plus 3pt minus 3pt } \noindent  \DndFontTableTitle{1st Level} \nopagebreak }

\begin{description}
\item \textit{Doom of Poor Fortune}
\item \textit{Elf Shot}
\item \textit{Pratfall}
\end{description}

{\par \vspace { 9pt plus 3pt minus 3pt } \noindent  \DndFontTableTitle{2nd Level} \nopagebreak }

\begin{description}
\item \textit{Drayfn's Blunted Blade}
\item \textit{Leiloch's Irritating Kazoo}
\item \textit{Ominous Winds}
\item \textit{Ugly Duckling}
\end{description}

{\par \vspace { 9pt plus 3pt minus 3pt } \noindent  \DndFontTableTitle{3rd Level} \nopagebreak }

\begin{description}
\item \textit{Faerie Toast}
\item \textit{Lost}
\item \textit{Shadow Portal}
\end{description}

{\par \vspace { 9pt plus 3pt minus 3pt } \noindent  \DndFontTableTitle{4th Level} \nopagebreak }

\begin{description}
\item \textit{Drayfn's Curse of Incompetence}
\item \textit{Fire Dance}
\item \textit{Polychromatic Bubble}
\item \textit{Spider Song}
\end{description}

{\par \vspace { 9pt plus 3pt minus 3pt } \noindent  \DndFontTableTitle{5th Level} \nopagebreak }

\begin{description}
\item \textit{Feast of Flesh}
\item \textit{Hero of Fable}
\item \textit{Leiloch's Interminable Yarn}
\item \textit{Lost and Wandering}
\end{description}

{\par \vspace { 9pt plus 3pt minus 3pt } \noindent  \DndFontTableTitle{7th Level} \nopagebreak }

\begin{description}
\item \textit{Charnel Banquet}
\item \textit{Leiloch's Arduous Shuffle}
\end{description}

{\par \vspace { 9pt plus 3pt minus 3pt } \noindent  \DndFontTableTitle{8th Level} \nopagebreak }

\begin{description}
\item \textit{Doom of False Friends}
\end{description}

\end{multicols}

\end{description}

\subsection{Cleric Spells}

\begin{description}
\begin{multicols}{2}
{\par \vspace { 9pt plus 3pt minus 3pt } \noindent  \DndFontTableTitle{Cantrips (0 Level)} \nopagebreak }

\begin{description}
\item \textit{Bright Sparks}
\item \textit{Drayfn's Bane of Excellence}
\item \textit{Flowering}
\end{description}

{\par \vspace { 9pt plus 3pt minus 3pt } \noindent  \DndFontTableTitle{1st Level} \nopagebreak }

\begin{description}
\item \textit{Doom of Poor Fortune}
\item \textit{Emerald Goblet}
\item \textit{Fireflies}
\end{description}

{\par \vspace { 9pt plus 3pt minus 3pt } \noindent  \DndFontTableTitle{2nd Level} \nopagebreak }

\begin{description}
\item \textit{Conjure Ferryman}
\item \textit{Drayfn's Blunted Blade}
\item \textit{Ominous Winds}
\end{description}

{\par \vspace { 9pt plus 3pt minus 3pt } \noindent  \DndFontTableTitle{3rd Level} \nopagebreak }

\begin{description}
\item \textit{Deep Roots of the Moon}
\item \textit{Faerie Toast}
\item \textit{Lost}
\item \textit{Orros Mark of Fate}
\end{description}

{\par \vspace { 9pt plus 3pt minus 3pt } \noindent  \DndFontTableTitle{4th Level} \nopagebreak }

\begin{description}
\item \textit{Drayfn's Curse of Incompetence}
\item \textit{Knife of Fate}
\item \textit{Moonlight Sending}
\item \textit{Portho's Portal}
\item \textit{Storm Door}
\end{description}

{\par \vspace { 9pt plus 3pt minus 3pt } \noindent  \DndFontTableTitle{5th Level} \nopagebreak }

\begin{description}
\item \textit{Hero of Fable}
\item \textit{Krail's Rot}
\item \textit{Lost and Wandering}
\end{description}

{\par \vspace { 9pt plus 3pt minus 3pt } \noindent  \DndFontTableTitle{6th Level} \nopagebreak }

\begin{description}
\item \textit{Grim Harvest}
\end{description}

{\par \vspace { 9pt plus 3pt minus 3pt } \noindent  \DndFontTableTitle{7th Level} \nopagebreak }

\begin{description}
\item \textit{Doom of Summer Years}
\item \textit{Ebon Tide}
\end{description}

{\par \vspace { 9pt plus 3pt minus 3pt } \noindent  \DndFontTableTitle{8th Level} \nopagebreak }

\begin{description}
\item \textit{Child of Light and Darkness}
\end{description}

\end{multicols}

\end{description}

\subsection{Druid Spells}

\begin{description}
\begin{multicols}{2}
{\par \vspace { 9pt plus 3pt minus 3pt } \noindent  \DndFontTableTitle{Cantrips (0 Level)} \nopagebreak }

\begin{description}
\item \textit{Drizzle}
\item \textit{Flowering}
\end{description}

{\par \vspace { 9pt plus 3pt minus 3pt } \noindent  \DndFontTableTitle{1st Level} \nopagebreak }

\begin{description}
\item \textit{Fireflies}
\end{description}

{\par \vspace { 9pt plus 3pt minus 3pt } \noindent  \DndFontTableTitle{2nd Level} \nopagebreak }

\begin{description}
\item \textit{Hibernation}
\item \textit{Krail's Maggot}
\item \textit{Krail's Rupture}
\item \textit{Ominous Winds}
\end{description}

{\par \vspace { 9pt plus 3pt minus 3pt } \noindent  \DndFontTableTitle{3rd Level} \nopagebreak }

\begin{description}
\item \textit{Bitter Wind}
\item \textit{Deep Roots of the Moon}
\item \textit{Faerie Toast}
\item \textit{Lost}
\end{description}

{\par \vspace { 9pt plus 3pt minus 3pt } \noindent  \DndFontTableTitle{4th Level} \nopagebreak }

\begin{description}
\item \textit{Drayfn's Curse of Incompetence}
\item \textit{Knife of Fate}
\item \textit{Lunar Transfer}
\item \textit{Moonlight Sending}
\item \textit{Storm Door}
\end{description}

{\par \vspace { 9pt plus 3pt minus 3pt } \noindent  \DndFontTableTitle{5th Level} \nopagebreak }

\begin{description}
\item \textit{Feast of Flesh}
\item \textit{Krail's Rot}
\item \textit{Lost and Wandering}
\item \textit{Radiant Beacon}
\end{description}

{\par \vspace { 9pt plus 3pt minus 3pt } \noindent  \DndFontTableTitle{6th Level} \nopagebreak }

\begin{description}
\item \textit{Conjure Giant}
\item \textit{Doom of Stacked Stones}
\item \textit{Grim Harvest}
\end{description}

{\par \vspace { 9pt plus 3pt minus 3pt } \noindent  \DndFontTableTitle{7th Level} \nopagebreak }

\begin{description}
\item \textit{Ebon Tide}
\end{description}

\end{multicols}

\end{description}

\subsection{Paladin Spells}

\begin{description}
\begin{multicols}{2}
{\par \vspace { 9pt plus 3pt minus 3pt } \noindent  \DndFontTableTitle{1st Level} \nopagebreak }

\begin{description}
\item \textit{Blade of Blood and Bone}
\item \textit{Doom of Poor Fortune}
\item \textit{Emerald Goblet}
\end{description}

{\par \vspace { 9pt plus 3pt minus 3pt } \noindent  \DndFontTableTitle{2nd Level} \nopagebreak }

\begin{description}
\item \textit{Drayfn's Blunted Blade}
\item \textit{Krail's Rupture}
\end{description}

{\par \vspace { 9pt plus 3pt minus 3pt } \noindent  \DndFontTableTitle{3rd Level} \nopagebreak }

\begin{description}
\item \textit{Faerie Toast}
\end{description}

{\par \vspace { 9pt plus 3pt minus 3pt } \noindent  \DndFontTableTitle{4th Level} \nopagebreak }

\begin{description}
\item \textit{Drayfn's Curse of Incompetence}
\end{description}

{\par \vspace { 9pt plus 3pt minus 3pt } \noindent  \DndFontTableTitle{5th Level} \nopagebreak }

\begin{description}
\item \textit{Hero of Fable}
\end{description}

\end{multicols}

\end{description}

\subsection{Ranger Spells}

\begin{description}
\begin{multicols}{2}
{\par \vspace { 9pt plus 3pt minus 3pt } \noindent  \DndFontTableTitle{1st Level} \nopagebreak }

\begin{description}
\item \textit{Blade of Blood and Bone}
\end{description}

{\par \vspace { 9pt plus 3pt minus 3pt } \noindent  \DndFontTableTitle{2nd Level} \nopagebreak }

\begin{description}
\item \textit{Drayfn's Blunted Blade}
\item \textit{Hibernation}
\item \textit{Krail's Maggot}
\item \textit{Krail's Rupture}
\end{description}

{\par \vspace { 9pt plus 3pt minus 3pt } \noindent  \DndFontTableTitle{3rd Level} \nopagebreak }

\begin{description}
\item \textit{Faerie Toast}
\item \textit{Lost}
\end{description}

{\par \vspace { 9pt plus 3pt minus 3pt } \noindent  \DndFontTableTitle{5th Level} \nopagebreak }

\begin{description}
\item \textit{Lost and Wandering}
\end{description}

\end{multicols}

\end{description}

\subsection{Sorcerer Spells}

\begin{description}
\begin{multicols}{2}
{\par \vspace { 9pt plus 3pt minus 3pt } \noindent  \DndFontTableTitle{Cantrips (0 Level)} \nopagebreak }

\begin{description}
\item \textit{Bright Sparks}
\item \textit{Drayfn's Bane of Excellence}
\end{description}

{\par \vspace { 9pt plus 3pt minus 3pt } \noindent  \DndFontTableTitle{1st Level} \nopagebreak }

\begin{description}
\item \textit{Doom of Poor Fortune}
\item \textit{Elf Shot}
\item \textit{Gloaming}
\item \textit{Pratfall}
\end{description}

{\par \vspace { 9pt plus 3pt minus 3pt } \noindent  \DndFontTableTitle{2nd Level} \nopagebreak }

\begin{description}
\item \textit{Drayfn's Blunted Blade}
\item \textit{Krail's Rupture}
\item \textit{Ominous Winds}
\item \textit{Shadow Adaptation}
\item \textit{Ugly Duckling}
\end{description}

{\par \vspace { 9pt plus 3pt minus 3pt } \noindent  \DndFontTableTitle{3rd Level} \nopagebreak }

\begin{description}
\item \textit{Bitter Wind}
\item \textit{Orros Mark of Fate}
\item \textit{Shadow Portal}
\end{description}

{\par \vspace { 9pt plus 3pt minus 3pt } \noindent  \DndFontTableTitle{4th Level} \nopagebreak }

\begin{description}
\item \textit{Drayfn's Curse of Incompetence}
\item \textit{Fire Dance}
\item \textit{Knife of Fate}
\item \textit{Lunar Transfer}
\item \textit{Moonlight Sending}
\item \textit{Polychromatic Bubble}
\item \textit{Portho's Portal}
\item \textit{Spider Song}
\end{description}

{\par \vspace { 9pt plus 3pt minus 3pt } \noindent  \DndFontTableTitle{5th Level} \nopagebreak }

\begin{description}
\item \textit{Feast of Flesh}
\item \textit{Krail's Rot}
\item \textit{Radiant Beacon}
\end{description}

{\par \vspace { 9pt plus 3pt minus 3pt } \noindent  \DndFontTableTitle{6th Level} \nopagebreak }

\begin{description}
\item \textit{Doom of Stacked Stones}
\end{description}

{\par \vspace { 9pt plus 3pt minus 3pt } \noindent  \DndFontTableTitle{7th Level} \nopagebreak }

\begin{description}
\item \textit{Charnel Banquet}
\item \textit{Doom of Summer Years}
\end{description}

{\par \vspace { 9pt plus 3pt minus 3pt } \noindent  \DndFontTableTitle{8th Level} \nopagebreak }

\begin{description}
\item \textit{Child of Light and Darkness}
\item \textit{Doom of False Friends}
\end{description}

\end{multicols}

\end{description}

\subsection{Warlock Spells}

\begin{description}
\begin{multicols}{2}
{\par \vspace { 9pt plus 3pt minus 3pt } \noindent  \DndFontTableTitle{Cantrips (0 Level)} \nopagebreak }

\begin{description}
\item \textit{Drayfn's Bane of Excellence}
\item \textit{Obfuscate Object}
\end{description}

{\par \vspace { 9pt plus 3pt minus 3pt } \noindent  \DndFontTableTitle{1st Level} \nopagebreak }

\begin{description}
\item \textit{Blade of Blood and Bone}
\item \textit{Gloaming}
\item \textit{Pratfall}
\end{description}

{\par \vspace { 9pt plus 3pt minus 3pt } \noindent  \DndFontTableTitle{2nd Level} \nopagebreak }

\begin{description}
\item \textit{Drayfn's Blunted Blade}
\item \textit{Krail's Maggot}
\item \textit{Krail's Rupture}
\item \textit{Leiloch's Irritating Kazoo}
\item \textit{Shadow Adaptation}
\end{description}

{\par \vspace { 9pt plus 3pt minus 3pt } \noindent  \DndFontTableTitle{3rd Level} \nopagebreak }

\begin{description}
\item \textit{Grim Shadows}
\item \textit{Shadow Portal}
\end{description}

{\par \vspace { 9pt plus 3pt minus 3pt } \noindent  \DndFontTableTitle{4th Level} \nopagebreak }

\begin{description}
\item \textit{Drayfn's Curse of Incompetence}
\item \textit{Fire Dance}
\item \textit{Spider Song}
\end{description}

{\par \vspace { 9pt plus 3pt minus 3pt } \noindent  \DndFontTableTitle{5th Level} \nopagebreak }

\begin{description}
\item \textit{Feast of Flesh}
\item \textit{Leiloch's Interminable Yarn}
\end{description}

{\par \vspace { 9pt plus 3pt minus 3pt } \noindent  \DndFontTableTitle{6th Level} \nopagebreak }

\begin{description}
\item \textit{Grim Harvest}
\end{description}

{\par \vspace { 9pt plus 3pt minus 3pt } \noindent  \DndFontTableTitle{7th Level} \nopagebreak }

\begin{description}
\item \textit{Leiloch's Arduous Shuffle}
\end{description}

\end{multicols}

\end{description}

\subsection{Wizard Spells}

\begin{description}
\begin{multicols}{2}
{\par \vspace { 9pt plus 3pt minus 3pt } \noindent  \DndFontTableTitle{Cantrips (0 Level)} \nopagebreak }

\begin{description}
\item \textit{Bright Sparks}
\item \textit{Drayfn's Bane of Excellence}
\item \textit{Drizzle}
\end{description}

{\par \vspace { 9pt plus 3pt minus 3pt } \noindent  \DndFontTableTitle{1st Level} \nopagebreak }

\begin{description}
\item \textit{Doom of Poor Fortune}
\item \textit{Elf Shot}
\item \textit{Emerald Goblet}
\item \textit{Gloaming}
\item \textit{Pratfall}
\end{description}

{\par \vspace { 9pt plus 3pt minus 3pt } \noindent  \DndFontTableTitle{2nd Level} \nopagebreak }

\begin{description}
\item \textit{Conjure Ferryman}
\item \textit{Drayfn's Blunted Blade}
\item \textit{Krail's Rupture}
\item \textit{Ominous Winds}
\item \textit{Shadow Adaptation}
\item \textit{Ugly Duckling}
\end{description}

{\par \vspace { 9pt plus 3pt minus 3pt } \noindent  \DndFontTableTitle{3rd Level} \nopagebreak }

\begin{description}
\item \textit{Bitter Wind}
\item \textit{Grim Shadows}
\item \textit{Orros Mark of Fate}
\item \textit{Shadow Portal}
\end{description}

{\par \vspace { 9pt plus 3pt minus 3pt } \noindent  \DndFontTableTitle{4th Level} \nopagebreak }

\begin{description}
\item \textit{Drayfn's Curse of Incompetence}
\item \textit{Fire Dance}
\item \textit{Knife of Fate}
\item \textit{Lunar Transfer}
\item \textit{Moonlight Sending}
\item \textit{Polychromatic Bubble}
\item \textit{Portho's Portal}
\item \textit{Spider Song}
\item \textit{Storm Door}
\end{description}

{\par \vspace { 9pt plus 3pt minus 3pt } \noindent  \DndFontTableTitle{5th Level} \nopagebreak }

\begin{description}
\item \textit{Feast of Flesh}
\item \textit{Krail's Rot}
\item \textit{Radiant Beacon}
\end{description}

{\par \vspace { 9pt plus 3pt minus 3pt } \noindent  \DndFontTableTitle{6th Level} \nopagebreak }

\begin{description}
\item \textit{Conjure Giant}
\item \textit{Doom of Stacked Stones}
\item \textit{Grim Harvest}
\end{description}

{\par \vspace { 9pt plus 3pt minus 3pt } \noindent  \DndFontTableTitle{7th Level} \nopagebreak }

\begin{description}
\item \textit{Charnel Banquet}
\item \textit{Doom of Summer Years}
\item \textit{Ebon Tide}
\end{description}

{\par \vspace { 9pt plus 3pt minus 3pt } \noindent  \DndFontTableTitle{8th Level} \nopagebreak }

\begin{description}
\item \textit{Child of Light and Darkness}
\item \textit{Doom of False Friends}
\end{description}

\end{multicols}

\end{description}

\section{Spell Descriptions}
The spells presented below are relatively common in the Shadow Realm and rather less frequently found in the bright lands. The spells of the Schools of Illumination and of Shadow (see \textit{Deep Magic}) are also widely known in the realm. Spells included in those traditions are marked here with either the \textbf{illumination} or \textbf{shadow} keywords.

Two new schools of magic are introduced as well: that of fey magic and of weather magic. These are, as the names suggest, spells that originated with fey spellcasters and weather mages. They are marked here with either the \textbf{fey} or \textbf{weather} keywords.


\begin{itemize}
\begin{multicols}{3}
\item \textit{Bitter Wind}
\item \textit{Blade of Blood and Bone}
\item \textit{Bright Sparks}
\item \textit{Charnel Banquet}
\item \textit{Child of Light and Darkness}
\item \textit{Conjure Ferryman}
\item \textit{Conjure Giant}
\item \textit{Deep Roots of the Moon}
\item \textit{Doom of False Friends}
\item \textit{Doom of Poor Fortune}
\item \textit{Doom of Stacked Stones}
\item \textit{Doom of Summer Years}
\item \textit{Drayfn's Bane of Excellence}
\item \textit{Drayfn's Blunted Blade}
\item \textit{Drayfn's Curse of Incompetence}
\item \textit{Drizzle}
\item \textit{Ebon Tide}
\item \textit{Elf Shot}
\item \textit{Emerald Goblet}
\item \textit{Faerie Toast}
\item \textit{Feast of Flesh}
\item \textit{Fire Dance}
\item \textit{Fireflies}
\item \textit{Flowering}
\item \textit{Gloaming}
\item \textit{Grim Harvest}
\item \textit{Grim Shadows}
\item \textit{Hero of Fable}
\item \textit{Hibernation}
\item \textit{Knife of Fate}
\item \textit{Krail's Maggot}
\item \textit{Krail's Rot}
\item \textit{Krail's Rupture}
\item \textit{Leiloch's Arduous Shuffle}
\item \textit{Leiloch's Interminable Yarn}
\item \textit{Leiloch's Irritating Kazoo}
\item \textit{Lost}
\item \textit{Lost and Wandering}
\item \textit{Lunar Transfer}
\item \textit{Moonlight Sending}
\item \textit{Obfuscate Object}
\item \textit{Ominous Winds}
\item \textit{Orros Mark of Fate}
\item \textit{Polychromatic Bubble}
\item \textit{Portho's Portal}
\item \textit{Pratfall}
\item \textit{Radiant Beacon}
\item \textit{Shadow Adaptation}
\item \textit{Shadow Portal}
\item \textit{Spider Song}
\item \textit{Storm Door}
\item \textit{Ugly Duckling}
\end{multicols}

\end{itemize}

\chapter{5. The Nature of Shadow}
The foremost questions from many travelers to the Shadow Realm have to do with its shadow matter, its malleable nature, and how humans and other intelligent species survive in such a strange place. To natives, these questions are foolish, for the laws of nature may be different than they are on other planes or in the mortal realm, but the Shadow Realm is hospitable to those who take time to learn its ways. This chapter provides that foundation.

While it is known for its occasionally warped, echoed locations with parallels in the mortal world, the Shadow Realm also harbors twisted plants, animals, monsters, conditions of travel, and even strange laws of nature worth closer inspection. Turn your gaze to the mutable lands of shadow, and pray you have a keen eye to see what marvels they contain.

\section{Moons, Stars, and Sky}
The Shadow Realm has no full day and night cycle, and instead, a perpetual twilight cloaks the realm. This dim, ambient illumination deepens into full darkness in the valleys and canyons that crisscross the land. And it isn't sunlight, so creatures that suffer in the sun, including vampires, endure it without difficulty.

The primary sources of light are the two moons of the Shadow Realm, Celene and Lindar, which most inhabitants of the realm know as the feminine and masculine faces of Hecate, or sometimes as Hecate and Kupkoresh or Lada. The two moons provide enough light for the realm's rather pale, thin crops and trees to grow, and both have linked phases. When Celene is full, so is Lindar, and when Lindar is new, so is Celene.

The stars are a handful of jewels, often cloaked in clouds, but when visible, they are just like the stars of the mortal world: chips of pure brilliance, red, blue, and diamond white. The constellations familiar to umbral peoples are the Dragon, the North Arrow, the Leaping Salmon, the Swan, the Hunter, and the Greater and Lesser Fox.

The sky in the Shadow Realm is often cloudy or misty, but when open from horizon to horizon, it is a richer palette than most: black and gray of course but with overtones of mauve, violet, and deep blue, velvety deep tones, sometimes even vivid orange and rich burgundy hues by the full moon. Its clouds are commonly gray or white and uncommonly silver or blue—perhaps even luminous.

\section{Animals and Encounters}
The plants and wildlife in the Shadow Realm survive as one would expect given the surrounding terrain. Trees and undergrowth fill the dark forests. Birds fly above and nest among twisted branches. Herd animals wander the plains. Fish swim through the black rivers and murky lakes. Most are entirely normal in their size and behavior.

Some creatures native to the Shadow Realm are a dark reflection of their counterparts in the mortal world, such as dark-furred squirrels and twilight-purple deer with black antlers. Some are clearly warped and tainted. The owl-like but human-mouthed stryx ghost-glides on silent wings through dark forests. Shadow fey nobles capture and tend flocks of steel-winged eala, and the corrupted, unicorn-like shadhavar wander the woods and race across the plains. White wolves, black bears, and even crimson-feathered hawks are common, thriving on a seemingly ready supply of pale field mice, the aforementioned black squirrels, and similarly shadow-hued prey.

Far more than the odd creatures still familiar from fields and rivers though, monsters are abundant within the Shadow Realm. The \textbf{gloomflower} (see \textit{Creature Codex}), \textbf{glass gator} (see \textit{Tome of Beasts 1}), and \textbf{unhatched} (see \textit{Creature Codex}) are all well-known threats, as are \textbf{light dragons} (see \textit{Creature Codex}) and black dragons, lantern dragonettes (see \textit{Tome of Beasts 1}), and \textbf{overshadows} (see \textit{Tome of Beasts 2}). (For Shadow Realm encounter tables, see \textit{Appendix}.)

\section{Food and Water}
Finding fresh water often presents a pressing challenge in the Shadow Realm since darkness flows through some waterways instead of water. This inky substance quenches the thirst of the denizens of the plane (including all those detailed in \textit{Chapter 2}), but visitors from the bright lands find it bitter and unpleasant. That said, this inky liquid can be consumed and provides refreshment, but it also carries the taint of shadow corruption (\textit{see below}). A few springs produce pure water, and sometimes a stream runs clear for a time before darkness fills it again. These precious resources are vital to those that haven't given themselves over to the plane's corruption, and they guard them viciously.

As for food, fruit-bearing plants yield their bounty (though their produce appears black and glossy), and animals can be hunted for meat. Whether fruit or meat, food gathered in the Shadow Realm tastes bland to visitors and often hints at something unpleasant, such as rot or mold. Ability checks made to forage for food and water by visitors to the realm are at disadvantage, yielding half the normal amount of edible food or potable water (though natives suffer no penalty). Brave or desperate characters can choose to forage as normal (without disadvantage and for a full yield), but any creature who consumes food gathered in this way suffers exposure to shadow corruption.

\subsection{Dark Water}
The black shadows that often pass for water in the Shadow Realm run swift and cold—so cold that no matter the surrounding terrain or climate, every stream and river and lake counts as frigid water. Worse, the spirits of things that died in or near the water constitute a separate hazard. A creature that starts its turn in dark water must succeed on a DC 12 Strength saving throw or be grappled and pulled 5 feet deeper into the water by the hungry spirits. The grappled condition lasts until the start of the creature's next turn. These ghosts long to drag every living thing they encounter into the inky depths, and even watercraft provide uncertain protection against these spirits.

A creature that drowns in dark water can't be returned to life by \textit{revivify} or \textit{raise dead}. More powerful magic, such as \textit{resurrection}, is necessary to free the creature's soul.

A few creatures of the Shadow Realm are naturally immune to the effects of dark water, including shadow goblins, shadow fey, river giants, and most fiends.

\subsubsection{Variant: Shadow-Touched Beasts}
Beasts in the Shadow Realm reflect the plane's unique nature. Ordinary beasts have darkvision out to 60 feet and are resistant to cold damage. At the GM's discretion, these traits fade away a week after the beast leaves the plane.

\begin{DndSidebar}[float=!b]{Ancient Animals Still Thrive}
Fauna such as mammoths, enormous elks, ground sloths, sabretooth tigers, and cave bears still thrive in the Shadow Realm. It seems that the dreams of the bearfolk keep them nearby as they are most common in the Moonlit Groves with occasional visits to dark forests and open grasslands where the sable elves and shadow fey certainly know how to hunt them.



Many of these are referred to as "elders," such as elder elk or elder mammoth. They are considered especially worthy prey and are sometimes protected by will-o'-wisps, witchlights, or other unusual guardians. Goblin tribes occasionally bring one to a court or to Fandeval as a curiosity and showpiece.



\end{DndSidebar}

\section{Methods for Entering Shadow}
Natives to the Shadow Realm need no help arriving, but the region is also popular with creatures from elven lands, from mortal worlds, and from the various hells and heavens, seeking souls lost in the Styx or for powerful magical weapons to recharge with the arcane forces of the realm.

The five principal methods for entering the realm are nodes, the Planar River, radiant wells, shadow roads, and spells. Each of these is further described below.

\subparagraph{Nodes and Portals}
Many visitors from outside the Shadow Realm enter via a magical gate, portal, or node of some kind. On their home plane, these locations appear as doors, gateways, arches, standing stones, or arcane circles, which must be activated by a magical item, seasonal arcane sacrifice, ritual magic, or similar charging energy. Once activated, they remain open for a short period, allowing multiple creatures to enter at the same time.

\subparagraph{Planar River Styx}
Travelers can and do float into the Shadow Realm from the hells, from the roots of Yggdrasil, or from other locations where the River Styx is known to flow. Charun's boatmen, various yugoloths, and even angelic guides have been known to point out the way to safe places to dock or beach.

\subparagraph{Radiant Wells}
Connected to the celestial planes, these glowing beacons provide a great deal of light, but the ability to use them to enter the Shadow Realm is limited primarily to angels, celestials, and good-aligned creatures generally. The fey, undead, and others consider them a hazard (see also \textit{Hazards of the Shadow Realm} below).

\subparagraph{Spells}
Both \textit{plane shift} and \textit{gate} spells can bring a party of adventurers to the Shadow Realm. Some scrolls are created, especially by powerful fey or ghoulish spellcasters, to bring new and entertaining visitors to the realm, and these are referred to as Shadow Portal Scroll (see \textit{Chapter 10}).

\subparagraph{Shadow Roads}
The most common way to enter the Shadow Realm is on the roads that connect the worlds. These shadow roads are obscured and difficult to open—for humans and other non-fey. For the shadow fey, elves, gnomes, and others, opening the shadow roads is much easier. Using them requires making a trip of several days to get from the Shadow Realm to the mortal lands or vice versa. More detail on these roads is found \textit{later in this chapter}.

\begin{DndSidebar}[float=!b]{The Nature of Shadow / from the Book of Ebon Tides}
While many wise souls pretend to know exactly what shadow matter is and how it is best distilled, shifted, changed, and locked into new forms, the truth is that the primal matter of the Shadow Realm is complex, often chaotic, and different from the simple elements of other realms, such as fire, water, earth, and air. It is malleable in the extreme, making it more susceptible to shaping by the untrained.



Primal shadow can be manipulated by a conjuror, illusionist, or druid—and sometimes, by the mere belief held by a prince or a fey child. In most cases, these changes give life to an object that is normally lifeless, such as a stone or a doll, purely through belief, and these changes last long enough for a birch tree to smile and bow, a cherry tree to blossom, a doll to wink, or a boulder to stroll down the lane.



In a few cases, animation and form can be provided that last longer and have greater effects, but in almost all cases, these require a practitioner of the magical arts.



\end{DndSidebar}

\section{Hazards of the Shadow Realm}
While most of the realm is ordinary earth and forests, several elements of the Shadow Realm are unique to it—or at least far more common than elsewhere. These hazards and landmarks are described below.

\subparagraph{Doomsand}
Similar in most ways to standard quicksand, doomsand follows the same rules for sinking and escaping. This ashen-gray slurry is infused with the nature of shadow, and it drains the resolve from those in its grasp. A creature that starts its turn in doomsand must succeed on a DC 12 Wisdom saving throw or become despondent (see \textit{Shadow Corruption} below). When the creature finishes a long rest, it can repeat the saving throw, ending the effect on a success. Doomsand usually appears in marshes and swamps, or in deserts as fine dust.

\subparagraph{Ebon Tides}
Ebon tides are enormous waves of shadow mire, of crackling shadow matter—a temporary wave of arcane power pulsing through a region and unmaking the land. When they affect a region, use the effects of the \textit{ebon tide} spell (see \textit{Chapter 4}).

\subparagraph{Hungry Gloom}
Some areas of the Shadow Realm—or places in the mortal world where the barrier to the Shadow Realm is thin—develop hungry gloom. An area of hungry gloom appears particularly dark and grim, swallowing light. An area of light that overlaps hungry gloom has its illumination reduced by one step (bright light becomes dim, dim light becomes darkness). Darkvision can't improve the ambient light conditions in the area. Shadow-touched creatures (see \textit{Shadow Corruption} below) can see in hungry gloom as if it were bright light.

\subparagraph{Leaden Clouds}
These dark, heavy clouds oppress the spirit of those beneath them, serving as a sort of psychic weight. Creatures under a leaden cloud of sorrow must succeed on a DC 12 Wisdom or Charisma saving throw to make an attack, cast a spell, or move at more than a standard walking pace (though they can still defend themselves normally). Certain items, such as \textit{Emerald Goblet} (see \textit{Chapter 10}), and spells, such as \textit{aid}, \textit{calm emotions}, \textit{enthrall}, and \textit{lesser restoration}, and even a barbarian's Rage ability negate the effects of leaden clouds.

\subparagraph{Radiant Wells}
This form of well resembles a beacon of light at the surface, though it swallows up anything thrown into it. A radiant well is a planar link to one of the celestial planes, and it sheds searing, bright daylight in a 60-foot radius around itself and up into the clouds. This makes them excellent landmarks for navigation, though their position can be shifted by ebon tides and terrain manipulation. Any creature entering a radiant well falls for 1 round and then is ejected upward at relatively high speed from another well within 2d20 miles (for evil creatures, this increases to 1d100 miles). The force of being ejected from the well propels a Small or smaller creature up about 30 feet before falling to the ground beside the well (potentially taking normal falling damage). Medium and Large creatures are pushed up about 15 feet, and Huge or larger creatures simply land at the edge of the well.

Particularly large or active radiant wells sometimes inflict radiant damage on evil-aligned creatures. These more active wells are sometimes used as nodes or portals to celestial planes by fey wizards and clerics.

\subparagraph{Shadow Mire}
Usually found along the shores of ponds, lakes, and streams, shadow mire is a concentration of pure shadow matter that has eaten away the soil, roots, and plants of a place and replaced them with a thin layer of viscous purple-black gloop. This resembles dark water and is indistinguishable from water at night—leaves and reeds may float on it, and it reflects moonlight normally. However, stepping into shadow mire generates an immediate, oversized wave that radiates from the creature entering the mire, knocking down all in its path. Any creatures within 30 feet, which are smaller than the creature stepping into the mire, must succeed on a DC 13 Strength (Athletics) or Dexterity (Acrobatics) saving throw or be knocked prone by the wave.

The creature in the center of the mire will be covered by the shadow mire as it rolls back to its point of origin, like the splash of an enormous raindrop. The creature must make a DC 13 Constitution saving throw or be knocked prone and unconscious.

Some wizards can distill or conjure pits of shadow mire (see spells in \textit{Chapter 4} and magic items in \textit{Chapter 10}) and thus reshape a riverbed or undermine a bridge, for instance.

\subparagraph{Witchlights and Will-o'-Wisps}
The Shadow Realm is filled with lesser lights than that of the moons and stars, wandering along with typically inscrutable aims. A white glow shines in the distance, offering a path or a sign of habitation—or simply zipping through the deep woods. These are the tiny witchlights (see \textit{Tome of Beasts 1}) and the larger will-o'-wisps, both of which create their own light. The witchlights are tiny shards of crystal, constructed creatures believed to have been made as friends, servants, and companions to the Ancients of shadow. Only relatively powerful wizards know how to make them, but many roam the Shadow Realm on their own.

Witchlights pursue their own purposes: guarding dangerous sites, helping travelers, seeking out old runestones, or gathering at certain seasons to dance and blink in unison. Others tag along with strangers for hours, perhaps driven by curiosity. Their appearance is well known, but their goals are often quite obscure. Some are friendly while others guide travelers astray, and a few of them serve as familiars to wizards and sorcerers among the shadow fey and shadow goblins.

Will-o'-wisps, their larger cousins, are generally malign and capable of more damage in combat. They lead travelers into doomsand or shadow mire and even use their attacks to sever boat lines or rope bridges.

\section{Shadow Corruption}
The Shadow Realm threads its tendrils into those places and creatures it can reach, tainting those that strike bargains with dark powers, linger too long in places infused with shadow, or even simply eat the wrong food in the Shadow Realm. This applies primarily to visitors, as the realm's inhabitants are wary and used to its dangers, though sometimes they also fall prey to them.

Creatures corrupted by shadow grow distracted and withdrawn, shunning light in all its forms. At the most severe levels, the corrupted creature gives itself wholly to the Shadow Realm. Constructs, fiends, shadow fey, shadow goblins, and undead aren't susceptible to shadow corruption. Constructs (other than gearforged or similar construct-like creatures) lack souls, and darkness of another sort already claims the essential nature of fiends and undead. Shadow fey and shadow goblins have lived in the realm for so long, their nature has already been changed by it.

Shadow corruption is measured in six levels. An effect can give a creature one or more levels of shadow corruption, as specified in the effect's description detailed in the following table.

\begin{DndTable}[header=Table 5-1: Shadow Corruption]{cX}

Level & Effect\\
1 & You have disadvantage on Wisdom and Charisma checks made against non-shadow creatures.\\
2 & You gain darkvision out to 30 feet or increases existing darkvision by 30 feet.\\
3 & You have disadvantage on Wisdom (Perception) checks and attack rolls made while in bright light.\\
4 & You have disadvantage on saving throws made while in bright light.\\
5 & You take 2d6 radiant damage when starting your turn in sunlight.\\
6 & You become a shadow thrall (\textit{see condition below}).\\
\end{DndTable}

If a corrupted creature suffers another effect that causes shadow corruption, its current level of shadow corruption increases by the amount specified in the effect's description. A creature suffers the effect of its current level of shadow corruption as well as all lower levels.

An effect that removes shadow corruption reduces its level, as specified in the effect's description, with all shadow corruption effects ending if all levels of a creature's shadow corruption level are removed.

For each week spent in the Shadow Realm, a susceptible creature must make a DC 10 Charisma saving throw. On a failure, the creature gains one level of shadow corruption. On a success, the creature resists corruption for the time being, but the DC of future saves made for this reason increase by 1. The DC returns to 10 when the creature gains one or more levels of shadow corruption or when it spends at least 1 week outside the Shadow Realm.

Shadow magic is dangerous for mortal creatures to wield since it carries the risk of shadow corruption. Whenever any humanoid that isn't an umbral human, bearfolk, darakhul, shadow fey, or shadow goblin casts any shadow magic spell, then the next time it finishes a long rest, it must make a Charisma saving throw with a DC of 10 + the highest level of shadow magic spell cast. On a failure, the creature gains one level of shadow corruption. Gnomes, elfmarked, and elves make the saving throw with advantage.

Shadow corruption can be removed in three ways:


\begin{itemize}
\item A corrupted creature that spends 1 week per current level of shadow corruption outside the Shadow Realm reduces its level of shadow corruption by one. If the creature casts or is affected by a shadow magic spell during this time, the recovery time starts over.
\item A \textit{dispel evil and good} spell cast on a creature reduces its shadow corruption level by one.
\item An \textit{emerald goblet} can remove one level of shadow corruption each time it is used.
\end{itemize}

\subsection{Shadow-Touched}
A creature that is native to the Shadow Realm or that possesses at least one level of shadow corruption is considered shadow-touched.

\subsubsection{New Condition: Despondent}
A despondent creature gains the following flaw: "I find it difficult to see how anything could turn out well no matter how hard I try. Why even bother?" In addition, they are directly affected as follows:


\begin{itemize}
\item The creature has disadvantage on ability checks.
\item The creature has advantage on saving throws against being charmed and disadvantage on all other saving throws.
\item The \textit{calm emotions} spell suppresses the despondent condition for its duration. \textit{Lesser restoration} or more powerful magic ends the condition.
\end{itemize}

\subsubsection{New Condition: Shadow Thrall}
A shadow thrall is considered a denizen of the Shadow Realm and won't willingly leave:


\begin{itemize}
\item If forced to leave the Shadow Realm, the creature becomes despondent (see condition above) until it returns.
\item The thrall's form becomes mottled with shifting patches of inky blackness. It automatically fails all Wisdom (Insight) and Charisma checks (except for Intimidation checks) against a non-shadow creature.
\item The thrall is considered charmed by any fey or shadow fey native to the Shadow Realm who has an Intelligence or Charisma score higher than 14. This charmed effect can't be removed as long as the creature is a shadow thrall. However, a \textit{remove curse} or \textit{lesser restoration} spell suppresses the effect for 24 hours.
\end{itemize}

\section{Traveling on Shadow Roads}
The network of shadow roads grows, withers, and changes over the years. However, the experience of traveling on them is usually similar, no matter who their maker might be. (For more on making shadow roads, see \textit{Chapter 1}.) These roads are built by magic, protected by the fey, bearfolk, and other road wardens, and provide swifter, more reliable travel than venturing over open land. In the Shadow Realm, open grasslands, forests, and any other natural terrain is subject to change and warpage, so the use of roads provides one of the few certain ways of getting from point \textit{A} to point \textit{B}. Though while more certain, they are still perilous, and travel comes with a number of possible dangers.

\subsection{Shadow Road Wardens}
The various roads through the Shadow Realm are often quick and easy, especially when the owners of the lands nearby do not cast any charms or hexes to slow a visitor's progress or to hinder enemies. But at other times, bandits, enchantments, and natural or supernatural obstacles prevent swift passage, and when these interfere with the regular travel of wagons, pilgrims, merchants, or noble processions through the realm, the road wardens are dispatched from the nearest court or city to clear things up.

The major chapterhouses of the road wardens stand in the \textit{Courts of Mice and Ravens} and the \textit{Court of Roses} as well as in \textit{Corremel}, \textit{Dalliance}, \textit{Fandeval}, and \textit{Merrymead}. The shades and gnomes do not provide road wardens near Soriglass or along the more distant reaches of the western woods, but the darakhul legions perform a similar function within the rifts in the name of the emperor.

Shadow road wardens are often veterans of courtly life who have been given the task after they tire of the daily rituals, gatherings, and petty slights of fey status-seeking. These are often quite deadly warriors, enchanters, or rogues, known and feared by bandits, satarre, vampires, and void cultists for their skill at sniffing out bandits, road-cutters, and various other plotters against the Old Roads.

Road wardens are notorious tale-spinners but also delightful in the depth of their understanding of shadow roads, tides, and matters arcane.

\begin{DndTable}[header=Table 5-2: Six Famous Road Wardens]{cX}

d6 & Result\\
1 & Ambrosus Fairhand, a sable elf general who tired of life in the Court of Golden Oaks, leads a detachment of sable elves and erina along the northern banks of the Styx. He is known for chopping off the hands of bandits, road-cutters, and other miscreants found on the roads—always pronouncing it, "Sentence performed in the name of Queen Valda."\\
2 & Hillar Vandegast, a shadow fey archmage who travels on a nightmare as a steed and who keeps a small company of young shadow fey squires and apprentices. He guards the eastern Queen's Wood and transforms criminals into toads, sparrows, and sometimes weasels.\\
3 & Gillyan the Sailor, a shadow goblin once of Fandeval, saw something in the Sable Sea and vowed never to sail again. Known for using radiant magic against his foes, it's said he once won and lost a \textit{crown of infinite midnight} (see \textit{Chapter 10}) in a card game.\\
4 & Kei Greenleaf, a gnomish druid, serves the bearfolk and protects the trails and paths of the north. Said to work with erina and alseid, though these are rarely visible when Kei questions a traveler.\\
5 & Toress Bonespitter, a darakhul captain assigned to service outside the great ghoul rift, enjoys human and fey company but despises the bearfolk. Her crew of darakhul and human rangers protect the rather new ghoulish roads around the rim of the rift.\\
6 & Garwen the Pure, a bearfolk paladin who travels as a humble but relentless defender of light and justice, is a road warden who lets nothing slide. Smugglers, necromancers, thieves, and simple gamblers may all be required to pay a toll—or may be carried off bodily to the court of justice at Merrymead.\\
\end{DndTable}

\begin{DndSidebar}[float=!b]{Highway Robbery in the Shadow Realm}
The frequent visitors to the Shadow Realm make it almost irresistible to attempt highway robbery. A group of enormous bearfolk, a horrifying company of darakhul soldiers, or a heavily armed shadow fey stepping out of the shadows around a single traveler or even a small group will often demand a toll or to "deliver up a parcel to the road gods." While most travelers from the bright lands are willing and able to pay a small fee, sometimes the highway robbers get much more than they bargained for. In these cases, they often scatter among the difficult terrain on either side of the road (forest, tall plains grasses, scrublands, or hills). Most non-native travelers are reluctant to leave the road, and bandits take advantage of this.



\end{DndSidebar}

\subsection{Travel Times on the Shadow Roads}
Traveling the hundreds of miles along a shadow road is fairly swift but not instantaneous, and the better-kept and better-warded roads of the shadow fey, bearfolk, and elves (though in some mild disrepair) are far swifter than roads elsewhere.

For instance, it takes 1d3 days to pass from any major city with fey or gnomish influence in the mortal world to the major elven or fey courts in the Shadow Realm. Travel from regions without elves, fey, or gnomes often takes a bit longer, on the order of 1d6 days, and the known and functional roads are fewer and less well kept. Travel along a shadow road in a region of little or no magic, such as wastelands, mage-blasted deserts, or the like, can take 1d12 days. For enormous distances, such as across seas and oceans or across continents, travel time can be up to 2d12 days along a shadow road—still much shorter than the months or years such trips normally take.

For creatures from the Shadow Realm, the roads are common enough. For creatures from elsewhere, make an Intelligence (Arcana) check at a −5 penalty for non-elves and consult the following table to determine what (if anything) they know about them.

\begin{DndTable}[header=Table 5-3: Shadow Road Lore]{cX}

DC & Result\\
1–5 & It is not known how the shadow roads work.\\
6–10 & Shadow roads connect two points through magical travel. A spell or passphrase is often required to open one. Elves, shadow fey, and similar peoples seem to have an easier time on them than anyone else.\\
11–15 & The start and end of a shadow road are set in advance, and the doors to them are never open long.\\
16–20 & The destination of a shadow road can be changed but only by fey. Otherwise, new destinations require obscure incantations or arcane tinkering.\\
21+ & Some shadow roads are guarded, and others lead in only one direction.\\
\end{DndTable}

\subsection{Sights and Encounters on Shadow Roads}
Those from outside the Shadow Realm who travel along the shadow roads for adventure find the plane's roads can be an overgrown chaos of hedge mazes, twisting passages, hollows, and delves—each road may vary, but all are notorious for lacking wide vistas or straight paths. Travelers may leave the path and wander at will, but once they attempt to go home, it takes them 1d3 days to find their way. Many have never returned since those who step off the paths sacrifice any protection offered by the fey that built the road or by warding stones and waymarks placed by other hands. Shadow creatures, fiends, and celestials all use these roads for purposes of their own and dislike competition.

In game terms, the shadow roads are filled with dangerous encounters. More powerful adventurers will attract more powerful creatures. Fey and undead are common as well as beasts of the night, marauding giants, and various other shadow creatures. Additionally, the shadow fey guard some roads with their black stone keeps, from which they send their minions forth to prey on the unwary—or at least, to collect tolls in the form of memory philters (see \textit{Chapter 10}), scrolls, or silver.

Travelers using these arcane highways do so at their own risk. (See \textit{Appendix} for suitable encounters.)

\section{Navigating the Black River}
The Lord of the Rivers is Charun the Boatman, and his followers are those who most often travel it, safely and swiftly, just as the followers of the shadow fey pass most swiftly through their forests and the darakhul most quickly through their dark canyons. While the rivers seem in many respects similar to rivers in mortal lands, they also loop and wander across the planes in ways that are difficult to predict. And on the twisting paths of the Shadow Realm, sometimes traveling downriver leads one to places that are surely upriver if one goes by the roads. At other times, one river melds into another. This is especially true when one travels by shadow goblin barge or with the help of a conjured ferryman (see \textit{conjure ferryman} spell in \textit{Chapter 4}).

For normal heroes and travelers without access to a magical ferry, travel by boat is easy enough downstream, providing a steady rate of about 5 miles per hour with the current—but travel is often confoundingly less direct than expected with strange tributaries flowing into the river, other boats seen in the perpetual mists, and various rapids, whirlpools, and backwater oxbows that lead nowhere.

Encounters on the river are common, usually with shadow goblins, river giants, or water fey, and more rarely with fiends or incorporeal undead. Characters must make a DC 15 Intelligence (Arcana) check for each day of navigation. On a success, the party has traveled as expected, and on a roll of 20 they have arrived after a single day on the water. On a failure, they have looped or dawdled or otherwise made far less progress, limited to 10 miles for the day. On a roll of 1 or 2, roll 1d12 and consult the following table.

\begin{DndTable}[header=Table 5-4: Navigating the Styx]{cX}

d12 & Result\\
1 & Boat lands exactly where it embarked in the morning.\\
2 & Ghosts plague the boat all day, annoying passengers with shouts and noxious mists and obscuring sight.\\
3 & Some coins or minor possessions (such as rope, amulet, or favorite socks) go missing.\\
4 & A member of the party gains a level of shadow corruption (see \textit{Shadow Corruption} above).\\
5 & A familiar or companion animal takes ill and must be carried off the boat.\\
6 & A member of the party becomes partly transparent, their bones now visible.\\
7 & The ferry slides into another plane of the GM's choosing, such as one of the hells, the celestial planes, or a mortal land.\\
8 & Everyone aboard the boat is drained by travel and gains a level of exhaustion.\\
9 & The life-draining power of the river inflicts 2d12 necrotic damage to each passenger each hour.\\
10 & The boat enters dark water (see \textit{Dark Water} above), and a \textbf{water elemental} attacks.\\
11 & The boat travels without effort, guided by unseen forces, and all its passengers gain the benefits of a long rest.\\
12 & One passenger has a vision of the afterlife and may have a premonition of danger, learn a secret from a ghost, or gain 2d12 \textit{temporary hit points} from understanding the power of the river.\\
\end{DndTable}

\begin{DndSidebar}[float=!b]{The Waters of the Styx}
The River Styx is a planar waterway found flowing through many planes, and it varies in its color, inhabitants, and even name from realm to realm. Always though, it carries the souls of the dead to the afterlife and binds the many planes of the multiverse together.



In the Shadow Realm, its course is rarely constant as the land itself shifts every few years. However, the creatures that make their living from it know how it runs and how to avoid its dangers. River giants, shadow goblins, and shadow river lords are all relatively common along its banks as are other creatures. (See \textit{Appendix} for suitable encounters.)



\end{DndSidebar}

\section{Lost Lands and Shifting Destinations}
A place as malleable and prone to revision as the Shadow Realm inevitably finds that some places are cut off from the shadow roads, unmoored in the currents, and seen intermittently at best. The driftwood villages are one example, and the city of \textit{Oshragora} another, though its movements are fairly steady. Others were once well known but are now quite difficult to find: the Radiant Isle, the Land of Apple Blossoms, the Keep of the Shadow Drakes, and the Paletree Rise out on the Whispering Plain. All are well-known legends to umbral humans, shadow fey, and others, though none are especially well fixed on a map. A few are briefly noted here as legendary locales for adventure.

\subsection{Keep of the Shadow Drakes}
This set of black towers resembles a tangled grove of trees more than any fey architecture. It was raised by Nonderil, an elder shadow drake, as a gift to his mate, Violanda, and it became a place of draconic power as well as a school for mages of the traditions of shadow and illumination. At various times, it has been seen on the Whispering Plains, as an island in the Shade Sea, and even as a dark cloud castle floating over the hills. In all cases, dozens of drakes, dragonettes, and dragons nest there because they are the only ones capable of opening its gates and doorways. Rumors run rampant that the keep of the shadow drakes is filled with the bones of a thousand fey or with the grimoires of a hundred mages or with the wealth of generations of gnome and human traders. Some even believe that it commands a radiant well that is somehow harnessed to provide arcane power (just as one might tap a barrel of wine). A few claim that the whole building is inscribed in Draconic with the history of the Ancients, who (in this telling) were shape-shifting dragons.

While bones and books are found there, more often, the Keep of the Shadow Drakes is a nesting ground for drakes, protected by their elders, and the dragons and drakes of the Shadow Realm gather there until their wyrmlings hatch and then depart as soon as their progeny can fly, leaving the place to some other type of draconic flock looking for a perfect nesting site. Eggshells litter the courtyards, lending concrete evidence of this use by dragons in breeding season.

\subsection{Land of Apple Blossoms}
Swarming with bees and fragrant with apple blossoms in spring, the Land of Apple Blossoms is sometimes described as an island in the sea and at other times as a set of five hills and seven villages, each inhabited by fair, kind-hearted humans and bearfolk who settled it long ago and wisely decided never to connect their realm to that of the fey.

The Land of Apple Blossoms is the subject of songs and sagas, and many of them feature its treasures: golden bees and honey of immortality, apples whose taste reverses aging and cures the sick, even berry bushes hung with jewels the size of a human thumb, and springs that run with sweet water, pure as snow, that teem with bright-red salmon each autumn. Some of this is surely poets at work, but the existence of a circle of apple druids or the guidance of Halla Berlesfand, a great druid of wisdom and kindness, are themes that often recur in tales of the Land of Apple Blossoms. The humans there are said to grow tall and fair, untroubled by hunger or plague.

The great druid of the Land of Apple Blossoms can open a magical moon portal to the city of Fandeval, which appears once every few years. On those occasions, a few of its young people take its apples, pears, and other fruits to market, exchanging them for Fandeval's salt, spices, and other rare items.

\subsection{Paletree Rise}
This drifting village was built around a circular grove of oaks and a single stone, sometimes thought to be one of the failed stones that Hecate placed in the waters before the Founding. (Indeed, the legends say it floated well and gathered earth and plants to it, though it failed to root in one place.) The village itself is seen in a dozen places each year, drifting faster than most villages, and some of the faithful of Hecate consider it an especially propitious site, worth seeking out for the devout. Others believe that the "pale tree" that gives the village its name is a sign that the place is quite dangerous, unstable, and best avoided. Chaos cultists and the satarre claim the village might be a good place to unbind and unravel magic for this reason.

\subsection{The Radiant Isle}
A secluded island located somewhere in the middle of a black lake in the south of the realm regularly draws dozens or even hundreds of witchlights to its shores. The island can be seen from the far side of the lake thanks to the pale glow that emanates from the place. Brave souls who use magic or take a boat across the lake find a building composed of solid light nestled in an otherwise desolate rocky outcropping. Low and sprawling, the structure has a blue central dome and several white marble wings, radiating from that hub. Witchlights flit across the lake and into the building, gliding through the opaque, glowing walls.

For creatures that can't pass through the light walls, a normal entrance opens into a library, but instead of books and scrolls, sheets of multicolored light are stacked atop one another in orderly fashion. Each sheet of light contains recorded knowledge. The witchlights move about the archive, a few at a time, to add some of their own light to what's already there.

In addition, brass-and-silver library automatons move among the stacks, scanning the library and its visitors. Visitors that convince the automatons to help can find strange lore gathered from across the realm, even the multiverse, stored in the archive. The library is said to connect to other storehouses of wisdom on other planes through doors limned in strange alphabets.

The \textbf{liosalfar}—elementals composed of pure, colored light—are frequent visitors to the Radiant Isle. They have a great love of patterns and a hatred of any creature that seeks to disrupt the patterns of "color" they identify as the world around them. They can offer invaluable insight on locating a specific piece of lore, but any who threaten the archive or the knowledge stored within inspire their wrath.

\begin{DndSidebar}[float=!b]{Shadow Landmarks and Natural Anchors}
While shadow fey, bearfolk, and others raise buildings, castles, and entire cities on the strength of foundation stones that preserve structures during times of flux and tide, other elements of the Shadow Realm also resist the currents and motions of ebon tides. Shadow landmarks and natural anchors perform the same function, and as elements of the plane, they do not require arcane power to create. Instead, these features of the land are notable as geographic features: white cliffs, enormous house-sized boulders out on the plains, a beach of pure white sand, a set of purple dunes, or a particularly stunning mountain peak or outcropping.



These sites are preserved over years without any help from arcanists or fey. Instead, the memory of the Shadow Realm and its inhabitants themselves maintains them. Most are associated with names, legends, or associated figures from history. The older and more visited the site, the more resistant it becomes to damage by black water, harsh radiance, or similar natural effects. Natives of the plane call them "natural anchors," places too distinctive to be washed away or to sink silently beneath other terrain.



\end{DndSidebar}

\chapter{6. Fey Courts and Servitors}
The Shadow Realm has always been most closely associated with the fey, who used it as a source of magic and magical transportation long before they learned to make warded foundation stones and to settle it. Their courts resemble those of the elves of Silendora, the Summer Lands, in their overall size and organization but with a darker cast and often with far more gnomes, goblins, and humans. The shadow fey themselves have always been relatively few in number, though their promises of power and their ability to enchant and control mortal rulers have given them the ability to pull entire cities into shadow.

\section{What is a Fey Court?}
A fey court is roughly similar to a human fiefdom: a land ruled by a noble or monarch, that ruler embodying its people and attracting like-minded fey to their court. A small court, for instance, might style itself the Baron's Court of Millrun and rule over a handful of villages along a stream. A large court, such as the \textit{Court of Night and Magic}, might have lesser courts that pay it fealty, be ruled by a queen, and command entire cities, forests, and small armies.

A dozen major courts and several minor ones are described here. More certainly exist, and new ones could be founded by any fey-blooded inhabitants of the Shadow Realm who lay claim to existing lands and castles or who can master the ebon tides to make their own lands and castles from nothing.

Any fey can found a court, though only the most powerful or charismatic seem to have success in growing them from a backwater named Frog Manor to a wildly popular Court of Emeralds. Shadow fey, shadow goblins, lunar elves, sable elves, wyrd gnomes, quicksteps, living stars, and even trollkin (and sometimes half-elves) have all ruled fey courts. Bearfolk, darakhul, halflings, alseid, ravenfolk, satarre, shades, and humans are sometimes members of a court, but they are not eligible to found or rule one.

Rules for styling oneself a title and naming the court itself are arcane even by fey standards. A new court, for instance, might take a "properly modest" name such as Lilypad Manor, the Court of Cherries, or the Court of Three Bridges, ruled by a lord, lady, or keeper. In time, the court might grow and become the Court of Ruby Blossoms or the Court of the Silent River, ruled by a princess, marquess, or count. Such names can be proposed yet fail to become accepted by the members of the court if a ruler judges their own advancement too generously. When other fey refuse to acknowledge the new placename, it may easily become a mark against an "upstart court," and the court's popularity may wither.

\begin{DndSidebar}[float=!b]{Changelings at Court}
Many fey courts have a tradition of hosting changelings, which they see as an honorable tradition of fostering the children of various non-fey races among themselves. (To the parents of the kidnapped children who become changelings, the matter is somewhat different.) These human children are stolen away as infants, taught Elvish, Umbral, and other refined languages, and their demeanor and personalities change to match that of their shadow fey, gnomish, or sable elvish parents. (Note that the fey do not choose to adopt dragonborn, dwarves, kobolds, tieflings, or trollkin as changelings, and they consider the various animal folk such as alseids, catfolk, ravenfolk, or subek as half-fey already. Kidnapped halflings generally become courtly servants.)



Changeling player characters can use the changeling variant race for umbral humans (see \textit{Chapter 2}), which provides rules governing their abilities.



\end{DndSidebar}

\section{Corremel, the City of Lanterns}
\begin{DndSidebar}[float=!b]{Corremel, the City of Lanterns}
\textbf{Ruler}: \textbf{Hander Svenk, the Black Prince} (see \textit{Courts of the Shadow Fey})



\textbf{Important Personages}: Ander Wyth, Master of the Market Stalls (LE \textbf{shadow fey knight of the road}, see \textit{Creature Codex}); Aten Aerlich, First Blade of the City Watch (LN \textbf{shadow fey duelist}, see \textit{Tome of Beasts 1}); Bort, seneschal to the Black Prince (LE \textbf{shadow goblin scribe}, see \textit{Chapter 9}); Gunrik Stoll (LE \textbf{dwarf necromancer}, see \textit{Creature Codex}); Henrik Stoll (LE \textbf{dream eater}, see \textit{Tome of Beasts 1}); Marti Dwerik, Head Lamp of the Lamplighter's Guild (LN elf \textbf{mage}); Masika Gamall, Mistress of the Pluming Lantern (CG werelion \textbf{war priest}, see \textit{Creature Codex}); Oyrepple Flinders (LE devilbound gnomish prince, see \textit{Tome of Beasts 1}); Rodrig Goldseam (N dwarf \textbf{bandit lord}, see \textit{Tome of Beasts 1}); Shiallee of the Carriage House (N human \textbf{mage}); Smiling Aymag, First Tormentor of Scourge Circle (NE \textbf{shadow fey enchantress}, see \textit{Tome of Beasts 1}); Tamrana Lanzt, clerk to the prince (N \textbf{shadow fey ambassador}, see \textit{Creature Codex}); Tristelia Caernefille, Second Blade of the City Watch (LE \textbf{shadow fey forest hunter}, see \textit{Tome of Beasts 1});



Twylla Yoop, Mistress of the Sheltered Market (LE \textbf{shadow fey knight of the road}, see \textit{Creature Codex})



\textbf{Population}: 50,000 (23,000 shadow fey, 15,000 shadow goblins, and 12,000 umbral humans)



\textbf{Great Gods}: \textit{Sarastra} (patron), \textit{Laughing Loki}, Mammon, Ninkash, and Rava



\textbf{Trade Goods}: Lanterns, papyrus, ale, figs, pork, fish, melons, oranges, lemons, and goats (in order of importance)





\end{DndSidebar}

A bright light on a dark river, the city of Corremel shines like a beacon from a thousand black-iron lanterns. Ruled by the Black Prince of the shadow fey, Hander Svenk, this well-lit city is their sparkling stronghold and the source of their wealth and power. Its foundations are enormous, both in its palaces and in its charcoal-cobbled streets, and its watch is vigilant and suspicious of troublemakers.

A busy metropolis, it is sometimes referred to as Corremel-in-Shadow (but usually only by those of the mortal realm). Corremel is situated on the banks of the River Lethe near two shadow roads to the mortal lands, making the City of Lanterns a distinct trade hub that sees people of shadow mingling with all varieties of mortals. This mash of cultures has left Corremel with a distinct amalgam in architectural style and in its dialect of Umbral, using more human loan-words.

The city's bridge is an important passage for crossing from the forests and courts of the Queen's Wood down to the plains, to the Saltwood, to the Shade Sea, and to the trading opportunities of \textit{Fandeval}. Visitors from throughout the Shadow Realm come for trade goods, for particular services from arcanists, spies, loremasters, and assassins, and for the various joys and pleasures that the fey offer to visitors, for the right price.

\begin{DndSidebar}[float=!b]{Corremel's Connections to Mortal Lands}
Much of the wealth and power of Corremel comes from its twin shadow roads, which connect the city to two other cities of your choice. In the Midgard setting, these cities are Zobeck, the Crossroads City, and Corremel-in-Nuria, a desert city along a crocodile-infested river.



The first of these is a small city of clockwork magic, kobold miners, and trade near the ancient Margreve Forest and Ironcrag Mountains and described in \textit{Zobeck: Clockwork City}. Any temperate trade or river town could stand in for Zobeck in your campaign as needed.



With the latter, the city enjoys a shadow road to the human City of Ale, called Corremel-in-Nuria and described in the \textit{Southlands Worldbook}. If you are not using the Southlands in your campaign, any other sun-filled, human city will do as the "mirror of Corremel."



The idea is to give the shadow fey a connection to places unlike the Shadow Realm and to a road for human foodstuffs, textiles, and other trade goods to arrive in the hands of the shadow fey.



\end{DndSidebar}

\subsection{Prominent Locations}
\subparagraph{1. Horn Gate Inn}
Many mortal visitors to the city find this lively tavern, blaring with music and conversation, to be their best first stop. Located on the northern edge of town, this large, horn-colored mansion with a slate-gray roof is bedecked with numerous gables and turrets. The stone exterior sports numerous balconies with detailed balustrades while the inside features hand-carved trim, hardwood floors, a carved wooden staircase, and massive stained-glass windows. Each bedroom has a stone fireplace and small-but-comfortable bunks or a bed large enough for two. The entertaining hall is an expansive room with a worn wooden floor and numerous wooden tables and chairs.

The music inside the Horn Gate Inn varies from night to night, but it is always merry. Playing for the inn's patrons is an honor granted only to the most talented of musicians. Locals often visit the inn to dance a jig with friends and visitors; those from elsewhere may find fey forwardness a bit off-putting. The Horn Gate Inn deals exclusively in the currency of the mortal realm, rather than any minted in the Shadow Realm, though its owners will accept other goods in exchange for lodging. Food, also available for purchase with coin, is fairly typical of the mortal realm and free of the Shadow Realm's corruption.

Harried Harald, a disheveled, hardworking human (\textbf{commoner}) in his late fifties, tends the bar, pouring drinks and sharing information with patrons. Tamikula, a fair-haired half-elf with sharp green eyes and soft cheeks, waits the tables. Her patience with demanding customers is thin, but she goes out of her way to assist polite travelers. The inn is run by Rodrig Goldseam, a portly, middle-aged dwarf who reports interesting findings to the city's Second Blade, Tristelia Caernefille.

The inn's most opulent room, called the Whispering Room by the staff, is haunted by the spirit of Loris Treik, a maiden who was seduced by Gunrik Stoll and murdered in the room by Stoll's jealous human mistress, Hiernanda Kitza. Loris's ghost rarely appears in the room, preferring to haunt the dreams of those who sleep in it (see \textbf{Whispering Room} sidebar).

\begin{DndSidebar}[float=!b]{Whispering Room}
A creature that takes a long rest in this room is wracked by vivid dreams, which rapidly alternate from a warm embrace and the sound of whispered endearments to skeletal hands wrapped around the dreamer's throat and the sound of hissed curses. At the end of the long rest, the creature gains no benefit from the rest, takes 10 (3d6) psychic damage, and gains one level of exhaustion. A successful DC 20 Wisdom saving throw avoids the psychic damage but not the other effects.



To end the haunting and bring rest to the ghost, Loris's murderer, Hiernanda Kitza (LE human \textbf{archmage}), now a powerful magic user, must spend a full night in the room—this may result in her death or a transfer of the whisper-curse, the constant sound of accusing whispers surrounding Kitza. Ending the curse of the Whispering Room earns Rodrig Goldseam's eternal gratitude, and he offers free lodging to the heroes for the rest of his life.



\end{DndSidebar}

\subparagraph{2. The Pluming Lantern Pleasure Salon}
This establishment, called the Stick and Dome by locals, is a large, square building of turquoise marble. A few small windows sit on each side, and two small doors stand at its front and back. A crystal dome atop the building casts broken rainbows and provides much of the light inside the otherwise solid structure.

A single narrow tower juts up from one corner. Capped with a pointed roof, the almost featureless structure has a single balcony that rings its girth.

Layers of perfume and incense hang in the air and stick to any cloth inside the pleasure salon. Plush carpets line its floors, and tapestries depicting lewd images of fey consorting with bestial humanoids adorn the walls. The colors, once vibrant, have been bleached by time and the Shadow Realm.

A public bath costs 1 sp. A private bath, which can accommodate up to eight Medium creatures, costs 1 gp per attendee. A secret bath, named for its enchantments that prevent magical or mundane eavesdropping, costs 100 gp per attendee. A secret bath can accommodate up to ten Medium creatures and is a popular place to make deals despite the high cost.

Those seeking to exercise more carnal desires can find nearly any consensual pleasure at the Pluming Lantern, barring congress with children or the dead. The rates for such activities range from 10 gp for 30 minutes with one of the skilled courtesans to several thousand gold for such complex joys as being immersed completely in the daydreams of young lovers while being whipped by gossamer-strand-wielding sprites. Payment at the Pluming Lantern must be made in mortal coin, and the establishment does not trade services for goods or other services.

The Pluming Lantern is run by Masika Gamall, a priestess of Bastet in exile from her home in Per-Bastet. Fleeing that city's ruler, Meskhenit, for her hand in directing the slaughter of a group of slavers, Masika made her way to Corremel-in-Shadow. She was unaware that the slavers were spies returning with urgent news regarding an impending dragonborn incursion against the City of Cats. Meskhenit is furious with Masika's actions and will pay handsomely for information that leads to her capture.

\subparagraph{3. The Sheltered Market}
This massive, ancient obsidian building is surrounded by fluted columns that support a gently sloped roof. Inside the Sheltered Market, friezes decorate the walls and ceiling, depicting diplomatic scenes of the shadow fey interacting with others. Most of the goods that make their way into Corremel are bought and sold here.

Unlike the market stalls in other parts of the city, the Sheltered Market is almost staid. Spaces are uniformly sized lots with the more prosperous merchants taking several adjacent lots to display their wares. Mistress Twylla Yoop, the owner of the Sheltered Market, runs it with greedy efficiency, levying fines against any creature who disobeys the edicts clearly posted around the market. When disputes arise, she rules in favor of whomever pays her the largest bribe, and she often levies a steep tariff on the loser for wasting her valuable time. Twylla is a confidante and sometimes lover of Gunrik Stoll, co-leader of the Stoll Cartel, and she turns a blind eye to most cartel activity in her demesne.

Those who cannot pay their fines at the Sheltered Market are held for humiliation at Scourge Circle. Detainment usually lasts 1 day for every 50 gp in unpaid fines, though punishments vary at Twylla's discretion. She is more lenient toward musicians and much harsher toward kobolds and goblins.

\begin{DndSidebar}[float=!b]{Going to the Sheltered Market}
In theory, one can find nearly any non-artifact item at the Sheltered Market, though availability and price vary:




\begin{itemize}
\item A buyer can find any mundane, non-consumable item in 1d3 hours. A successful DC 10 Charisma (Deception, Intimidation, or Persuasion) check reduces this to 1 hour.
\item A buyer can find an exotic item or a common or uncommon magic item with 1d6 hours of talking and searching and with a successful DC 20 Intelligence (Investigation) or Charisma (Deception, Intimidation, or Persuasion) check.
\item A customer can find a rare magical item with 4d6 hours of talking and searching and with a successful DC 25 Intelligence (Investigation) or Charisma (Deception, Intimidation, or Persuasion) check. Vendors selling these items rarely accept coin for them, preferring barter.
\item Very rare or unique magical items might not be available, at the GM's discretion, and a customer looking for such an item usually must endure a dangerous task in exchange.
\end{itemize}



\end{DndSidebar}

\subparagraph{4. Prince Hander Svenk's Palace}
The Black Prince's palace is a series of six amethyst towers, each topped by a statue of a former city ruler. Each tower is attached to a larger, central tower, which is topped with a conical-pyramid roof and flying buttresses.

The lower floors of the central tower are open to the public, and visitors often tour the tower's ancient and exquisite collections of art, trophies, murals, friezes, and wall hangings. The tower's upper floors are closed to the public and house the Black Prince and his numerous concubines and shadow fey servants. Excepting his seneschal, Bort, the Black Prince's hundreds of shadow goblin servants live in two basement levels below the central tower.

On occasion, the Black Prince's younger sister, a pale shadow fey named Rosali Svenk, wanders the floors of the central tower, unnerving staff and visitors alike. No matter her location in the tower, its occupants always hear the subtle clinking of the tiny animal skulls tied to her gnarled staff as if she is walking just behind them. For more information on Rosali Svenk, the Pale Witch, see \textit{Courts of the Shadow Fey}.

The six adjoining towers are cleverly designed prisons, housing the few creatures who have truly pricked the Black Prince's ire. Three of the towers are occupied and hold the following prisoners:


\begin{description}
\item \textbf{\textit{Nuallia Hillier}.} The Black Prince's first lover was imprisoned when she spurned him. Nuallia can move through her beautifully appointed tower freely, but she is treated with indifference by any guards or servants who attend to her needs.
\item \textbf{\textit{Ceres}.} This lunar elf assassin came close to publicly ending the Black Prince's life. The tower is full of greased platforms and illusory warriors who often beat the prisoner.
\item \textbf{\textit{Tig}.} This shadow fey critic made the mistake of publishing a satire, depicting the Black Prince as a kind-hearted fairy. His tower prison resounds with mocking laughter at all hours. Every time he moves to another room, an audible parody of his sentencing hearing plays throughout the tower.
\end{description}

\subparagraph{5. Firefade Park}
Occupying the northern banks of the River Lethe, this park is filled with maple and oak trees separated by narrow paths of gray cobblestone. The entire park is suffused with a dim orange-yellow illumination, and the park's central clearing is never empty as visitors and residents alike come to enjoy the view of the park's large fountain.

A majestic fountain, made of the same white marble as the River Bridge, depicts a bear's head in profile with a human head caught in its jaws. The human head in turn has a frog in its mouth, and the frog has a carp by the tail. The carp's body is carved as if the frog caught it mid-leap, and its mouth emits a constant stream of pure, clear water. The fountain is guarded by six well-armed shadow fey guardians (see \textit{Tome of Beasts 1}). Those wishing to drink at the fountain must pay the clerk, Tamrana Lanzt, before they can drink the water. The clerk works on behalf of Prince Hander Svenk and accepts silver coin as well as more esoteric payments for use of the fountain.

Fire damage is halved while in the park, and creatures caught attempting to burn the trees are detained for trial at the Black Prince's leisure. Those caught drinking the water without sufficient payment are also detained for trial.

\begin{DndSidebar}[float=!b]{Firefade Park Fountain}
The crystal-clear water flowing from this fountain comes directly from the Elemental Plane of Water. It is pure, cold, and delicious. A creature that drinks this water removes one level of exhaustion and gains \textit{temporary hit points} equal to its CR or level. Once a creature has drunk from the fountain, it must finish a long rest before it can benefit from the fountain's water again.



\end{DndSidebar}

\subparagraph{6. Shady Lane}
While the bright lanterns of the city illuminate most of the streets and buildings, less savory individuals have a knack for discovering darkened nooks and crannies. Shady Lane, located off Guild Avenue, is a dark, narrow alley. Lanterns are dimmer here, and even magical light is muted, much to the consternation of the city's lamplighters. The light in Shady Lane is never brighter than dim light. On rare occasions, creatures who have stumbled into Shady Lane exit confused onto another road in the city—one which is not actually connected to the lane.

\subparagraph{7. Lamplighter's Guildhall}
In the Shadow Realm's most well-lit city, the Lamplighter's Guildhall, unsurprisingly, sees reasonable traffic. This ostentatious mansion has multiple gables and turrets, with a fine exterior of wood painted bright blue with gold trim. While the building has no balconies, a large, covered porch runs along one side and over the main entrance. The guildhall has many windows, and its tallest turrets are four stories high.

The current head of the Lamplighter's Guildhall is Head Lamp Marti Dwerik, a ruthless careerist who detects and exterminates potential rivals before they can formalize any sort of opposition. Marti is a controlling guildmaster who demands her dozens of underlings report directly to her throughout the day before returning to their task of ensuring the lanterns about the city remain lit.

Most of the lights in the city are magical, and the Lamplighter's Guild includes numerous magically talented individuals. Dozens of guild members with little to no magical ability, referred to as "skips" behind their backs, light, douse, and maintain the non-magical lights in Corremel. These guildmembers also report any broken or counterspelled magical lights to the Illuminator's Committee.

In recent years, the Lamplighter's Guild has been infiltrated by doppelgangers working for the Stoll Cartel. As new doppelganger agents become available, Henrik Stoll sends them to murder a member of the guild and assume that member's life and personality. It is the Stoll brothers' hope that soon their agents can rise up and butcher the guild's leadership. Once this occurs, the brothers plan to assume control of the guild and exploit homes and businesses for preferential illumination. Unknown to the Stoll brothers, most of the doppelgangers actually enjoy their new lives and aren't too keen on replacing the guild's current leadership.

\subparagraph{8. The River Bridge}
Easily the most notable feature of the city, the white marble bridge spans the black River Lethe and is featured on the city's coat of arms. It provides one of the few ways to safely cross the specter-haunted river, making it one of the busiest and most traveled locations in the entire city. The bridge is truly ancient, and its creation predates the founding of Corremel. Lunar elves and devotees of Hecate claim that the goddess herself set the bridge's foundation stones, and small candle-lit shrines sometimes spring up to the goddess on either bank.

Each of the bridge's supports features a carving of an unfortunate soul recently lost to the waters of the River Lethe. The expressions of the graven figures are always both horrified and haunting. The images are ever-changing as the Lethe quickly claims the unwary who fall into its dark waters. The Lanternkeepers regularly patrol the bridge.

\subparagraph{9. River Lethe}
One of the major waterways in the Shadow Realm, the River Lethe is a dangerous torrent of inky shadow, running through Corremel. Like all the waters in this realm, the River Lethe is frigidly cold and haunted by the creatures who perished in or near it. Many a soul has attempted to wade or swim across this river only to be consumed by it or its hungry spirits, making safe passage via ferry or the River Bridge so attractive.

A shadow road stretches along the riverbank inside the city, connecting Corremel to the Crossroad region near the Free City of Zobeck. Shadow goblin guides sometimes offer visitors assistance in opening the shadow road's portal to begin the journey to the mortal world, though in most such cases they merely open a portal and do not participate in any further travel.

\begin{DndSidebar}[float=!b]{The Law in Corremel}
All visitors to the city should be aware of some of its stranger laws.




\begin{itemize}
\item To speak ill of the \textbf{Queen of Night and Magic} or the Black Prince is treason.
\item Vandalizing an image of any member of the court is grounds for permanent expulsion from the city.
\item Damaging another's reputation will be punished as severely as if the offender had harmed the victim's body—possibly more so. (And yet, gossip is a local pastime.)
\item An offender has no right to representation at trial. Guilt or innocence can be determined in the offender's absence.
\item Ignorance is no excuse for acting ignorant.
\item The Black Prince has the final word and is exempt from all laws.
\end{itemize}



\end{DndSidebar}

\subparagraph{10. Stoll's House}
Like a dead flower in a flawless garden, this small dwelling stands out from the beautiful buildings in the southwest quarter. Peeling paint, crumbling roof tiles, and other evidence of neglect make this run-down peasant's abode appear to belong more in the Tumbles than anywhere else.

However, Stoll's House is a facade. The extensive network of catacombs below it houses the network of spies, assassins, and other criminals Gunrik and Henrik Stoll employ. Officially, the group is called the Queensman's Guild, though most locals simply call it the Stoll Cartel. The guild is responsible for most organized crime in Corremel, excepting slavery, which is run by Oyrepple Flinders (see \textbf{Joyful Park} below). The guild's inability to break into the slavery racket is a sore subject with the Stoll brothers, and they would richly reward any group or individual who could help them dismantle Flinders's organization.

The Queensman's Guild is largely comprised of human, shadow fey, and doppelganger agents with a few halflings, tieflings, shadow goblins, and half-elves. They do not employ darakhul, lunar or sable elves, bearfolk, or wyrd gnomes. The organization has 117 members with other specialists contracted as needed.

The Stoll brothers appear to be cantonal dwarves, and indeed Gunrik is what he appears to be. Henrik, however, was dispatched and replaced by a \textbf{dream eater} (see \textit{Tome of Beasts 1}) about a decade ago. Since Henrik's ideals were largely the same as his replacement's, the change has gone unnoticed.

\subparagraph{11. Joyful Park}
Joyful Park sits on the southern banks of the River Lethe. Countless visitors to Corremel stop for a rest here, enjoying the park's proximity to the bridge and its many street entertainers, performing at the small amphitheater. Very few guests are aware of the park's secret: a subterranean slave market sits beneath Joyful Park, where slavers auction off unfortunates to the highest bidder. The varied assortment of bidders makes it difficult to determine a slave's fate. Some purchasers work their new property to death while others prefer to house their recent acquisition in comfort, merely withdrawing a memory at a time until there is nothing left.

The slave market is overseen by Oyrepple Flinders, a gnome who secretly serves Nislev the Tidy, a \textbf{young blue dragon}. Oyrepple stole the dragon's egg from a smuggler who escaped into the Shadow Realm from a desert kingdom over 70 years ago. He cultivated the dragon's greed and vanity over the decades and now serves the powerful creature in the shadows of Joyful Park. Nislev's ultimate goal is to assemble a hoard that contains a living member of every sentient race, and he encouraged the gnome to open the slave market a few decades ago to further this goal. Oyrepple's minions, a seemingly endless number of \textbf{dark folk} (see \textit{Creature Codex}), murder anyone who causes trouble for him, Nislev, or the slave market.

\subparagraph{12. Market Stalls}
This collection of carts, makeshift structures, mats piled high with goods, hawkers, and wandering peddlers is a riot of colors, scents, and sounds as suits the ground where the majority of the foodstuff in Corremel is bought and sold. Vendors from the mortal world, the Shadow Realm, and many other planes fill the stalls, trying to earn a living in the City of Lanterns. They learn quickly there is no point getting too comfortable in a location as they are unlikely to have the same spot on their next visit to the city.

The Market Stalls are overseen by Ander Wyth, a charmingly wry shadow fey who treats all vendors equally. When disputes between customers and vendors are brought to him, Ander generally rules in the most ironic fashion possible, ensuring neither party is truly satisfied.

\begin{DndSidebar}[float=!b]{Going to the Market Stalls}
Given enough time, a market visitor can find almost any type of food or drink, prepared or raw, in the Market Stalls:




\begin{itemize}
\item A visitor can find common foodstuff, including dishes from all nations of the mortal realm, in less than an hour.
\item A visitor seeking more exotic food and drink must spend 1d4 hours searching the chaotic jumble of stalls and succeed on a DC 15 Intelligence (Investigation) or Charisma (Deception, Intimidation, or Persuasion) check.
\item A visitor wanting a rare or illicit substance must spend 2d4 hours searching and succeed on a DC 20 Intelligence (Investigation) or Charisma (Deception, Intimidation, or Persuasion) check. At the GM's discretion, some substances may not be found at the Market Stalls.
\end{itemize}



\end{DndSidebar}

\subparagraph{13. The Tumbles}
Most buildings in the Tumbles are crumbling or rotten, and some homes are missing walls, stairs, or large sections of roof. While they remain standing, pieces of houses have been known to shift and fall on unsuspecting passersby. Inhabitants of the Tumbles are often desperate for food, water, or medicine, and most residents of Corremel avoid the area.

The fabric of the Shadow Realm is fragile here, and invisible portals to other planes form frequently. These portals are short-lived, eventually collapsing and tearing small portions of the nearby area away with them. Those adventurous or desperate enough to brave the area wander aimlessly among its buildings, looking for an escape to somewhere else. A few enterprising locals sometimes hire their services out as guides. A small shrine to Laughing Loki stands here, though residents say it can be found only by those who aren't looking for it.

\begin{DndSidebar}[float=!b]{Planar Travel Via the Tumbles}
For every hour a character spends wandering through the Tumbles, there is a 1 in 20 chance they encounter a portal. A character must succeed on a Wisdom (Perception) check to detect and avoid the portal. A character must succeed on an Intelligence (Arcana) check to discern the portal's destination. The character has advantage on this check if they know the \textit{gate} or \textit{teleportation circle} spells (or similar magic from \textit{Chapter 4}).



Use the following table to determine the destination of the portal, its DC, and its noticeable characteristics. Choose a plane from the cosmology you prefer that matches the themes listed.



\end{DndSidebar}

\begin{DndTable}[header=Table 6-1: Portals in the Tumbles]{cXXX}

d20 & Destination Plane & Arcana DC & Characteristics\\
1–2 & Heavens or hells & 13 & The portal is slightly warmer than its surroundings and smells of ozone (heavens) or sulfur (hells).\\
3 & Astral or the Void & 13 & A chill hangs in the air around the portal.\\
4–7 & Pocket dimension (a hub of trade) & 15 & The sharp scent of spices and the sound of thousands of whispering voices waft through the air near the portal.\\
8 & Lawful, clockwork plane & 14 & The air around the portal smells and tastes of iron.\\
9 & Elvish realm or plane & 14 & Faint motes of floating light dance around the portal.\\
10 & Plane of war & 13 & The scent of blood hangs in the air around the portal, and the sounds of clashing blades and screams echo from it.\\
11–16 & Random Shadow Realm location & 16 & Shadows thicken around the portal, and the air feels oily.\\
17–20 & Random mortal world location & 17 & Colors near the portal are brighter, and the air smells fresh.\\
\end{DndTable}

\subparagraph{14. Scourge Circle}
This small square is tiled in marble and contains dozens of cages, platforms, stocks, pillories, and other devices used to punish criminals. At any given time, upward of 20 creatures are suffering punishment of some sort.

While painful punishments and executions are not unknown in Corremel, the most common form of punishment is humiliation. The First Tormentor, Smiling Aymag, is particularly inventive in devising humiliations to inflict upon the prisoners in her care. She is currently fond of having prisoners stripped, shaven, coated in honey, rolled in breadcrumbs, and posed improbably for the viewing public. All creatures known to have been punished here have −1 to their Status. (For more on Status, see \textit{Chapter 1}.)

Prisoners are housed in sparse cellhouses located directly north of Scourge Circle.

\subparagraph{15. Shiallee's Carriage House}
This expansive lot has ample space to park carriages, and the staff makes minor and major carriage repairs. Travelers find this a worthwhile stop as the workers of the carriage house are knowledgeable and meticulous.

An immense white-brick mansion dominates the beautifully maintained lot, and it welcomes all visitors to Corremel—for 5 gp of course. This fee often limits the carriage house's clientele to denizens of the Shadow Realm and the more successful merchants of the mortal world. The carriage house, owned by the elusive and rarely seen Shiallee, is run by the numerous people under her employ. It has 20 guestrooms, a small bathhouse, two dining rooms, a small ballroom, a large study, and a games room. All of the guest rooms are on the upper floor while the shared rooms are located on the main level. (A separate wing belonging to Shiallee herself is closed to the public.)

The bathhouse can house no more than six people at a time, but admission, and a bar of scented soap, is included with a stay at Shiallee's. The large study houses an extensive collection of books from various planes on a broad range of topics, from angelic plumage to horn structures in demons to the customs of the shoth. The games room features a billiards table, darts, chess, assorted cards and dice, and other games. Guests dance each evening within the ballroom, a beautifully detailed space with carved wooden panels and exquisite glass lanterns, floating near the high ceiling. A quintet of shadow fey serves as the house musicians, nicknamed the Handsome Five.

Coin and other goods are rarely accepted as payment in the mansion, though they may be accepted for work on a carriage. Shiallee prefers payment in favors, services, or information and often sells such things to others. If it happened in the Shadow Realm, chances are Shiallee has heard of it or knows someone who has.

\subparagraph{16. Lanternkeepers' Barracks}
This large communal home for the city watch is built of gray stone with a peaked black roof. Members of the watch are known as Lanternkeepers for the distinctive lanterns they carry as their badges of office, though locals just call them Keepers.

Martial equipment litters many of the rooms, cleverly scattered to allow the Lanternkeepers easy access to weaponry while hindering those unfamiliar with the layout. First Blade Aten Aerlich and Second Blade Tristelia Caernefille run the barracks with strict adherence to time and duties. A small building near the gates of the barracks houses an on-duty Lanternkeeper who is meant to assist the general public. In truth, most Keepers at the gatehouse post despise it and pointedly ignore visitors. The gates to the barracks themselves are kept closed, except during changes of the watch.

Several hundred Lanternkeepers serve in Corremel, with a minimum of twelve five-person patrols active in the city at any given time. An average patrol consists of a \textbf{city watch captain} (see \textit{Tome of Beasts 1}) and four guards, though a \textbf{mage} or \textbf{priest} of Sarastra may join any given patrol. Each Lanternkeeper carries a distinctive oval lantern made of silver that glows with magical light. These lanterns serve as the badges of office for each Lanternkeeper, and it is a crime punishable by death for a non-Lanternkeeper to possess one. If a lantern is ever lost or stolen, the entire barracks comes to life, striving en masse to retrieve the lantern.

\subparagraph{17. House of Dream Birds}
This square home surrounds a closed courtyard filled with lemon trees, flowering vines, and many small birds with white plumage. They are all tended by Ainoress Gabboneth, a sable elf matron who is the priestess of this strange temple to Gytellisor. Devotees come to feed the birds, to write gloomy poetry or joyous love notes, and to drink wine or smoke requiem in the evenings. The place seems more a fashionable public house than a temple, though each day the high priestess of the dreaming god and her acolytes perform the rites and sacrifice a bird on the god's small basalt altar in the courtyard. A scriptorium in the temple copies out the god's sacred verse and (for a fee) anything else that either the faithful or merchants, lovers, or spies ask them to write.

\begin{DndSidebar}[float=!b]{Saving the City of Lanterns / from the Book of Ebon Tides}
Alas, Corremel, where magic rose to such heights along its arcane river that all its beauty was swept into shadow! Many shelter there, stepping between the bright world and the shadow world to bring all the bounty of the mortals to the fey. It often seems a place waiting to return to the sunlit world. A small gathering of mages there, calling themselves the Bright Sparks, dares to entertain the notion that Corremel could fall back out of shadow and become a mortal land again if only the attention of the right old gods were brought to bear.



\end{DndSidebar}

\section{Court of Golden Oaks}
\begin{DndSidebar}[float=!b]{Court of Golden Oaks}
\textbf{Ruler}: Her Kind Majesty and Queen of the Sable Elves, Valda Limineska, Keeper of the Golden Grove, Mistress of Laughter, Mother to Dawn, Her Resplendent Glory Undimmed (CG sable elf \textbf{archmage})



\textbf{Important Personages}: Ainu-Olla the Swift‑Handed, Captain of the City Watch (NG sable elf \textbf{city watch captain}, see \textit{Tome of Beasts 1}); Bowyer to the Court Vilhopelli the Generous (LE \textbf{shadow fey knight of the road}, see \textit{Creature Codex}); Chamberlain and Court Sorcerer Lumi-Livia the Kind (NG \textbf{sable elf hierophant}, see \textit{Chapter 9}); Chief Librarian and Royal Gardener Drossilda of the Golden Eye (LG bearfolk \textbf{apostle} of Bengta, see \textit{Tome of Beasts 3}); Chief Magistrate and Arcane Proctor Justee Arnosh the Seventh (NE shadow goblin \textbf{archmage}); Ekkitta‑rup‑Wisewhisker, Keeper of the Golden Acorn (CG \textbf{ratatosk warlord}, see \textit{Creature Codex}); Grumbor Bluebelly of the Second Circle (N bearfolk \textbf{druid}); Keeper of the Tree, Maariko-Tirhi (NE human \textbf{veteran})



\textbf{Population}: 9,000 (6,000 sable elves, 1,000 shadow goblins, 900 wyrd gnomes, 600 bearfolk, 300 umbral humans, 150 half-elves, and 50 ratatosk)



\textbf{Great Gods}: Valeros (patron), \textit{Baccho}, \textit{Bengta}, \textit{Charun}, and \textit{Sarastra}



\textbf{Trade Goods}: Furs, silver, amber, salt fish, oak, wool, and tin (in order of importance)





\end{DndSidebar}

This court is the strongest holding of the sable elves in the Shadow Realm, ruled by the bright spark that is Her Resplendent Glory Undimmed, Queen Valda, whose joy and sternness are fused into a perfect golden-haired combination of careful planning, shrewd insight, and arcane power. As a mistress of the arcane, Valda has outwitted and ensorcelled suitors, foes, and assassins, and as a queen, she has won the hearts of her own people since she ascended the throne and truly began to rule as a young sable elf.

In her days as a bride and young royal, most elves of her court considered Queen Valda a charming, well-educated ornament rather than a serious ruler or head of state. However, when her husband, the Crow Prince, abandoned the court to seek wisdom from Yggdrasil, she was perforce required to suddenly leave her books, calligraphy, and arcane musings, and she took command to keep the fortress of the sable elves from falling to its enemies. She has taken little time to think of where the Crow Prince is and what he might be up to on his quest, for the queen has enemies aplenty among the shadow fey, the satarre, and the gnomes. The court's position at a planar nexus means that outside forces both strengthen it (when they seek to trade and pay handsomely for goods from far afield) and undermine it (when they seek to seize its position for themselves).

The court's patron god is Valeros, patron of the elves, and both the queen and her people value martial strength, arcane power, and personal courage. Archery is a popular pastime, and the court bowyer, Vilhopelli, produces astonishing, sometimes-enchanted bows, arrows, and fletching, using materials from the Golden Oak. However, Vilhopelli always demands great favors or other magical items in exchange—his sobriquet, "the Generous," a bit of an insult as he is notoriously tight-fisted.

\subparagraph{The Golden Branch of Yggdrasil}
The strength of the Court of Golden Oaks is closely linked to its position within and around a World Tree, the Golden Oak, and its various paths and mystical acorns. The Golden Oak itself towers as most World Trees do, and the portal to its branches is treated as any other gate or road out of the court: watched and respected as the Golden Oak connects easily to Grenstad, the city of the ratatosk on Yggdrasil, as well as to the Uppermarket of a pocket plane devoted to trade, and Wotan's Tree in the snowy northlands of the mortal world. Ratatosks make up a small but loudly chittering minority in the court itself, and many of the boldest serve as guides and heralds to other realms.

Merchants and travelers who come through the Golden Oak vary widely and include celestials, humans, ratatosks, tieflings, satarre, and even fiends and fey lords. For some reason, almost all are granted passage so long as they swear to remain peaceful within the court and its various fiefdoms, villages, and holdings. Those who refuse are turned back with wand and sword if necessary.

\subparagraph{Sweet Field of Victory}
The Court of Golden Oaks is loved by humans, shadow fey, and others who are primarily visitors for the competitions, tourneys, and feats of strength that Queen Valda sponsors. As is tradition in many elven courts, Queen Valda declares the contests, called the Tournaments of Oak and Ash, when the season seems right. She invites all comers to the court to prove their ability as archers, as knights, or as wizards. Sometimes these include particular tests such as jousts, melees, target shooting, shooting on the wing, battle magic, counterspell duels, feats of horsemanship, drawing shadow power, or tilting at a golden ring. At other times, she may include a test of poetry or dancing, merely to show that not all a knight's prowess comes from strength of arms. The contests are the ones the queen declares, and the winners are given prizes and win renown both at the court itself and among other elves, shadow fey, gnomes, and the like.

\subparagraph{College of the Acorn}
While outsiders prize the tourneys and the spectacle, the queen and her advisors spend more time concentrating on studious matters and raising new battle mages for the Company of Golden Wands, the court's arcane soldiers. These troops are trained by Arcane Proctor Justee Arnosh, a shadow goblin wizard of great power but (unusually for a shadow goblin) a very humble attitude. He delights in serving Queen Valda and bringing glory to the court and in teaching young humans, sable elves, and shadow goblins the rudiments of battlefield magic. Ratatosk rune wizards and bearfolk druids sometimes also contribute to the curriculum.

Both Court Sorcerer Lumi-Livia the Kind and the golden-furred bearfolk Chief Librarian Drossilda serve as instructors, though much of the work of repetition and testing is done by various sable elf apprentices and underlings. The College of the Acorn also maintains a tradition of asking all visiting druids, wizards, and sorcerers to address the students on a topic of interest.

\begin{DndSidebar}[float=!b]{Adventures at the Court of the Golden Oak}
While the court is generally peaceful and civilized, it occasionally attracts dangerous visitors:




\begin{itemize}
\item A fiend enters the court via the gate in the Golden Oak and demands an audience in the name of one of the dukes of hell. Surprisingly, the queen agrees to meet with it but asks the PCs to provide additional sword arms.
\item A druid of the gardens is found as little more than a dry husk among the peonies. A \textbf{vampire} or \textbf{wraith} has drawn all life from the body and now haunts the gardens each night.
\item A courtier asks a PC spellcaster to provide a lecture on an arcane or druidic topic. During the lesson at the college, the speaker is booed, insulted, and heckled by an old enemy who is now a student here.
\item A messenger sent out to the Golden Oak has not returned from Grenstad, city of the ratatosk. The queen and the captain of the guard alike fear a gang of \textbf{void cultists} (see \textit{Creature Codex}) or simply bandits may have established themselves near the portal. The sable elves have not been able to find their stronghold, and Watch Captain Ainu-Olla is now offering a reward.
\end{itemize}



\end{DndSidebar}

\section{Court of Mice and Ravens}
\begin{DndSidebar}[float=!b]{Court of Mice and Ravens}
\textbf{Ruler}: His Infinite Wisdom and Sly Omnipotence, His Jocund Majesty King \textbf{Reynard}, the Fox Lord (see \textit{Creature Codex}), the First Keeper of Ravens, Master of the Roads, Lord of the Blue Pine Forest and High Prelate of Laughing Loki



\textbf{Important Personages}: Chief Scout and Master of War Grumfell Skarsgaar (LN bearfolk \textbf{field commander}, see \textit{Tome of Beasts 3}); Court Jester Illahellan Silverbells (CG sable elf \textbf{thief lord}, see \textit{Creature Codex}); Darresh‑kip-Bristletooth, Swift Flying Lord of Branches (CG \textbf{ratatosk warlord}, see \textit{Creature Codex}); Haberdasher Royal to the Court Tii-kaw Tokkaw (CE ravenfolk \textbf{archmage}); Master of Protocol Winsy Cooperson (LN courtfolk halfling \textbf{assassin}); Most Watchful Guardian of the Gate Unwina the Cook (NG erina \textbf{city watch captain}); Noble Trickster and High Prelate of Santerr Marossa Carashtella Fimbulfinger (CN shadow goblin \textbf{first servant} of Santerr Marossa, see \textit{Tome of Beasts 3})



\textbf{Population}: 8,000 (3,000 shadow goblins, 2,000 sable elves, 800 wyrd gnomes, 600 erina, 500 courtfolk halflings, 400 ravenfolk, 400 ratfolk, 150 bearfolk, 100 alseid, and 50 ratatosk)



\textbf{Great Gods}: \textit{Laughing Loki} (patron), \textit{Baccho}, \textit{Charun}, \textit{Santerr Marossa}, and \textit{Sarastra}



\textbf{Trade Goods}: Furs, dried meat, maple wood, dyed leathers, hats, ironwork, and musical instruments (in order of importance)





\end{DndSidebar}

Ruled—or at least sponsored—by the fey lord Reynard, the Court of Mice and Ravens is an amalgam of trickery, clever plots, and alchemy, mixed with impulsive magic and with a bent toward the mystical or humorous. While it may seem like chaos on the surface, the erinas at the bridge gates manage to keep away most threats to the island-city itself, and the Haberdasher Royal seems to always be preternaturally aware of things, either through divination or his powerful network of ratfolk and other furred and feathered spies. The city is also known to have a shadow road connection to the court of the Mouse King.

\subparagraph{Scoundrels Galore}
Rogues and bards adore their patron lord, and this is why the place is often called the Rogue's Court among the fey. The variety of inhabitants, courtiers, officious minions, and seal-toting officials is high—as is the number of fraudsters, charlatans, jongleurs, fire-eaters, runaway apprentices, and disinherited noble children. Every inhabitant has a story, and most seek just a little help to be restored to their rightful title, shape, inheritance, or lands. Alongside a small but lively ghetto of ratfolk, jokingly called the Pantry, a mixture of sable elf rogues, human bards, courtfolk halflings, fey owls, bearfolk, alseid, and erina are all relatively common.

\subparagraph{Lost Goods and Trickery}
The Court of Mice and Ravens has a thriving trade in lost goods and somewhat shady potions. Outrageous Philter of Lust or Philter of Enchanted Fortune (see memory philters in \textit{Chapter 10}) are readily available here—as is the highly addictive requiem (see Requiem, Drug of the Afterlife sidebar in \textit{Chapter 7}). The wyrd gnomes grow a few rose golems (though not as many as at the \textit{Court of Pale Roses}), and gnomish pranks involving exploding frogs, luminously painted bats, and magical balloons and fireworks are relatively frequent. Items purloined anywhere in the Shadow Realm can sometimes be found here—or at the very least the well-informed master of protocol or the extremely keen-eyed master of the gate may have noted it passing in or out of the court by boat or via one of its many bridges.

\subparagraph{Stepping-Stone River}
The court itself is built on a series of islands on a river, deep woods to one side. Its gardens are lush, and small boats and barges connect the various noble houses on the islands: some have bright towers, others feature large docks and a constant hum of activity, and a few are silent and a bit ominous with treelined shores, few lights, and a bad reputation.

\section{Court of Midnight Teeth}
\begin{DndSidebar}[float=!b]{Court of Midnight Teeth}
\textbf{Ruler}: Countess Phylomara Gladrienne, Countess of Flesh and Secrets, \textbf{Mistress of Midnight Teeth}, Mistress of the Blue Barbers, Patron of Gnomes, The Divine Gossip, Baroness of Gloom, Conjuror Supreme, Lady of Hungers, and Keeper of the True Mirror (see \textit{Tome of Beasts 2})



\textbf{Important Personages}: Court Mage Trissmena Istovash, the Invisible Illusionist (LE gnome \textbf{archmage}); General Inspector Kubberolus Gingercrick (LE wyrd gnome \textbf{field commander}, see \textit{Tome of Beasts 3}); Guildmaster of the Blue Barbers Garrolo Goldenlock (NE wyrd gnome \textbf{bandit captain}); Ingo Ingomarrick, Royal Mirror Polisher (LN erina \textbf{gladiator}); Keeper of Ravens Visilius Barbatus (N human \textbf{keeper of ravens}, see \textit{Chapter 9}); Librarian‑in-Chief Melora Sadrienn (LN gnome \textbf{assassin})



\textbf{Population}: 5,600 (2,300 wyrd gnomes, 1,500 umbral humans, 800 shadow fey, 500 quicksteps, 300 erina, 100 sable elves, and 100 ravenfolk)



\textbf{Great Gods}: \textit{Hecate} (patron), \textit{Ailuros}, \textit{Alquam}, \textit{Charun}, and \textit{Santerr Illosi}



\textbf{Trade Goods}: Mirrors, scroll paper, silk cloth, fine wines, spiced sausages, illusion seeds, leatherwork, and woolens (in order of importance)





\end{DndSidebar}

Gnomes are peculiar creatures, all pointed teeth, narcissistic tendencies, blue hair or—in the case of the Countess of the Court of Midnight Teeth—all three. Countess Phylomara Gladrienne, an influential and loyal vassal to the queen and king of the \textit{Court of Night and Magic} (and sometimes the \textit{Court of the Moonlit King}), is a refined and fashionable gnome bedecked in fine silks and gleaming jewels that compliment her azure hair. She styles her locks elaborately and never quite the same way twice, setting trends among her courtiers. Beneath this veneer of civility and style however, blood and secrets await. Manicured lawns, incredibly lifelike statues, and bubbling fountains filled with inky midnight waters surround her home. String music drifts on the perpetual gloom, beckoning travelers to this oasis of lantern-dappled beauty.

\subparagraph{Lady of the Blue Barbers}
The gnome courtiers of Midnight Teeth who attend the mistress in her home also serve as her agents abroad. They call themselves the Blue Barbers after their universally blue hair and their desire to bring style to all they meet. They serve as valets, barbers, art dealers, decorators, and stylists of all kinds. The Blue Barbers are polite, gregarious, and knowledgeable on a variety of topics, which they eagerly share while plying their trade.

Their leader and guildmaster is Garrolo Goldenlock, and he and his associates are various stripes of bards, bandit captains, peddlers, and spy masters. Their work keeps the gnomes well-informed and their foes nervous. Sabotage, blackmail, and assassination are all part of their extremely broad and extremely ruthless portfolio.

\subparagraph{Infiltration of Mirrors}
While Countess Phylomara is a great user of mirror portals, scrying, and divination, one of her foes, the Duchess Shelesorra, the Demon of Lies and Cruelty of a minor hell, considers the mirror realms her own domain and sees the gnome countess's use of mirrors as an affront to her claim. Phylomara pays close attention to warding her mirrors and even employs a stubborn, keen-eyed erina named Ingo to observe, polish, and report on any unusual sightings in her mirrors. All the same, demonic visitations and even attacks are not unknown, all of them originating from Sooleleed, the Mirror Hell of Lies (see Warlock 16: Eleven Hells for more details).

\subsection{Halls of the Blue Barbers}
The halls of the Blue Barbers comprise a tower of white marble and a maze of manor-like wings, tunnels, and two-story banquet halls thrown together and remixed in a style that can charitably be called ad hoc. There's so little logic to it that guides are often required to reach even the most common destinations: the guildmaster's audience chamber, the purser's office, the scriptorium, and the mess hall. Adding to the confusion are the dozens of Blue Barbers who flit about with seeming unerring sense of where they are and where they are headed.

In truth, the whole place is meant to befuddle and confound visitors, leaving them somewhat exhausted before they arrive to speak to any guild member of importance. Dead-end staircases, doors leading nowhere (but which a Blue Barber may have just exited from), secret passages, rotating fireplaces, or chimneys containing hidden stairs are all elements of the halls.

The most likely destinations for visitors are often locked and marked with a sign reading "Closed for Cleaning" or simply "Closed." Breaking into a closed chamber is grounds for expulsion or even imprisonment, though in many cases only the primary entrance is closed. A secondary, hidden, or servant's entrance may still be perfectly serviceable.

One of the stranger elements of the halls is the shrine to the Demon Lord of Night, Alquam, which can be seen primarily as an azure dome from outside the building. Choral praise to the darkness is rendered there by devotees who call themselves the Alabar Chorus and who include a half-dozen rather famous gnomish bards and jesters. It's unclear just how sincere this group is, for they claim it is little more than a friendly gathering of wine, song, and macabre storytelling.

\subsection{Tower of the Keening Winds}
A prison for rebels, nobles, traitors, and unfaithful royal spouses, the Tower of the Keening Winds is made of gray-black stone and heavily warded to reveal invisible or ethereal infiltrators (who appear as glittering outlines in gold or silver while on the premises). The guards include a dozen mages and knights, led by the Commander of the Keening Tower, Kubberolus Gingercrick, and supported by two dozen trusted and carefully trained gnomish servants.

The tower's prisoners include both the serious and the ridiculous, including the ghost of Keening Kangovar (LE \textbf{ghost}), a swindler of the treasury; His Majestic Radiance Chancellor Umberwick (LN \textbf{shadow fey ambassador}, see \textit{Creature Codex}), a royal pretender to the throne of the Golden Oaks; the heretical Seer of the Twin Light, called Starfire (NE shadow fey \textbf{priest} of Sarastra); the warlord Neeka Nook (NE shadow goblin \textbf{bandit lord}, see \textit{Tome of Beasts 1}); and various arcane frauds, charlatans, and dispensers of bright, shining lies, including Yssandros the Clever (CE human \textbf{mage}), Arbecca Tollshi (LE shadow goblin \textbf{moonshadow catcher}, see \textit{Chapter 9}), and the arboreal, blood magic wielder Redfur (NE \textbf{bearfolk druid}, see \textit{Chapter 9}).

The tower's most famous royal prisoner is surely Mistress of Pearls, Baroness of the Purple Pavilion, Marina Nikla (LE gnome \textbf{thief lord}, see \textit{Creature Codex}) a longtime companion and courtesan of Countess Phylomara, who abandoned the queen for a lover among the gnomes of the \textit{Court of Pale Roses} for a time. So enraged and heartbroken was the countess that, when Nikla returned for a diplomatic visit, she seized her and locked her up. Relations with the Court of Pale Roses have been icy ever since.

\section{Court of Night and Magic}
\begin{DndSidebar}[float=!b]{Court of Night and Magic}
\textbf{Ruler}: Sarastra Aestruum, the \textbf{Queen of Night and Magic} (see \textit{Tome of Beasts 1})



\textbf{Important Personages}: Akyishigal, Demon Lord of Cockroaches (see \textit{Tome of Beasts 1}); Analissa, Verousha, and Lyssalyn, the Ladies of Nightbrook (NE \textbf{night hag} priests of Hecate); \textbf{Hander Svenk, the Black Prince} (see \textit{Courts of the Shadow Fey}); Kaedrin Blackwing, \textbf{Commander of Horn and Gold} (see \textit{Tome of Beasts 2}); the \textbf{Lord of the Hunt} (see \textit{Tome of Beasts 1}); Revich, the Blind Seer (see \textit{Courts of the Shadow Fey}); Yleira the Swift, Baroness of the Sable Court (LE \textbf{shadow fey forest hunter}, see \textit{Tome of Beasts 1})



\textbf{Population}: 660,000 (100,000 shadow fey, 300,000 shadow goblins, and 260,000 various fey, supplicants, and traders)



\textbf{Major Enclaves}: Corremel, the City of Lanterns, population 50,000 (see \textit{Corremel, the City of Lanterns} above); Hunt's Retreat, population 7,700 (4,000 shadow fey, 3,000 shadow goblins, 500 wyrd gnomes, and 200 umbral humans); Dalliance, population 5,300 (see \textit{Dalliance} below); Wormwood, population 3,500 (1,800 shadow fey, 900 shadow goblins, and 800 wyrd gnomes); Tower of Horn and Gold, population 600 (see \textit{Tower of Horn and Gold} below); Court of Night and Magic, population 300 (150 shadow fey, 100 shadow goblins, and 50 supplicants to the court);



Nightbrook Court, population 100 (see \textit{Nightbrook} below); Sable Court, population 100 (shadow fey)



\textbf{Great Gods}: \textit{Sarastra} (patron), \textit{Anu-Akma}, \textit{Charun}, \textit{the Hunter}, and \textit{Laughing Loki}



\textbf{Trade Goods}: Secrets, fey wine, metalwork, jewelry, memories, illumination magic, and shadow magic (in order of importance)





\end{DndSidebar}

Ruled by Sarastra, the Queen of Night and Magic, this is among the largest and best organized of the shadow fey courts with a central position astride several shadow roads and with towers and wardings to control those roads. The Court of Night and Magic also has a long association with mortals as well as with a number of powerful demons and devils, including connections with Akyishigal, Lord of Cockroaches, and his roachlings and with Alquam, Lord of Night.

\subparagraph{Highborn Rulers}
The true rulers of the shadow fey reside in the upper portions of the courts, in the Winter Palace and the Royal Halls. The Winter Palace functions as the traditional seat of the Moonlit King, though he only walks its halls when the strange seasons of fey rule turn and his Winter Court ascends—a turn of events that has not been seen in centuries. It has been many years since the seasons last turned, and in all that time, the Moonlit King was banished to the faraway Tower of the Moon. Rumors run rampant now that the Moonlit King has left the Tower, and whispers circulate through the halls that he's taken a demon lover or learned secrets of radiant magic from Hecate herself. His new cloud court drives his wife, the Queen of Night and Magic, to distraction.

Whatever the truth of these whispers, the courts agree that the queen now seeks to replace her wayward king. The court functions in his absence with consorts and courtesans, servants, and the usual petitioners hanging about the entrance to the Royal Halls. The Gray Ladies dwell here, three ancient shadow fey crones with an affinity for spiders and weaving. They spend their days spinning silk and whispering to one another, searching for individuals they can use to bend fate to their whims.

Ultimate power over shadow fey society rests with Her Transcendent Majesty Sarastra Aestruum, the Queen of Night and Magic. She dwells in the Royal Halls along with the upper crust of shadow fey society. Only the most respected supplicants receive invitations to enter the Royal Halls for audience, and of these, most receive only a sliver of the queen's attention. They make their obeisance, ask a favor, and quickly get handed off to court functionaries.

A lucky few, however, manage to catch the queen's interest. No matter the apparent circumstance, this always furthers the queen's designs or appeases a momentary whim. Her Majesty's direct attention is never taken lightly, and sudden alliances or rivalries spontaneously spring up around visitors so "blessed." The Royal Halls also host honored guests, usually envoys from other fey courts.

The courts also contain a strange pit called the Black Well of Night. A cavern at the bottom of this well houses Akyishigal, the Demon Lord of Roaches, and his roachling servants, creatures of filth and evil. The demon lord chafes under his current servitude to the queen and seeks potential allies to turn the tide against his "hostess" and win his freedom.

\begin{DndSidebar}[float=!b]{The Nature of Shadow / from the Book of Ebon Tides}
My various arcanists, priests, and jesters have long debated what makes the Shadow Realm malleable and unique. My own thought is that it is both an echo of the worlds of mortals and a place whose people understand that echoes have substance—and light and radiance are every bit as powerful as elemental earth, air, fire, and water. My realm is one made of elemental light and elemental darkness, and this makes it entirely different than realms of crass clay or raging fire. Our truths and our lands can be bent to serve my vision for the elves.



\end{DndSidebar}

\subsection{Hunt's Retreat}
Nestled at the midpoint of a hunting path between the road bisecting the forest and the cliffs of the courts of the shadow fey stands a stone lodge. This is Hunt's Retreat, a haven for the hunters and trappers who work the Blackwood and the Forest of the Firebirds. It serves as shelter to the weary traveler or lost soul wandering the woods, but it carries dangers of its own. The Black Prince lodges here during his hunting excursions, and the powerful fey lord of the Wild Hunt, the Lord of the Hunt, camps with his retinue just outside the retreat before sounding the horn and charging into the forest.

\subsection{Sable Court}
The Sable Court is a small forest court, nestled in tree branches high above the dark Sable Forest floor, built into and around the trunks of great shadow oaks. While it resembles a high elven forest dwelling, the Sable Court reminds the local inhabitants (particularly the Ladies of Nightbrook Court) of their place as vassals of the queen. Appointed by the Queen of Night and Magic, Baroness Yleira the Swift oversees the shadow fey and the hounds of this court. The Sable Court can be reached only by flight, aerial magic, a dangerous climb in the sight-line—and bowshot—of the forest sentries, or the series of lifts manned at all times by shadow fey guardians.

The Sable Court serves as the base of operations and training ground for the shadow fey forest hunters. Those shadow fey unfit for life in the courts or in the more civilized settlements often seek to join the hunters. Deadly assassins, forest guides, and iron-tough trailblazers all rise from the ranks of the forest hunters.

\subsection{Wormwood}
In the northernmost quarter of the shadow fey domains stands the gate town of Wormwood. It guards the road south into the dark woods that leads to the courts of the shadow fey and eventually to Corremel, keeping away the haunts and ghouls of the plains. A formidable wall with stout shadow oak gates at the north and south encircles the small settlement. The inhabitants of Wormwood are all battle-tested soldiers with a few spies and assassins thrown into the mix.

Wormwood also serves as a bulwark against the umbral vampires and their minions from Oshragora, the City Fallen into Shadow (see \textit{Chapter 7}). They regard travelers with scrutiny and heavy suspicion and take a pitch-black view of anyone carrying relics from the umbral vampire city without express leave of the Queen of Night and Magic or another high-ranking shadow fey noble.

\section{Court of One Million Stars}
\begin{DndSidebar}[float=!b]{Court of One Million Stars}
\textbf{Ruler}: Prince Valendan (N male \textbf{selang}, see \textit{Tome of Beasts}, \textbf{archmage})



\textbf{Important Personages}: Cylentha the Silver Mistress (LG \textbf{deva}); First Consort Shemani (CE \textbf{living star}, see \textit{Creature Codex}); Phaerliggath, the dragon of Void Reach (NE \textbf{adult void dragon}, see \textit{Tome of Beasts 1}); Radiant Chancellor Ectius, Ambassador to the Stars (N lunar elf \textbf{radiant lord}, see \textit{Chapter 9}); Second Consort Tremanca (CG \textbf{shadow fey enchantress}, see \textit{Tome of Beasts 1}); Verskas Bright Eye (NG \textbf{sprite})



\textbf{Population}: 1,850 (900 shadow fey, 300 courtfolk halflings, 250 witchlights, 200 celestials, 100 shadow goblins, 70 umbral humans, and 30 starry courtiers, or living stars)



\textbf{Great Gods}: \textit{Gytellisor} (patron), \textit{Black Goat of the Woods}, Lada, \textit{Sarastra}, and \textit{Umbeserno}



\textbf{Trade Goods}: Star metals, liquid fire, stardust, void essences, memory philters, celestial silks and satins, and stellar radiances (in order of importance)





\end{DndSidebar}

Stars and moon rarely shine directly over the Shadow Realm. Though often hidden by its dreary sky, stars still occasionally gleam in all their glory, and one in particular shines bright. Above the cloud cover rests the most enigmatic of dark fey gatherings: the Court of One Million Stars. Delicate spires of silver and pearl spiral up from great discs spun of dreamstuff and starlight. Clusters of these floating platforms and towers gather, bound by gracefully arching bridges forged from moonbeams woven with silver.

Prince Valendan, a noble selang wizard, rules over the Starry Court with an aloof manner and quiet wisdom. The prince makes it a point to meet every visitor to the court at least once, though few receive a second audience. His strange appearance repels some, but he is far less bound to chaos or evil than most selang.

\subsection{The Starry Court}
The Starry Court serves as the heart of the Court of One Million Stars, a magnificent spire of pearl and silver, spiraling like an unicorn's horn from the court's center and audience chamber. This seven-sided structure of white stone features seven gates, one on each of the building's sides, that lead into a great hall. Courtiers cluster around the wings of the chamber with those of increasing station closer to the center. When present, Prince Valendan sits in the precise center, resting upon a tripod seat of silver and silk.

The spire rising toward the stars houses the prince's personal chambers as well as those of his consorts and the other royal figures of the Starry Court. Balconies dot the outer surface of the spire, allowing a breathtaking view of both the stars above and of the surface of the Shadow Realm below.

A series of magnificent gardens surround the court, home to a wide array of plant life from the Shadow Realm as well from the mortal world. Walkways of crushed pearl trace paths through the central gardens and the seven outer gardens. Floating lanterns shining with pastel light illuminate the walkways, and each outer garden anchors a moonlight bridge, leading to another section of the courts. Notably, one of the bridges is blackened and shattered, ending a few yards from where it begins, destroyed long ago by devastating magic.

Fey from all walks of life visit the Starry Court to study illumination magic, the stars, and celestial movements. Dark pixie scholars can spend decades tracking the movement of a single star, meticulously recording every hint of prophecy in anticipation of a momentous event—the starfall.

\subsection{When the Stars Walk the World}
The stars are not just points of light burning in an uncaring sky. They are living things of power and unfathomable knowledge. When momentous events shake the world, a star might deign to leave its bower of black velvet and fall to the mortal world or to the Shadow Realm. When it lands, it becomes a creature of luminous beauty and impenetrable mystery.

Once in physical form, these starry emissaries appear cloaked and hooded, shrouded in a luminous aura of starlight. They tend to be aloof, and on the rare occasions when they exchange words with a supplicant, understanding isn't guaranteed. The emissaries speak in a strange manner, often reusing terms, but with vastly different contexts and meanings from one part of a statement to the next. Magic such as \textit{tongues} and \textit{comprehend languages} can help but never provides a perfectly clear meaning.

Even more unsettling, the courtiers believe that when such a sojourn ends, the star never leaves unaccompanied. Who or what vanishes with the star varies with each visit and isn't always apparent. Such starfalls are rare, but at least one emissary from the stars remains as a permanent fixture within the Court of One Million Stars: First Consort Shemani.

\subsection{The Heavens Seek Answers}
The heavens have taken a special interest in the Court of One Million Stars. Cylentha the Silver Mistress descended to the Shadow Realm nearly 200 years ago and has yet to return home. Many believe that she is searching for something. She advises Prince Valendan who keeps her counsel quite private. Certain courtiers would offer a treasure trove of stardust diamonds and prophecies of future glory to understand why the celestial spends so much time at a court filled with fey decadence and void followers.

When Cylentha isn't offering advice to the prince, she spends most of her time in her observatory. This unadorned platform holds no gardens nor decorations of any kind. Seven blue-white balls of light, like miniature stars, surround the observatory, which appears as a round building surmounted by a dome of iridescent pearl. Inside, the dome allows a perfect view of any star, every far-flung celestial cloud or nebula, and the farthest reaches of the utterdark, perhaps even offering a glimpse into the Void.

\subsection{Whispers from the Void}
The Starry Court reaches toward the stars, but that necessitates drawing close to the cold, vast emptiness between the heavenly lights. In that endless nothing, ancient things stir and whisper their secrets. A careless or foolish investigator can tap into those whispers and glean their secrets, but such knowledge of course has its price. The Void seeks to spread its influence just as the Shadow Realm has, but it has a more difficult time finding purchase in the world. Despite the Void's destructive influence, Prince Valendan hasn't forbidden study of its magic yet, though he takes an interest in anyone who seeks it.

One of the more recent and momentous visitors to the Court of One Million Stars, the void dragon Phaerliggath, is from the cold dark beyond the stars. This hoary old wyrm has long drifted within the utterdark, and like all of his ilk, he has seen far too much of it. The dragon's mind has very much frayed. Phaerliggath claimed an entire spire of the Starry Court when he arrived and rarely emerges from his commandeered lair. The spire, once magnificent, now crumbles and falls into decay and is called Voidreach. Phaerliggath entertains visitors from time to time, but only those guests truly captivated by the Void ever visit a second time. Phaerliggath seems disturbingly willing to share the forbidden secrets of the Void with those brave or unhinged enough to seek him out.

Lady Cylentha makes occasional visits to the dragon's tower, and investigates anyone who speaks with the dragon. She is polite during such inquiries, but won't take no as an answer to her questioning.

\subsection{The Spiral-Downs}
While enclaves of the Court of One Million Stars dot the skies above the Shadow Realm, the largest concentration of starry courtiers live in a group of spires known as the Spiral-Downs. Connected to the Starry Court by a moon bridge, the Spiral-Downs are connected platforms that hang in the sky, joined by moonlight bridge staircases, forming a spiral of ascending status and power. The lowest and largest platform houses the largest residential spires where courtiers of the lowest station reside. Each ascending platform is slightly smaller, more ornate, and houses fewer residents than the platform below it. The loftiest platform boasts verdant gardens, ethereal music, delicate spires spun of starlight and dreams, and the elite of the Starry Court.

\subsection{Rising Above}
Visiting the Court of One Million Stars is no easy task. Unseen paths lead to it from the surface, paths nearly impossible to decipher without an invitation or a guide. Occasionally a witchlight chooses to lead a pilgrim from the Shadow Realm up to the stars, leaving a trail of glittering light that somehow remains solid during the climb.

Shadow fey wizards speak in hushed tones of star bridges that appear only in certain locations of the Shadow Realm and only during certain sidereal alignments. Rare glimpses of the stars through the realm's gloomy sky can offer such an illuminated mind enough insight to calculate where and when such a bridge will open next.

Each path to the Court of One Million Stars ends at a gatehouse where warriors of the court guard against hostile incursions from the surface. The most formidable of these gatehouses, Vigil, stands unblinking over the only path leading from the surface to the Starry Court. Vigil appears as a hovering disc of tightly fitted tiles. A gate of spun silver, seemingly delicate but nearly indestructible, bars both the entrance from the surface and the bridge to the Starry Court. Twisted spires like horns defend the disc's edge, and a bronze dome in the center provides shelter for the wardens. The gate wardens, the fiercest warriors among the court, closely question any newcomers.

\section{Court of Pale Roses}
\begin{DndSidebar}[float=!b]{Court of Pale Roses}
\textbf{Ruler}: King Frulio Cornutus, Duke of Roses, Count of Spiders, Lord of Illusions, Lord Mayor of Lurgestown, Baron of the Void Immaculate (LN shadow fey \textbf{first servant} of Hecate, see \textit{Tome of Beasts 3})



\textbf{Important Personages}: Alchemist-in-chief Olmed Carroso (N human \textbf{mage}); Chancellor of the Royal College Alchemical Igo Natterbutt (NG \textbf{gnomish distiller}, see \textit{Chapter 9}); General Flavestor Murena (LE wyrd gnome \textbf{field commander}, see \textit{Tome of Beasts 3}); High Priest Vares the Silent (CN shadow fey \textbf{apostle} of Laughing Loki, see \textit{Tome of Beasts 3}); Keeper of Roses, Aiyesh bin Teleros (N human \textbf{druid}); Lady Vindesol (CG high elf \textbf{mage}); Mistress of the Cloisters Ellec Ish-Vardes (NE satarre \textbf{first servant} of Vardesain, see \textit{Tome of Beasts 3})



\textbf{Population}: 2,450 (1,300 wyrd gnomes, 600 umbral humans, 300 shadow fey, 100 quicksteps, 50 tieflings, 50 bearfolk, and 50 satarre)



\textbf{Great Gods}: \textit{Hecate} (patron), \textit{Ailuros}, \textit{Baccho}, \textit{Charun}, and \textit{Laughing Loki}



\textbf{Trade Goods}: Flowers, weavings, ale, glasswork, and ceramics (in order of importance)





\end{DndSidebar}

Home to an enormous community of wyrd gnomes, the Court of Pale Roses is young by fey standards, steeped in magic and open to a degree of civil relations between shadow fey, wyrd gnomes, and umbral humans. With a tolerant ruler and open doors, the place is home to many, with occasional visits by tiefling arcanists, quickstep duelists, and even bearfolk druids who share the secrets of their groves with this fey court but no other. Its ruler is King Frulio Cornutus, a young and vigorous shadow fey whose obsession with arcane workings and enchantments brings him many followers—and those who pretend to some arcane knowledge but whose true purpose is spying for weakness at the court.

\subparagraph{Gardens of Illusion}
Illusion seeds (see \textit{Chapter 10}) grow in great profusion in the court's namesake gardens, watered by chill, dark buckets of half-frozen water draw up from the Somering Well, a planar spring with great powers of fertility and growth. This connecti90on runs through a root of Yggdrasil somewhere underneath the court itself, and only the shadow fey and a few trusted gnomes know where to find it. The gardens themselves are warded by magic and watched by a group of shadow fey guards sworn to the king, called the Rose Guard.

\subparagraph{The Young King's Fancies}
The young King Frulio is wooing the puckish Lady Vindesol of the Linden Tree, an elf maid from the Summer Lands with a love of birds and fine embroidery. He is also a frequent visitor to the halls of the cat-goddess Ailuros, and some say he is Ailuros's favored king among the fey. He has commissioned prayers, scents, and holidays for the cat goddess.

King Frulio enjoys putting on great festivals combined with shows of strength, arranged by General Murena, a wyrd gnome with an eye for a keen warrior. The king lavishes praise on his alchemists and gardeners alike, who both add to the court's growing reputation, and he has founded the Royal College Alchemical, a place for the study of natural philosophy with an emphasis on botany, enchantments, and philters and potions.

\subparagraph{Cultish Fervor}
From time to time, the Court of Pale Roses seems overtaken with some new fashion: clothes, a particular cantrip, a fashionable memory philter (see \textit{Chapter 10}). At the moment, the cult of the Void, as proclaimed by the satarre priestess Ellec Ish-Vardes, is in vogue, and it is considered elegant to at least pretend to know what these doom-ridden creatures rattle on about. Most sable elves seem to find it amusing more than revelatory, but the satarre has brought some tiefling underpriests to aid in ministering to the shadow fey and wyrd gnomes. Her goal is to establish a direct link to the cold waters of the lower branches of Yggdrasil or the River Lethe and in time bring the king under her sway (and the court with him). She believes that corrupting enough of the fey courts will lead to an enormous Great Tide, which will sweep away all settlements and indeed all land in the realm, leaving formless chaos and pure Void.

\section{Court of Sparrows}
\begin{DndSidebar}[float=!b]{Court of Sparrows}
\textbf{Ruler}: The Undefeated King Marrow of the Court of Sparrows, Bondsman of King Reynard, and Commander of the Greenwood Company (NE trollkin \textbf{field commander}, see \textit{Tome of Beasts 3})



\textbf{Important Personages}: Commander of Foot Sventonna the Quick (LN bearfolk \textbf{gladiator}); Lady Castellan Grinholda (CN \textbf{trollkin shaman}, see \textit{Creature Codex}); Master of Protocol Hithbert Humbervale (LN courtfolk halfling \textbf{merchant captain}, see \textit{Tome of Beasts 3}); Master of Wands Ardinaldo the Clever (CE human \textbf{battle mage}, see \textit{Creature Codex}); Royal Keeper of Ducks Mardovac the Foul (CG trollkin \textbf{bandit lord}, see \textit{Tome of Beasts}); Soothsayer to His Majesty and Court Healer Ploofi Flatfoot (CE shadow goblin \textbf{apostle} of Santerr Marossa, see \textit{Tome of Beasts 3})



\textbf{Population}: 800 (300 courtfolk halflings, 200 bearfolk, 130 trollkin, 100 shadow goblins, and 70 umbral humans)



\textbf{Great Gods}: \textit{Laughing Loki} (patron), \textit{Bengta}, \textit{the Hunter}, \textit{Santerr Marossa}, and \textit{Umbeserno}



\textbf{Trade Goods}: Duck eggs, duck meat, chickens, goose feathers, fletching, fine yew wood, small woodwork, and salt meat (in order of importance)





\end{DndSidebar}

A minor court owing fealty to the much more extensive \textit{Court of Mice and Ravens}, the Court of Sparrows is a tan stone stronghold that has taken a recent turn. It began as a whimsical or at least polite court, founded by courtfolk halflings breeding sparrows, ducks, geese, and chickens for royal banquets—a place of jesters, farts, and jugglers—and visitors came for the cuisine and entertainment. The castellan was Traverty Peaseblossom, a courtfolk halfling bard widely loved by both his people and his poultry.

\subparagraph{Change of Rulers}
A few years ago though, the court was stormed and overtaken by King Marrow, a trollkin warlord with a company of intimidating and ruthless followers, including trollkin, bearfolk, and a few umbral humans. They've made it their home base and retained the halflings, and their arrows are now fletched with extremely fine goose feathers. Prince Peaseblossom fled via a secret passageway somewhere in the castle, and he may now be a wandering minstrel at another court entirely.

The place's change of ownership was not marked (as would be traditional) by a change of name to, say, the Court of Trolls or Court of Hammers. Being not quite bright enough to rename the place has worked out well for King Marrow though, and so long as the eggs and ducks and other foodstuffs flow to the nearby courts of the major fey, no one seems to mind who is in charge.

King Marrow also raids the shadow roads and hunts in the Queen's Forest, all the while eating a great many eggs. The halflings now in the role of servants sometimes creep away to seek service elsewhere, but for many of them, having a great big trollkin warlord "in charge" saves them from a lot of tedious diplomacy and bootlicking of the shadow fey snobs, sable elf visitors, and greedy umbral human merchants.

\subparagraph{Duck Eggs and Diplomacy}
The lady castellan, the master of wands, and the commander of foot are Marrow's captains of trollkin soldiers, trollkin wands, and bearfolk heavy infantry, respectively. However, a considerable degree of power is shifting to the Royal Keeper of Ducks, Mardovac the Foul, who happens to be Marrow's son and heir. Mardovac is charged with wringing every ducat out of the courtfolk halflings and their flocks of ducks and geese, and he seems to have a greater knack for that than thievery and ambuscade. So while the elder captains thrive on banditry, Mardovac keeps amassing coins and power among the trollkin and halflings. Meanwhile, Hithbert, the master of protocol, sends useful diplomatic notes to King Reynard.

The fey lord of the Court of Mice and Ravens finds the entire situation vastly amusing, though if a shipment of fresh duck ever arrives at his halls late, he will surely come to pay a visit to King Marrow in person, as "one ruler to another," and perhaps install his own, more trustworthy vassal. Or he might just stay a while—until the eggs run out entirely.

\section{Court of the Moonlit King}
\begin{DndSidebar}[float=!b]{Court of the Moonlit King}
\textbf{Ruler}: Ludomir Imbrium the XVI, the \textbf{Moonlit King} (see \textit{Tome of Beasts 1})



\textbf{Important Personages}: His Radiance, Popplan Lusteron, the Minotaur (LN minotaur \textbf{archmage})



\textbf{Population}: 2,200 (1,000 shadow fey, 400 minotaurs, 300 bearfolk, 300 erina, 100 lunar elves, 50 quicksteps, and 50 wyrd gnomes)



\textbf{Great Gods}: \textit{Hecate} (patron), \textit{Ailuros}, \textit{Bengta}, \textit{Charun}, and Valeros



\textbf{Trade Goods}: Silver jewelry, moonstones, sweet moonlit wine, lunar magic, and illumination magic (in order of importance)





\end{DndSidebar}

This court was once part of a dual crown with Ludomir as its king and Sarastra his wife and co-ruler. However, with the Moonlit King's descent into strange magic and times of pure fantasy unlinked to the facts of the Shadow Realm, most of his followers drifted away—and the crown with them. All that remains of the Court of the Moonlit King is a set of names, spells, philters, titles, a handful of servitors, including lunar devils (see \textit{Tome of Beasts 1}) and \textbf{quicksteps} (see \textit{Creature Codex}). His old lone tower and labyrinth still stand, and various small groups (or shards) of followers and friends remain loyal, including hundreds of followers among the bearfolk, lunar elves, minotaurs, and erina who seek to restore the king's title and cast down Sarastra, so the Moonlit King may shine once more. Currently, this means they are building a new and highly magical court in the skies, the first major new court in some time.

\subparagraph{The Labyrinth of Moonlight}
The Court of the Moonlit King was once found in the King's Tower at the center of a frequently changing maze, which still exists and is inhabited by eye golems, lantern dragonettes, witchlights, shadow goblins, and more. However, because this maze is surrounded by the territory of the Queen of Night and Magic, it has been largely abandoned as a seat of rulership. The Moonlit King maintains an embassy there to landbound fey but has proclaimed his new domain is "the sky and the moon and all the territory they touch."

\subparagraph{A Wandering Court}
The Court of the Moonlit King is currently nomadic, wandering over vast distances and often keeping its people together by the most tenuous of messengers, code phrases, and meeting points. The lesser lunar devils (see \textit{Chapter 9}) sometimes carry small fey, quicksteps, and others from these cloud fortresses to the earth and forests, but just as often, they use luminous staircases at radiant wells to ascend or descend, or they ride griffons or use flying spells and Cloak of the Raven (see \textit{Chapter 10}) to travel from earth to sky. The Court of the Moonlit King is one of high magic and mystery, and clearly the restored King is gathering power and increasing his number of followers month by month.

\subparagraph{Alliance with the Minotaurs}
With the Court of the Moonlit King now mobile and active, striking out for new sources of power and bending radiant wells into luminous cloud palaces, it has become easier for his ambassadors and loyal followers to gather allies. The strongest of these are the minotaurs of the mortal world, who revere the moon and her champions. The cloud-palace of the Moonlit King is a place best known as a rumor more than as a place that many have visited. The king's newfound power from radiant wells intrigues all shadow mages and those who study illumination magic, but its details are scarce.

\section{Court of Winter's Love}
\begin{DndSidebar}[float=!b]{Court of Winter's Love}
\textbf{Ruler}: Queen Kalaslurr (LE \textbf{ghost}), Lady of Ash and Sorrow, Keening Mother of the North, Mistress of Snows, Keeper of the Winter Gate, Friend to the Seven Hands of Ice



\textbf{Important Personages}: Chief Duelist and Champion to the Queen, Illebor Circon (LE \textbf{quickstep}, see \textit{Creature Codex}); High Priestess Silloon Honeyeater (CN shadow fey \textbf{priest} of Bengta); Keeper of the Scrolls Tikkaw Tokkaw (N \textbf{ravenfolk doom croaker}, see \textit{Tome of Beasts 1}); Lord of Illusions and Keeper of Song Umberell Tumwater (NG bearfolk \textbf{merchant captain}, see \textit{Tome of Beasts 3}); Master of the Hunt Arros Freor (NE human \textbf{keeper of hounds}, see \textit{Chapter 9}); Most Devout Steward and Head Chamberlain Porthan Villeros (LE shade \textbf{apostle} of Charun, see \textit{Tome of Beasts 3})



\textbf{Population}: 1,350 (400 shadow fey, 350 quicksteps, 200 wyrd gnomes, 200 umbral humans, 50 snow maidens, 50 bearfolk, 50 ice bogies, and 50 shades)



\textbf{Great Gods}: \textit{Hecate} (patron), \textit{Ailuros}, \textit{Alquam}, \textit{Baccho}, and \textit{Charun}



\textbf{Trade Goods}: Ice, sunstones, fine furs, and woolen yarn (in order of importance)





\end{DndSidebar}

The coldest and most desolate of the fey courts is the Court of Winter's Love, ruled by Kalaslurr, the Lady of Ash and Sorrow, a ghost whose touch and voice destroy most mortals who meet her. The court itself is home to ice bogies and snow maidens, various ice demons, and a thriving community of quicksteps, who all seem immune to the fey lady's song and terrifying presence. Most of the court's inhabitants stay far away from its ruler—other than the devoted shades of Kalaslurr, her ghostly court of undead courtiers, all of whom are immune to her deadly presence. The head chamberlain, Porthan, is also the chief priest of Charun and one of the few voices that the queen considers when deciding where to send her demons next.

\subparagraph{The Singing Tower}
The tall white tower of the court is unmistakable, and at most hours of the day, a slow, low howling can be heard from it. These are the wails and moans of Queen Kalaslurr, whose voice draws in and slowly devours any visitors she finds entertaining. Rumors claim that these are sometimes old lovers and at other times foes, rivals, or simply those whose appearance she finds pleasing. She is easy to gain audience with, but her attention is difficult to survive as the distant moan is a song of terror when heard within the tower.

The queen's servants among the wyrd gnomes are often made deaf by her song, or they (so it is said) destroy their own eardrums to better serve her. The only servant who has no fear of the queen's dire singing is the bearfolk bard Umberell, for he has long since learned the knack of a counter song. He will teach this to any bard or talented singer who seems trustworthy and honest.

\subparagraph{The Demon's Hall}
A black stone binds seven ice demons who serve Queen Kalaslurr without hesitation, tearing apart foes, fetching gentian flowers from alpine meadows, or threatening travelers for tolls or information. They are her hands, bringing her interesting news, a talented cook, or a worthy morsel of mountain goat. The demons are named First Hand, Second Hand, All Thumbs, Wicked Frost, Silver Eye, Little Cook, and Wanderer.

\subparagraph{Frost Bogie Workshop}
The smallest and strangest of the inhabitants of the court are the \textbf{ice bogies} (see \textit{Tome of Beasts 2}), small snow fey who obey any order given to them by anyone but who positively worship the ice demons and the keeper of scrolls and keeper of song, as well as the queen, the chamberlain, and the queen's champion Illebor. When not fetching ale or collecting pine nuts from the forest, they often carve small wooden trinkets, card and spindle woolen yarn, or clean and tan furs of fox, rabbit, ermine, and beaver.

\section{Dalliance}
\begin{DndSidebar}[float=!b]{Dalliance}
\textbf{Ruler}: Lady of Whims and Fancies, Her Divine Debauchery, Baroness Lupina Valdross (LE \textbf{shadow fey enchantress}, see \textit{Tome of Beasts 1}), Matron of the Red Door



\textbf{Important Personages}: Lady of Small Mercies, Viscountess Ellana Hibiscalla (CG human \textbf{spy}); Master of Coin Tristoffos Illemartos (LE \textbf{shadow fey knight of the road}, see \textit{Creature Codex}); Master of Shades Lord Correto Anders (LE human \textbf{mage}); Matron of the Cat's Paw, Meli Hawthorne (N wyrd gnome \textbf{merchant captain}, see \textit{Tome of Beasts 3})



\textbf{Population}: 5,300 (1,900 shadow goblins, 1,500 shadow fey, 1,300 umbral humans, 400 wyrd gnomes, 180 quicksteps, and 20 shades)



\textbf{Great Gods}: \textit{Baccholon} (patron), \textit{Ailuros}, \textit{Charun}, \textit{Hecate}, and \textit{Santerr Illosi}



\textbf{Trade Goods}: Furs, timber, requiem, wine, and brandy (in order of importance)





\end{DndSidebar}

Where the Nightbrook spills into the River Lethe, Dalliance sits like a fat but happy leech. Its people make their living as traders for forest dwellers and trappers and, for the \textit{Court of Night and Magic} and the \textit{Court of Roses}, as woodsmen and occasional hunters. The water fey from Nightbrook Court sometimes visit with dark philters in hand, and it is sometimes host to a goblin troupe of actors or jugglers. However, the small town's real wealth comes from trading pleasure, games of chance, and debauched entertainments to anyone willing to pay.

Three pleasure houses—the River House, the Red Door, and the Cat's Paw—are fine, stone-walled manors that cater to anyone and everyone, from the most straightforward hunter seeking wine and friendly companions to the dissolute gnome princes who demand exotic entertainments. Any bit of intoxication and any form of company is available if one asks the right house and offers a fat coin purse. This is a well-known secret, but Dalliance is famous for keeping its secrets very well hidden. Whatever goes on behind the walls and doors does not leave town.

\subparagraph{Shades and Requiem}
Master of Shades, Lord Correto Anders, keeps a string of ghosts, haunts, and shades of various kinds working for him to bring requiem from Soriglass to Dalliance. The addictive, soul-bending substance is combined with other pleasures of the flesh in ways best left to the completely debauched customers of the small River House where it is offered.

\subparagraph{Blackmail and Extortion}
While most of Dalliance is honest in its whoring and pleasure seeking, Master of Coin Tristoffos is a spymaster and a blackmailer of umbral humans, shadow fey, and shadow goblins alike. He pays extremely well for gossip, for news, and for well-founded speculation. He is most often found questioning the staff at the Red Door, which leads some to believe his information is first and foremost for Baroness Lupina's ear. Very few who share information with Tristoffos are willing to share information about a client's personal secrets though, only about their political aims, their religious quirks, or their business ventures. (Which is certainly still enough to keep Tristoffos in business.)

\subparagraph{The River Spirits}
The waters near Dalliance are quite active with water fey, including river spirits (see \textit{Chapter 9}), \textbf{shadow river lords} (see \textit{Creature Codex}), and \textbf{rusalka} (see \textit{Tome of Beasts}). They are led by a \textbf{lorelei} (see \textit{Tome of Beasts 1}) in love with various fey at various times, who sings in the brightest moonlight and welcomes gifts from suitors as well as songs and poems from bards and passing travelers. In return for song and elegance, she provides safe passage along the Dalliance stretch of the river as far as \textit{Corremel}.

\section{Langmire's Court}
\begin{DndSidebar}[float=!b]{Langmire's Court}
\textbf{Ruler}: Lord of Misrule Oshlan the Chef (NE halfling \textbf{bandit lord}, see \textit{Tome of Beasts})



\textbf{Important Personages}: Keeper of Cards Phildra VII (LE tiefling \textbf{mage}); Keeper of the Dice Naldos Even-Handed (LN human \textbf{bandit captain}); Sister Fox (CN shadow fey \textbf{priest} of Loki)



\textbf{Population}: 3,800 (1,300 umbral humans, 800 shadow fey, 600 shadow goblins, 600 halflings, and 500 wyrd gnomes)



\textbf{Great Gods}: \textit{Laughing Loki} (patron), \textit{Ailuros}, \textit{Charun}, \textit{Hecate}, and \textit{Santerr Illosi}



\textbf{Trade Goods}: Carved bone dice, inscribed books and cards, inks, hides, gold ingots, herbs and potions, woolen and linen cloth, and leather goods (in order of importance)





\end{DndSidebar}

A minor court popular among young fey gentlemen, Langmire's Court is full of gamblers, secret lovers, and indolent wastrels, very loosely owing fealty to the \textit{Court of Winter's Love}. The court itself has dozens of small outbuildings, but three great halls are at its heart. Gnomes, halflings, and shadow goblins all serve as cooks, maids, and house servants in a constant bustle of servile attention—and they also take most of the profits at the end of the day.

\subparagraph{Long Hall}
While the whole sometimes resembles a bacchanalian festival at night, by day, Long Hall is a bright, smelly place with flaking black paint on the walls and a crew of hardworking halflings clearing away the mess and preparing some extravagant display of flowers, scenery for a musical play, or some similar entertainments. The gambling here begins at moonrise and continues into the morning with vast sums of coin, magic, and even land and entire mansions frequently wagered. The Lord of Misrule, Oshlan the Chef, prepares the entertainments, hires singers, brings in flocks of trained birds, and generally makes Long Hall a wonderland for the wealthy and foolish to spend their coin. Rumor has it that Oshlan owns the entire estate of several young shadow fey nobles and holds the keys to mansions in \textit{Corremel}, \textit{Merrymead}, and \textit{Fandeval} as well.

\subparagraph{Loki's Golden Hall}
Perhaps the largest shrine to Laughing Loki found anywhere, this somewhat garish wooden temple features gilded pillars, silver nets, and a statue of the god himself in a pose of mirth and good spirits. The statue changes its poses from time to time in a series of minor miracles (or the shaping work of the priests and jokers who maintain the shrine—some believe they have a series of five or more statues displayed at various times). A small number of priests here serve as entertainers and gossips as often as they minister to the spirits of those who have lost everything gambling with cards or dice, and their green and golden cloaks are a common sight in the Long Hall where they are often asked to bless dice or pray with those who hazard enormous sums in Loki's name.

\subparagraph{Servant's Hall}
This is home to the wyrd gnomes, halflings, and shadow goblins who truly keep Langmire running. Family life, weaving, cooking, and instruction for the young of the court are all part of Servant's Hall. Visitors are rarely brought in. The main doorway is guarded by two gargoyles instructed to stop anyone taller than a shadow goblin from entering.

\section{Nightbrook}
\begin{DndSidebar}[float=!b]{Nightbrook}
\textbf{Rulers}: Ladies of Nightbrook—Analissa, Verousha, and Lyssalyn (NE night hag priests of Hecate)



\textbf{Important Personages}: Master Shipwright Regunn Vorbeck (LN human \textbf{merchant captain}, see \textit{Tome of Beasts 3})



\textbf{Population}: 100 (50 shadow fey and 50 shadow goblins)



\textbf{Great Gods}: \textit{Charun} (patron), \textit{Alquam}, \textit{Black Goat of the Woods}, \textit{Hecate}, and \textit{Santerr Illosi}



\textbf{Trade Goods}: Boats and timber, herbs and potions, woolen and linen cloth, memory philters, and leather goods (in order of importance)





\end{DndSidebar}

A minor court owing fealty to the \textit{Court of Night and Magic}, Nightbrook Court occupies a hunched building of stone and black wood in the Sable Forest. The structure resembles a temple more than a seat of nobility, and strange glyphs adorn its walls and domed roof. It stands at the headwaters of the Nightbrook, a small black waterway that eventually joins the River Lethe. The shadow fey of Nightbrook Court seem strangely at home on the Shadow Realm's waters. They live and work along the Nightbrook, offering passage on small boats or rafts from one shore to the other or even downriver to another settlement—for a price, of course. The vengeful spirits that swarm in the black rivers seem indifferent to them, though none outside the court's confines know why.

Three sisters rule Nightbrook Court: shadow fey women with sour expressions and back-curving horns. Their wiry, long black hair bears a streak of white in their otherwise midnight tresses. Despite their perpetual glowers, the Ladies of Nightbrook—Analissa, Verousha, and Lyssalyn—welcome guests with surprisingly warm hospitality, particularly mortal guests. The sisters are, in truth, a coven of night hags. They work tirelessly to locate black-hearted people whose dreams they can haunt, hounding the hapless victims to death, so they can steal their evil souls. They bring these souls to the headwaters of the Nightbrook, and in a dark ritual that requires a memory philter (see \textit{Chapter 10}) holding emotions of loss, longing, rage, or bitterness, they twist the souls into hungry shades. Once infused with those hollow emotions, the ladies pour the souls out of their soul bags into the river to join the others. These souls draw the memories from those they drown and carry them back to the sisters.

Good-hearted visitors to Nightbrook Court have no fear of being murdered for their souls, but their memories remain at risk. The ladies require a constant supply of appropriate memories to create their spectral spies. If they can bargain or purchase such memories from mortal visitors, so much the better, but they aren't above helping circumstances along to create the memories they require. Perhaps a husband chokes to death on his sumptuous (and cursed or poisoned) dinner in the court's dining hall, and the sisters offer to soothe his widow's anguish by taking the memory from her.

\section{Tower of Horn and Gold}
\begin{DndSidebar}[float=!b]{Tower of Horn and Gold}
\textbf{Ruler}: Kaedrin Blackwing, \textbf{Commander of Horn and Gold} (see \textit{Tome of Beasts 2})



\textbf{Important Personages}: Nest Matron and Eldest Weaver Cuilliarna (\textbf{elder shadow drake}, see \textit{Tome of Beasts 1}), Sir Parnath the Golden, Captain of the Knights of Horn (LE shadow fey \textbf{wyvern knight}, see \textit{Creature Codex})



\textbf{Population}: 600 (400 shadow goblins, 150 shadow fey, 50 shadow drakes)



\textbf{Great Gods}: \textit{Sarastra} (patron), \textit{Anu-Akma}, \textit{Gytellisor}, \textit{the Hunter}, and \textit{Santerr Marossa}



\textbf{Trade Goods}: Shadow gold





\end{DndSidebar}

High in the peaks of the Mistcall Mountains west of the courts of the shadow fey, where gray snow falls like ash from a pyre, an ancient tower broods—the Tower of Horn and Gold. So high its apex scrapes the perpetual cloud cover, the tower's black stone is hewn from the mountains, though it takes its name from the cladding adorning it. Rows of great, curving horns jut from the tower's surface, giving it a wicked, serrated silhouette. Beneath these foreboding trophies, shadow-mined gold bands the stone. The metal, a wavering blend of gleaming yellow and pitch-black, gives the tower a striated appearance that both catches and absorbs any stray beam of light.

The tower represents one of the oldest remnants of the shadow fey's first incursion into darkness in their desperate bid to remain powerful. While it commands one of the best vantage points in all the Shadow Realm, the Tower of Horn and Gold is no mere watchtower. It protects a cleft in the mountains, leading to tunnels that have never seen a single mote of light. Shadow drakes have long lived and died in these caverns, and they maintain a loose, tenuous alliance with the fey occupying the tower. In exchange for the drakes' defense of the tower and occasional service as messengers or mounts to the high-ranking knights, the shadow fey have given the drakes leave to hunt in the surrounding lands. Most important, the shadow fey mine the tunnels for the elusive shadow gold veins that run through them. The drakes take a generous share of the yield.

Shadow gold is in many ways the same as the gold mined in mortal lands, but it holds the essence of shadow within it. Most of the goblins here mine and smelt the metal without pause or holiday, for the shadow fey and drakes require it (while the others maintain the tower and prepare the food, clothes, and tools for the goblin miners and fey knights). In its raw form, shadow gold appears as a mottled mix of gold and black. When the metal is cast or forged and then polished, its color and luster become indistinguishable from normal gold. Only when the metal is molten, nicked, bent, or scratched is its true nature revealed. Any damage to the metal shows the shadow within.

The shadow fey mint coins of shadow gold and trade it generously to mortals in payment or bribery. It is said that with each shadow gold coin one accepts from the fey, one loses a bit more of one's will and soul to them. Greed remains a powerful force, and shadow fey diplomats, emissaries, and spies liberally buy whatever and whoever they need with the tainted coins.

\section{Witch Queen's Court}
\begin{DndSidebar}[float=!b]{Witch Queen's Court}
\textbf{Ruler}: Nicnevin, \textbf{Queen of Witches} (see \textit{Tome of Beasts 1}) and Lady of the Dolorous Mountain, Mistress of the Moonsilver Ring, Keeper of the Cauldron, and patron of hags, witches, and hedge wizards



\textbf{Important Personages}: Cauldron Maker and Master Smith Clagrieda Ironborn (LN \textbf{dwarven ringmage}, see \textit{Tome of Beasts 1}); First of the Night Hags, Utarna Terrosi (NE \textbf{night hag}); Potion Mistress Jingalla Portana (NE human \textbf{mage})



\textbf{Population}: 2,100 (600 umbral humans, 500 shadow fey, 400 shadow goblins, 250 dwarves, 100 alseid, 100 ravenfolk, 100 quicksteps, and 50 various hags)



\textbf{Great Gods}: \textit{Hecate} (patron), \textit{Ailuros}, \textit{Alquam}, \textit{Baccholon}, and \textit{Kupkoresh}



\textbf{Trade Goods}: Copper, herbs, fine woven cloth, wool, crystals, iron, and dyes (in order of importance)





\end{DndSidebar}

The Court of Nicnevin, more often called the Court of the Witch Queen, stands high on a windswept mountain and is a place of somewhat dangerous repute. While the court itself is in a sheltered valley in the mountain's foothills, Nicnevin's abode is somewhere under the mountain itself, where she spends the nights of a full moon. The court's strange halls include not just beautiful, spiral-pillared dancing chambers and hearths filled with herbal, steam-laden cauldrons but also active mines, smelters and chimneys, and forges for copper, silver, and iron. The Witch Queen's interests are many and varied, and her followers include a surprisingly large number of umbral humans and bearfolk as well as fey.

\subparagraph{Long Halls and Eternal Cauldrons}
The Court of the Witch Queen is built with mountain storms and snows in mind, with enormous timbers holding up a slate roof and extending deep into the mountains. Some of its long halls resemble stave churches of the mortal Northlands. Others are more akin to the onion domes of the east or the simple turf huts of the plains. In each great hall boils a cauldron, usually for a stew but occasionally for magic working, potions, and the distillation of souls and essences. A pair of elderly women, human or alseid, is often tasked with stirring, heating, and maintaining the cauldron for a period of time, always taking turns to deliver a fine meal or a well-turned potion to one of the hags or shadow fey who decide each cauldron's next purpose.

\subparagraph{Refuge for Humans and Dwarves}
The fey ruler herself depends on umbral humans and dwarves far more than most fey lords and ladies, and she welcomes hags, alseids, and many others to her mountain eyrie. In return, they show great loyalty to Nicnevin, often creating works of magic for her court. The Dolorous Mountain is known as one of the few active iron mines in the region, and access to cold iron weapons and armor gives the court an important advantage in dealing with other fey courts. Even the iron arrows that the Witch Queen's followers are known for are quite valuable to assassins in \textit{Corremel} or duelists in \textit{Dalliance}.

\subparagraph{Crystal and Copper, Herbals and Charms}
The Witch Court is a practical one compared to many. Nicnevin seeks followers to dig out valuable crystals, ores, and stone from the mountain (dwarves and umbral humans), to gather herbs and roots from the river valleys (ravenfolk and alseids and shadow goblins), and to enchant potions, charms, and trinkets that strengthen and sustain her followers.

\begin{DndSidebar}[float=!b]{Fallen Courts}
Over the eons of fey rule in the Shadow Realm, many rulers have fallen into oblivion, living on only in songs and legends. Others built such fine halls and castles that they were rebuilt and transformed into newer halls when a new dynasty ascended to power. And a few were abandoned, though their halls remain as sites where dark magic is worked, where curses still linger, or where furtive plotters gather against some foe. These fallen courts are well-remembered and still visited from time to time.



\subparagraph{Blackwolf's Court}
A notorious court for warmongers, tournaments, duels, and deeds of prowess, Blackwolf's Court was one of the few fey courts ruled by a quickstep, Prince Blackwolf, who gathered a group of powerful shadow fey, a horde of shadow goblins, and a particularly talented human wizard named Onri Farromir, the Mage Prince. These masters of nocturnal combat and strategic dueling overthrew the Nightbrook Court and ruled it independently of the \textit{Court of Night and Magic} for more than 25 years, teaching sword skills and banditry to anyone who swore allegiance. Blackwolf was eventually cut down by Reynard, the Fox Lord of the \textit{Court of Mice and Ravens}, who kept \textit{Nightbrook} for a time, until the tides drifted too far, and he gave it then as a gift to Sarastra.



\subparagraph{Botethar's Court}
A court of the shadow goblins under King Botethar, it is remembered both for its astonishing lunar towers and astrological hexes as well as for its employment of dwarven masons, rarely seen in the Shadow Realm and likely brought to build Botethar's tallest and most ambitious temple spires to Hecate.



\subparagraph{Court of Drizzling Fire}
This was a court of dancers who wove enchantment through movement and conjured both cold fire and impossibly potent enchantments in this school of patterned, dual-casting ritual magic. All its secrets were lost when the site was long ago overrun by the darakhul. The remaining elements became the foundation stones of Gloomhaven. The last rulers of the Court of Drizzling Fire were Her Grace Tommika Telanni, a shadow fey countess, and her consort, His Grace Ondros Paraselli. Some believe the pair became darakhul themselves and rule again in the \textit{Twilight Empire}.



\subparagraph{Court of Knives}
One of the stranger fallen courts, the Court of Knives was founded by an esoteric sect of dwarven ringmages and grew into a place of smithing and wizardry—where quickstep duelists paid great sums for enchanted moon steel, where shadow fey enchanters sought out particularly well-made rings and bracers, and where ravenfolk warriors bought axes that could cleave stone or steel alike. The king and master smith was a dwarf named Olbert the Lame, a twisted and thin-bearded fellow who some said was a son of Volund, the smith-god. The court fell into ruin when Olbert refused to sell a Vorpal Sword to a prince of the shadow fey, despite being offered a king's ransom. Assassins killed Olbert not long after, and his followers fled a plague of wraiths and darakhul.



\subparagraph{Court of Silver Words}
Once a court of sable elves devoted to spoken magic, poetry, and history, the Court of Silver Words was a beacon of learning and justly famous for a scriptorium that generated both magical and mundane books and scrolls. Many of these volumes still exist, though many others are lost. The \textit{Book of Ebon Tides} itself was said to have been written here by the last ruler of the Court of Silver Words, the young Princess Tephanora, who was murdered when the book was stolen by shadow goblin rogues said to have taken the volume to \textit{Fandeval}—though it has traveled far afield since. The well-carved marble and limestone buildings of the court still stand, and occasional visitors leave offerings of honey, flowers, candles, flags, stones, and small written notes at a memorial statue of the princess.



\end{DndSidebar}

\chapter{7. Realms Beyond the Courts}
While the fey courts are among the best-known settlements of the Shadow Realm, other races have their own strongholds, cities, and even nations here. The most significant of these are the strange goblin city of Fandeval, the ghost city of \textit{Soriglass}, and the halls of the bearfolk at \textit{Merrymead}. The darker lands include the \textit{Twilight Empire} of the darakhul, the reeking pits of the roachlings on the Whispering Plains, and \textit{Oshragora}, the City Fallen into Shadow. Each is described here.

\section{Fandeval, City of Goblins}
\begin{DndSidebar}[float=!b]{Fandeval, City of Goblins}
\textbf{Ruler}: Flibbarn Barkwater, His Majesty the King of All Goblins and Master of the Moons, Duke of the Golden Field, Baron of the Black Water Fleet, Lord Mayor of Fandeval, Potentate of All Tides and Elements, Lord of the Cat Legion, and Protector of the Blue Chambers (LE shadow goblin \textbf{moonshadow catcher}, see \textit{Chapter 9})



\textbf{Important Personages}: Dour Dillos Zanabbe, Bargemaster of the Guild (LN halfling \textbf{bandit lord}, see \textit{Tome of Beasts 1}); Grandfather Clay Nethersmoke, master of the docks (LE shadow goblin \textbf{thief lord}, see \textit{Creature Codex}); Her Luminance Salarra Mortre, high priestess of Hecate the Silver Goddess (NE human \textbf{first servant} of Hecate, see \textit{Tome of Beasts 3}); Master Shipwright Adalbarr Lunkvoll, Guildmaster of the Shipwright's Guild (LN human \textbf{merchant captain}, see \textit{Tome of Beasts 3}); Royal Keeper of Cats Jennai Petticlaw (NE shadow goblin \textbf{druid}); Royal Netmaker Kotteran the Quick (LN \textbf{shadow goblin scribe}, see \textit{Chapter 9}); Ullestro of the Vale, Duke of the Lunar Light (CN sable elf \textbf{mage}); Vespilla Fleabite, Luminous Princess and Consort to King Flibbarn (LN shadow goblin enchantress, use \textbf{shadow fey enchantress}, see \textit{Tome of Beasts 1});



\textbf{Population}: 40,000 (20,000 umbral humans, 17,000 shadow goblins, 1,000 sable elves, 1,000 catfolk, and 1,000 shadow fey)



\textbf{Great Gods}: \textit{Santerr Marossa} (patron), \textit{Anu-Akma}, \textit{Charun}, \textit{Hecate}, and \textit{Kupkoresh}



\textbf{Trade Goods}: Fish, beer, grain, white oak ships, wool, salt, moonstones, and leather (in order of importance)





\end{DndSidebar}

The shadow goblin and umbral human population of the Shadow Realm is large in number but often subservient to shadow fey or darakhul rule. Not so in Fandeval, the City of Goblins, a realm of strange customs and traditions dating back to the establishment of the free city's goblin rulers, including its moon towers, gangsters, and scavengers at the mouth of the River Styx.

In the case of Fandeval, the city was granted to the Barkwater soothsayers for services rendered to the Crow Prince, a fey lord who has since abandoned his own court—but his gift is still remembered fondly by the shadow goblins. In gratitude, they still carry his messages or tokens to Queen Valda in the Court of Golden Oaks (see \textit{Chapter 6}).

\begin{DndSidebar}[float=!b]{The Crow Prince \& the Goblins}
While he ruled, the Crow Prince shared treasures, fiefdoms, arcane lore, and more with the ravenfolk, shadow goblins, and umbral humans who encouraged and supported his work in illusion arcana, in ancient shadow roads, and in lunar prophecies and moonlight weaving. The shadow goblins were extremely valuable in the last of these elements, and the Crow Prince rewarded their generosity with his own. The quest he now pursues may involve the goblins in some form. Occasionally he visits Fandeval to consult its soothsayers and scholars.



\end{DndSidebar}

Small towers throughout the city are used to shout out the hours of moonrise, midday, and moonset as well as to watch for signs of fire in the largely wooden city. The shadow goblins get their timber from the nearby forest, and while most think of the city as riverine or swampy, it is just as much a producer of barrels and ships and of carts for the shadow roads.

\subparagraph{Waterfront}
Fandeval itself lies on the Black River Styx, and it makes what it can of those waters, in boatbuilding and in the very active fishmongers' market and costermongers' stalls for its freshwater oysters. The shadow goblin fortunetellers often use fish guts or oyster shells in their soothsaying, so their stands also clutter the waterfront, between reeking pubs and well-watched warehouses. Despite the wider reputation of shadow goblins as lazy and thieving, among themselves, they take care to watch over their friends and neighbors.

The largest house in the waterfront district is Grandfather's House, for the waterfront of Fandeval is controlled by Grandfather Clay Nethersmoke, a particularly rotund shadow goblin. His network of corrupt lieutenants, dock bosses, watchmen—even judges, mooncatchers, and guildmasters—enforces the practice of skimming a touch of silver from every shipment on every barge. While he is a loyal supporter of His Majesty, Grandfather Clay collects a toll from the sale of most goods, whether brought through the city gates or made within the city, including everything from common fish and beer to memory philters, nets, and other enchanted goods.

\begin{DndSidebar}[float=!b]{Shores of the Great and Endless Sea}
The dark and roiling waters of the sea, churning beyond the marshes, are dangerous, and the shadow fey, shadow goblins, and umbral humans all know that the ebon tides are strong and destructive here, often creating reefs or shoals or towering waves. Fishing is a dangerous proposition, and most fisherfolk stick to the marshes, the river, or the tidal flats. The small Saltworker Tribe of shadow goblins uses kettles and timber from the Salt Wood to boil away water, creating bushels of a fine, dusty, flavorful salt that is sold in Fandeval and throughout the Shadow Realm. But no trade extends far on the Shade Sea. The rumors of islands somewhere beyond the horizon, created by archmagi or called up from the depths by great priests of Charun, have never borne out. The sea is the limit and boundary of all land that the umbral peoples know, and their ships always keep to the coastline.



\end{DndSidebar}

\subparagraph{Cat Legion}
Among those who patrol the city are a group of quite notably quick and lithe catfolk, the descendants of visitors from a distant desert land who choose to serve the goblin king, Flibbarn, as an elite guard and fighting corps. The Cat Legion is fond of keeping its privileges as the finest warriors and sharpest-eyed bodyguards to His Majesty, and they also serve as special messengers and agents of the king's will against the shadow fey, the ghouls, and others. Grandfather Clay seems to have some arrangement with the legion, for his shops and warehouses are rarely disturbed while those of his enemies are sometimes shut down and illegal goods confiscated.

\subparagraph{Wading Giants}
Three river giants (or Styx giants, see \textit{Chapter 9}) make their homes on the eastern shore of Fandeval and provide transport both into the Sodgrass Marshes and across the river from near the Styx Tower to the Goblin Plains to the east. These giants, nicknamed Memory, Sorrow, and Loss, use glistening Mooncatcher's Net (see \textit{Chapter 10}), made only by the mooncatcher goblins and priests of Hecate in Fandeval, to pull souls and memories from the water. Once they are netted and bottled, these are traded to the mooncatchers and wizards who transform them into memory philters (see \textit{Chapter 10}), healing draughts, and sometimes into Moonsteel Weapon (see \textit{Chapter 10}).

\subparagraph{Moon Towers and Moonlight Masses}
The city of Fandeval is famed for its moon towers, which are each topped by a globular dome that lights up at night. Built as shrines to Hecate—and as acts of ostentatious devotion by the wealthy—these towers also house the Lady of Darkness's priests. The brightest moon tower is connected to the Cathedral of the Dark Moons, which holds private services for initiates of the goddess's inner mysteries at each of the new moons and full moons (both moons reach their greatest and least brightness in synchrony). Her Luminance Salarra Mortre leads these ceremonies, and her various blessings and sacrifices include the Transformation of Moonlight, which turns lesser fey into shadow goblins. Other ceremonies are said to turn silver into memory philters or to weave nets for the Styx giants.

\section{Groves of the Wardens}
The bearfolk groves of the wardens are druidic sites and training grounds for bearfolk shadow gnawers. Several have grown into small villages or settlements, but their primary function remains as training grounds for bearfolk wardens and guardians against shadow. The primary groves are briefly described here.

\subparagraph{Eagle Grove}
Up in the hills by a stream's headwater stands a single stone temple to Umbeserno, a pilgrimage site for the bearfolk God of Dreams. His priests maintain the grove and teach shadow gnawers and others how to protect themselves against shadow corruption. They also sell amulets and holy draughts of mead to aid in winter's sleep. These sales and pilgrimages end in late fall when the temple priests retire to their own slumbers until spring.

\subparagraph{Hawthorn Grove}
Favored by birch sirens and river giants (see \textit{Chapter 9} for both) as well as treants and various ogres, unicorns, and alseids, the Hawthorn Grove remains the most senior school and seminary for young druids of all kinds in the Shadow Realm. As many as two dozen initiates receive tutelage here at any given time under the auspices of Lady Rowan (CG \textbf{bearfolk druid}, see \textit{Chapter 9}).

\subparagraph{Leaping Salmon Grove}
A small bearfolk settlement with a rickety wooden bridge over a narrow point of the Styx, Leaping Salmon Grove was once purely a druidic circle but has since grown into a small town of almost 1,000 bearfolk, wyrd gnomes, and erina, most of them fisherfolk. They smoke and preserve much of their catch for trade.

\subparagraph{Lost Grove}
Far from the larger bearfolk realms, this grove once stood near the Whispering Plain. A ritual gone awry at a time of great flux in the ebon tides broke it loose from the wider bearfolk realm, and it drifted in a matter of months across the Whispering Plain and finally came to rest near the Silver Hills. It is still a potent stone circle of power, ably led by a fierce defender, Lord Istler (LG bearfolk \textbf{knight of shadows}, see \textit{Chapter 9}), and his company of shadow gnawers and two lunar elf archers, nicknamed Folly and Flight.

\subparagraph{Moonstone Grove}
Known for its druid school and its enormous standing stones, tall trees, and a moon tower, Moonstone Grove is too quiet to interest most fey and too tranquil for most goblins. The chief druid and village elder is Franelle Elderberry, a human woman who has spent her life among the bearfolk. She left the fey courts and retains a deep dislike of the shadow fey and shadow goblins generally.

\subparagraph{Sleeper's Grove}
High in the mountains, this deeply scented grove of pines is inaccessible in winter. The druids here are skilled in winter magic and know many avian forms for shapeshifting, leading some jesters to refer to them as the "owlbearfolk." It is the site of dozens of tombs for honored bearfolk, especially royalty and archdruids. It stands high in the mountains with a few year-round guardians and caretakers, led by Chief Tomb Warden Tromble Dorran (LN bearfolk \textbf{gladiator}). Small bands of ravenfolk visit from time to time, perhaps to speak with the dead. The guardians frown on this but do not prevent it.

\subparagraph{Three Hives Grove}
A center of honey production, as its name implies, this grove is a ramshackle collection of bearfolk cabins and a wooden palisade—a sleepy timber town. It also features an arcane well of power, a radiant well that somehow emits moonlight and arcane power that the local druids harness to turn into moonstone, moonsteel, and other useful substances of pure, solidified radiance. Lunar devils and followers of the Moonlit King have visited repeatedly, perhaps to trade arcane secrets with the Beekeeper in Chief, Lord Aparombo (CG \textbf{bearfolk thunderstomper}, see \textit{Chapter 9}).

\subparagraph{Veiled Grove}
Barely connected to the capital of Solvheim by a narrow track, the Veiled Grove is a sanctuary and holy ground where any outlaw or exile may find shelter, be they bearfolk, shadow fey, or a shade from \textit{Soriglass}. Its trees are enormous pines of a kind not often seen, their trunks reddish and their leaves tending to dark-plum shades. The druid of the Veiled Grove is Rumbella Nuuneva, a grandmotherly bearfolk who seems able to pierce illusions and sense the truth of corrupted hearts. A small spring in the Veiled Grove is used to cleanse shadow corruption from those who fall prey to it, attracting occasional pilgrims and wayward adventurers.

\begin{DndSidebar}[float=!b]{The Coming of the Bearfolk / from the Book of Ebon Tides}
The bearfolk came to the land of shadows in two great waves: first were the druids and their servants, and then the warriors, chewing shadows and destroying the darakhul, the vampires, and others. Their forests here grow quickly, and their roads are strengthened with radiance, but the bearfolk remain vulnerable in the sleepy winter months. Over the centuries, their presence has gone from a few wanderers and stone huts to a northerly realm of its own. Everywhere they settle, the snows are thicker and the forests a bit healthier. Their druids can remove shadow contamination from even the worst afflicted, and their warriors are a match for anyone.



\end{DndSidebar}

\section{Merrymead, City of Revels}
\begin{DndSidebar}[float=!b]{Merrymead, City of Revels}
\textbf{Rulers}: His Sonorous Majesty King Zaim Gundervolk (CG \textbf{bearfolk druid}, see \textit{Chapter 9}) and Her Gilded Majesty Queen Hollibee (LG bearfolk \textbf{apostle} of Bengta, see \textit{Tome of Beasts 3})



\textbf{Important Personages}: Double Time Karlo von Karlose (CN \textbf{quickstep}, see \textit{Creature Codex}); Ferryboat's Keeper Linne Mallemele, high priestess of Charun (NE human \textbf{apostle} of Charun, see \textit{Tome of Beasts 3}); Fragrant Servant Sayzanette Belleflor, high priestess of Ailuros (CG human \textbf{priest} of Ailuros); Her Solemnance Snuffling Tilda (NG bearfolk \textbf{first servant} of Umbeserno, see \textit{Tome of Beasts 3}); Lady of the Forest Alossa Tergehar (LN human \textbf{druid}); Master of the Hives Ingmore Meadbeard (NG \textbf{bearfolk druid}, see \textit{Chapter 9})



\textbf{Population}: 11,000 (6,500 bearfolk, 2,000 umbral humans, 1,600 shadow goblins, 500 darakhul, and 400 wyrd gnomes)



\textbf{Great Gods}: \textit{Bengta} (patron), \textit{Ailuros}, \textit{Charun}, \textit{Kupkoresh}, and \textit{Umbeserno}



\textbf{Trade Goods}: Honey, wheat, oats, ironwork, maple and pine wood, and pork (in order of importance)





\end{DndSidebar}

A city of joy and light and strength is unusual in any kingdom, but in the dark landscape of the Shadow Realm, Merrymead seems a mirage. Built of great oak timber daubed with clay and painted in bright, warding patterns of flowering vines, every rooftop studded with beehives, and wandering droves of pigs grown fat from acorns, Merrymead is a lively point in a shadowy forest, and it enjoys frequent trade as it sits at the crossroads between \textit{Fandeval}, the \textit{Moonlit Glades}, and the forests of the shadow fey. It is an especially seasonal sort of city, having a much faster, even frenetic pace in high summer and fall, compared to the sleepy winter and the difficult days of spring.

The city's role as a crossroads extends beyond the roads within the Shadow Realm. One of its roads leads back to the mortal world, to a sophisticated, well-ruled human city with a large temple to a goddess of cats, hunters, and perfumes.

\subparagraph{The Honey Fair}
Each year at high summer, the King and Queen of Merrymead open a trade fair, one of the few in the Shadow Realm. It is open to everyone who can obey its rules: fair speech, fair bargains, and no brawling outside the Royal Ring. Many deals are struck for great vats of honey and mead, for lengths of bearfolk timber and goblin moonsteel, and even for strange treasures dug up from \textit{Soriglass} or brought down from the stars. The trade fair sometimes even carries on into mortal lands when a bear king of the mortal world opens a shadow road, as the Honey Fair ends, and extends an invitation to travel the Honey Road to bearfolk lands in the mortal world.

\begin{DndSidebar}[float=!b]{A Road to Midgard}
The forest paths to the northeast of Merrymead contain a shadow road to Midgard's Western Wastes. Those many layers of fern, bracken, and undergrowth eventually lead directly into the temple of Ailuros in Bourgund, and worshippers there are sometimes surprised by the sudden appearance of bearfolk.



In addition, each year at the honey fair, the bearfolk king of Björnrike or the bearfolk queen of the Arbonesse opens the Honey Road for a week. They then extend an invitation for bearfolk of the realm to come to their sunny lands, strengthening the bonds between the bearfolk of the Shadow came from, long ago.



\end{DndSidebar}

\subparagraph{The Slumbering Winter}
The city of Merrymead is still populated in winter months, but its bearfolk do very little but sleep, dream, and wake for brief periods to see the snow. During the slumbering months, the regency is granted to a trusted human, such as the high priestess of Charun, the boatwoman Linne Mallemele, or else the high priest of Umbeserno, Snuffling Tilda, Keeper of the Snoring Temple—who, as all priestesses of the Lord of Sleep, enters the winter sleep late and wakes early. One of the priestesses watches over the bears and keeps the city's business moving, albeit at a sedate pace until spring. Their primary goal is to provision the city well for the spring hunger.

\subparagraph{The Spring Hunger}
Come the spring and the snowmelt, Merrymead stirs and wakes itself, extremely thin and hungering. Darakhul in particular seem to understand the bearfolk hunger in a way that gnomes and humans don't quite, and their merchants always forge through the late snows to bring fish, spring lamb, potatoes, and other provisions to Merrymead. Vast amounts of food are devoured in just a few weeks as the bearfolk awaken family by family, and the price of spring provender is double that of the rest of the year. In short order, the festivals of Bengta and Charun are held, the shrine of Umbeserno holds the rite of Waking the Last Sleeper, and the dream season ends with bearfolk returning to the hives, fishponds, groves, and fields.

\section{Moonlit Glades}
\begin{DndSidebar}[float=!b]{Moonlit Glades}
\textbf{Ruler}: Gulfwyr Moonrage, Chief of the Moonlit Glades (CG werebear \textbf{bearfolk chieftain}, see \textit{Creature Codex})



\textbf{Important Personages}: Elder Whitepaw (CG \textbf{bearfolk druid}, see \textit{Chapter 9}); Hierophant Ernalda Berlasdottir (NG \textbf{bearfolk druid}, see \textit{Chapter 9}); Kordan Wodenson, reaver chieftain (LN \textbf{wolf reaver dwarf}, see \textit{Tome of Beasts 1})



\textbf{Population}: 27,000 (21,000 bearfolk, 3,500 umbral humans, 1,500 dwarves, 500 half-elves, 400 shadow fey, and 100 various werebears, alseid, and others)



\textbf{Capital}: Solvheim, population 12,000




\begin{description}
\item[Other Settlements:]
Eagle Grove, Hawthorn Grove, Leaping Salmon Grove, Sleepers Grove, Three Hives Grove, Veiled Grove
\end{description}



\textbf{Great Gods}: \textit{Bengta} (patron), \textit{Charun} (called Hod), Freyr and Freyja (Yarila and Porevit), Thor, and \textit{Umbeserno}



\textbf{Trade Goods}: Medicinal herbs, pure food and water, silver, moonlight honey (in order of importance)





\end{DndSidebar}

The Moonlit Glades is the most prominent ray of light and hope in the Shadow Realm. Here the bearfolk dwell and tend their small patch of home against the darkness, and their numbers grow steadily from year to year. While the bearfolk are natives of the Shadow Realm, they remain as vulnerable to the plane's corruptive influence as any visitor from the mortal world. That generations have lived and died within the shadows without falling to its temptation represents a minor miracle—and a sign that the bearfolk understand shadow's temptations better than most and that their druids and gods protect them.

\subsection{The Moonswept North}
The Moonlit Glades nestle deep in a black pine forest of the Shadow Realm's far north, corresponding loosely to the lands of ice and snow in the mortal world. The moon shines its face upon the glades as if blessing the bearfolk in their desperate struggle. Shadow roads to \textit{Soriglass} and the fey courts offer the most reliable and timely path of entry into the glades, and the bearfolk have wrested control of these paths from the shadow fey. Stone and silver cairns mark the borders of the Moonlit Glades, and moonlight always bathes these markers. Bearfolk druids conduct rituals at these stones, renewing the magic that holds shadow's touch at bay. Each bearfolk settlement similarly features a standing stone or cairn at its center, creating another layer of protection for its residents.

Given their relative safety from the more powerful denizens of the Shadow Realm, the bearfolk have nurtured their claim of the dark land. The Moonlit Glades have secured numerous pure resources here, and the bearfolk fight furiously to maintain them.

Waterways in the Moonlit Glades run clear, glittering like rivers of molten silver beneath the moonlit sky. The spirits that haunt rivers and streams in the rest of the Shadow Realm grow quiet and peaceful within the bearfolk's borders, allowing travelers to drink unmolested. The bearfolk tribes maintain a strong druidic tradition as well as the worship of their own gods, Bengta and Umbeserno, and the Green Gods in their masks as Freyr and Freyja. These divine servants spend most of their time tending to agriculture, tending to their fruits and vegetables free of shadow corruption. Even the beasts are more like their material counterparts here, and their meat is wholesome and nourishing. Few things rouse the ire of the bearfolk more quickly or more violently than poachers.

Settlements in the Moonlit Glades are small and dispersed, but each remains in close contact with the others. Runners, hunters, druids, and priests make regular rounds between the villages and the capital settlement of Solvheim, carrying news, delivering requests for aid, and reinforcing close ties with the surrounding tribes. Originally the population of the Moonlit Glades was purely bearfolk, but over the years, it has adapted to include humans, elves, and others. Other races from the Kingdom of the Bear have come to the Shadow Realm, initially to maintain contact between the shadow colony and their kin in the mortal world, but eventually many integrated into the existing bearfolk communities. Humans are the most prominent of these newcomer races, but a contingent of reaver dwarves who venerate bear spirit totems found their way here and decided to settle. The dwarves represent an invaluable asset to the Moonlit Glades since they are among the few seafarers willing to sail on the fog-shrouded black waters of Frost Gulf. Even these stalwart sailors don't venture far from land, but their crafts allow the bearfolk and their allies unprecedented, swift access to other areas of the Shadow Realm, when the shadow roads become too dangerous or don't lead to the proper destination. The Rivers Styx and Lethe are far safer for the moonlit reavers to navigate though, and they occasionally raid darakhul, goblin, and shadow fey villages along the waterways.

Some renegade shadow fey call the Moonlit Glades home, though they must go to extraordinary lengths to prove their honorable intentions and win the bearfolk's trust. Half-elves as well as a fair number of alseid and werebears are scattered throughout the population. The Moonlit Glades have a calming effect on werebears, and even during the full moon, these lycanthropes find their more destructive urges curbed. Werebears often seek the way to the glades, so they can live out their lives in righteous service with less fear of losing themselves to the curse.

\begin{DndSidebar}[float=!b]{The Bear Kingdom of Midgard}
A proud people, the bearfolk originally hailed from Björnrike, the Bear Kingdom in the Northlands of Midgard. The ruler of Björnrike, Mesikämmen the Bear King, foresaw a terrible doom growing in the Shadow Realm that threatened all his people. He gathered his most trusted warriors and druids and charged them with an impossible task—to enter the Shadow Realm and fight back the corruption at its source. His brave subjects set forth on a forbidden shadow road without hesitation and never looked back.



\end{DndSidebar}

\subsection{Ties that Bind: Rulers of the Glades}
Despite their differences, the peoples of the glades consider each other to be family. A diarchy consisting of the chieftain and the hierophant rules the nation. The chieftain oversees martial defense and civil government, serving as the highest authority to settle disputes among tribal leaders. The hierophant represents the spiritual center of the Moonlit Glades and is the most trusted counselor on magical matters, including for the elevation and ordinance of new druids and clerics and for consultation on matters of shadow corruption, the flow of the ebon tides, the management of boundary stones, and similar issues.

\subparagraph{The Chieftain}
The current chieftain, Gulfwyr Moonrage, is a massive grizzlehide bearfolk werebear whose pitch-black fur is deeply frosted with silver. A terrifying opponent in battle, Gulfwyr earns his deed name hundreds of times over, channeling his rage into the destruction of his foes and the protection of the Moonlit Glades. His restraint regarding the werebear curse borders on the stuff of legends, and he has never infected another creature. Outside of battle, Gulfwyr is wise and fair, and many already consider him a living ancestor. He brooks no foolishness when deciding disputes and is quick to punish a party bringing frivolous suits. Living in the depths of the Shadow Realm, there is always more at stake than spite or selfishness.

\subparagraph{The Hierophant}
Ernalda Berlasdottir, the hierophant, provides the spiritual and magical counterpart to Gulfwyr's battle prowess and civic dictates. A slender but tough purifier bearfolk, she decorates her long, red-gold fur with braids, beads, feathers, and silver charms. The most powerful druid in the Moonlit Glades, she takes great care to learn of any shadow incursion into the relative purity of her home. She provides the single greatest source of healing and restorative lore in the entire realm: there is no affliction that she can't identify, if not relieve. She devotes her time to maintaining the moonlit wards that protect her home and pursuing cases of shadow corruption.

\subparagraph{The Wanderer}
A strange sight in the holdings of the Moonlit Glades, a venerable bearfolk with snow-white fur, called Elder Whitepaw, calls no settlement nor village home. Instead, he wanders from town to town. The shaman's demeanor can be strange and sometimes off-putting, for he converses at random, answering unheard questions or shifting with no warning or preamble from a discussion with present individuals to one with figments of his imagination or beings otherwise imperceptible to any but himself.

Despite his strange habits, Whitepaw is the only bearfolk other than the chieftain with a direct line of contact to Hierophant Ernalda. Though his wandering seems random, Whitepaw is deeply in tune with the healing and awakening spirits of nature within the cleansed Moonlit Glades. Elder Whitepaw has a deep knowledge of the spread of shadow corruption and the influence of dark entities, such as the umbral vampires, cultists of the Black Goat of the Woods, and the undead, but communicating that knowledge is often difficult. Whitepaw often stumbles into agents of the Lantern Bearers or a visiting band of adventurers and offers them a task or request of great import. Individuals who perform admirably earn Whitepaw's favor and the goodwill of the bearfolk.

\subsection{Against the Darkness}
When the bearfolk first walked the shadow roads, they had only the faintest idea of what awaited them or even what foe they sought to thwart. The peril they faced was deceptively minimal at first: the occasional shadow or unquiet spirit, predatory beasts that sought their first taste of uncorrupted flesh. These dangers were faced without difficulty, insulated from the most insidious creatures of shadow by their remote location. All too soon, they learned the true danger of the Shadow Realm remains one's own imperfect nature. Their numbers began to feel the force of corruption, and division struck their ranks. Melancholy gave way to paranoia, and that gave way to conflict. As the corruption spread, those who proved more resistant tried in vain to help their fellows, which only added fuel to the fire. Before long, the corrupted bearfolk turned on their own kin, and the blood spilled drew the attention of the umbral vampires. One day they were simply there, arising from dark ruins the bearfolk had never seen before. The warriors fought valiantly, but with their ranks reduced by spreading corruption, they knew they were doomed.

It was then that the druids' last desperate plan succeeded, and hope kindled anew. Using the last drops of pure water brought from their home along with silver mined in the Shadow Realm, the druids drove back the darkness empowering the invaders. Through this victory, the bearfolk discovered how shadow corruption—coupled with acts of betrayal, murder, and the resulting despair—bored holes through the veil between the Shadow Realm and the mortal world. The umbral vampires slipped through these widening cracks, their timeless nature drawing them to places where shadow bleeds into the world. In the light that shone from the heavens, the bearfolk saw the truth that led them to this place and resolved to oppose the spread of shadow from within. Bolstered by their new discovery, the bearfolk rallied and drove the wretched creatures back into their ruins. These strange slivers of the City Fallen into Shadow melted back into the darkness as moonlight overtook them.

The bearfolk's knowledge brings them into conflict with powerful forces within the Shadow Realm, namely the darakhul of the \textit{Twilight Empire} and the umbral vampires of \textit{Oshragora}. Each of these has a terrible knack for finding rifts and cracks to exploit, and the bearfolk know they pose a risk of large-scale shadow corruption pushing into their homeland here as well as their mother lands in the mortal world.

\subsection{Allies in the Light}
The Moonlit Glades form a solid bastion of hope in the Shadow Realm. The nation is protected from shadow corruption and from the worst depredations of shadow creatures because of its moonlit wards. The bearfolk understand how pervasive shadow is. Bearfolk and their allies are occasionally corrupted, but they hold the key to long-term survival: they can perform a ritual passed down from the first settlers of the Moonlit Glades that can scour shadow corruption with cleansing moonlight.

The druids regularly conduct rituals to maintain the moonlit wards around their borders and in the center of each settlement. Because they can remove shadow's touch, the Moonlit Glades became the center of the network now known as the Lantern Bearers. Though the bearfolk never teach their purification rituals to outsiders for fear that forces of shadow will find a way to corrupt or circumvent them, they are willing to use their gifts to aid outsiders struggling against shadow corruption. The relative safety and purity of the Moonlit Glades is recognized throughout the Shadow Realm, and the hidden allies of light know the place may be their only hope in desperate times.

A few far-flung settlements of bearfolk and their allies have been established far from the Moonlit Glades. Each revolves around a standing stone inlaid with silver or a stone-and-silver cairn as in the glades, holding back the shadow. Usually away from well-traveled areas, these outposts always have a specific purpose. Most often they observe an area where the barrier between the Shadow Realm and the mortal world has worn thin. They watch to ensure that incursions from the Shadow Realm can't push through these eroded doorways. These villages and small forts also serve as a sanctuary or last-ditch salvation for friendly travelers lost in the darkness.

Guides from the Moonlit Glades are fascinated by the strange glowing crystals that float through the Shadow Realm—the witchlights. The bearfolk and their ilk aren't entirely sure what to make of these constructs, but they have seen them time and again lead travelers out of danger—and sometimes to exactly what the lost souls need to survive. They don't fully trust the witchlights though since they resemble will-o'-wisps. Although there are no confirmed reports of witchlights harming or misleading travelers, they are suspicious enough to warrant healthy caution.

\section{Oshragora, the City Fallen into Shadow}
\begin{DndSidebar}[float=!b]{Oshragora, the City Fallen into Shadow}
Oshragora, the City Fallen into Shadow



\textbf{Ruler}: Lord of Oshragora and Master of the Violet River, Keeper of the Sanguine Shrine, and Seneschal of the Salamander City, Iktorros Nalgevar (CE \textbf{umbral vampire}, see \textit{Tome of Beasts 1})



\textbf{Important Personages}: Elenya the Proud, Lady of the Castle Darros, First Spawn of Iktorros (NE \textbf{umbral vampire}, see \textit{Tome of Beasts 1}); His Excellence Cilestros Margevan, high priest of the Hunter (NE \textbf{lich}); Treposs Palendulo (LE human \textbf{archmage})



\textbf{Population}: 600 (100 dark folk, 100 umbral humans, 100 shadow goblins, 50 darakhul, 50 derro, 50 satarre, 50 shades, 50 selang, and 50 umbral vampires)



\textbf{Great Gods}: \textit{Alquam} (patron), \textit{Goat of the Woods}, \textit{Hunter}, \textit{Santerr Illosi}, and \textit{Vardesain}



\textbf{Trade Goods}: Statuary, trained gargoyles and gremlins, amulets, wolf pelts, and fish (in order of importance)





\end{DndSidebar}

Once men and women, the umbral vampires now exist as twisted fiends sustained by the power of shadow and hungry for the life they once held. Not undead but instead fully creatures of the Shadow Realm, umbral vampires resemble pale, gaunt humanoids with long, pointed ears and sharp features—darkness weeping like smoke from their eyes, mouth, and nose—and their touch drains the life and vitality from their victims. They drift through the City Fallen into Shadow, a crumbling ruin twisted through space and time.

\begin{DndSidebar}[float=!b]{Using the Oshragora 8-Section Map}
The City Fallen into Shadow is splintered along a fourfold axis, dividing the city into eight segments around the Storm Tower at the center. When using the map, you can either keep segments \textit{A}, \textit{B}, \textit{C}, and so forth in clockwise sequence up to \textit{H}, or you can shuffle the segments as PCs explore it. The blue barrier lines mark the map but are rarely visible in the city itself: these are where the splinters move, grind against each other, or rotate to shuffle the city's geometry. The Storm Tower is always the center and often visible from any other section.



\end{DndSidebar}

\subparagraph{Albent Stadium}
Athletic feats and prisoner blood sports are held here to amuse the vampire overlords.

\subparagraph{Burning Towers}
The city has at least two towers wrapped in perpetual flames, usually a deep orange although sometimes purple—sometimes it is three or four such towers, when visual or temporal distortions double them. These towers were built over radiant wells that were tampered with, and lines of arcane power provide a constant stream of phlogiston to burn without any other fuel. The umbral vampires avoid them, but fire elementals sometimes dance in the flames, and some believe a shadow road to the City of Brass once began in one of these towers.

\subparagraph{Castle Darros}
The castle and the nearby barracks are home to Lady Elenya and her handmaidens, a group of wraiths, as well as her court of shadow goblins and human thralls.

\subparagraph{City Hall}
A fine, half-timbered building that is one of the few structures in very good repair, this hall is the hub of commerce, industry, and plotting for the creatures that dare to visit the city. Its basement is the Ratskellar, the only tavern in town, and the fare is notoriously poor. The specialty is fried purpure, a sort of lilac-colored fish native to the city's river. Lord Nalgevar maintains an audience chamber somewhere in the dungeons beneath the hall.

\subparagraph{Dusk Guildhall}
The vampires that often visit their umbral cousins founded a guild here, devoted to the spread of the city's people and power. As a strange sort of explorer's society, they exchange information and pursue targets among the darakhul, humans, and shadow fey. Most members are dhampir, vampires, vampire spawn, or shades. Some of the Dusk Guildhall members seek a close alliance with \textit{Soriglass}.

\subparagraph{Storm Tower}
This tall, Gothic tower with its many gargoyles, bloodthirsty vines, and elaborate balconies is home to the city's umbral vampires, and features eight large bronze doors at its base. As the home to many of the city's commoners and the most habitable of its buildings, the Storm Tower is a useful crossroads between all the splinters of the city, no matter how they twist in space. Lightning and flickering witchfire often illuminate the tower and azza gremlins and spire walkers (see \textit{Tome of Beasts 1} for both) are often visible. The city's rulers are rarely found here though, instead favoring Castle Darros and City Hall.

The ghouls of the \textit{Twilight Empire} believe an ancient artifact is buried somewhere beneath the tower and frequently send spies or scouts to attempt to recover it (thus far, without success).

\subparagraph{Violet River}
This tributary of the River Styx runs though Oshragora silently and largely without hindrance. Sometimes, the river seems to flow into a portal or void, and at other times, it appears seemingly from nowhere in the middle of a city street.

\begin{DndSidebar}[float=!b]{How to Use Temporal Distortions in the City}
The degree of chaos and weirdness visible on the streets and among the ruins is up to the GM. To add some level of strangeness to exploration, roll on the \textbf{Temporal Distortions} table every time creatures pass from one section (or splinter) of the city to another, ending any prior effect and using the new one. Creatures that spend more than a day in Oshragora slowly see less and less of such distortions as the city seems to acknowledge them as natives within a week.



\end{DndSidebar}

\begin{DndTable}[header=Table 7-1: Temporal Distortions]{cX}

d8 & Temporal Effect\\
1 & Perception slows, causing everything that moves to become blurry and indistinct. All creatures in the splinter are affected by the \textit{blur} spell.\\
2 & Everyone in the splinter hops between two different timelines. Each living creature is affected by the \textit{blink} spell. Affected creatures are safe from attacks when they jump to the other timeline.\\
3 & Time is inconstant in the splinter. Each creature that rolls an even result on its initiative roll gains the effects of the \textit{haste} spell. Each creature that rolls an odd result on its initiative roll gains the effects of the \textit{slow} spell.\\
4 & Each creature in the splinter gains a slightly out-of-sync double. Each time a creature targets you with an attack, it must roll a d20. On a result of 1–10, the attack hits your double. On a result of 11–20 the attack hits you. Your double's AC and hit points are equal to yours at the beginning of combat. Your double can't move or take actions independently of you.\\
5 & Time jumps backward in short bursts. At the end of each of its turns, each creature regains 10 hit points, up to its maximum hit points. On the first round after it dies, each creature that is dead must roll a death saving throw. On a success, it regains 10 hit points and resumes living as though it had never been dead.\\
6–8 & No temporal effect.\\
\end{DndTable}

\subsection{Sundered History}
The greatest cautionary tale of the Shadow Realm is undoubtedly that of the umbral vampires. Once residents of a thriving metropolis, the rulers of this long-forgotten city were powerful spellcasters whose might put to shame even the greatest sorcerers of this current age. They were wise and just and used their power to preserve the glory of their people. And their decline began subtly.

In an effort to prevent the inevitable march of time, the great archmages twisted it and diverted its flow. The entire population of the great city freed themselves of the shackles of age and death. For generations, they existed in this state, developing even greater magic, fine art, and architecture to rival the palaces of the gods. But for all their wisdom and might, they were blind to the danger set before them. As the river of time in the world around the city flowed on, it began to erode the spells holding it at bay. Eventually the spells failed, and the city drowned in all those many years it had denied. In the rush of uncontrolled magic and rampantly fluctuating time, the entire city vanished from the mortal world. Only the barest traces remain because the entire city was pulled into the darkness between worlds. The people of the once-great city felt the ravages of shadow corruption, but they were trapped in chaotic temporal storms. They spent an eternity in that closed loop of time, devoured by shadow, reborn and consumed again and again until finally the timestorm abated.

When the winds of history settled, Oshragora, the City Fallen into Shadow, was born. Great, towering ruins appeared in the landscape of the Shadow Realm. Cobbled streets cracked and twisted, following new paths dictated by shadow's distortions. Buildings that once touched the sky now stand in a frozen jumble of collapse, held in unstable pockets of temporal stasis. These pockets fail, unwind, and reassert themselves at random. Bits of the city crumble or reassemble before an observer's eyes.

The city defies conventional understanding of geography and space, even among other places in the Shadow Realm. While the rest of the plane shifts over the course of months or years, the City Fallen into Shadow seems to change without a predictable pattern or schedule.

\subsection{Inscrutable and Uncountable}
Oshragora isn't sufficiently organized to be called a city-state, much less a nation. An unknown number of umbral vampires lurk in its crumbled buildings and shadow-choked alleys. Despite repeated attempts, no scouting force or explorer has ever made an accurate estimate of the city's cursed population. Umbral vampires simply slip between the cracks of reality as if they're still shifting through time and space. One half-mad shadow fey hunter raved that she fought and slew the same umbral vampire for weeks while lost in a forest that continually coiled back on itself.

The umbral vampires' motivations remain obscure—or at least prone to chaotic shifts. On occasion, one might be willing to speak with visitors rather than simply devouring their essence, but these conversations always prove frustrating. The vampires mention events out of chronological order and make references to people and places no sage can identify. In spite of this, they give the impression that they understand far more than they're able to effectively communicate. Shadow fey emissaries are certain that a ruling council still exists somewhere in the city, but the place's unstable nature has made it impossible to find. The Queen of Night and Magic has proclaimed that she will trade an emperor's ransom in treasure and secrets for a reliable connection to the umbral vampires' dark council—or at least to Lord Nalgevar.

Despite the difficulty of operating within the city and communicating with the vampires, other creatures exist there. Undead shadows thrive in and around the city (and in the splinters of the city that appear elsewhere), drifting in the wake of umbral vampires. For some reason, the shadows don't attack the umbral vampires, suggesting either some mastery the umbral vampires hold over the shadows or at least a perceived kinship. Ghouls from the \textit{Twilight Empire} also make frequent trips to the city. Through covert means and interrogation, the shadow fey believe that the Twilight Emperor of the ghouls searches for some lost artifact hidden in the city. Whatever this artifact is, it likely has power over time, and it might even have been a part of the magic that doomed the city in the first place.

Other stranger visitors appear in and around the city and its splinters, such as shadow mages and eonic drifters (see \textit{Tome of Beasts 1}) who step between the past and future as one might cross into the next room, enigmatic travelers of the aeons. This is hardly surprising given the city's tumultuous relationship with time. Due to the city's unique nature and potent magic, travelers from all realms and times find themselves drawn to it.

\subsection{Ever Reaching}
The other denizens of the Shadow Realm shun the umbral vampires and drive them back with blade and spell whenever their paths cross. As a result, the umbral vampires can't readily make use of shadow roads or even those spots where the planes connect to the Shadow Realm. Though much reduced from their former glory, they remain powerful and devilishly cunning. Perhaps their most disturbing talent is the ability to find weak spots in the barrier between the realms. This shadow reflection of the mortal world, where despair, blood, and darkness tear rifts between the worlds, draws the vampires, and their presence seems to hasten the opening of doorways to the Shadow Realm. Thankfully, the umbral vampires seem compelled to remain close to the ruins of Oshragora, so they exist in a limited area.

Just as the city seems to warp the land around it, or at least to move in some fashion, so too does it occasionally shed pieces of itself. These splinters of the city, sometimes as little as a single ruined building and other times an entire neighborhood or district, simply appear elsewhere. As with the other distortion effects of the city, a splinter might emerge miles away and remain for days or weeks at a time. The splinters drag umbral vampires and other denizens along with them, depending on the size of the splinter. A single building might contain one umbral vampire and a host of shadows. A district, on the other hand, might release an onslaught of a dozen or more umbral vampires on the surrounding area.

The bearfolk from the \textit{Moonlit Glades} have kept a watchful eye on Oshragora after a few disturbing appearances of splinters within their domain. They grow increasingly worried with each vampire encounter but don't have enough useful information to plan a suitable defense. So far, they've contained the umbral vampires in a dozen splinters until the splinters vanished, but their lack of arcane knowledge hampers their efforts beyond pure reaction. They're quite eager to learn how the splinters function and how to keep them at bay.

The most recent addition to those scrutinizing the city includes the \textit{Court of One Million Stars}. These strange fey drift through the ruins while enshrouded in starlight, which pains the shadows and umbral vampires. They study the city, but they meticulously avoid the areas of distorted time. Only one such creature has engaged in the briefest exchange while about its task. After obliterating a pack of umbral vampires with an incandescent wave of starfire, the luminous creature whispered to the only surviving adventurer, "The Void calls to the Void. It must not be allowed to answer." The adventurer recounted his story to Revich, the blind angel in the \textit{Court of Night and Magic}, before hurling himself from the top of the angel's tower.

\section{Soriglass, City of Shades}
\begin{DndSidebar}[float=!b]{Soriglass, City of Shades}
\textbf{Ruler}: Tomarrich the Howler, His Spectral Majesty and High Lord of Soriglass, the Keeper of the Final Portal, Warden of the Violet Marshes, Lord Protector of the Boneyard Pass (CE \textbf{wraith})



\textbf{Important Personages}: Happy Fillbec, Bargemaster of the Guild (NG riverfolk halfling \textbf{merchant captain}, see \textit{Tome of Beasts 3}); Her Luminance Ingalla Vert, pantheist priestess of Hecate and Santerr Illosi (NE darakhul \textbf{first servant} of Santerr Illosi, see \textit{Tome of Beasts 3}); Valish the Starcatcher (LE satarre \textbf{void speaker}, see \textit{Creature Codex})



\textbf{Population}: 2,000 (1,000 umbral humans, 600 shades, 200 darakhul, 100 shadow goblins, 50 satarre, and 50 halflings)



\textbf{Great Gods}: \textit{Anu-Akma} (patron), \textit{Charun}, \textit{Hecate}, \textit{Santerr Illosi}, and \textit{Vardesain}



\textbf{Trade Goods}: Papyrus, purple dyes, silver, moondust, requiem, and statuary (in order of importance)





\end{DndSidebar}

Built in aeons past along the banks of the River Styx by the sable elves, the city of Soriglass is an enormous mausoleum—a ruin more famous for its crypts, ghosts, and shades than for any particular power, art, or wealth. Its inhabitants collect tomb-dust to make the drug called requiem (see \textit{Requiem, Drug of the Afterlife} sidebar), and visitors often come to explore one of the two main dungeon complexes of the region: the Cloisters of the Outer Void, largely held by ghouls, satarre, and undead; and the Pit of the Squamous Fiend, said to connect directly to the lower layers of the Eleven Hells, inhabited by enough fiends and aberrations to make that seem likely. Both are death traps, brimming with fiends and undead, with connections to various deadly sorcerers, frothing cultists, and satarre doomsayers, and the two locales have been ruled by a succession of lich lords, minor demon lords, and evil high priests. They are too strong for the local king of the shades to root out, but fortunately they spend most of their time and energy fighting each other rather than honest travelers or the citizens of the city.

\begin{DndTable}[header=Table 7-2: 10 Ghosts of Soriglass]{cX}

d10 & Result\\
1 & Old King Usen. A human lord of a long-forgotten land, who asks visitors to find the bones plundered from his tomb.\\
2 & The Pale Innkeeper. A human woman known for serving fresh ale that turns to blood or salt water.\\
3 & Cursing Marbella. An elven woman who cannot be seen but who swears in Elvish. She seeks vengeance against her killer.\\
4 & The Watchman. A human city guard whose shouting could wake anyone but who can never grab a prisoner. He is looking to find and catch a thief and earn his rest.\\
5 & The Drowned Man. Tangled in kelp and bits of net, he walks the streets, perpetually leaving a damp trail. If given a small sacrifice (spilled drink, bit of smoke or incense), he tells tales of the dark, ebon sea and the riverboats... and of a treasure he buried with his crew, long ago.\\
6 & The Moaning Joker. A gnomish woman who always seeks to fill a ghostly goblet for a toast and who weeps that she can no longer enjoy water, wine, nor mead. Tells terrible jokes.\\
7 & The Empty Princess. An eyeless princess surrounded by a nimbus of burning void energy. Within her gown and sleeves is a tapestry of beetles, flies, and scorpions. If caught in a \textit{mooncatcher's net}, these insects yield a rich purple dye.\\
8 & The Gray Knight. A silent, armored knight with helm and sword who stands watch over a tomb each night. Those who seek to pass into the tomb must defeat him for three nights in a row to lay his spirit to rest.\\
9 & The Gnawing Ogre. A fat, spectral ogre who chews on any meat or bones left in the graveyard. Some say he devours the living as well and can only rest if he eats his fill.\\
10 & Ingortellek, Giant of the Ancients. A red-haired giant who speaks in strangely stilted tones and who seeks to return down the River Styx to a final rest—if only someone would perform the burial rites.\\
\end{DndTable}

\subparagraph{Bright Barge Dock}
While Soriglass has no bridge across the Styx, as does the \textit{Court of Golden Oaks}, it does have a small community of halflings, Styx-loving satarre, and umbral humans who build, load, and work barges up to Vysanth—a pilgrimage town devoted to Charun—and down the river as far as the goblin city of \textit{Fandeval}. The mistress of the Bargeworker's Guild, Happy Fillbec, is a strong, no-nonsense leader of the barge captains and crews, and she brooks absolutely no foolishness when it comes to cargo and payments. Visitors to Soriglass often come by boat rather than by road, and her work securing safe, effective barge traffic on the Styx is one reason why.

\subparagraph{Castle of Frozen Tears}
King Tomarrich the Howler rules a small court of shades, minor phantoms, and necromancers who use the rich array of tombs and ghosts here as raw materials in the performance of dark magic. He dreams of gathering the darakhul to his side, though in practice he is very much a minor vassal to Emperor Vilmos Marquering of the \textit{Twilight Empire}, which is always nearby and in his thoughts. The castle itself is made of finely worked blue-white stone, luminous over the plains or from across the Styx.

\subparagraph{Low Temple of the Ghouls, High Temple of Hecate}
While Charun has several thriving shrines and priests in Soriglass, the most active temple is the Great Vault of Vardesain and Santerr Illosi, where ghoul and darakhul pilgrims come to pay homage to a religious relic: the Bone Cloak of Vardesain, said to be made from the fingerbones of more than 400 darakhul and one bone from an avatar of Vardesain himself. This relic is brought out by the priests on holy days, and darakhul wrapped in it are said to gain relief from the endless hunger and to be cured of various ailments of the undead. Her Luminance Ingalla Vert is the keeper of the cloak while the high priests of Vardesain from both Vysanth and Ossean often make the trek to visit the relic themselves, but neither priest cares to spend much time at the Great Vault.

The Cathedral of the Luminous Moon was once the largest temple to Hecate in the Shadow Realm. With the long decay of Soriglass though, its stained-glass windows are empty sockets in a weathered shell. High Priestess Vert leads other darakhul priestesses and several human adepts and initiates in the rites of the moon, and all are welcome in the black pews and at the full moon blessings of weapons, armor, livestock, and wands. Rumor has it that the chancel is enchanted to ward against evildoers, protecting the altar and its platinum dagger and chalice. The cathedral's real treasures are kept underground and include a lock of Hecate's silvery hair and a bottled piece of her sacred shadow, both relics of the faith. They are brought out for pilgrims to see each full moon and new moon.

\subparagraph{Sixty Cyclopean Tombs}
Outside the city of Soriglass but always moving in its wake when the ground shifts and the land ebbs or flows, the Sixty Tombs are a cluster of mausoleums and pillared shrines with occasional obelisks. Their script it not related to modern Umbral but to an ancient dialect, and the bones found within them are more akin to those of ogres or trollkin than to humans or fey: heavy, tall, and of an imposing scale. Hidden passages, buried tombs, and major undead are always rumored to live among the Sixty Tombs, though few know quite where to find them. Necromancers and priests of Charun often visit to investigate their mysteries.

\begin{DndSidebar}[float=!b]{Requiem, Drug of the Afterlife}
Requiem is a potent drug that allows its users to speak with the dead when smoked. Made from death's head mushrooms, the cremated ashes of sentient beings, and various other unsavory components, it gives its users visions of the dead but is highly addictive. It comes in two forms: a muddy substance called clay and a refined powder called bliss.



Requiem Clay



\textbf{Price per Dose}: 750 gp



When you smoke requiem clay, you summon the spirit of a single deceased person that you personally knew in life. The spirit's image is visible, and its voice can be heard in whispers, but it cannot touch you. You can ask the spirit up to five questions, as the \textit{speak with dead} spell. The spirit knows only what it knew in life, including the languages it knew. Answers are usually brief, cryptic, or repetitive, and the spirit is under no compulsion to offer a truthful answer if you are hostile to it or if it recognizes you as an enemy. The spirit can't learn new information, doesn't comprehend anything that has happened since it died, and can't speculate about future events.



Afterward, you suffer 3 (1d6) poison damage per question asked and must succeed on a Constitution saving throw (DC 10 + 1 per question asked) or become addicted to the drug. If you become addicted, you suffer one level of exhaustion 1d4 days after you last smoked the drug and gain a further level of exhaustion for each week that requiem is not smoked. Levels of exhaustion gained through requiem use do not reduce as normal after taking a long rest. Smoking requiem again eliminates levels of exhaustion gained from abstaining.



At the end of each week, you can make a DC 13 Constitution saving throw. Two consecutive weekly successful saving throws are necessary to break the addiction through abstinence. Alternatively, the addiction can be broken with a \textit{lesser restoration} or \textit{heal} spell.



Requiem Bliss



\textbf{Price per Dose}: 2,500 gp



When you smoke requiem bliss, you experience feelings of euphoria. You summon the spirit of a single deceased person whose name you know. The spirit takes on the physical characteristics it had in life and might touch you, though it cannot harm you. You can ask the spirit questions as the \textit{speak with dead} spell with the following exceptions: you can ask up to ten questions, and the spirit cannot lie to you. In addition, you are simultaneously granted the benefit of the \textit{contact other plane} spell (without taking psychic damage or going insane) as the summoned spirit consults with others and relates information about a single question that would otherwise be unknown to it.



You suffer 3 (1d6) poison damage per question asked and must succeed on a Constitution saving throw (DC 12 + 1 per question asked) or become addicted to the drug.



If you become addicted, you suffer two levels of exhaustion 1d4 days after you last smoked the drug and gain a further level of exhaustion for each week that requiem is not smoked. Levels of exhaustion gained through requiem use do not reduce as normal after taking a long rest. Smoking requiem again eliminates levels of exhaustion gained from abstaining.



At the end of each week, you can make a DC 15 Constitution saving throw. Three consecutive weekly successful saving throws are necessary to break the addiction through abstinence. Alternatively, the addiction can be broken with a \textit{lesser restoration} or \textit{heal} spell.



\end{DndSidebar}

\section{Tenebrous Plain}
A plain of tall, luscious grass downstream from the Forest of the Firebirds stretches out around a tributary of the River Lethe. Gnarled trees and skeletal, leafless shrubs contrast sharply with the serene greenspace. The shadows cast by the plants, both living and dead, fall in strange directions with some stretching farther than others.

A fast-moving tributary flows through the plain, winding its way around clumps of vegetation. This offshoot of the River Lethe is fickle, alternating between flowing with the same inky blackness as most water sources in the Shadow Realm and running as clear and cool as the most refreshing stream in elven lands, based on the whim of \textbf{Oma Rattenfanger} (see \textit{Chapter 9}), the Keeper of the Plain.

It is difficult to say what causes her to show such generosity, but tales tell of countless mortals who make offerings in the Tenebrous Plain, eager for the untainted water. On some occasions, the inky blackness disappears from the riverbed, leaving behind perfectly drinkable water. However, tales also tell of ungrateful young adventurers greedily guzzling the clear water only to have it swirl with black motes while they're drinking it. Some say lost souls who have been mistreated by those they revered are likely to find aid from Oma Rattenfanger, but others say such individuals are never seen again, recruited into her eternal service.

\subsection{Hidden Dangers}
While most creatures of the Shadow Realm avoid Oma Rattenfanger's domain, there are still plenty of dangers awaiting those who travel the Tenebrous Plain. Oma Rattenfanger is the keeper not just of the plain itself but of its inhabitants too, including the \textbf{orphans of the black} (see \textit{Tome of Beasts 2}). These feral, childlike creatures were mistreated youths who seek revenge on those in positions of authority. While orphans of the black are not often loyal creatures, those from the Tenebrous Plain are fiercely devoted to Oma Rattenfanger. Anyone attempting to harm her must contend with these savage souls.

Oma Rattenfanger has granted sanctuary to other creatures seeking shelter from the mortal world or from other portions of the realm. Shadow-touched plants such as \textbf{shadow blights} and \textbf{razorleafs} (see \textit{Creature Codex} for both) find comfort in Oma's domain. Others, such as a grove of weeping treants (see \textit{Tome of Beasts 1}), have been on the Tenebrous Plain as long as Oma and watch over both the keeper and the plain. A single \textbf{feyward tree} (see \textit{Tome of Beasts 1}) roams the plain, protecting the area from fey creatures looking to trespass. It doesn't harm Oma or those in her care though, a mystery the tree's victims do not have a chance to solve.

\begin{DndSidebar}[float=!b]{Pits of the Roachlings}
While many regions of the Shadow Realm resemble the mortal world, some areas adrift in the shadows are much stranger than that. In particular, the Pits of the Roachlings are an entirely bizarre form of terrain made by congealing, spinning, and reforming the raw shadow into what can best be described as demonic termite mounds or perhaps hollow dunes and curved pits. A roachling pit of this kind is usually found on the plains where it appears as a gentle hollow or dip in the landscape with an enticing body of water at the bottom surrounded by reeds or trees. The water is indeed pure rainwater, drinkable and crisp—but the gentle slope around that water is often quite slippery, making it difficult to leave the area around the water. Creatures must make a strong effort to climb back up the slope, requiring a DC 15 Strength (Athletics) check.



At other times, the roachling pit may resemble a set of stones or a small mesa (again, usually with a pure water spring) that sinks into the shadowy plain when approached. Creatures that stay near the water overnight or for a long rest slowly sink with the stones and the pool of water, and on waking, may find themselves 30 or 40 feet below the level of the surroundings. The walls of the pit are commonly pockmarked with niches, tunnels, and small dwelling spaces, often inhabited.



In either case, the nearby ground is riddled with roachling (see \textit{Tome of Beasts 1}) tunnels—and various lesser demons and devils such as \textbf{rattok} (see \textit{Creature Codex}), quasits, imps, and lemures may make their homes there as well. True to their heritage, roachlings can indeed survive almost anywhere, including in the Shadow Realm. Their communities rely on trapping various game (which roachlings can devour without danger of corruption) and on a limited amount of trade with the more civilized races.



The primary foe of the roachlings are darakhul, who find roachling flesh entirely palatable. The Twilight Empire organizes frequent "roach hunts" on the Whispering Plains and elsewhere to feed its expanding population.



\end{DndSidebar}

\section{Twilight Empire}
\begin{DndSidebar}[float=!b]{Twilight Empire}
\textbf{Ruler}: Emperor Vilmos Marquering, the Black Fang (NE \textbf{darakhul general}, see \textit{Chapter 9})



\textbf{Important Personages}: Duchess Angvyr Ssetha the Lady of Chains, Slave Mistress of Chaingard (LE \textbf{darakhul captain}, see \textit{Tome of Beasts 2}); Duchess Mikalea Soulreaper, Lorekeeper of Ossean (NE darakhul necrophage, use \textbf{necromancer}, see \textit{Creature Codex}); Duke Borag the Executioner, Warlord of Gallwheor (NE \textbf{darakhul captain}, see \textit{Tome of Beasts 2}); Duke Eloghar Vorghesht, Regent of Evernight (NE \textbf{darakhul high priest} of Vardesain, see \textit{Creature Codex}); Duke Wierdunn Bonehand, Warden of Blackstone (LE \textbf{darakhul captain}, see \textit{Tome of Beasts 2}); First Captain Ingelmarr Rothand, Steward of Bargol's Watch (LE darakhul \textbf{gladiator}); Her Witchlight Resplendence Tiberessa Vert (LE darakhul \textbf{priest} of Santerr Illosi); His Carmine Eminence and High Patriarch of Zhurahk, Devros Perghallen (LE darakhul \textbf{priest} of Vardesain); Kassimir Valengurd the Confessor, adviser to the Black Iron Throne (NE darakhul necrophage, use \textbf{archmage}); Protector of Souls Sen-Nefer (LN darakhul \textbf{priest} of Anu-Akma)



\textbf{Population}: 65,000 (30,000 slaves [various bearfolk, elves, umbral humans, shadow goblins, shadow fey], 10,000 lesser ghouls and ghasts, 7,000 darakhul, 6,000 imperial ghouls, and 6,000 shadow goblins), plus Zhurakh houses one legion of 6,000 active soldiers (90 percent of which is lesser ghouls and ghasts)



\textbf{Capital}: Zhurakh, City of Bone and Iron, population 15,000 (6,000 slaves, 2,500 lesser ghouls and ghasts and 500 darakhul, plus one legion of 6,000 ghouls)



\textbf{Enclaves}: Chaingard, population 2,300 (1,200 slaves, 870 lesser ghouls and ghasts, 200 imperial ghouls, and 30 darakhul); Gallwheor, population 1,000 (600 slaves, 250 lesser ghouls, 100 iron ghouls, and 50 darakhul); Blackstone, population 750 (450 lesser ghouls and ghasts, 200 slaves, and 100 imperial darakhul); Evernight, population 700 (350 lesser ghouls and ghasts, 200 slaves, 100 darakhul, and 50 priests of Vardesain); Ossean, population 600 (400 slaves, 100 lesser ghouls and ghasts, 50 necrophagi, and 50 darakhul sages and scholars); Bargol's Watch, population 350 (200 slaves, 100 ghouls, and 50 darakhul)



\textbf{Great Gods}: \textit{Vardesain} (patron), \textit{Anu-Akma}, Mavros, \textit{Santerr Illosi}, and \textit{Sarastra}



\textbf{Trade Goods}: Iron, weapons, armor, gemstones, and relics from the City Fallen into Shadow (in order of importance)





\end{DndSidebar}

The Empire of the Ghouls exists in the Crossroads region of Midgard and beyond, tunneling through the mortal world. But the darakhul are also active in the Shadow Realm, and their empire here is, if anything, even more dangerous, as it relies on alliances with shadow dragons and a cave dragon lich, Dreadwing.

Their ability to walk shadow roads and to visit Evermaw, the planar realm of the undead, also makes the Shadow Realm's darakhul exceptionally dangerous. The realm is ruled by an emperor who considers Nicoforus his "imperial cousin," and the two nations seem to exchange ambassadors, trade in goods, and conduct various supporting activities to aid one another as might two somewhat distant nations in any other realm.

The primary difference is that the ghouls in the Shadow Realm have the ability to feed on the strange trees of the Crackling Forests (which means they hunt and devour humanoids somewhat less often)—plus they are fewer in number here compared to other undead realms or even in most human lands, where the undead are a larger presence.

To the north and west of the Blackwood where the shadow fey dwell lies a great pit in the Shadow Realm's landscape. A series of deep depressions gouge the dark earth, diving deep into thick shadows and darkened tunnels. In these Black Iron Depths thrives an exiled arm of the Ghoul Empire known as the Twilight Empire. These darkness-touched ghouls thrive in the shadow of the Ironcrags, searching for a way to breach the barrier between darkness and light. Their undead bodies absorb the essence of shadow through the flesh of their meals (shadow fey and other living creatures of the realm), and every ghoul of the Twilight Empire bears the mark of that shadow—their teeth appear pitch black.

Driven out of the Ghoul Empire after a failed coup, renegade ghouls fled deep into the Underworld. Starving and broken, they found a shadow road in a lightless cavern filled with twisted abominations. The ensuing battle and bloodshed briefly activated the road, and the commander of the traitorous legion led his troops into the Shadow Realm.

The road led the ragged ghouls into the shadows. The ghouls collapsed the tunnel behind them, sealing the shadow road and blocking pursuit. Safe from further reprisal, they set about building a new empire.

\begin{DndSidebar}[float=!b]{Crackling Forests}
Originally found in the Dry Lands of Evermaw, these alien trees are not green and woody but rather cartilaginous masses of splintered bone with leathery leaves, and they exude an air of menace and move rather more vigorously than normal trees (though not enough to fight off a ghoul or woodsman). A rather small number of these trees now grow in the Shadow Realm, imported from Evermaw, but getting more of them to survive and flower is extremely difficult, for the ghouls devour the saplings too quickly to amass a proper forest.



When cut or broken, the trees themselves bleed a black milk that undead find nourishing and pleasant (though living creatures describe it variously as rancid, foul, or rotten). Ghouls, vampires, and other corporeal undead delight in cracking open the hairy bark and gorging themselves on the thick syrup. For such creatures, devouring a pint of black milk is as restorative as devouring a living creature.



\end{DndSidebar}

\subparagraph{Bargol's Watch}
This small, fortified manor house and its outbuildings protect a small village of tanners and bonecutters and serves as the unofficial embassy between the darakhul and the fey—and occasionally between the darakhul and the umbral vampires of \textit{Oshragora}, when that strange city drifts close by. The steward of the watch is a hard case, a former leader of the legions' shock troops who was retired here after he lost his right hand to fey magic: Captain Ingelmarr Rothand despises the shadow fey, gnomes, anything light and cheerful. He keeps the holding in fine condition, marches and drills his small garrison frequently, and generally tolerates no nonsense and encourages no untoward laughter, carousing, or joy of any kind among his troops. He spends most of his time in the Great Tower, which stands 70 feet tall and commands an excellent view of the chasm nearby.

The troops mostly hate him except for his lieutenant, Valmira Brandorf, a warlock who spends most of her time conjuring spirits, imps, and animal spies to gather information from the manor's visitors. Bearfolk merchants, fey road wardens, even gnomish distillers all occasionally visit the house to trade goods with the darakhul or trade information about the position of Oshragora or the movements of elder shadow drakes.

\subparagraph{Blackstone}
Home to the Twilight Empire's shock troops, the Ravenous Legion, Blackstone broods over one of the few roads into the Black Iron Depths, guarded by the gate-fortress Gloomhold. The warden of Blackstone, Duke Wierdunn Bonehand, scours the populace for ghouls strong and cruel enough to join the Ravenous Guard, a group of a hundred imperial ghouls. Blackstone is also home to the Temple of the Faceless Warden, home to the faith of Anu-Akma among the darakhul. Its chief priest is the Protector of Souls Sen-Nefer.

\subparagraph{Chaingard}
The largest settlement on the surface, just on the edge of the Black Iron Depths, Chaingard stands as the foreboding main gateway into the Twilight Empire. This city serves as the receiving point for all slave traffic. Duchess Angvyr Ssetha, the Lady of Chains, oversees the slave supply with meticulous efficiency. She enjoys greeting high-ranking emissaries to the empire in a haughty, sneering manner. The Shackled Legion guards the walls and gates, directing the flow of slaves with brutal efficiency. Despite its title as a legion, it is a motley group of hunters, slavers, and guards more familiar with cruelty than a battlefield.

\subparagraph{Evernight}
Hewn from the black stone and iron ore of the depths, Evernight sits on the edge of a dark lake. Ruled by Duke Eloghar Vorghesht, high priest of Vardesain, Evernight is the center of faith for the empire. A grand cathedral to the ghoul god dominates the city's center, its black altar stained darker with the blood and flesh of countless sacrifices. The Midnight Company of 100 elite darakhul defends the city, boasting warrior clerics and dark reflections of paladins among its ranks.

\subparagraph{Gallwheor}
Whereas Blackstone produces terrifying shock troops, Gallwheor provides the finest tactical minds in the Twilight Empire. Duke Borag the Executioner, Warlord of Gallwheor, never lost a battle during his time in the Ghoul Empire or since their exile into the Shadow Realm. The Headsmen are a group of assassins and stealthy killers known for their deadly efficient strikes to the heart of any opposing force. Rumor has it they have slain bearfolk nobles, fey ladies, and gnomish princes who spoke against the emperor.

\subparagraph{Kaerim's Deep}
Widely believed to be the work of the ghouls, this deep canyon is not known to be inhabited by the darakhul, bearfolk, or others. Occasionally pale giants or cave dragons are seen on its narrow paths and tunnels, but the region is avoided by most.

\subparagraph{Ossean}
The first stop after the steep and perilous Clanking Descent from Chaingard, down the sheer cliffs of the Black Iron Depths, Ossean receives first pick of any new shipment of slaves. Amid the bone-clad buildings, the Twilight Empire's necrophages labor and study under the direction of Duchess Mikalea Soulreaper, Lorekeeper of Ossean. Her apprentices and loyal necrophages hone their practice and ply the secrets from the flesh and souls of captives, guarded by the Bone Company's arcane and martial might. Its ranks swell not just with stout ghouls but also with dark wizards, darakhul necrophages, and eldritch knights who blend arcane magic with the strength of steel.

In addition, the settlement hosts a powerful priestess of Santerr Illosi, Her Witchlight Resplendence Tiberessa Vert. She has for decades built a repository of scrolls, tomes, and grimoires, chronicling the ghoul expansion and the unique fusion of shadow magic with necromancy (see also \textit{Chapter 4}). She has also befriended nearly 100 witchlights, who seem to enjoy her company and sometimes carry messages for her in some form of luminous code. A dozen or more always accompany her like a veil of lights.

\subparagraph{Zhurakh, City of Bone and Iron}
The deepest city of the Black Iron Depths, Zhurakh serves as the seat of the Twilight Emperor and as the unbeating heart of the empire. The Twilight Emperor learned from his failed coup, and his secret police saturate every level of the city, ever watchful for signs of dissent. For this reason, would-be conspirators avoid the City of Bone and Iron completely.

Nearby, the Endless Falls take a seemingly endless plunge into the black depths of the Shadow Realm. This is the Twilight Emperor's favorite place to deal with traitors and criminals, personally casting them into the depths. The most elite twilight soldiers, the Blackmaw Legion, stand ready to enforce the emperor's will.

In addition, Zhurakh is the seat of the Carmine Temple of Vardesain and the seat of power of the red-handed master of the faith, His Carmine Eminence and High Patriarch of Zhurahk, Devros Perghallen. The temple itself is carved of disquieting red stone that seems to ripple and move in torchlight, the walls resembling a bloody abattoir. The place is a source of power for the entire priesthood, carved with sacred writings and with a new testament of corruption, praising the lord of hunger and confirming his blessings on Emperor Vilmos. The temple is the seat and first house of the inquisitors of Vardesain, who enforce the feeding laws and keep a watchful eye on the darakhul and their servants alike.

\subsection{Emperor in Exile}
Vilmos Marquering, called the Black Fang, once commanded the Iron Legion in the White City in the mortal lands. His coup was well planned and brilliantly executed. He turned key conspirators to his cause, several close to the true emperor, and those he couldn't turn he killed and replaced with loyal cat's-paws. When he struck, Vilmos learned just how deep the emperor's secrets ran and how woefully inadequate all those careful preparations truly were. He barely escaped, and the most powerful warriors of his legion were reduced to dust by the emperor's magic. With what must have been Vardesain's intervention, Vilmos and his surviving allies fled to the Shadow Realm.

The Twilight Emperor now rules from his Black Iron Throne in the heart of the Black Iron Depths. Vilmos reforged the battered remnants of his coup into a fighting force and anointed his five most trusted lieutenants as his dukes. Each duke assumed command of a city within the Depths and sought to raise a twilight legion to defend that stronghold in the Twilight Emperor's name. The Black Fang is wary of even these trusted advisers, for he remembers that betrayal arises from those closest to power—and thus he keeps these half-legions weak enough that they cannot challenge him. Further, their commanders' devotion is enforced through magic, by a black-mithral bracer recovered from a strange ruin near the Depths. If any of his dukes turn on him, they'll swiftly find themselves at the points of their own legion's weapons.

\begin{DndSidebar}[float=!b]{The Coming of the Ghouls / from the Book of Ebon Tides}
The ghouls came to the shadows already half-corrupted, and they took to the shaping of the land with uncommon speed and tireless efficiency. Their necrophages quickly mastered the art of fixing land into a permanent form, and their use of rune-marked bone wands to break the forests, fields, and walls of other inhabitants of shadow makes them much-feared. Their numbers remain relatively small compared to the shadow fey, goblins, humans, and bearfolk, but their arcane aggression and wild ambition make them extremely dangerous. Their gods have already carved a bloody home for them from the land, and it seems to threaten all their neighbors.



\end{DndSidebar}

\subsection{Dealings with Darkness}
The Twilight Empire has grown swiftly. They make regular raids on the living: communities of fey, the \textit{Moonlit Glades}, goblin settlements, and any other breathing creatures they can find. The twilight ghouls manage their food sources efficiently, and they rarely suffer a ravenous beggar ghoul to exist among them.

Living in the Shadow Realm means greater difficulty in finding sources of living flesh, so the twilight ghouls forged peaceful contact with the shadow fey courts. Since the shadow fey control the vast majority of shadow roads, they can arrange passage for ghoul raiding parties to the mortal world and ensure their safe return with their precious cargo of "food." In return, the Twilight Emperor assigns legionnaires to help maintain shadow roads that pass through particularly dangerous sections of the realm as well as those known to draw unwanted trespassers. The Queen of Night and Magic found the sudden appearance of a ghoul empire intriguing and amuses herself by keeping tabs on their doings. Lately, she has grown concerned at the increasing ghoul activity in and around the \textit{City Fallen into Shadow}. Their activities somehow elude her most powerful magic and her canniest spies. She contents herself with the strange magical relics the ghouls unearth through trade and favor for the time being, but that won't last long. Her Majesty isn't known for patience once her curiosity is aroused.

Ever since the chance discovery of the Twilight Emperor's magical bracer within a splinter of the City Fallen into Shadow, the Black Fang's ambitions have kindled anew. He hasn't forgotten his humiliating defeat and still dreams of one day returning to the mortal world and seizing power in the Underworld. The umbral vampires and their city suggest a means. The Black Fang charged his necrophages to learn everything they could about the umbral vampires. The cost has been steep, for gleaning secrets from the vampires has lost several a ghoul wizard their sanity, but it has been worthwhile. The Twilight Emperor now knows of an artifact orrery hidden deep within the City Fallen into Shadow that can give its wielder mastery over time and space, said to be hidden in the Storm Tower of Oshragora. He currently lacks the forces to strike into the city, so he sends scouting parties to search for signs of the so-called \textit{Eye of Veles}.

\subsection{Raids to the Mortal World}
The Shadow Realm ghouls thrive in the Black Iron Depths, but that doesn't mean they are all content. Many harbor deep resentment because of their exile, and they channel that sullen anger in different directions. The cadres of lesser ghouls who blame the Twilight Emperor for his failure and for dragging them into exile don't have the stomach to take direct action against their new empire, but these pockets of resistance are starting to organize. The hidden traitors know the road home is a long one, but they're working to recruit allies of higher station and with greater knowledge who share their discontent. If they can find a legion officer or a necrophage wizard who understands what they've lost, they might discover a way to return to the Ghoul Empire.

The loyal ghouls, on the other hand, know that just on the other side of an intangible but nearly impenetrable barrier awaits a nation of flesh and blood unaware of the danger that essentially walks among them. These ghouls hunger for dwarf flesh and have turned their efforts to finding paths through the veil between the Shadow Realm and Midgard, seeking to open the floodgates for the slaves, food, and death waiting in the Ironcrag Cantons and the Crossroads region of the mortal world.

\begin{DndSidebar}[float=!b]{Connections to Evermaw: The Road of Bones}
In recent times, the twilight ghouls have made an effort to forge paths to the Evermaw—a planar realm of the undead ruled by Mot, god of the undead, and Vardesain, god of hunger. As part of this, a black path called the Road of Bones has been laid from the palace in Zhurakh, which requires blood magic to open and leads directly to Evermaw. Darakhul, satarre, vampires, fiends, and others find this an extremely useful and direct method of crossing from plane to plane, and the emperor collects a toll for using the road. It is unclear how the necrophages and priests of Vardesain built it, and the shades of \textit{Soriglass} are mad with curiosity to understand how it was formed. The Road of Bones brings stability, trade, and income to the empire from both sides.



\end{DndSidebar}

\chapter{8. Umbral Pantheon}
The gods of the Shadow Realm are a strange mix of those of the fey and those of the dark and with a few borrowed gods from the bearfolk, from the Nurians, and from the darker divine realms. The gods rarely found elsewhere are those gods of the shadow goblins and the bearfolk, whose worship is less common in the mortal world and whose mysteries in shadow are more powerful and more useful in the Shadow Realm's twilight forests.

Some of them offer access to the Keeper and Shadow domains (see \textit{Chapter 3}).

\section{Great Gods of the Umbral Pantheon}
Hecate, Ailuros, Anu-Akma, Bengta, and Charun are the most widely worshipped and recognized gods of the Shadow Realm. Even if they are not worshipped in a particular settlement or nation, they are known by name, and their many temples, groves, and worshippers outstrip both the lesser godlings and the dark gods.

\begin{DndSidebar}[float=!b]{Gods and Masks}
Similar to the deities of Midgard, the gods of the Shadow Realm are not individuals but rather archetypes or incarnations of universal forces. They aren't people in the same sense that individual mortals are. Though they all arise from the strength of their priests and the cosmos, the same god can look and act entirely differently from one region to another. The beliefs of the faithful similarly vary from place to place, even among those who purport to worship the same deity. Many gods of the Shadow Realm go by multiple names and even switch gender and appearance when it suits them.



As a result, the gods are unknowable and mysterious, and their faiths embody shifting channels of power. Savants believe that only five or six gods might truly thrive in a city, perhaps only three in a town, and one in a village. Shrines devoted to more than a few gods rarely prosper—the priests of these supernumerary gods find that their ability to heal or to intercede with the faithful through divine magic is limited: their power is constrained to a single domain, their prayers go unheard, and their sacrifices are of no consequence. But the heart has room for many gods. Faith in the Shadow Realm is not a matter of choosing a single god but of choosing the right god for a particular need or occasion. In practical terms, this means that pantheist priests (see the \textit{Midgard Heroes Handbook} or the \textit{Southlands Player's Guide}) are sometimes found in the Shadow Realm.



The gods seem determined to hide their true identities from worshippers, and sometimes they pretend to be other gods entirely. This sort of deception is common: Ailuros might be Bastet and Bengta might be Isis. Although malleable and changeable, each of the gods of the Shadow Realm retains a core identity that the god's faithful recognize, even if some of the details might change. The true god, they say, casts different shadows for different worshippers.



\end{DndSidebar}


\begin{itemize}
\begin{multicols}{3}
\item Ailuros
\item Anu-Akma
\item Bengta the Bear Maiden
\item Charun
\item Hecate
\end{multicols}

\end{itemize}

\section{Godlings of the Shadow Realm}
These seven demigods and saints of shadow are not well-known outside the Shadow Realm, and their followers are relatively few compared to major figures such as Charun or Hecate.


\begin{itemize}
\begin{multicols}{3}
\item Baccholon
\item Gytellisor
\item Kupkoresh
\item The Old Roads
\item Santerr Illosi, Our Lady of Blood
\item Santerr Marossa, Our Lady of Roses
\item Umbeserno
\end{multicols}

\end{itemize}

\section{Dark Gods of the Shadow Realm}
The Shadow Realm has more than its share of night terrors, demon lords, and dark gods as well as godlings exiled or banished from brighter realms. A few of them are familiar from the mortal world of Midgard, and a few are similar but not quite the same when seen in the shadows. The five most important dark gods are described here.

\begin{DndSidebar}[float=!b]{The True Dark Gods / from the Book of Ebon Tides}
While the dark gods of the umbral peoples are dangerous, other forces do swim in the shadow matter and hide beneath the hills. The two best known of these are thought to be gods of the Ancients: Nammerkah the Devourer, a sort of tunneling horror thought to be a god of war or the patron of the vampires of Oshragora, and Friselka, the Seal Maiden, whose cold touch is the hand of death and a bringer of new life, sometimes associated with Charun or Lada. Their statues and images are commonly found in tombs of the Ancients. No one knows their prayers or rites, and even their shapes are softened by time.



\end{DndSidebar}


\begin{itemize}
\begin{multicols}{3}
\item Alquam
\item Black Goat of the Woods
\item The Hunter
\item Laughing Loki
\item Vardesain
\end{multicols}

\end{itemize}

\chapter{9. Monsters and NPCs}
The peoples and monsters of the Shadow Realm are numerous and varied. It's difficult to know what to expect. A small sampling of those creatures is included here.

\begin{DndSidebar}[float=!b]{Monsters of Fantasy}
Every fantasy world is filled with a wide variety of antagonists, both monsters and NPCs. Here you'll find new monsters, like the \textit{birch siren} and the \textit{river giant}, along with new NPCs, like the \textit{bearfolk druid} and the \textit{sable elf hierophant}.



Each character and creature has statistics that provide a guideline for a typical encounter, but not every individual of that type need exhibit the traits listed, be it in alignment, traits, actions, or something else. If you want to create monsters and NPCs different from what's listed, go for it!



\end{DndSidebar}


\begin{itemize}
\begin{multicols}{3}
\item \textbf{Bearfolk Druid}
\item \textbf{Bearfolk Thunderstomper}
\item \textbf{Birch Siren}
\item \textbf{Darakhul General}
\item \textbf{Gnomish Distiller}
\item \textbf{Keeper of Hounds}
\item \textbf{Keeper of Ravens}
\item \textbf{Knight of Shadows}
\item \textbf{Lesser Lunar Devil}
\item \textbf{Memory Thief}
\item \textbf{Molefolk}
\item \textbf{Moonshadow Catcher}
\item \textbf{Oma Rattenfanger}
\item \textbf{Radiant Lord}
\item \textbf{River Giant}
\item \textbf{Styx Giant}
\item \textbf{River Spirit}
\item \textbf{Rose Golem}
\item \textbf{Sable Elf}
\item \textbf{Sable Elf Hierophant}
\item \textbf{Shadow Fey Bandit}
\item \textbf{Shadow Goblin Chieftain}
\item \textbf{Shadow Goblin Scribe}
\item \textbf{Shadowspider Swarm}
\item \textbf{Umbral Tailor}
\item \textbf{Wandering Pond}
\end{multicols}

\end{itemize}

\chapter{10. Magic Items and Trickery}
The Shadow Realm is known as a place where magic items, magic tricks, and the peculiar memory philters and illusion seeds are common. Other magical items found in the Shadow Realm are focused on navigating or surviving the dangers of the plane or are derived from the traditions of its goblin, fey, bearfolk, and human inhabitants.

\section{Random Tables for Magic Items}
Items marked with an asterisk (*) are from the \textit{core rules} and are included here because they are found often in the Shadow Realm.

\begin{DndTable}[header=Table 10-1: Apparel]{cX}

d20 & Item\\
1 & \textit{Bearskin cloak}\\
2 & \textit{Boots of dancing}\\
3 & \textit{Boots of elvenkind}*\\
5 & \textit{Boots of shadow walking}\\
4 & \textit{Boots of striding and springing}*\\
6 & \textit{Cloak of elvenkind}*\\
7 & \textit{Cloak of splendor}\\
8 & \textit{Cloak of the bat}*\\
9 & \textit{Cloak of the raven}\\
10 & \textit{Earrings of eclipse}\\
11 & \textit{Emerald goblet}\\
12 & \textit{Gloves of the walking shade}\\
13 & \textit{Glowstone helmet}\\
14 & \textit{Sealskin cloak}\\
15 & \textit{Styx boots}\\
16 & \textit{Trickster's boots}\\
17 & \textit{Whispering cloak}\\
18 & \textit{Winged boots}*\\
19 & \textit{Wolfskin cloak}\\
20 & GM's choice from the list\\
\end{DndTable}

\begin{DndTable}[header=Table 10-2: Armor]{cX}

d20 & Item\\
1 & \textit{Armor of the Black River}\\
2 & \textit{Armor of the white rose}\\
3 & \textit{Ebony beetle armor}\\
4 & \textit{Elven chain}*\\
5 & \textit{Ghost armor}\\
6 & \textit{Glamoured studded leather}*\\
7 & \textit{Midnight armor}\\
8 & \textit{Mithral armor}*\\
9 & \textit{Moonbeam chain mail}\\
10 & \textit{Moonstone banded armor}\\
11 & \textit{Rattling armor}\\
12 & \textit{Road warden shield}\\
13 & \textit{Serpent hide}\\
14 & \textit{Shield of blinding}\\
15 & \textit{Shield of crows}\\
16 & \textit{Shield of dawn's herald}\\
17 & \textit{Shield of otters}\\
18 & \textit{Shield of swans}\\
19 & \textit{Umbral armor}\\
20 & GM's choice from the list\\
\end{DndTable}

\begin{DndTable}[header=Table 10-3: Illusion Seeds]{cX}

d20 & Item\\
1 & \textit{Ale seed}\\
2 & \textit{Bonfire seed}\\
3 & \textit{Bridge seed}\\
4 & \textit{Dragon seed}\\
5 & \textit{Forest seed}\\
6 & \textit{Knightly seed}\\
7 & \textit{Messenger owl seed}\\
8 & \textit{Orchard seed}\\
9 & \textit{Rain and thunder seed}\\
10 & \textit{Road seed}\\
11 & \textit{Stallion seed}\\
12 & \textit{Tavern seed}\\
13–20 & GM's choice from this list\\
\end{DndTable}

\begin{DndTable}[header=Table 10-4: Memory Philters]{cX}

d10 & Item\\
1 & Philter of Boundless Joy\\
2 & Philter of Deep Despair\\
3 & Philter of Enchanted Fortune\\
4 & Philter of Innocence\\
5 & Philter of Loving Kindness\\
6 & Philter of Lust\\
7 & Philter of Mercy\\
8 & Philter of Righteous Anger\\
9 & Philter of Waking Dreams\\
10 & GM's choice from this list\\
\end{DndTable}

\begin{DndTable}[header=Table 10-5: Potions, Scrolls, and Wondrous Items]{cX}

d20 & Item\\
1 & Apparel (see \textit{Table 10-1})\\
2 & \textit{Bag of beans}*\\
3 & \textit{Book of Ebon Tides}\\
4 & \textit{Collapsible mountain}\\
5 & \textit{Crown of infinite midnight}\\
6 & \textit{Deck of illusions}*\\
7 & Figurine of Wondrous Power*\\
8 & Illusion seed (see \textit{Table 10-3})\\
9 & \textit{Hecate's Lantern}\\
10 & \textit{Instant door}\\
11 & \textit{Labyrinth rope}\\
12 & Memory philter (see \textit{Table 10-4})\\
13 & \textit{Moonseeker's lantern of safe returns}\\
14 & \textit{Pipes of haunting}*\\
15 & \textit{Shadow portal scroll}\\
16 & \textit{Unerring dowsing rod}\\
17 & \textit{Wine of the starry void}\\
18 & \textit{Wine of the summer court}\\
19 & \textit{Wine of the winter court}\\
20 & GM's choice from this list\\
\end{DndTable}

\begin{DndTable}[header=Table 10-6: Rings, Rods, Staves, and Wands]{cX}

d20 & Item\\
1 & \textit{Ring of dancing}\\
2 & \textit{Ring of daylight}\\
3 & \textit{Ring of invisibility}*\\
4 & \textit{Ring of mind shielding}*\\
5 & \textit{Ring of shooting stars}*\\
6 & \textit{Ring of the ghouls}\\
7 & \textit{Ring of winter's kiss}\\
8 & \textit{Rod of alertness}*\\
9 & \textit{Staff of charming}*\\
10 & \textit{Staff of frost}*\\
11 & \textit{Staff of swarming insects}*\\
12 & \textit{Staff of withering}*\\
13 & \textit{Staff of radiance}\\
14 & \textit{Staff of the void}\\
15 & \textit{Wand of binding}*\\
16 & \textit{Wand of fear}*\\
17 & \textit{Wand of shadows}\\
18 & \textit{Wand of tides}\\
19 & \textit{Wand of wonder}*\\
20 & GM's choice from the list\\
\end{DndTable}

\begin{DndTable}[header=Table 10-7: Weapons]{cX}

d20 & Item\\
1 & \textit{Blade of grass}\\
2 & \textit{Dagger of venom}*\\
3 & \textit{Dancing sword}*\\
4 & \textit{Fey blade}\\
5 & \textit{Frostglow weapon}\\
6 & \textit{Hecate's dagger}\\
7 & \textit{Lunar lance}\\
8 & \textit{Mace of terror}*\\
9 & \textit{Mooncatcher's net}\\
10 & \textit{Moonsteel weapon}\\
11 & \textit{Nine lives stealer}*\\
12 & \textit{Oathbow}*\\
13 & \textit{Quickstep weapon}\\
14 & \textit{Scepter weapon}\\
15 & \textit{Scimitar of speed}*\\
16 & \textit{Sword of sharpness}*\\
17 & \textit{Sword of wounding}*\\
18 & \textit{Vicious weapon}*\\
19 & \textit{Witchlight scimitar}\\
20 & GM's choice from the list\\
\end{DndTable}

\section{Apparel}
Magical apparel in the Shadow Realm focuses on making a powerful impression and frequently on shapeshifting. These items are described here.


\begin{itemize}
\begin{multicols}{3}
\item \textit{Bearskin Cloak}
\item \textit{Boots of Dancing}
\item \textit{Boots of Shadow Walking}
\item \textit{Cloak of Splendor}
\item \textit{Cloak of the Raven}
\item \textit{Earrings of Eclipse}
\item \textit{Emerald Goblet}
\item \textit{Gloves of the Walking Shade}
\item \textit{Glowstone Helmet}
\item \textit{Sealskin Cloak}
\item \textit{Styx Boots}
\item \textit{Trickster's Boots}
\item \textit{Whispering Cloak}
\item \textit{Wolfskin Cloak}
\end{multicols}

\end{itemize}

\section{Armor}
Some items of magical armor and weaponry were either invented in the Shadow Realm or are found there more commonly than elsewhere.


\begin{itemize}
\begin{multicols}{3}
\item \textit{Armor of the Black River}
\item \textit{Armor of the White Rose}
\item \textit{Ebony Beetle Armor}
\item \textit{Ghost Armor}
\item \textit{Midnight Armor}
\item \textit{Moonbeam Chain Mail}
\item \textit{Moonstone Banded Armor}
\item \textit{Rattling Armor}
\item \textit{Road Warden Shield}
\item \textit{Serpent Hide}
\item \textit{Shield of Blinding}
\item \textit{Shield of Crows}
\item \textit{Shield of Dawn's Herald}
\item \textit{Shield of Otters}
\item \textit{Shield of Swans}
\item \textit{Umbral Armor}
\end{multicols}

\end{itemize}

\section{Illusion Seeds}
Illusion seeds are wondrous items first created in the Shadow Realm and reliant on a supply of shadow matter to power their effects—each such seed generates an illusion and those illusions can take quite powerful forms. While illusion seeds are fairly permanent in the Shadow Realm, where there is shadow matter to empower and sustain the illusions, they are notably less permanent anywhere else. In the mortal world and on other planes, the effects of illusion seeds last for 2d12 hours and then crumble away in a small burst of glittering light and shadows.

To use an illusion seed, its owner merely needs to throw it 5 feet or more from themselves (to avoid being caught in the magical effect). Once expended, their effects are permanent (in the Shadow Realm), though they can be unraveled using \textit{dispel magic} (DC 13).


\begin{itemize}
\begin{multicols}{3}
\item \textit{Ale Seed}
\item \textit{Bonfire Seed}
\item \textit{Bridge Seed}
\item \textit{Dragon Seed}
\item \textit{Forest Seed}
\item \textit{Knightly Seed}
\item \textit{Messenger Owl Seed}
\item \textit{Orchard Seed}
\item \textit{Rain and Thunder Seed}
\item \textit{Road Seed}
\item \textit{Stallion Seed}
\item \textit{Tavern Seed}
\end{multicols}

\end{itemize}

\section{Memory Philters}
First distilled by gnomes in the Court of Midnight Teeth, memory philters have since gained a large following among the fey. Each such philter contains a distilled emotion, specific memories, or a particular moment in time. The sable elves, shadow fey, shadow goblins, and wyrd gnomes all prize these ersatz memories, and the longer-lived fey in particular prize the memories of the young as a way of reliving fresh, youthful emotion. Some of these philters have game effects while others are prized more for their value as status items, entertainment, or simply for the smug sense of having stolen passionate moments from mortals.

All memory philters superficially resemble potions: they are contained in flasks or small vials, often quite beautifully shaped or painted to hint at the contents. Their ability to transfer memories makes them also useful to a GM to dispense clues or story elements to players (as in a memory that reveals a secret door, a plot point, or a murderer's identity).


\begin{itemize}
\begin{multicols}{3}
\item \textit{Philter of Boundless Joy}
\item \textit{Philter of Deep Despair}
\item \textit{Philter of Enchanted Fortune}
\item \textit{Philter of Innocence}
\item \textit{Philter of Loving Kindness}
\item \textit{Philter of Lust}
\item \textit{Philter of Mercy}
\item \textit{Philter of Righteous Anger}
\item \textit{Philter of Waking Dreams}
\end{multicols}

\end{itemize}

\section{Potions, Scrolls, and Wondrous Items}
The items listed here are either created exclusively in the Shadow Realm or simply more common there than elsewhere.


\begin{itemize}
\begin{multicols}{3}
\item \textit{Book of Ebon Tides}
\item \textit{Collapsible Mountain}
\item \textit{Crown of Infinite Midnight}
\item \textit{Hecate's Lantern}
\item \textit{Instant Door}
\item \textit{Labyrinth Rope}
\item \textit{Moonseeker's Lantern of Safe Returns}
\item \textit{Shadow Portal Scroll}
\item \textit{Unerring Dowsing Rod}
\item \textit{Wine of the Starry Void}
\item \textit{Wine of the Summer Court}
\item \textit{Wine of the Winter Court}
\end{multicols}

\end{itemize}

\section{Rings, Rods, Staves, and Wands}
Items with a substantial degree of work invested in their creation—rings, rods, staves, and wands—are very popular as status items among the shadow fey, shadow goblins, and umbral humans, even for those who cannot use their powers.


\begin{itemize}
\begin{multicols}{3}
\item \textit{Ring of Dancing}
\item \textit{Ring of Daylight}
\item \textit{Ring of the Ghouls}
\item \textit{Ring of Winter's Kiss}
\item \textit{Staff of Radiance}
\item \textit{Staff of the Void}
\item \textit{Wand of Shadows}
\item \textit{Wand of Tides}
\end{multicols}

\end{itemize}

\section{Weapons}
These weapons are all found more often in the Shadow Realm than in the mortal world, and many were first created there. Racial weapons of the bearfolk, fey, and ghouls are all included.


\begin{itemize}
\begin{multicols}{3}
\item \textit{Blade of Grass}
\item \textit{Fey Blade}
\item \textit{Frostglow Weapon}
\item \textit{Hecate's Dagger}
\item \textit{Lunar Lance}
\item \textit{Mooncatcher's Net}
\item \textit{Moonsteel Weapon}
\item \textit{Quickstep Weapon}
\item \textit{Scepter Weapon}
\item \textit{Witchlight Scimitar}
\end{multicols}

\end{itemize}

\appendix
\chapter{Appendix: Life in the Shadow Realm}
The Shadow Realm is a grand, strange domain of fey magics and hungry shadows.

\section{Magic of the Fey}
Items listed here are magical but rarely in a way that is useful in combat. They are meant to enhance the user's appearance, provide entertainment, or display status. The owner of one or more such items gains +1 Status while visiting a fey court.

\begin{DndTable}[header=Lesser Umbral Items of the Shadow Fey]{cX}

d20 & Umbral Items\\
1 & Dancing Shadow Amulet. Wearing this amulet makes your shadow dance and caper in amusing ways.\\
2 & Shadow Umbrella. This parasol floats near you, protecting you from direct sunlight or from rain.\\
3 & Floating Hat. This is a normal hat that always floats a handspan above your head. It tips itself courteously to ladies and gentlemen.\\
4 & Animated Goblet. This fine silver goblet floats above the table and will go refill itself from a cask or spring when ordered. It comes when called.\\
5 & Animated Cutlery. This fine silver—a knife, fork, and spoon set—will dance for you, create ringing noise (as before a toast), and even carve a roast or stir soup or tea for you.\\
6 & Glittering Shadow Powder. This powder adds luminous green, blue, or purple sparkles to your shadow.\\
7 & Glimmering Hair Comb. While this comb remains in your locks, your hair shines and reflects light almost as bright as sunlight on water.\\
8 & Flawless Raiment. These clothes are always clean, dry, fashionable, and well-tailored.\\
9 & Purifying Canteen. The water poured from this canteen is always free of shadow corruption.\\
10 & Umbral Incense. This incense draws shadows together in a 5-foot radius and burns for up to 15 minutes. Even in blazing sunlight, the conditions more closely resemble dim light.\\
11 & Dawn Alarm. When the moons rise (or the sun, in the mortal realm or Summer Lands), this pale-white stone glows bright yellow, and a songbird sings.\\
12 & Gilded Shadow. Your shadow is not black and gray but gold and silver. Your Dexterity (Stealth) checks are all at disadvantage.\\
13 & Sunlight Candle. This candle burns with the full light of the sun in a 10-foot radius for 2 hours.\\
14 & Phantasmal Tapestry. A traditional tapestry, often showing a unicorn, hippogriff, dragon, hunting party, or similar scene—but animated and often staring at the viewer. Depending on the subject, rabbits twitch, birds wheel about, waves crash on the shore, and wind sways the woven trees.\\
15 & Perfumed Bath Tent. This tent resembles a bedroll or small rucksack. When shaken, it becomes a 10-by-10-foot tent complete with a wooden tub of hot water, fragrant pine or rose scents, and towels. It is far too thin to provide much protection against wind, cold, or rain, but it makes a fine way to remove the mud of roads and fields.\\
16 & Everblooming Vase. Flowers placed in this vase never fade or wither.\\
17 & Spider Silk Rope. With a sticky end as effective as a grappling hook and a weight of mere ounces, this 50-foot rope assists climbers with foot loops and can fit into a small pocket.\\
18 & Watchful Bat. This small bat statue takes on the coloration of its perch and can watch for a particular person, creature, or event. When that person or event is witnessed, it finds you and tell you about it.\\
19 & Icy Wind Fan. With a cooling breeze, the icy wind fan provides comfort on summer days.\\
20 & GM's Choice from the list.\\
\end{DndTable}

\begin{DndTable}[header=Greater Umbral Items of the Shadow Fey]{cX}

d20 & Umbral Items\\
1 & Moon Helmet. This helmet is made entirely of transparent metal, so your vision is not obstructed.\\
2 & Ebony Scabbard. A scabbard of purest darkness, it helps you gather shadows about yourself. You may attempt to hide while in plain sight once per day.\\
3 & Animated Banner. This banner shows a moving, ferocious form of your heraldic crest, such as a lion rampant or a dragon breathing fire.\\
4 & Umbral Palanquin. This palanquin seats one person and is carried by invisible forces. You need never walk anywhere.\\
5 & Ephemeral Tent. This tent appears fully set when called, complete with brazier of coals, a pair of comfortable beds, and a small table set with nourishing food. It provides limited protection against bitter cold, storm winds, and heavy rain. Poor conditions may shred it.\\
6 & Shadow Reins. These black reins with silver bells are hardly visible, but they are highly effective in transmitting your wishes to a mount. All Wisdom (Animal Handling) checks when dealing with a mount fitted with shadow reins have advantage.\\
7 & Mooncatcher's Delight. This small amulet always draws moonlight to you. You'll never stand in darkness.\\
8 & Phantasmal Feast. A deep-red tablecloth woven of shadows, it creates enough nutritious food for two creatures when placed on a table.\\
9 & Dog Shadow. Your shadow resembles a faithful dog, which is watchful and barks when creatures approach. You cannot be surprised.\\
10 & Sunlight Lantern. This hooded lantern burns with the full light of the sun in a 30-foot radius with dim light for an additional 30 feet. It burns for 6 hours on enchanted oil and can be damped to a 5-foot radius as needed.\\
11 & Umbral Hunting Hounds. Two shadowy hounds that come when called. You have advantage on Wisdom (Survival) checks to follow a trail using scent.\\
12 & Shadow Fanfare. A set of silver trumpets announces your arrival every time you enter a room.\\
13 & Spider Silk Boat. The size of a handkerchief, it becomes a small, two-person rowboat when shaken out. It must be rowed normally but can be returned to kerchief size once you reach shore.\\
14 & Ghost Jar. This jar can hold one ghost or phantom, which answers questions or moans hauntingly from time to time. If a second question is asked on any given day, a successful DC 14 Charisma (Persuasion) check is required, or the ghost departs.\\
15 & Whispering Pebble. A small stone that whispers nonsense without pause. If left alone in a room, it always sounds as if a conversation is happening, but the gist of it is never clear. If addressed, the stone falls silent.\\
16 & Illusory Tower. This gleaming, transparent stone tower contains a single set of spiral stairs up to a height of about 50 feet. At the top, a flat platform and battlements provide an excellent view in all directions. The tower provides no cover or protection but is useful for seeing into the distance.\\
17 & Evergroomed Kit. This pouch is made of oiled leather and holds various grooming products: such as shampoo, hair oil, a comb, a shaving razor, pressed powder, and a hand mirror. On spending 1 minute grooming yourself with the items from this kit, you are clean and fresh, hair perfectly styled, makeup party ready.\\
18 & Feline Ring. This silver band creates an adorable illusory kitten. The kitten appears from behind a chair, someone's leg, or other cover and frolics and tumbles about. If approached, it flees and disappears once it's behind \textit{total cover}.\\
19 & Spectral Kestrel. This translucent bird of prey perches on your shoulder or wrist, keen-eyed and waiting. It can follow simple commands (such as fly, return, fetch, and harry) and has advantage on Wisdom (Perception) checks.\\
20 & GM's Choice from the list.\\
\end{DndTable}

\section{Trickery and Whimsy}
The items in this section are tricks, art, and illusions done for show or to discomfit a rival without injury. These are performed with cantrip-level magic for the most part, though some merely use oils, scents, or ventriloquism. The fey gifts and glamours are charms or knacks sometimes bestowed as boons by fey lords and ladies to their followers or favorites. These gifts can be rescinded from those who fall out of favor.

\begin{DndTable}[header=Tricks and Pranks of the Gnomes and Shadow Goblins]{cX}

d12 & Tricks and Pranks\\
1 & Frosty Frown. Your facial expression is frozen in a scowl for at least 1 hour. A DC 11 Constitution saving throw negates this.\\
2 & Hidden Door. An important chamber's doorway is hidden behind an illusion. Courtiers seem to pass through the wall when they enter or leave this way.\\
3 & Living Statue. A statue speaks and moves, mocking visitors or commenting on their choice of apparel.\\
4 & Irritated Guard. Words are put in a guard's mouth, such as "You're not allowed here" or "Her majesty wishes you to leave the court at once." These are lies, but the guard is either paralyzed or silenced and cannot correct the statement verbally.\\
5 & Sewer Stench of Garlic. A particular set of boots or a cloak is infused with a rancid smell. Magic or several washings are required to get it out.\\
6 & Rotting Garments. Your clothes are made to look moth-eaten, worn, patched, or otherwise slipshod.\\
7 & Sliding Carpet. A carpet is infused with grease or oil that is not visible. Stepping on it requires a DC 14 Dexterity saving throw, or you fall prone (no damage except to dignity). When placed on a staircase, 1d6 falling damage may apply.\\
8 & Enormous Belch. A huge belch seems to come from a genteel fey or human.\\
9 & Stampede! A herd of sheep, goats, oxen, or cattle seems to be rushing down the hallway. Those who stand in its way are not injured, but goblins laugh at anyone who is startled.\\
10 & Tiny Thieves. Trained crows, ravens, or even squirrels take food or small objects (such as buttons, laces, or hair pins) away from a creature and scamper off to the nearest tree or tower with their haul.\\
11 & Storm Cloud. A small rain cloud follows you, drizzling on you (and only you) and leaving you damp.\\
12 & GM's Choice from the list.\\
\end{DndTable}

\begin{DndTable}[header=Gifts and Glamours of the Fey]{cX}

d12 & Gifts and Glamours\\
1 & Scent of Lemons and Roses. You smell of roses and lemons, even in the dead of winter or other unlikely times.\\
2 & Flower Clothing. Your clothes are embroidered with flowers that sway in an unfelt breeze and that blossom, close, and turn to seed in a swift-flowing cycle of a few minutes.\\
3 & Cloud of Kind Lights. You are surrounded by a cloud of wispy lights that always display you to advantage.\\
4 & Deep Eyes like Pools. Your eyes are entrancing. If you succeed with a Wisdom (Insight) check, you can always catch another creature's eyes within 30 feet.\\
5 & Entrancing Voice. Your voice is sweet or resonant; if you succeed with a Charisma (Performance) check, you can always gather the attention of creatures within 30 feet.\\
6 & Artist's Eye. You can describe any creature in detail or sketch it with uncanny accuracy.\\
7 & A Fey Kiss. If a creature permits you to kiss it, your lips leave a lasting memory. Creatures you have kissed always remember you fondly, even if they never see you again.\\
8 & Beguiling Caper. You dance with grace and beauty. At any party or festival where you dance, others stop to watch you.\\
9 & Elegant Speech. Your rhetoric is sublime. You make Charisma (Persuasion) checks with advantage when first meeting a creature you share a language with.\\
10 & Swift Rider. You can spur any mount to extraordinary speed. Your mount's Dash actions always add both your own speed and the mount's speed.\\
11 & Follow Me. When you crook a finger at another creature, it often feels inspired to follow you. This applies as much to hounds and horses as it does to humanoids, fey, and giants. Sometimes, however, the charm fails; if \textit{Counterspell} by an arcanist or resisted by a powerful fey, you may be compelled to move toward the creature that you sought to draw close to yourself.\\
12 & GM's Choice from the list.\\
\end{DndTable}

\begin{DndTable}[header=Beloved Illusions of the Fey Courts]{cX}

d12 & Illusory Effects\\
1 & Enchanted Forest. A great hall or audience chamber resembles a lush forest of fragrant trees, ferns, and similar growth. Trees may hide doorways or balconies. It is difficult to tell what is a pillar and what is a tree.\\
2 & Moonlit Hallway. An indoor hallway is lit by clear beams of moonlight.\\
3 & Starry Carpet. The carpet resembles an open void filled with stars, comets, and moonlight. The carpet is disquieting to walk on.\\
4 & Distant Towers. A set of tall, shining towers is visible from a balcony or window, but it is impossible to visit them. They are illusions near the horizon.\\
5 & Silent Parade. A favorite of gnomes and fey children, a silent parade involves elves, goblins, bearfolk and erina, all marching and singing, shouting, playing trumpets, or pounding drums—but completely silent. Humans tend to find it eerie rather than hilarious.\\
6 & Living Tapestry. Though similar to a phantasmal tapestry (see \textit{Lesser Umbral Items of the Shadow Fey} above), these feature unicorns, hippogriffs, lions, or other animals that can exit the tapestry, wandering within 120 feet of it, fully alive.\\
7 & Speaking Doors. A door with a metal knocker speaks and argues with those seeking to pass through it.\\
8 & Infinite Hallway. The hallway seems to lead to a golden doorway, but walking down it never brings you any closer to the end.\\
9 & Bed Monsters. Growling and hissing can be heard beneath a particular bed. When investigated, the space is empty.\\
10 & Beehives. A section of window seems to be a transparent beehive with the bees themselves always visible.\\
11 & Phantom Orchestra. Light chamber music drifts on the air without any obvious source, as if played by some small group always just out of sight.\\
12 & GM's Choice from the list.\\
\end{DndTable}

\section{Life in Shadows}
The peoples of the Shadow Realm have adapted and thrived in a wide assortment of ways, stamping their own personalities and dreams into the realm.

\begin{DndTable}[header=Fashions of the Courts]{cX}

d12 & Apparel\\
1 & Shadow Silk Cloak. Woven out of darkness itself, a shadow silk cloak is a simple yet impressive accessory. This silky garment floats about you as if in a breeze and evaporates when removed, leaving behind nothing but the clasp until it is donned again.\\
2 & Gossamer Veil. This veil is made from strands of gossamer and starlight, light as air and delicate as a dream.\\
3 & Redcap. Found often in the shadow goblin courts, a redcap is a hat of any style dyed a deep, blood red. Those who hold to tradition dye their cap by soaking it in the fresh blood of a slain foe on a regular basis, but that custom has largely fallen out of favor—though the darakhul have taken up the tradition.\\
4 & Heavens' Weave Clothing. Clothes made from heavens' weave fabric are black or dark blue and shimmer with constellations and star patterns of the mortal world.\\
5 & Glass Gatorskin Accessories. Once tanned and dyed, the translucent skin of the glass gator becomes partially opaque, lending it a fascinating, crystal-like appearance, and it has the added benefit of being completely waterproof. Boots and gloves made of this material are especially popular with the shadow fey of all social classes, though only the nobility can readily afford them.\\
6 & Sun-Stitched Raiment. Created and prized by the bearfolk of the Shadow Realm, these clothes are stitched over with luminous golden thread. The cloth is warm to the touch, as if laid out in the sun, and has a bolstering effect on your emotional state.\\
7 & Moon Sable Fur Mantelet. Commonly considered a more feminine accessory, this short cape is made of the pure-white fur of the moon sable and is worn more for style rather than utility.\\
8 & Living Tattoo. While not technically alive, these intricate tattoos are infused with illusion magic, allowing them to move and shift over your skin, and they can even travel from one part of the body to another based on your whims.\\
9 & Sackcloth Clothing. Resembling prison garb, sackcloth clothing is a sign that you have been exiled from a court or city. Often a punishment in lieu of death, offending parties are stripped of belongings, dressed in sackcloth, and run out of town. Some wear the uncomfortable apparel as a form of atonement for their transgressions.\\
10 & Moonglow Jewelry. The glowing gems in these elaborate accessories are made of moonlight, captured and cut into dazzling jewels.\\
11 & Starlight Coronet. Tiny, shimmering motes of light encircle your head or sit nestled within your hair.\\
12 & GM's choice from this list.\\
\end{DndTable}

\begin{DndTable}[header=Beloved Pets of the Shadows]{cX}

d12 & Strange Creatures\\
1 & Shimmer Serpent. This snake is 3–5 feet long but remains slender and delicate. The serpent is a solid, glossy black except for its belly, which features a shimmering, spangled pattern that resembles a starry night sky.\\
2 & Moon Sable. A sleek, keen-eyed musteline with pure-white fur and the ability to convey simple emotions via empathic projection.\\
3 & Umbraling. This piece of living shadow can be as small as a mouse or as large as a cat, and it takes on the form of whatever animal is most pleasing to its owner, though it has no features beyond a black silhouette.\\
4 & Weal Bug. This jewel-green, cricket-like insect is often carried in a tiny, elaborate cage tied to the owner's belt or scabbard or even carried by a servant. It is said that its chirp draws good fortune to its owner and that it goes silent to warn of ill luck nearby.\\
5 & Wolpertinger. A tiny creature, the wolpertinger (see \textit{Creature Codex}) resembles a fanged rabbit with wings and horns. There are many breeds, from raven-black wolpertingers with a single spiral horn to the more common mottled-brown, deer-antlered variety.\\
6 & Midnight Dragonette. Also known as the twilight dragonling, this tiny, dragon-like creature has an uncanny intelligence and a fondness for shiny trinkets. Its coloration tends toward blacks, greys, purples, and deep blues with the occasional hint of the deep reds of sunset. Its pupils are shaped like starbursts, and it breathes wreaths of purple vapor when happy.\\
7 & Umbral Hunting Hound Puppies. A pair of shadowy puppies that come when called. They wrestle with each other, chase shadows, and are otherwise cute and cuddly.\\
8 & Pocket Pig. This diminutive blue piglet is small enough to fit in a pouch or backpack. It is prized by some households for its ability to sniff out rare culinary ingredients.\\
9 & Flutterby. This feline creature resembles a common housecat with large butterfly wings. Though only able to fly for short distances, it is a prized pet in many noble fey households.\\
10 & Aetherfish. This shimmering fish "swims" in the air around its owner, its long, fan-like fins rippling gently in an unseen current. It occasionally takes on a ghostly appearance or momentarily disappears altogether as it passes too close to the border of the ethereal plane.\\
11 & Lamplark. This diminutive songbird has soft gray-brown feathers and two distinctive black crescents just beneath its eyes. Belying its somewhat lackluster appearance, the lamplark sings beautiful and elaborate songs but only at night. When it sings, the lamplark glows like a golden beacon, shedding bright light in a 15-foot radius.\\
12 & Snapdragon. A curious, sentient plant, the snapdragon has a long stem, broad leaves, and a vivid blossom that comes in a variety of colors. It is content to remain potted and indoors, though it can turn grumpy and brown if it outgrows its pot or if the pot is too plain for its tastes. An unhappy snapdragon has the tendency to bite.\\
\end{DndTable}

\begin{DndTable}[header=Foods of the Shadow Realm]{cX}

d12 & Dishes and Delicacies\\
1 & Twilight Venison Stew. This hearty stew is made from the meat of the deer that wander the Shadow Realm, cooked in a rich, thick gravy with white and purple root vegetables. It is enjoyed by all classes and creatures within the Realm.\\
2 & Air and Darkness Chiffon Cake. This fluffy, airy cake is flavored with a dark, bitter chocolate and topped with powdered sugar and black walnuts.\\
3 & Fried Scuttles. Scuttles are a white beetle found feasting on deadfall in shadowy forests. Once fried, they become tender yet their shells crispy, and they're often served with a side of stinging nettle sauce for dipping. This dish is favored by the shadow goblins, erina, and some alseid.\\
4 & Sautéed Midnight Morels. Conical black mushrooms cooked in butter and topped with shadow salts. Raw, this fungus is toxic, but cooked properly, it takes on a nutty, earthy flavor.\\
5 & Golden Nightjar. Dusk-singing whip-poor-will roasted in a crust of saffron. Each chef adds additional spices to create a custom version of the dish, and the requisite saffron turns the skin and meat of the cooked bird a deep yellow.\\
6 & Dreamingwine. A sweet dessert wine, this deep-red digestif is made with grapes that have been exposed to noble rot and the daydreams of young lovers. A small family of wyrd gnomes have begun developing a new vintage, made with white grapes and the bitter regrets of one dying of old age.\\
7 & Candied Twilight Pears. This purple fruit is sliced thin and simmered in honey. It's served warm with a side of clotted cream. Consuming it conveys a temporary sense of mild euphoria.\\
8 & Stygian Swill. A cheap and heady spirit made by distilling wheat or other grain with waters from the Styx or Lethe. It is highly potent and extremely bitter, and those not native to the Shadow Realm often find they wake the next day with memories that do not belong to them.\\
9 & Whisper Blossom Pods. The long, finger-shaped seed pods of the whisper blossom are a versatile and delicious side dish. They can be battered and fried, sautéed in butter, or even added to stews and soups to add a dash of spice and melancholy.\\
10 & Glass Gator Sweetbreads. The breaded and fried glands and organs of the glass gator are a particular delicacy in many noble households. The transparent nature of the creature and its meat make it difficult to hunt and butcher, making this a costly offering.\\
11 & Steamed Shadowspider Legs. While the legs of the shadowspider matron are black and spiney, the meat within is light and flakey when steamed. This dish is often served with clarified butter seasoned with black truffle.\\
12 & GM's choice from this list.\\
\end{DndTable}

\begin{DndTable}[header=Cantrips of the Fey]{cX}

d12 & Cantrip Effects\\
1 & Eldritch Blast. Your \textit{eldritch blast} briefly leaves behind a lingering trail of illusory flower petals, falling leaves, or butterflies in the color of your choice.\\
2 & Light. Though your \textit{light} spell provides the same illumination as always, it appears as a purple-black glow casting off dark rainbows.\\
3 & Spare the Dying. When you cast \textit{spare the dying}, shadows pour from your mouth and fingers and sink into your target's chest.\\
4 & Resistance. When you cast \textit{resistance}, your target is surrounded by iridescent, sparkling motes that remain until the spell ends.\\
5 & Mage Hand. Your \textit{mage hand} appears as a knobby, twisted hag's hand.\\
6 & Message. When you cast \textit{message}, your voice arrives at the target's ear, discordant and disconcerting, with whispered echoes and strange voices layered among your own.\\
7 & Produce Flame. When you cast \textit{produce flame}, the fire appears as black, blue, pink, or purple.\\
8 & Mending. When you cast \textit{mending}, dozens of tiny, shimmering spectral beetles envelop the damaged area, dissipating into glittering smoke when the spell is complete.\\
9 & Dancing Lights. Your \textit{dancing lights} appear as swarms of blue and purple fireflies, spectral glowing goblin skulls, or tiny illusory moons.\\
10 & Shillelagh. When you cast \textit{shillelagh}, the wood of your staff or club becomes twisted and black and glowing, wicked-looking thorns emerge from it.\\
11 & Guidance. When you cast \textit{guidance}, your target's eyes are briefly framed by glowing purple or smoky black butterfly wings.\\
12 & GM's choice from this list.\\
\end{DndTable}

\begin{DndTable}[header=Omens of the Shadows]{cX}

d12 & Omens\\
1 & Raven Feather. A single black raven feather in your path means death will find you soon.\\
2 & Dancing Stars. Occasionally, a star may reorient itself in the heavens, moving from one point in the sky to another. This indicates shifting allegiances—either one that is already in place or one that should be strongly considered.\\
3 & Backward-Flowing River. If you stumble upon a portion of the Styx or Lethe that is flowing in the wrong direction, the spirits are unquiet or displeased with your current course of action.\\
4 & Rotten Fruit. If you bite into a piece of fruit and find it rotten, someone close to you will soon betray you.\\
5 & Colored Flame. If your candle or fire burns blue, the time is ripe for you to act. If it burns black, beware a blade in the dark. If it burns purple, you'll have a visitor soon. If it burns green, good news is on its way.\\
6 & Cracked Cauldron. A crack appearing in your cauldron or other dishware means you'll soon meet your love.\\
7 & Stryx Shadow. A stryx shadow crossing over you means your current path leads to tragedy.\\
8 & Phantom Knock. If you hear a knocking but find no one there, a known secret must be kept at all costs.\\
9 & Thick Mist. A heavy mist or fog means the spirits are trying to convey a message. It is the ideal time to attempt to contact the dead.\\
10 & Spilled Wine. If someone spills their wine at dinner, they will die within two nights. If you spill someone else's wine, their death will be your fault.\\
11 & Spare Tooth. If you discover an unexpected tooth—in your food, on the ground, or even in your own mouth—good luck is with you.\\
12 & GM's choice from this list.\\
\end{DndTable}

\begin{DndTable}[header=Mounts in the Shadows]{cX}

d12 & Mounts (Large unless noted otherwise)\\
1 & Lesser Unicorn. Also known as a false unicorn, this creature resembles a slender white horse with split hooves and a single shimmering horn on its forehead. It lacks the intelligence and magical abilities of a true unicorn but can use an action to teleport up to 60 feet when in moonlight. Its speed is 50 feet.\\
2 & Valravn. This mount has the body of a wolf with the head and wings of a raven. It has a speed (including flying) of 50 feet.\\
3 & Horned Frog. This black and purple amphibian has two unicorn-like horns that protrude just above its eyes. With a speed (including swimming) of 40 feet, a long jump of 30 feet, and a high jump of 20 feet, this mount is especially advantageous in areas of uneven terrain.\\
4 & Racing Snail. This large snail's spiraling conical shell comes in a variety of hues, and the top breeders pride themselves on creating new and unique color combinations. The racing snail has a speed (including climbing) of 50 feet and can climb difficult surfaces, including upside down on ceilings, without needing to make an ability check.\\
5 & Giant White Rook. This white bird has dark eyes and a keen intelligence. It has a speed of 10 feet and a fly speed of 60 feet. It can follow simple instructions and can mimic short phrases of six words or less.\\
6 & Sylvan Hare. Also known as a sylvie, this Medium rabbit has feather-soft fur—in black, dark gray, or white—red eyes, and overlong ears. It has a speed of 40 feet, a long jump of 30 feet, and a high jump of 20 feet, and it has advantage on Wisdom (Perception) checks related to sound. Halflings in the Shadow Realm prefer sylvies as their mounts.\\
7 & Aquileona. A white and gold lion with a thick mane and broad, eagle-like wings, the aquileona is a fierce and terrifying mount. It has a speed (including flying) of 50 feet.\\
8 & Marbled Widow. A mottled black and white spider, the marbled widow has slender legs, a bulbous abdomen, and paralyzing venom (DC 15 Constitution saving throw or be paralyzed for 1 minute). Though its speed is only 40 feet, it can climb difficult surfaces, including upside down on ceilings, without needing to make an ability check, and it has advantage on Dexterity (Stealth) rolls.\\
9 & Onyx Moa. This Medium flightless bird has jet-black feathers, wicked talons, and a broad, razor-sharp beak. With a speed of 60 feet, it's able to outpace many creatures and is the favored mount of the Shadow Realm's goblins and gnomes.\\
10 & Reitenhund. A sleek black and brown canine, the reitenhund is the ideal mount for hunters and trackers. It moves quickly, with a speed of 60 feet, and has advantage on Wisdom (Survival) checks when tracking prey.\\
11 & Kalong. A fuzzy black bat with a foxlike face, the kalong is awkward and slow on the ground but particularly quick and agile in the air. It has a speed of 10 feet and a fly speed of 60 feet, and it has blindsight to 60 feet.\\
12 & GM's choice from this list.\\
\end{DndTable}

\begin{DndTable}[header=Strange Servants of the Shadows]{cX}

d12 & Servants\\
1 & Knitter. Appearing as a humanoid or fey woman with overlarge solid-black eyes, a knitter has two extra sets of arms that she normally keeps hidden beneath her clothes. Knitters are highly proficient in dexterous tasks, including sewing, mending, weaving, and crochet, but her extra arms come in handy for all sorts of jobs.\\
2 & Shadow Attendant. This shadowy, humanoid form coalesces out of the darkness when called and returns to the shadows when dismissed. It can perform simple tasks, like tending to a fire, cleaning, serving drinks, and carrying written messages.\\
3 & Chatterbox. A 2-foot-tall, white-skinned humanoid with a too-wide mouth and vacant eyes, the chatterbox makes a magical record of any conversation it hears and can repeat it back verbatim later. They often replace scribes for nobles who need to constantly jot down lists or reminders—and the opportunities for spies are endless.\\
4 & Pyreflies. These Tiny insects flit from room to room, kindling fires in hearths and lighting candles as darkness falls.\\
5 & Nanny Goat. This goatlike humanoid has horizontal, rectangular pupils, two small, curved horns, and is quick and agile on their hooved feet. They're often employed as an assistant to nursemaids and governesses to shadow fey children.\\
6 & House Brownie. Tiny and shy, the house brownie is very rarely seen, but the fruits of its labor manifest as mysteriously completed household tasks. A fireplace cleaned, floors scrubbed, and pots and pans sparkling in the drying rack are all evidence of a house brownie. While it only requires payment of milk and honey, those who go out of their way to befriend a house brownie find their kindness repaid in unexpected ways.\\
7 & Skullery Maid. This Small skeleton walks on two legs but does not resemble any of the known humanoids in the Shadow Realm. Its head is a horse's skull, and its eyes burn with white fire. It cleans, mends, and does simple chores in exchange for the ashes from the hearth, which it consumes with gusto.\\
8 & Teacup Fairy. So named because they're small enough to sleep in a teacup, this Tiny, winged fey cultivates tea and other herbs in the kitchen garden. They bite anyone who interferes with their work, gets in their way, or otherwise annoys them.\\
9 & Shadow Watchman. This guard blends neatly with the darkness, standing at its post atop ramparts, parapets, or even outside a chamber door, alert for any encroaching danger. It requires no sleep or food, making it the ideal guard.\\
10 & Umbral Butler. Punctual, tidy, and unfailingly polite, the umbral butler is the ideal head servant. If the other servants resent working beneath a creature made entirely of smoke and shadow, they never show it.\\
11 & Silent Squire. A humanoid with white skin, hair, and eyes, the silent squire is perfectly attentive to a knight's every need, ready with shield, sword, or lance. Just don't look too deeply into their eyes—lest they look back.\\
12 & GM's choice from this list.\\
\end{DndTable}

\begin{DndTable}[header=Secrets of the Courts]{cX}

d12 & NPC Secrets\\
1 & Lightweight. They cannot hold their liquor, and they spill others' secrets when drunk.\\
2 & Dalliance. They are carrying on a secret affair with someone scandalous.\\
3 & Blackmailer. They know someone else's secret and are blackmailing them to keep it quiet.\\
4 & Indebted. They are not nearly as wealthy as they appear to be. They are heavily in debt to unscrupulous moneylenders or are pawning off old family heirlooms in order to afford the newer fashions.\\
5 & Parricide. They killed a family member, either through an arranged killing or by their own hand.\\
6 & Imposter. They are not who they say they are.\\
7 & Conspirator. They are part of a plot to unseat the ruling monarch of their court.\\
8 & Allergy. They are deathly allergic to a common food.\\
9 & Sticky Fingers. They steal things. A lot. It could be small trinkets, large amounts of coin, or even priceless heirlooms.\\
10 & Star-Crossed. They are in love with someone they can never be with.\\
11 & Illusory Appearance. They heavily modify their appearance with illusions and enchantments, far past what is considered socially acceptable.\\
12 & GM's choice from this list.\\
\end{DndTable}

\begin{DndTable}[header=Tresures and Trinkets of the Shadows]{cX}

d12 & Treasures and Trinkets\\
1 & Memento Mori Music Box. This palm-sized music box is made of etched silver. When wound, it plays an eerie and entrancing melody. The inside of the lid holds a lock of black hair set behind glass.\\
2 & Gemstone-Studded Fan. This fan is made from teal ostrich feathers, held together at the base by gilded ribs adorned with onyx cabochons.\\
3 & Bloodletting Brooch. This tiny and extremely sharp dagger is disguised as a fashionable accessory, hidden within the design of a silver and garnet brooch. It can easily be carried into an event where weapons are frowned upon.\\
4 & Pressed Memories. A 3-inch-square frame holds five different colored blossoms pressed between glass. Each blossom is an important memory, removed from the owner's mind and preserved—and promptly forgotten.\\
5 & Everlight. This handheld taper burns with a warm golden glow and never gets any shorter.\\
6 & Book of Secrets. A small journal bound in strange black leather, its pages are filled with fine writing in elegant calligraphy—and it's completely encoded. It holds the secrets of dozens of luminaries and other important beings but requires the cypher to translate, and it resists magical means to bypass the code.\\
7 & Gilded Hummingbird Nest. This tiny accessory is created by harvesting a real hummingbird nest and coating it in gold, silver, or bronze—eggs and all. It is then worn in the hair or hung by a chain around the neck.\\
8 & Kintsugi Teapot. An elegant black stoneware teapot that has been shattered and repaired with gold.\\
9 & Porcelain Unicorn. This tiny, delicate statuette bears an uncanny resemblance to the real thing, having perfectly rendered features.\\
10 & Shadowsphere Pendent. This marble-sized glass sphere holds a tiny, shifting wisp of shadow or black smoke. The wisp is constantly moving and changing.\\
11 & Luminous Shard. This crystalline shard is 1 foot long and sharp as a knife. It glows softly whenever music is played in its vicinity.\\
12 & GM's choice from this list.\\
\end{DndTable}

\begin{DndTable}[header=Lore of the Bearfolk Druids]{cX}

d12 & Lore\\
1 & The Cracked Oak Circle. This bearfolk grove is built around a massive, broken oak tree in the center of their settlement. Green and alive when the bearfolk settled here, the oak was split by dark lightning some decades ago from a darakhul sorcerer locked in battle with the grove's elder. The sorcerer ultimately escaped, the elder died, and the tree never recovered. The bearfolk who train here are educated on the ways of the darakhul and their undead, all the better to destroy them.\\
2 & The Banshee-Haunted Grove. This tiny grove rests on the edge of an ancient battlefield. The site was initially chosen in order to watch for any undead, and when a banshee arose, she was immediately destroyed by the druids. However, she rose again the next night and was put down again—only to rise once more the night after that, and the night after that. After more than a week of this cycle, the bearfolk realized she had no intention of harming anyone within the grove if left alone, and they have instead taken to studying her pathways through the forests and fields in the hopes of discovering her origins and putting her to rest.\\
3 & The Circle of Song. The natural stone henge of the Circle of Song is a curiosity. Carved by the wind whistling through the hillsides, on a gusty day, the stones of the grove sing as air passes through dozens of holes and passageways. Its strange geography also seems to be a focus for magic, and the druids of the Circle of Song can channel their abilities and shape the natural flow of the Shadow Realm around them with ease. They have learned to harmonize their own singing with the songs of the stones, creating a unique form of ritual.\\
4 & The Moonshadow Grove. The druids here are attuned to the shadows and can pass between the pools of darkness with ease. Recently, however, something has infested their grove's shadows, and those who attempt to step between them do not always reappear.\\
5 & The Moonfire Circle. In the center of this druid grove stands a brazier filled with cold-burning fire. The flame gives off a hollow white light and no warmth. It is harmless to living creatures from the mortal world but burns undead and those native to the Shadow Realm. Unfortunately, due to its nature, the druids have been unable to get the flame to grow or spread—it does not ignite wood, paper, or other flammables like true fire does. If they can learn how to manipulate it, it would be a powerful weapon indeed. So far, they have managed to keep its existence a secret from other Shadow Realm denizens, for surely such a tool would never be allowed to remain in bearfolk hands if the shadow fey had any say over it. But it is only a matter of time until word gets out and armies descend upon them.\\
6 & The Grove of the Three Rivers. This grove is built in and around a fork of three small tributaries that branch off from the Styx and empty into three small pools. The bearfolk tend to these shadowy waters, attempting to commune with the souls trapped within. They have met with moderate success and have thus far "rescued" roughly 20 souls, encouraging them to pass on in some cases and sending them off as Stygian shades to discover the manner of their death and the secrets of their life in others. The druids here call themselves shade whisperers and find that violence is not the only answer to the evils of the Shadow Realm.\\
7 & The Bent Tree Grove. The druids of this grove have created an intricate, living housing system created by coaxing the trees to grow in certain designs and directions. Though they are patient when it comes to home construction, the druids are merciless in their training of shadow gnawers and have no hesitation when it comes to cutting out the rot of shadow corruption wherever they find it.\\
8 & The Circle of the Seven Sisters. Established by seven sisters who arrived in the Shadow Realm from the Arbonesse Forest, the Circle of the Seven Sisters has a tenuous connection to the mortal world. A pool of clear, fresh water stands in the center of the grove. The pool is fiercely guarded, both as a source of untainted water and as a pathway between the realms. Creatures who submerge themselves completely in the pool and swim directly toward the bottom for 1 minute emerge in the mortal world in Riverbend within the Arbonesse. Following the pathway back into the Shadow Realm is trickier however since the way through the river is only open during certain phases of the moon.\\
9 & The Hearthollow Grove. Once a grove with many students and teachers, the Hearthollow Grove is now all but abandoned. Only the caretaker, Umalla, remains, and no bearfolk travel here anymore. Recently, even communication with Umalla has gone silent, and none of the bearfolk know the reason for the loss of this stronghold.\\
10 & The Blackwater Grove. Nestled on the banks of a massive lake, the druids of the Blackwater Grove have pulled something curious coffin from the shadowy waters. The box is made of heavy iron and etched with runes, and though it seems old, it shows no signs of rust or other degradation. Still, something is within, and the bearfolk here argue about whether they should open it, guard it, or return it to the lake. So far, Ellar Dannos, the bearfolk leader of this grove, has been content to study it in quiet contemplation for endless hours.\\
11 & The Circle of the Gloomflower. Once a massive grove of gloomflowers, a large swath of forest was cleansed by bearfolk druids who established the Circle of the Gloomflower. Though a latent psychic scream still pierces the air periodically, an echo of a memory of darker times in the land, this grove has known relative peace for some time now. Not long ago, however, the patrolling druids discovered some strange silver seeds that buffeted them with corrupting visions when touched. The gloomflower grove may be returning—either through natural means or by someone's hand.\\
12 & GM's choice from this list.\\
\end{DndTable}

\chapter{Appendix: Shadow Realm Encounters}
You never can quite know what you'll run across in the Shadow Realms. Walk softly.

\section{Forest Encounters}
\begin{DndTable}[header=Encounters in the Forests (CR 1-4)]{cX}

d12 & Encounter\\
1 & The ruins of a crumbling stone palace stand before you. Partially collapsed walls hide whatever treasures may lie within. PCs who indulge their curiosity and enter the ruins are attacked by two cobbleswarms (see \textit{Tome of Beasts 1}), masquerading as the stone floor.\\
2 & A gruesome scene of violence lies in the clearing just ahead. Several corpses picked entirely clean lay on the ground, but one in the clearing center is covered in dozens of black butterflies. The butterflies are a \textbf{death butterfly swarm} (see \textit{Tome of Beasts 1}) that attacks if the corpses are disturbed.\\
3 & Painful groans alert the PCs to a \textbf{sable elf} (see \textit{Chapter 9}) bound to a tree. The \textbf{lunarchidna} (see \textit{Tome of Beasts 2}) that wrapped the elf waits above in the tree to ambush any rescuers.\\
4 & A small fire alerts you to a camp occupied by eight shadow fey forest hunters (see \textit{Tome of Beasts 1}). The fey are happy to share their fire with polite travelers and warn them that an unusual amount of lycanthropes have been seen in this region of the forest.\\
5 & A \textbf{giant elk} bearing a \textbf{far darrig} (see \textit{Tome of Beasts 1}) bounds onto the path. The far darrig is under orders to guard this path and demands the PCs turn back. If the PCs do not comply, both the elk and its rider attack.\\
6 & The PC with the highest passive Perception hears a tiny voice call out from the ground. The voice belongs to an enchanted snail who would like to be placed on a tall branch just above. A character who helps the snail gains the benefits of the \textit{bless} spell for 24 hours. A character who harms or kills the snail has the penalty of a \textit{bane} spell for 24 hours.\\
7 & The PCs step into a grove covered in booby traps left by a tribe of moss lurkers (see \textit{Tome of Beasts 1}).\\
8 & A bedraggled \textbf{veteran} stumbles onto the path and begs the PCs to help end its torment. The veteran slew a bird while hunting for food and has become the target of a vengeful \textbf{pombero} (see \textit{Tome of Beasts 1}).\\
9 & A clearing in the forest admits the shining silver light of the moons above. In the center of the clearing stands a solitary \textbf{gloomflower} (see \textit{Creature Codex}), illuminated by a beam of moonlight.\\
10 & A \textbf{shadow blight} (see \textit{Creature Codex}) has found its way to this part of the forest, searching for dead trees to infect with its taint. It stands motionless next to one such tree, waiting for the PCs to pass.\\
11 & A lumbering bear comes into view. It raises its head, and its deer-like antlers and burning red eyes mark it as more than a normal bear. The \textbf{spawn of Chernobog} (see \textit{Creature Codex}) roars and attacks.\\
12 & A \textbf{wolpertinger} (see \textit{Creature Codex}) sits in the party's path. It hops once, then twice, off the path and into the underbrush. If the party follows, it leads them to their destination in half the time. If the party attacks, it flees.\\
\end{DndTable}

\begin{DndTable}[header=Encounters in the Forests (CR 5-9)]{cX}

d12 & Encounter\\
1 & A trio of deathwisps (see \textit{Tome of Beasts 1}) begins to stalk the party. These wisps are born of three murdered shadow fey who mistakenly believe the PCs are the adventurers responsible for their death.\\
2 & The party is approached by a \textbf{bearfolk druid} (see \textit{Chapter 9}), tracking a creature corrupted by shadow. The druid offers the party information or a safe place to shelter for the night in exchange for assistance in hunting the beast.\\
3 & A chorus of barks heralds the arrival of two hounds of the night (see \textit{Tome of Beasts 1}) followed by a \textbf{keeper of hounds} (see \textit{Chapter 9}). The keeper is training the dogs to hunt humanoids, and they are delighted to sick the dogs on such vulnerable prey.\\
4 & Haunting rustling carries on the wind as two birch sirens (see \textit{Chapter 9}) have flanked the party, hoping to ensnare them with enchantment.\\
5 & The corpse of an enormous bear lies in the clearing ahead, ensnared by a bear trap. The trap is a \textbf{gnarljak} (see \textit{Tome of Beasts 1}), hoping to prey on any creature inspecting the corpse.\\
6 & A chorus of voices leads you to an overgrown black chapel. Inside, four owl harpies (see \textit{Tome of Beasts 1}) sing hymns praising \textbf{Alquam} (see \textit{Tome of Beasts}), the demon lord of night.\\
7 & The PCs are ambushed by a pair of \textbf{malphas} (see \textit{Tome of Beasts 1}), assassins ordered to kill the party at the behest of their mysterious shadow fey lord.\\
8 & The path ahead is filled with black webbing. A successful DC 18 Wisdom (Survival) check reveals these webs aren't from a spider, and if the DC is exceeded by 5 or more, it is identified as webbing from a \textbf{somberweave} (see \textit{Tome of Beasts 2}).\\
9 & The PCs hear sobbing cries for help. If they investigate, they discover a shadow fey trapped in the cage atop the head of a lumbering \textbf{chuhaister} (see \textit{Creature Codex}). The giant entered the Shadow Realm to destroy the fey on their home turf and has been harrying nearby villages.\\
10 & A pile of armor and other belongings can be seen through the trees. The items lie beneath a \textbf{goliath longlegs} (see \textit{Creature Codex}) that attacks with its Paralytic Web when the PCs are within range.\\
11 & An alseid grovekeeper (see \textit{Tome of Beasts 1}) and three \textbf{alseid} (see \textit{Tome of Beasts 1}) circle the party and demand to know why they are in their woods.\\
12 & Three \textbf{myling} approach the party for help in finding their final rest. They lead the PCs to a dense copse of corpse-white trees, where they attempt to drag the PCs into the ground.\\
\end{DndTable}

\begin{DndTable}[header=Encounters in the Forests (CR 10+)]{cX}

d12 & Encounter\\
1 & Profane chanting comes from the grove just beyond where five druids conduct an evil ritual under the guidance of a \textbf{lunar devil} (see \textit{Tome of Beasts 1}). If the devil is destroyed, the corrupting enchantments on the druids are broken, and they are horrified by their actions.\\
2 & A wounded berstuc (see \textit{Tome of Beasts 1}) stumbles onto the path with two treants following closely behind. The berstuc pretends to be a treant, running from the shadow-corrupted treants, and begs the party to help fight off the pursuers. If the PCs fall for this trick and the berstuc survives, it immediately turns its attack on the party.\\
3 & Baying hounds and rallying horn blasts ring through the night as the \textbf{Lord of the Hunt} (see \textit{Tome of Beasts 1}) passes through the wood. Every humanoid that can hear the horn must succeed on a DC 15 Charisma saving throw or be compelled to follow the lord through the night.\\
4 & The party happens upon a \textbf{mavka} (see \textit{Tome of Beasts 1}), setting fire to a huge swath of forest. She gives the party one chance to walk away, no questions asked, but attacks if they interfere.\\
5 & The earth groans as something massive moves far below your feet, sudden tremors rock the ground, and ancient trees fall into a widening chasm in the earth. The PCs must succeed on a DC 18 Strength (Athletics) check to outrun the fissure, or they fall 100 feet into the dark.\\
6 & The sounds of tinkling music and laughter draw you toward a garden party on the lawns of a glittering marble estate. The shadow fey and gnomish nobles in attendance are delighted to introduce inquiring PCs to the host, the \textbf{Mistress of Midnight Teeth} (see \textit{Tome of Beasts 2}).\\
7 & You catch glimpses of dark figures just beyond the tree line. PCs who take a closer look discover the figures are overgrown stone statues that look remarkably like the characters.\\
8 & Tinkling music comes from above as a breeze stirs through the trees. You see hundreds of chimes made from humanoid teeth are hanging from the tall branches.\\
9 & An \textbf{ancient mandriano} (see \textit{Creature Codex}) and three \textbf{mandriano} (see \textit{Creature Codex}) lurk near a forgotten burial ground. When the PCs pass nearby, they attack. When one of the PCs dies or the ancient mandriano drops to 25 hp, it uses its Call the Dead ability to raise the PC's corpse or one of the bodies in the burial ground.\\
10 & A \textbf{bearfolk chieftain} (see \textit{Creature Codex}) and four \textbf{bearfolk} (see \textit{Tome of Beasts 1}) charge into view, pursuing three umbral vampires (see \textit{Tome of Beasts 1}). If the PCs aid the bearfolk in destroying the vampires, he thanks them with gifts of fresh food and water and an open invitation to his grove.\\
11 & A \textbf{ravenfolk doom croaker} (see \textit{Tome of Beasts 1}) and two ravenfolk warriors (see \textit{Tome of Beasts 1}) make their way through the forest. They ask to accompany the PCs for a time and share a campfire before going their own way.\\
12 & A swarm of \textbf{nyctli} (see \textit{Tome of Beasts 2}) stalks the party for as long as it can without being spotted, spying for their mistress. If spotted, the swarm attacks.\\
\end{DndTable}

\section{Marshes and Rivers Encounters}
\begin{DndTable}[header=Encounters in the Marshes and Rivers (CR 1-4)]{cX}

d12 & Encounter\\
1 & A trio of shadow fey bandits (see \textit{Chapter 9}) are camped beside a ramshackle dock on the edge of a river. They offer to ferry the PCs across the river for a hefty sum of coin. If the party refuses, the bandits attack, planning to loot their corpses before pushing them into the river.\\
2 & A group of merchants (commoners) stands arguing in a cluster just in front of a stone bridge, crossing a black river. A \textbf{river spirit} (see \textit{Chapter 9}) has recently taken possession of the bridge and won't allow anyone to pass until a champion defeats it in single combat.\\
3 & The sound of flapping wings fills the air as three \textbf{stryx} (see \textit{Tome of Beasts 1}) land on the rocky shore ahead. The stryx are collecting information for the \textbf{Moonlit King} (see \textit{Tome of Beasts 1} and \textit{Chapter 6}), and they ask the PCs a series of odd questions before flying away.\\
4 & The PC with the highest passive Perception notices that two swarms of ravens have been following the party for the last hour. If the party does nothing, the ravens continue to watch the party while they remain in the area. If the party harms or interferes with the ravens, a \textbf{keeper of ravens} (see \textit{Chapter 9}) appears to dole out justice.\\
5 & One of the PCs has unwittingly caught the eye of a jealous shadow fey noble who sends a sluagh swarm (see \textit{Tome of Beasts 1}) to dispose of them.\\
6 & A ball of gentle blue light appears in the distance, parting the oppressive gloom. It patiently bobs in place as if waiting for you to follow. The light is a friendly \textbf{witchlight} (see \textit{Tome of Beasts 1}) that guides the party safely to their destination if followed.\\
7 & The PCs unwittingly catch the notice of a \textbf{woe siphon} (see \textit{Tome of Beasts 2}) that begins to stalk them. The party begins to experience an increasing series of annoyances and misfortunes until the woe siphon grows strong enough to attack or the PCs catch on to the presence of the creature.\\
8 & A \textbf{green hag} disguised as a lost traveler approaches the party asking for help. If the situation turns hostile, four \textbf{nyctli} (see \textit{Tome of Beasts 2}) spring out of hiding to help their mistress.\\
9 & Three \textbf{shadow skeletons} (see \textit{Creature Codex}) lie motionless beside the black waters of a shadowy river, waiting to ambush travelers.\\
10 & A hungry \textbf{lou caracohl} (see \textit{Creature Codex}) has extended its six tongues in a 60-foot-radius around itself. As the PCs pass near, they must make a DC 13 Dexterity saving throw or become restrained (as per the creature's Sticky Tongues ability). If all the PCs avoid the tongues, it attacks.\\
11 & A \textbf{glass gator} (see \textit{Tome of Beasts 1}) waits at the edge of the water, nearly invisible in the gloom.\\
12 & Three miremals (see \textit{Tome of Beasts 1}) offer to lead the party safely through the swamp but instead lure them deeper into the muck and mire until the PCs are lost and confused, and then they attack.\\
\end{DndTable}

\begin{DndTable}[header=Encounters in the Marshes and Rivers (CR 5-9)]{cX}

d12 & Encounter\\
1 & A dwarf stumbles out of a nearby bank of reeds and professes to have no memory of who they are or what they are doing here. This dwarf is one of many victims of a local \textbf{memory thief} (see \textit{Chapter 9}) who has been stalking victims and staging them to think they've fallen into memory-leeching waters of the realm.\\
2 & After miles of fetid bogs and brackish water, you spot a clear blue pond full of fresh water. The pond is a \textbf{wandering pond} (see \textit{Chapter 9}) disguised as the only drinkable water amidst the marsh.\\
3 & A \textbf{shadow fey enchantress} (see \textit{Tome of Beasts 1}) suddenly appears before you in a puff of smoke. The enchantress was recently banished from court and is interested in recruiting the PCs to help restore her position.\\
4 & Far in the distance, the party sees a flickering purple figure, depositing something on a marshy bank before it disappears. PCs who investigate the spot find a series of stones arranged to spell "Behind You." The \textbf{shadow beast} (see \textit{Tome of Beasts 1}) who left the message materializes to attack.\\
5 & Characters who succeed on a DC 15 Wisdom (Perception) check are alerted to the sound of squelching boots. Characters who fail this check are surprised as a group of three redcaps (see \textit{Tome of Beasts 1}) spring to attack from the reeds.\\
6 & Through the dense mist, you see a hunched figure, rooting through the mud. If the PCs attract any attention to themselves, the \textbf{faceless wanderer} (see \textit{Tome of Beasts 2}) moves to attack.\\
7 & Sticking out of the water, you spot a black \textbf{crystalline monolith} (see \textit{Tome of Beasts 2}) covered in bizarre symbols. Anyone who speaks Void Speech recognizes the symbols as part of a prophecy heralding the return of an elder god.\\
8 & You stumble upon a series of flat rocks laid out with a sumptuous feast. PCs who examine the scene notice a bright-yellow flower sprouted beside the table. The flower is a \textbf{swamp lily} (see \textit{Tome of Beasts 2}), and the feast is its poisonous trap.\\
9 & As the PCs travel near a river, the waters churn and erupt as a \textbf{shadow river lord} (see \textit{Creature Codex}) returns from the mortal world with a freshly reaped soul.\\
10 & A beautiful elven woman steps out of the shadows between trees. Her edges are hazy and indistinct, but her eyes are intense and beguiling. The \textbf{moon nymph} (see \textit{Creature Codex}) smiles and attacks.\\
11 & As the PCs travel through the marshland, flickering ghostlights and shimmering clouds can be seen, always in the distance. A \textbf{foxfire ooze} (see \textit{Creature Codex}) creeps closer to the party, using its natural camouflage to approach unseen.\\
12 & Partially submerged ruins jut from the mud and water. An \textbf{elder shadow drake} (see \textit{Tome of Beasts 1}) does not take kindly to the party intruding into its territory, and it attacks.\\
\end{DndTable}

\section{Plains, Hills, and Mountains Encounters}
\begin{DndTable}[header=Encounters in the Plains, Hills, and Mountains (CR 1-4)]{cX}

d12 & Encounter\\
1 & From its hiding place, a \textbf{domovoi} (see \textit{Tome of Beasts 1}) calls to the party, demanding that they pay a fine for trespassing on private property. The estate that once stood here is long gone, but the domovoi will not be pacified unless it receives something of value as payment.\\
2 & An irate \textbf{molefolk} (see \textit{Chapter 9}) pops out of the ground, demanding to know why the PCs are making such a racket right above their home.\\
3 & The PCs run into two \textbf{darakhul} (see \textit{Tome of Beasts 1}), hunting for fresh meat.\\
4 & A \textbf{shadhavar} (see \textit{Tome of Beasts 1}) has just escaped its cruel shadow fey master and tentatively approaches the party, hoping they can help it get to safety.\\
5 & You come upon a circle of tall stones arranged in a circle. Something sparkles in the grass at the center. The sparkling object in a shiny rock laying atop a \textbf{carnivorous sod} (see \textit{Tome of Beasts 2}) disguised as grass.\\
6 & Tucked amidst the rolling hills is a gated square surrounded by high stone walls. PCs who enter (the gate lock has long since rusted away) find the remains of a monastery garden kept by a lonely \textbf{dancing foliage} (see \textit{Tome of Beasts 2}).\\
7 & Glass marbles are hurled at you from twisted briars just beyond. An \textbf{orphan of the black} (see \textit{Tome of Beasts 2}) creeps through the brambles, hoping the party wants to play.\\
8 & Two shadow fey forest hunters (see \textit{Tome of Beasts 1}) come running down the hillside toward you, pursued by an oozing black \textbf{overshadow} (see \textit{Tome of Beasts 2}).\\
9 & A \textbf{dark father} (see \textit{Creature Codex}) has begun trailing the party. It is not hostile until attacked, and it instead watches them from a distance, though always within view.\\
10 & Six \textbf{bucca} (see \textit{Tome of Beasts 1}) emerge from a nearby cave to attack the passing PCs.\\
11 & A silent, cloaked figure approaches the party through the gloom. It does not answer any shouted inquiries, nor does it remove its hood. The figure is a \textbf{mindrot thrall} (see \textit{Tome of Beasts 1}) that continues toward the group until it is close enough to engulf most of the party with its Acid Breath.\\
12 & A massive, horned beetle rushes toward the party. The beetle is actually a \textbf{skitterhaunt} (see \textit{Tome of Beasts 1}), infesting the carcass of a giant beetle.\\
\end{DndTable}

\begin{DndTable}[header=Encounters in the Plains, Hills, and Mountains (CR 5-9)]{cX}

d12 & Encounter\\
1 & A sudden explosion draws attention to a ramshackle building—and it is engulfed in flames. Two squabbling gnomish distillers (see \textit{Chapter 9}) run outside their laboratory and attack each other if the party does not intervene.\\
2 & A \textbf{knight of shadows} (see \textit{Chapter 9}) has grown bored while patrolling the region and approaches the party for some entertainment (torment).\\
3 & You nearly trip over a human sleeping in the grass. A \textbf{rose golem} (see \textit{Chapter 9}) guards the ruins of a nearby estate and has put several recent passersby to sleep.\\
4 & Raucous laughter and music drift down the hillside from a camp of 13 satyrs. The satyrs are happy to wine and dine the party but intend to slit their throats and rob any PCs who let down their guard.\\
5 & The PCs run into two iron ghouls (see \textit{Tome of Beasts 1}), hunting for fresh meat.\\
6 & On the hill just beyond, you can see the smashed remains of a massive stone gate. PCs who succeed on a DC 15 Intelligence (Arcana) check know that this gate was once a giant portal to the mortal world. Only a creature of immense power could have completely destroyed such a structure.\\
7 & Sharp hissing is the only warning given before the party is swarmed by two umbral vampires (see \textit{Tome of Beasts 1}).\\
8 & Sounds of clattering and pickaxes alert you to a mining operation just over the next ridge. PCs who investigate find a band of darakhul led by a \textbf{darakhul captain} (see \textit{Tome of Beasts 2}), forcing a band of enslaved commoners to dig out a collapsed cavern.\\
9 & A \textbf{shadow fey knight of the road} (see \textit{Creature Codex}), riding a \textbf{lou carcohl} (see \textit{Creature Codex}), blocks the PCs' path and demands payment for passage through her territory. If the PCs refuse, she and her mount attack.\\
10 & A \textbf{moon drake} (see \textit{Creature Codex}) soars overhead. If the party follows it, they lose sight of it at an ancient circle of standing stones.\\
11 & Six \textbf{flesh reavers} (see \textit{Creature Codex}) have managed to get free of their handlers and are hunting the hillsides for living flesh. They converge on the party as soon as they are spotted.\\
12 & A portion of the sod shifts and moves toward the party as a \textbf{corrupting ooze} (see \textit{Tome of Beasts 1}) slinks along the ground to attack.\\
\end{DndTable}

\section{Road Encounters}
\begin{DndTable}[header=Encounters on the Road (CR 1-4)]{cX}

d12 & Encounter\\
1 & A group of commoners led by a \textbf{bearfolk thunderstomper} (see \textit{Chapter 9}) approaches from farther down the road. The bearfolk is escorting this group of settlers to the Moonlit Glades and offers to bring the PCs along.\\
2 & A \textbf{moonshadow catcher} (see \textit{Chapter 9}) and their three \textbf{shadow goblin} minions (see \textit{Creature Codex}) ambush the PCs, attempting to incapacitated them and take them to a local fey looking for test subjects.\\
3 & Hooting and jeering comes from just ahead where a crowd of goblins watches two shadow goblin chieftains (see \textit{Chapter 9}) fight. If the PCs watch the action, the losing chieftain offers the party a fortune to help him kill the winning chieftain. If the party accepts the offer, four \textbf{shadow goblins} (see \textit{Creature Codex}) loyal to the winning chieftain join the fight.\\
4 & A \textbf{sable elf hierophant} (see \textit{Chapter 9}), escorted by two \textbf{sable elf} (see \textit{Chapter 9}) handmaidens, takes an interest in one of the characters. She extends an invitation for the PC to visit her at her patron's mansion.\\
5 & A rowdy pair of shadow fey duelists (see \textit{Tome of Beasts 1}), escorted by a \textbf{shadow fey guardian} (see \textit{Tome of Beasts 1}), waylay any travelers they meet along the road.\\
6 & The road before you is stained with slashes of black-crusted blood. The grisly trail leads to a flipped-over cart that has been torn apart. There are no signs of the cart's owners.\\
7 & A solitary ravenfolk (see \textit{Tome of Beasts 1}) dressed in black robes waves down the party. It politely asks for one of the PCs to sing a song. If the PCs oblige, the ravenfolk listens, whispers, "That's not the one," and hands the singer some gold for their time.\\
8 & A recently exiled \textbf{scrofin} (see \textit{Tome of Beasts 2}) sits drinking on the side of the road. It waves over passersbys, eager to share its bottle and discuss its troubles with fey politics.\\
9 & A piercing shriek fills the air and continues unabated until the PCs investigate or leave the area. If the PCs investigate, they find a young basket-carrying shadow fey unconscious beside an uprooted \textbf{great mandrake} (see \textit{Creature Codex}). The mandrake continues to use its Shriek ability and attacks any creature that come within 10 feet of it.\\
10 & An inebriated \textbf{dhampir} (see \textit{Creature Codex}) stumbles upon the party and begs their help in finding his signet ring, which he lost in a drunken scuffle with a shadow fey.\\
11 & A dirty and scarred book lies abandoned on the side of the road. It might be a spellbook but will require more study to decode. An inkling (see \textit{Creature Codex}) hides within, waiting until it is put in a bag with other books to explore its new surroundings and devour the ink on their pages.\\
12 & A quiet, bookish shadow fey recently passed away, and his \textbf{lantern dragonette} (see \textit{Tome of Beasts 1}) found itself on the streets. It approaches the party, asking for candles it can eat. If they're kind and feed it, it accompanies them for a time.\\
\end{DndTable}

\begin{DndTable}[header=Encounters on the Road (CR 5-9)]{cX}

d12 & Encounter\\
1 & The imposing figure of a \textbf{monolith champion} (see \textit{Tome of Beasts 1}) greets you as it paces in front of a shrine just off the road.\\
2 & A \textbf{fardorocha} (see \textit{Tome of Beasts 2}) approaches the party, bearing an invitation to a nearby fey estate.\\
3 & A \textbf{fey revenant} (see \textit{Tome of Beasts 2}) has been given leave to hunt humanoids as a birthday present from the Queen of Night and Magic. The PCs are selected as the revenant's first quarry.\\
4 & Just ahead, villagers crowd around a hastily constructed stage. The people look on in horror as a \textbf{jack of strings} (see \textit{Tome of Beasts 2}) uses their magic to force a young man to beat himself.\\
5 & A carriage with blackened windows pulls up alongside the PCs. Its occupants are two \textbf{lunarians} (see \textit{Tome of Beasts 2}), hoping to get directions to the closest \textbf{umbral tailor} (see \textit{Chapter 9}). The lunarians hope the tailor can cure them of their curse, but they welcome any interested PCs to join them on their journey.\\
6 & A gargantuan wall of thorns completely blocks the road. The PCs must leave the safety of the road to go around the thorns or confront the \textbf{thornheart guardian} (see \textit{Tome of Beasts 2}) that has created the blockage.\\
7 & After a long stretch of travel through the cold fog, you see a stone guardhouse built along the roadside. It appears to be empty and to be a perfect place to rest. If the PCs take a rest in the guardhouse, they are attacked mid-sleep by one of the walls, which is a \textbf{walled horror} (see \textit{Tome of Beasts 2}).\\
8 & Thunderous hoofbeats come from behind as two knights of shadow (see \textit{Chapter 9}) ride to destroy a trio of \textbf{werynax} (see \textit{Tome of Beasts 2}), causing havoc on the road up ahead. PCs who assist the knights in their hunt have the opportunity to garner favor with the courts.\\
9 & A \textbf{karakura} (see \textit{Creature Codex}) lurks in the shadows just off the road. As the PCs pass by, it uses its Shroud of Darkness ability on the last PC in the group. Since the PCs are already in the Shadow Realm, if the PC fails their Charisma saving throw, they are instead engulfed in an invulnerable cocoon of shadow until the effect ends.\\
10 & A \textbf{mirror hag} (see \textit{Tome of Beasts 1}) has begun plaguing a nearby shadow fey town, delighting in causing chaos among the self-absorbed fey nobles. The hag disguises herself as an elderly fey and targets the PC with the highest Charisma score with her Reconfiguring Curse.\\
11 & A miserable-looking dog approaches the party. The dog is a \textbf{hulking whelp} (see \textit{Tome of Beasts 1}) in search of friendship. If the PCs treat it kindly, it accompanies them as long as they allow it. If they harm or are otherwise cruel to it, it uses its Unleashed Emotion ability to grow in size and then attacks.\\
12 & A \textbf{kikimora} (see \textit{Tome of Beasts 1}) has recently been forcibly evicted from the home it was plaguing. It spots the party and begs them to help it find new lodgings, where it will cause misery and destruction.\\
\end{DndTable}

\begin{DndTable}[header=Encounters on the Road (CR 10+)]{cX}

d12 & Encounter\\
1 & The roads buckle beneath your feet as a flickering, gigantic figure phases into view. The figure is an irate \textbf{shadow giant} (see \textit{Tome of Beasts 2}), looking for any excuse to take out its frustration.\\
2 & The mists part to reveal an abandoned village just off the road. The villagers fell prey to an \textbf{incarnate gloom} (see \textit{Tome of Beasts 2}), eager to make the PCs its next meal.\\
3 & A disguised \textbf{luck leech} (see \textit{Tome of Beasts 2}) approaches the PCs, challenging them to a "friendly" game of chance. If the luck leech begins to lose, it chooses to attack the victor and rob them.\\
4 & A gleaming \textbf{repository} (see \textit{Tome of Beasts 2}) emerges from the mists, bearing an invitation for the PCs to dine with its \textbf{psychic vampire} (see \textit{Tome of Beasts 2}) master. If the party accepts, the repository leads them to the vampire's nearby subterranean estate.\\
5 & A team of deep gnomes has pulled off the side of the road to repair their broken \textbf{steam golem} (see \textit{Tome of Beasts 1}). The golem is a gift to the \textbf{Queen of Night and Magic} (see \textit{Tome of Beasts 1}), and the gnomes are desperate for help repairing it.\\
6 & Screaming comes from just ahead as a squad of four \textbf{darakhul} (see \textit{Tome of Beasts 1}) are being chased by a hungry \textbf{gug} (see \textit{Tome of Beasts 1}).\\
7 & Trails of smoke drift toward you from a burned-down village. Anyone who approaches sees a mountain of bodies blazing at its center. The village is actually an illusory haunt, and the bodies are a \textbf{corpse mound} (see \textit{Tome of Beasts 1}), which attacks anyone who comes close.\\
8 & A being of shining golden light suddenly falls from the sky and lands before you. The \textbf{Radiant Lord} (see \textit{Chapter 9}) introduces themselves and politely asks if the PCs have seen any locations in their travels that would be fit for constructing an observation tower.\\
9 & Hoofbeats can be heard as a \textbf{vampire patrician} (see \textit{Creature Codex}), his three \textbf{vampire spawn} escorts, and their mounts draw near to the PCs on the road. The vampire patrician will stop and exchange pleasantries with the party, for he's interested in their business in the Shadow Realm. The vampires are not hostile and only attack if attacked first.\\
10 & An enraged \textbf{shadow fey ambassador} (see \textit{Creature Codex}) and her \textbf{dream squire} (see \textit{Creature Codex}) stand at the side of the road, robbed of their goods. The ambassador hails the PCs and demands they pursue the three \textbf{shadow fey knights of the road} (see \textit{Creature Codex}) and retrieve their stolen belongings.\\
11 & A \textbf{dullahan} (see \textit{Tome of Beasts 1}) gallops down the road, preparing to cut down the party and any who stand in their way.\\
12 & A \textbf{fear smith} (see \textit{Tome of Beasts 1}), accompanied by a \textbf{shadow fey guardian} (see \textit{Tome of Beasts 1}), confronts the characters in a dark alley or lonely road.\\
\end{DndTable}

\chapter{Appendix: Shadow Realm Creatures by Terrain}
A non-exhaustive list of possible creatures that you may encounter in the Shadow Realm is provided here, organized by terrain and challenge rating and with options provided from the core rules and from \textit{Tome of Beasts 1} (TOB1), \textit{Creature Codex} (CC), and \textit{Tome of Beasts 2} (TOB2)—and from this volume, \textit{Book of Ebon Tides} (see \textit{Chapter 9}).

\section{Forests}

\begin{DndTable}{cX}

CR & Creature\\
0 & \textbf{Awakened Shrub}\\
0 & \textbf{Commoner}\\
0 & \textbf{Deer}\\
0 & \textbf{Owl}\\
1/8 & \textbf{Bandit}\\
1/8 & \textbf{Giant Weasel}\\
1/8 & \textbf{Kobold}\\
1/8 & Leonino (CC)\\
1/8 & \textbf{Mastiff}\\
1/8 & \textbf{Poisonous Snake}\\
1/8 & \textbf{Stirge}\\
1/8 & \textbf{Stryx} (TOB1)\\
1/4 & \textbf{Blink Dog}\\
1/4 & \textbf{Boar}\\
1/4 & \textbf{Elk}\\
1/4 & Erina Scrounger (TOB1)\\
1/4 & Forest Falcon (TOB2)\\
1/4 & \textbf{Giant Badger}\\
1/4 & \textbf{Giant Bat}\\
1/4 & \textbf{Giant Owl}\\
1/4 & \textbf{Giant Poisonous Snake}\\
1/4 & \textbf{Giant Wolf Spider}\\
1/4 & \textbf{Goblin}\\
1/4 & Grove Bear (TOB2)\\
1/4 & Nyctli (TOB2)\\
1/4 & \textbf{Pseudodragon}\\
1/4 & Shadow Fey (Elves) (TOB1)\\
1/4 & Sootwing Giant Moth (CC)\\
1/4 & \textbf{Sprite}\\
1/4 & \textbf{Suturefly} (TOB1)\\
1/4 & \textbf{Swarm of Ravens}\\
1/4 & \textbf{Treacle} (TOB1)\\
1/4 & \textbf{Witchlight} (TOB1)\\
1/4 & \textbf{Wolf}\\
1/4 & Wolpertinger (CC)\\
1/2 & \textbf{Alseid} (TOB1)\\
1/2 & \textbf{Black Bear}\\
1/2 & Foxin (CC)\\
1/2 & \textbf{Giant Wasp}\\
1/2 & \textbf{Hobgoblin}\\
1/2 & \textbf{Lantern Dragonette} (TOB1)\\
1/2 & \textbf{Ravenfolk Scout} (TOB1)\\
1/2 & \textbf{Sable Elf} (see \textbf{Chapter 9})\\
1/2 & \textbf{Satyr}\\
1/2 & \textbf{Swarm of Insects}\\
1/2 & Woodwose (CC)\\
1/2 & \textbf{Worg}\\
1 & Albino Death Weasel (CC)\\
1 & \textbf{Brown Bear}\\
1 & \textbf{Bugbear}\\
1 & Bulbous Violet (TOB2)\\
1 & Carbuncle (CC)\\
1 & \textbf{Child of the Briar} (TOB1)\\
1 & \textbf{Crimson Drake} (TOB1)\\
1 & \textbf{Dire Wolf}\\
1 & \textbf{Dogmole} (TOB1)\\
1 & \textbf{Dryad}\\
1 & \textbf{Erina Defender} (TOB1)\\
1 & \textbf{Giant Spider}\\
1 & \textbf{Giant Toad}\\
1 & \textbf{Harpy}\\
1 & \textbf{Leshy} (TOB1)\\
1 & \textbf{Molefolk} (see \textbf{Chapter 9})\\
1 & \textbf{Moss Lurker} (TOB1)\\
1 & Serpentine Lamia (CC)\\
1 & Shadow Blight (CC)\\
1 & \textbf{Shadow Fey Bandit} (see \textbf{Chapter 9})\\
1 & Shockwing Giant Moth (CC)\\
1 & Wirbeln Fungi (CC)\\
2 & \textbf{Ankheg}\\
2 & \textbf{Awakened Tree}\\
2 & \textbf{Bandit Captain}\\
2 & \textbf{Centaur}\\
2 & \textbf{Ettercap}\\
2 & Forest Drake (CC)\\
2 & \textbf{Giant Boar}\\
2 & \textbf{Giant Elk}\\
2 & Giant Vampire Bat (CC)\\
2 & \textbf{Green Dragon Wyrmling}\\
2 & \textbf{Grick}\\
2 & Lunarchidna, Lesser (TOB2)\\
2 & \textbf{Ogre}\\
2 & \textbf{Pegasus}\\
2 & Rageipede (CC)\\
2 & \textbf{Shadhavar} (TOB1)\\
2 & \textbf{Shadow Goblin Chieftain} (see \textbf{Chapter 9})\\
2 & Shadow Skeleton (CC)\\
2 & \textbf{Swarm of Poisonous Snakes}\\
2 & Wererat (Lycanthrope)\\
2 & White Stag (CC)\\
2 & \textbf{Will-o'-wisp}\\
2 & Wind Weasel (CC)\\
3 & Alseid Grovekeeper (TOB1)\\
3 & \textbf{Amphiptere} (TOB1)\\
3 & \textbf{Bagiennik} (TOB1)\\
3 & \textbf{Basilisk}\\
3 & \textbf{Bearfolk} (TOB1)\\
3 & \textbf{Duskthorn Dryad} (TOB1)\\
3 & Fang of the Great Wolf (Worg) (CC)\\
3 & \textbf{Far Darrig} (TOB1)\\
3 & Giant Albino Bat (CC)\\
3 & \textbf{Goat-Man} (TOB1)\\
3 & \textbf{Gold Dragon Wyrmling}\\
3 & \textbf{Green Hag}\\
3 & Korrigan (CC)\\
3 & \textbf{Millitaur} (TOB1)\\
3 & \textbf{Moonshadow Catcher} (see \textbf{Chapter 9})\\
3 & \textbf{Owlbear}\\
3 & Peluda Drake (CC)\\
3 & \textbf{Phase Spider}\\
3 & \textbf{Pombero} (TOB1)\\
3 & Prismatic Beetle Swarm (TOB1)\\
3 & \textbf{Ravenfolk Warrior} (TOB1)\\
3 & \textbf{Veteran}\\
3 & Vine Golem (TOB2)\\
3 & Werewolf (Lycanthrope)\\
4 & Amber Ooze (TOB2)\\
4 & Azeban (CC)\\
4 & Dancing Foliage (TOB2)\\
4 & \textbf{Death Butterfly Swarm} (TOB1)\\
4 & \textbf{Deathcap Myconid} (TOB1)\\
4 & \textbf{Domovoi} (TOB1)\\
4 & \textbf{Forest Marauder} (TOB1)\\
4 & Lunarchidna, Greater (TOB2)\\
4 & Overshadow (TOB2)\\
4 & \textbf{Sap Demon} (TOB1)\\
4 & Wereboar (Lycanthrope)\\
5 & \textbf{Asanbosam} (TOB1)\\
5 & Centaur Chieftain (CC)\\
5 & \textbf{Gorgon}\\
5 & \textbf{Lindwurm} (TOB1)\\
5 & Mandriano (CC)\\
5 & Moon Drake (CC)\\
5 & \textbf{Owl Harpy} (TOB1)\\
5 & Quickstep (CC)\\
5 & \textbf{Ravenala} (TOB1)\\
5 & \textbf{Ravenfolk Doom Croaker} (TOB1)\\
5 & Razorleaf (CC)\\
5 & \textbf{Sable Elf Hierophant} (see \textbf{Chapter 9})\\
5 & Serpentine Lamia Matriarch (CC)\\
5 & \textbf{Shadow Fey Forest Hunter} (TOB1)\\
5 & \textbf{Shambling Mound}\\
5 & \textbf{Troll}\\
5 & \textbf{Unicorn}\\
5 & \textbf{Vila} (TOB1)\\
5 & Werebear (Lycanthrope)\\
5 & Xiphus (CC)\\
5 & \textbf{Zmey Headling} (TOB1)\\
6 & \textbf{Bearfolk Druid} (see \textbf{Chapter 9})\\
6 & Bloom Hydra (TOB2)\\
6 & Child of Yggdrasil (CC)\\
6 & Death Butterfly Swarm, Greater (TOB1)\\
6 & Green Knight of the Woods (CC)\\
6 & Gulon (CC)\\
6 & \textbf{Likho} (TOB1)\\
6 & Piasa (CC)\\
6 & \textbf{Redcap} (TOB1)\\
6 & Transcendent Lunarchidna (TOB2)\\
6 & \textbf{Weeping Treant} (TOB1)\\
6 & Wolf Spirit Swarm (TOB1)\\
7 & Chuhaister (CC)\\
7 & Goliath Longlegs (CC)\\
7 & \textbf{Shadow Beast} (TOB1)\\
7 & Shadow Drake, Elder (TOB1)\\
7 & \textbf{Shadow Fey Enchantress} (TOB1)\\
7 & Ulnorya (TOB2)\\
8 & Arborcyte (CC)\\
8 & \textbf{Dragonleaf Tree} (TOB1)\\
8 & Green Dragon, Young\\
8 & Mandriano, Ancient (CC)\\
8 & \textbf{Savager} (TOB1)\\
8 & Vines of Nemthyr (CC)\\
9 & \textbf{Bukavac} (TOB1)\\
9 & Flame-Scourged Scion (CC)\\
9 & \textbf{Ghostwalk Spider} (TOB1)\\
9 & \textbf{Knight of Shadows} (see \textbf{Chapter 9})\\
9 & Kulmking (CC)\\
9 & Shadow Fey Ambassador (Elves) (CC)\\
9 & Shadow River Lord (CC)\\
9 & \textbf{Treant}\\
9 & Wickerman (CC)\\
9 & Wraith Bear (TOB2)\\
10 & \textbf{Birch Siren} (see \textbf{Chapter 9})\\
10 & Chameleon Hydra (TOB2)\\
10 & Colláis (TOB2)\\
10 & Foxfire Ooze (CC)\\
10 & Gold Dragon, Young\\
10 & Nyctli Swarm (TOB2)\\
11 & \textbf{Ychen Bannog} (TOB1)\\
12 & Derendian Moth Abomination (TOB2)\\
12 & Bear King (Fey Lord) (TOB1)\\
12 & Reynard (Fey Lord) (CC)\\
12 & Horned Serpent (CC)\\
14 & \textbf{Zmey} (TOB1)\\
15 & Green Dragon, Adult\\
16 & Berchta (Fey Lady) (CC)\\
\end{DndTable}

\section{Grasslands and Plains}

\begin{DndTable}{cX}

CR & Creature\\
0 & \textbf{Commoner}\\
0 & \textbf{Deer}\\
0 & \textbf{Eagle}\\
0 & \textbf{Goat}\\
0 & \textbf{Vulture}\\
1/8 & \textbf{Flying Snake}\\
1/8 & \textbf{Giant Weasel}\\
1/8 & \textbf{Kobold}\\
1/8 & \textbf{Poisonous Snake}\\
1/8 & \textbf{Pony}\\
1/8 & \textbf{Stirge}\\
1/4 & Alliumite (CC)\\
1/4 & \textbf{Boar}\\
1/4 & \textbf{Elk}\\
1/4 & Erina Scrounger (TOB1)\\
1/4 & \textbf{Giant Poisonous Snake}\\
1/4 & \textbf{Giant Wolf Spider}\\
1/4 & \textbf{Goblin}\\
1/4 & \textbf{Riding Horse}\\
1/4 & \textbf{Sable Elf} (see \textbf{Chapter 9})\\
1/4 & Shadow Fey (Elves) (TOB1)\\
1/4 & Sootwing Giant Moth (CC)\\
1/4 & \textbf{Witchlight} (TOB1)\\
1/4 & \textbf{Wolf}\\
1/2 & \textbf{Burrowling} (TOB1)\\
1/2 & \textbf{Cockatrice}\\
1/2 & \textbf{Giant Goat}\\
1/2 & \textbf{Giant Wasp}\\
1/2 & \textbf{Hobgoblin}\\
1/2 & \textbf{Swarm of Insects}\\
1/2 & \textbf{Worg}\\
1 & \textbf{Bugbear}\\
1 & \textbf{Erina Defender} (TOB1)\\
1 & \textbf{Giant Eagle}\\
1 & \textbf{Giant Vulture}\\
1 & \textbf{Hippogriff}\\
1 & Shockwing Giant Moth (CC)\\
1 & \textbf{Wind Dragon Wyrmling} (TOB1)\\
2 & \textbf{Ankheg}\\
2 & \textbf{Centaur}\\
2 & \textbf{Eala} (TOB1)\\
2 & \textbf{Giant Boar}\\
2 & \textbf{Giant Elk}\\
2 & Giant Vampire Bat (CC)\\
2 & \textbf{Griffon}\\
2 & \textbf{Ogre}\\
2 & \textbf{Pegasus}\\
2 & Roggenwolf (CC)\\
2 & Shadhavar (CC)\\
2 & \textbf{Shadow Goblin Chieftain} (see \textbf{Chapter 9})\\
3 & Darakhul (Ghoul) (TOB1)\\
3 & Fang of the Great Wolf (Worg) (CC)\\
3 & \textbf{Far Darrig} (TOB1)\\
3 & Ghost Boar (CC)\\
3 & Giant Albino Bat (CC)\\
3 & \textbf{Manticore}\\
3 & \textbf{Phase Spider}\\
3 & \textbf{Veteran}\\
4 & \textbf{Death Butterfly Swarm} (TOB1)\\
4 & \textbf{Frostveil} (TOB1)\\
4 & \textbf{Mammoth}\\
4 & \textbf{Mngwa} (TOB1)\\
4 & Overshadow (TOB2)\\
4 & \textbf{Serpopard} (TOB1)\\
4 & \textbf{Tusked Skyfish} (TOB1)\\
4 & Wereboar (Lycanthrope)\\
5 & \textbf{Bulette}\\
5 & Centaur Chieftain (CC)\\
5 & \textbf{Gorgon}\\
5 & \textbf{Hound of the Night} (TOB1)\\
5 & Sleipnir (CC)\\
5 & Three-Headed Cobra (CC)\\
5 & Tusked Crimson Ogre (CC)\\
5 & \textbf{Vapor Lynx} (TOB1)\\
5 & \textbf{Wandering Pond} (see \textbf{Chapter 9})\\
6 & \textbf{Bearfolk Druid} (see \textbf{Chapter 9})\\
6 & \textbf{Chimera}\\
6 & Death Butterfly Swarm, Greater (TOB1)\\
6 & Dracotaur (CC)\\
6 & Ghost Boar, Elder (CC)\\
6 & Werynax (TOB2)\\
6 & Wind Dragon, Young (TOB1)\\
6 & Wolf Spirit Swarm (TOB1)\\
6 & Zouyu (TOB2)\\
7 & Giant Sloth (CC)\\
8 & Arborcyte (CC)\\
8 & \textbf{Degenerate Titan} (TOB1)\\
8 & Neophron Demon (CC)\\
8 & Oliphaunt (CC)\\
8 & \textbf{Savager} (TOB1)\\
8 & Vines of Nemthyr (CC)\\
9 & Fulminar (CC)\\
9 & Kulmking (CC)\\
10 & Foxfire Ooze (CC)\\
10 & Gold Dragon, Young\\
12 & Horned Serpent (CC)\\
13 & Bonespitter (TOB2)\\
13 & Noth-Norren (TOB2)\\
17 & Gold Dragon, Adult\\
17 & Wind Dragon, Adult (TOB1)\\
20+ & Jotun Giant (TOB1)\\
\end{DndTable}

\section{Hills and Mountains}

\begin{DndTable}{cX}

CR & Creature\\
0 & \textbf{Eagle}\\
0 & \textbf{Goat}\\
0 & \textbf{Raven}\\
0 & \textbf{Vulture}\\
1/8 & \textbf{Bandit}\\
1/8 & \textbf{Blood Hawk}\\
1/8 & Giant Moth (CC)\\
1/8 & \textbf{Giant Weasel}\\
1/8 & \textbf{Guard}\\
1/8 & \textbf{Mule}\\
1/8 & \textbf{Shroud} (TOB1)\\
1/8 & \textbf{Stirge}\\
1/4 & \textbf{Azza Gremlin} (TOB1)\\
1/4 & \textbf{Boar}\\
1/4 & \textbf{Elk}\\
1/4 & Erina Scrounger (TOB1)\\
1/4 & \textbf{Giant Owl}\\
1/4 & \textbf{Giant Wolf Spider}\\
1/4 & \textbf{Goblin}\\
1/4 & Living Shade (CC)\\
1/4 & \textbf{Panther}\\
1/4 & \textbf{Pseudodragon}\\
1/4 & Shadow Fey (Elves) (TOB1)\\
1/4 & Skull Lantern (CC)\\
1/4 & Snow Cat (CC)\\
1/4 & Sootwing Giant Moth (CC)\\
1/4 & \textbf{Swarm of Bats}\\
1/4 & \textbf{Swarm of Ravens}\\
1/4 & \textbf{Witchlight} (TOB1)\\
1/4 & \textbf{Wolf}\\
1/2 & Alkonost (CC)\\
1/2 & \textbf{Bucca} (TOB1)\\
1/2 & \textbf{Fraughashar} (TOB1)\\
1/2 & \textbf{Giant Goat}\\
1/2 & \textbf{Hobgoblin}\\
1/2 & \textbf{Orc}\\
1/2 & \textbf{Ravenfolk Scout} (TOB1)\\
1/2 & \textbf{Skin bat} (TOB1)\\
1/2 & Storm Spirit (CC)\\
1/2 & Void Cultist (CC)\\
1/2 & \textbf{Worg}\\
1 & \textbf{Brown Bear}\\
1 & \textbf{Copper Dragon Wyrmling}\\
1 & Dark Servant (Dark Folk) (CC)\\
1 & \textbf{Dire Wolf}\\
1 & \textbf{Dogmole} (TOB1)\\
1 & \textbf{Erina Defender} (TOB1)\\
1 & Eye of the Gods Angel (CC)\\
1 & \textbf{Giant Eagle}\\
1 & \textbf{Harpy}\\
1 & \textbf{Hippogriff}\\
1 & Mountain Strider (TOB2)\\
1 & Phantom (CC)\\
1 & Serpentine Lamia (CC)\\
2 & \textbf{Beli} (TOB1)\\
2 & \textbf{Eala} (TOB1)\\
2 & Fear Liath (CC)\\
2 & \textbf{Giant Boar}\\
2 & \textbf{Giant Elk}\\
2 & Giant Vampire Bat (CC)\\
2 & \textbf{Griffon}\\
2 & \textbf{Kot Bayun} (TOB1)\\
2 & Lady in White (CC)\\
2 & Light Dragon Wyrmling (CC)\\
2 & \textbf{Noctiny} (TOB1)\\
2 & \textbf{Ogre}\\
2 & \textbf{Pegasus}\\
2 & \textbf{Peryton}\\
2 & \textbf{Saber-Toothed Tiger}\\
2 & Scitalis (CC)\\
2 & \textbf{Shadhavar} (TOB1)\\
2 & \textbf{Silver Dragon Wyrmling}\\
2 & Trollkin Grunt (CC)\\
2 & Vampire Thrall (CC)\\
2 & \textbf{Vile Barber} (TOB1)\\
2 & \textbf{Void Dragon Wyrmling} (TOB1)\\
2 & Wind Demon (CC)\\
2 & Wind Eater (CC)\\
3 & Abbanith Giant (TOB2)\\
3 & \textbf{Basilisk}\\
3 & Cave Drake (TOB2)\\
3 & Dark Eye (Dark Folk) (CC)\\
3 & \textbf{Elvish Veteran Archer} (TOB1)\\
3 & \textbf{Flame Dragon Wyrmling} (TOB1)\\
3 & Giant Albino Bat (CC)\\
3 & Gloomflower (CC)\\
3 & \textbf{Goat-Man} (TOB1)\\
3 & \textbf{Green Hag}\\
3 & \textbf{Hell Hound}\\
3 & \textbf{Manticore}\\
3 & \textbf{Phase Spider}\\
3 & \textbf{Ravenfolk Warrior} (TOB1)\\
3 & Skull Drake (CC)\\
3 & Sluagh Swarm (TOB1)\\
3 & \textbf{Spire Walker} (TOB1)\\
3 & Thursir Giant (TOB1)\\
3 & \textbf{Veteran}\\
3 & Werebat (Lycanthrope) (CC)\\
3 & Werewolf (Lycanthrope)\\
4 & \textbf{Ash Drake} (TOB1)\\
4 & \textbf{Bereginyas} (TOB1)\\
4 & \textbf{Carrion Beetle} (TOB1)\\
4 & \textbf{Ettin}\\
4 & \textbf{Firebird} (TOB1)\\
4 & \textbf{Flab Giant} (TOB1)\\
4 & \textbf{Frostveil} (TOB1)\\
4 & \textbf{Lich Hound} (TOB1)\\
4 & Overshadow (TOB2)\\
4 & \textbf{Red Dragon Wyrmling}\\
4 & Shadow Fey Guardian (Elves) (TOB1)\\
4 & \textbf{Trollkin Reaver} (TOB1)\\
4 & Trollkin Shaman (CC)\\
4 & Unhatched (CC)\\
4 & \textbf{Wereboar}\\
4 & Wind's Harp Devil (CC)\\
5 & \textbf{Air Elemental}\\
5 & \textbf{Aridni} (TOB1)\\
5 & \textbf{Black Knight Commander} (TOB1)\\
5 & \textbf{Bulette}\\
5 & \textbf{Corrupting Ooze} (TOB1)\\
5 & Dark Voice (Dark Folk) (CC)\\
5 & \textbf{Dogmole Juggernaut} (TOB1)\\
5 & Dream Wraith (CC)\\
5 & \textbf{Fellforged} (TOB1)\\
5 & \textbf{Gorgon}\\
5 & \textbf{Hill Giant}\\
5 & \textbf{Hound of the Night} (TOB1)\\
5 & \textbf{Hulking Whelp} (TOB1)\\
5 & Lamia Serpentine Matriarch (CC)\\
5 & \textbf{Lindwurm} (TOB1)\\
5 & \textbf{Ravenfolk Doom Croaker} (TOB1)\\
5 & \textbf{Sandman} (TOB1)\\
5 & Shadow Fey Forest Hunter (Elves) (TOB1)\\
5 & Shadow Fey Knight of the Road (Elves) (CC)\\
5 & Sleipnir (CC)\\
5 & \textbf{Troll}\\
5 & Werebear (Lycanthrope)\\
5 & Wyvern Knight (CC)\\
6 & Bearfolk Chieftain (CC)\\
6 & \textbf{Chimera}\\
6 & Corrupted Ogre Chieftain (TOB1)\\
6 & Crimson Mist (CC)\\
6 & Elf Deathsworn (CC)\\
6 & \textbf{Fate Eater}\\
6 & \textbf{Fext} (TOB1)\\
6 & Fey Drake (CC)\\
6 & \textbf{Ghost Knight} (TOB1)\\
6 & Haunted Giant (CC)\\
6 & \textbf{Malphas} (TOB1)\\
6 & \textbf{Mirror Hag} (TOB1)\\
6 & Mithral Dragon, Young (TOB1)\\
6 & Mountain Nymph (TOB2)\\
6 & Necromancer (CC)\\
6 & \textbf{Nichny} (TOB1)\\
6 & Pixiu (CC)\\
6 & \textbf{Redcap} (TOB1)\\
6 & \textbf{Rotting Wind} (TOB1)\\
6 & \textbf{Scheznyki} (TOB1)\\
6 & Shadow Fey Duelist (Elves) (TOB1)\\
6 & Wind Dragon, Young (TOB1)\\
6 & \textbf{Wyvern}\\
7 & Alabaster Tree (Planar Flora) (CC)\\
7 & Blood Mage (CC)\\
7 & Copper Dragon, Young\\
7 & \textbf{Deathwisp} (TOB1)\\
7 & Dhampir Commander (CC)\\
7 & Elf Enchanter (CC)\\
7 & Fire Dancer Swarm (TOB1)\\
7 & Giant Sloth (CC)\\
7 & \textbf{Red Hag} (TOB1)\\
7 & Sarsaok (TOB2)\\
7 & Shadow Fey Enchantress (Elves) (TOB1)\\
7 & \textbf{Spark} (TOB1)\\
7 & \textbf{Stone Giant}\\
7 & Two-Headed Eagle (CC)\\
7 & \textbf{Umbral Vampire} (TOB1)\\
7 & Vellso Demon (CC)\\
7 & Void Speaker (CC)\\
7 & War Wyvern (CC)\\
8 & \textbf{Ala} (TOB1)\\
8 & Ankou Soul Seeker (CC)\\
8 & \textbf{Blemmyes} (TOB1)\\
8 & Blood Giant (CC)\\
8 & \textbf{Bone Collective} (TOB1)\\
8 & \textbf{Chained Angel} (TOB1)\\
8 & \textbf{Frost Giant}\\
8 & Hoard Drake (TOB2)\\
8 & \textbf{Liosalfar} (TOB1)\\
8 & Nachzehrer (CC)\\
8 & \textbf{Qwyllion} (TOB1)\\
8 & Voidwracked Mage (CC)\\
9 & \textbf{Cloud Giant}\\
9 & \textbf{Fire Giant}\\
9 & Flame Dragon, Young (TOB1)\\
9 & Lich Hierophant (CC)\\
9 & Monarch Skeleton (CC)\\
9 & Silver Dragon, Young\\
9 & \textbf{Vine Troll Skeleton} (TOB1)\\
9 & Void Dragon, Young (TOB1)\\
10 & \textbf{Bone Swarm} (TOB1)\\
10 & Elf Alchemist Archer (CC)\\
10 & \textbf{Fear Smith} (TOB1)\\
10 & Grave Behemoth (CC)\\
10 & Nightgaunt (CC)\\
10 & Red Dragon, Young\\
11 & \textbf{Abominable Beauty} (TOB1)\\
11 & \textbf{Blood Hag} (TOB1)\\
11 & \textbf{Corpse Mound} (TOB1)\\
11 & \textbf{Dullahan} (TOB1)\\
11 & Ghost Dragon (CC)\\
11 & \textbf{Grim Jester} (TOB1)\\
11 & Mountain Dryad (TOB2)\\
11 & \textbf{Roc}\\
11 & Vampiric Knight (CC)\\
11 & Void Giant (CC)\\
11 & \textbf{Voidling} (TOB1)\\
11 & Wendigo (CC)\\
11 & \textbf{Ychen Bannog} (TOB1)\\
12 & \textbf{Ancient Titan} (TOB1)\\
12 & Bear King (Fey Lord) (TOB1)\\
12 & \textbf{Flutterflesh} (TOB1)\\
12 & \textbf{Skein Witch} (TOB1)\\
13 & \textbf{Mask Wight} (TOB1)\\
13 & \textbf{Storm Giant}\\
13 & \textbf{Stuhac} (TOB1)\\
13 & Vampire Warlock (Variant) (TOB)\\
14 & Cherufe (TOB2)\\
14 & Copper Dragon, Adult\\
14 & Void Dragon, Adult (TOB1)\\
15 & Commander of Horn and Gold (Fey Lord) (TOB2)\\
15 & Drake, Star (TOB1)\\
15 & \textbf{Purple Worm}\\
15 & \textbf{Slow Storm} (TOB1)\\
16 & Silver Dragon, Adult\\
16 & Snow Queen (Fey Lady) (TOB1)\\
17 & Moonlit King (Fey Lord) (TOB1)\\
17 & Queen of Witches (Fey Lady) (TOB1)\\
17 & Red Dragon, Adult\\
17 & Wind Dragon, Adult (TOB1)\\
CR 18 & Lord of the Hunt (Fey Lord) (TOB1)\\
\end{DndTable}

\section{Marshes}

\begin{DndTable}{cX}

CR & Creature\\
1/8 & \textbf{Giant Rat}\\
1/8 & \textbf{Kobold}\\
1/8 & Lone Sluagh (Sluagh Swarm) (TOB1)\\
1/8 & \textbf{Poisonous Snake}\\
1/8 & \textbf{Shroud} (TOB1)\\
1/8 & \textbf{Stirge}\\
1/4 & \textbf{Constrictor Snake}\\
1/4 & Corpse Thief (CC)\\
1/4 & Exploding Toad (CC)\\
1/4 & \textbf{Giant Frog}\\
1/4 & \textbf{Giant Lizard}\\
1/4 & \textbf{Giant Poisonous Snake}\\
1/4 & Living Shade (CC)\\
1/4 & \textbf{Mud Mephit}\\
1/4 & Otterfolk (TOB2)\\
1/4 & Resinous Frog (TOB2)\\
1/4 & Shadow Fey (Elves) (TOB1)\\
1/4 & Skull Lantern (CC)\\
1/4 & Swamp Adder (Snake) (TOB1)\\
1/4 & \textbf{Swarm of Rats}\\
1/4 & \textbf{Swarm of Ravens}\\
1/4 & \textbf{Witchlight} (TOB1)\\
1/2 & Cackling Skeleton (TOB2)\\
1/2 & Corrupted Pixie (TOB2)\\
1/2 & \textbf{Crocodile}\\
1/2 & Execrable Shrub (Planar Flora) (CC)\\
1/2 & Flesh Reaver (CC)\\
1/2 & \textbf{Fraughashar} (TOB1)\\
1/2 & Hallowed Reed (Planar Flora) (CC)\\
1/2 & Husk (TOB2)\\
1/2 & \textbf{Lizardfolk}\\
1/2 & \textbf{Orc}\\
1/2 & Shadow Goblin (CC)\\
1/2 & \textbf{Skin bat} (TOB1)\\
1/2 & \textbf{Swarm of Insects}\\
1/2 & Void Cultist (CC)\\
1 & Anophiloi (CC)\\
1 & \textbf{Boloti} (TOB1)\\
1 & Dark Servant (Dark Folk) (CC)\\
1 & \textbf{Ghoul}\\
1 & \textbf{Giant Spider}\\
1 & \textbf{Giant Toad}\\
1 & \textbf{Glass Gator} (TOB1)\\
1 & Phantom (CC)\\
1 & Sporous Crab (TOB2)\\
1 & Spurred Water Skate (TOB2)\\
1 & Trollkin Raider (TOB2)\\
1 & Wirbeln Fungi (CC)\\
2 & \textbf{Beli} (TOB1)\\
2 & \textbf{Black Dragon Wyrmling}\\
2 & Blood Zombie (CC)\\
2 & Dragon Skeleton (TOB2)\\
2 & \textbf{Eala} (TOB1)\\
2 & Egret Harpy (TOB2)\\
2 & \textbf{Gargoyle}\\
2 & \textbf{Ghast}\\
2 & \textbf{Giant Constrictor Snake}\\
2 & Lady in White (CC)\\
2 & Leech Swarm (TOB2)\\
2 & Light Dragon, Wyrmling (CC)\\
2 & Lou Carcolh\\
2 & \textbf{Moonshadow Catcher} (see \textbf{Chapter 9})\\
2 & \textbf{Ogre}\\
2 & \textbf{Putrid Haunt} (TOB1)\\
2 & Rotsam (TOB2)\\
2 & Scrag, Lesser (CC)\\
2 & \textbf{Shadhavar} (TOB1)\\
2 & \textbf{Shadow Goblin Chieftain} (see \textbf{Chapter 9})\\
2 & Shadow Skeleton (CC)\\
2 & Swampgas Bubble (TOB2)\\
2 & \textbf{Swarm of Poisonous Snakes}\\
2 & Trollkin Grunt (CC)\\
2 & Vampire Thrall (CC)\\
2 & \textbf{Vile Barber} (TOB1)\\
2 & \textbf{Will-o'-wisp}\\
2 & Wind Eater (CC)\\
2 & Yumerai (TOB2)\\
3 & \textbf{Bagiennik} (TOB1)\\
3 & Bloodsapper (TOB2)\\
3 & Carrier Mosquito (TOB2)\\
3 & Dark Eye (Dark Folk) (CC)\\
3 & Drake, Skull (CC)\\
3 & Gloomflower (CC)\\
3 & \textbf{Green Hag}\\
3 & Lazavik (TOB2)\\
3 & Mold Zombie (CC)\\
3 & Purple Slime (CC)\\
3 & Sluagh Swarm (TOB1)\\
3 & Swordbreaker Skeleton (TOB2)\\
3 & Viiret (TOB2)\\
3 & \textbf{Wight}\\
4 & \textbf{Bandit Lord} (TOB1)\\
4 & Battle Mage (CC)\\
4 & Brachyura Shambler (TOB2)\\
4 & \textbf{Firebird} (TOB1)\\
4 & \textbf{Lich Hound} (TOB1)\\
4 & Lunar Devil, Lesser (see \textbf{Chapter 9})\\
4 & Mari Lwyd (TOB2)\\
4 & Quickserpent (TOB2)\\
4 & Rattok Demon (CC)\\
4 & \textbf{Serpopard} (TOB1)\\
4 & Servant of the Unsated God (TOB2)\\
4 & Shadow Fey Guardian (Elves) (TOB1)\\
4 & \textbf{Trollkin Reaver} (TOB1)\\
4 & Trollkin Shaman (CC)\\
4 & Unhatched (CC)\\
4 & Venom Elemental (CC)\\
5 & \textbf{Aridni} (TOB1)\\
5 & Befouled Weird (TOB2)\\
5 & \textbf{Bouda} (TOB1)\\
5 & Cipactli (Demon) (CC)\\
5 & \textbf{Corrupting Ooze} (TOB1)\\
5 & Dark Voice (Dark Folk) (CC)\\
5 & Eleinomae (TOB1)\\
5 & \textbf{Fellforged} (TOB1)\\
5 & Fleshdreg (TOB2)\\
5 & \textbf{Giant Crocodile}\\
5 & Lord Zombie (CC)\\
5 & \textbf{Otyugh}\\
5 & Scarlet Ibis (TOB2)\\
5 & Shadow Fey Forest Hunter (Elves) (TOB1)\\
5 & Shadow Fey Knight of the Road (Elves) (CC)\\
5 & \textbf{Shambling Mound}\\
5 & Somberweave (TOB2)\\
5 & \textbf{Troll}\\
5 & \textbf{Wandering Pond} (see \textbf{Chapter 9})\\
5 & \textbf{Water Elemental}\\
5 & Willowhaunt (TOB2)\\
5 & \textbf{Zmey headling} (TOB1)\\
6 & Crimson Mist (CC)\\
6 & Doomspeaker (CC)\\
6 & Drake, Fey (CC)\\
6 & \textbf{Fext} (TOB1)\\
6 & \textbf{Ghost Knight} (TOB1)\\
6 & \textbf{Malphas} (TOB1)\\
6 & \textbf{Mirror Hag} (TOB1)\\
6 & Moonchild Naga (CC)\\
6 & Necromancer (CC)\\
6 & \textbf{Nichny} (TOB1)\\
6 & Ogrepede (TOB2)\\
6 & Primal Oozer (TOB2)\\
6 & \textbf{River Giant} (see \textbf{Chapter 9})\\
6 & Rotsam Swarm (TOB2)\\
6 & \textbf{Rotting Wind} (TOB1)\\
6 & Shadow Fey Duelist (Elves) (TOB1)\\
6 & \textbf{Spectral Guardian} (TOB1)\\
6 & Swamp Naga (TOB2)\\
6 & Tar Ooze (TOB2)\\
7 & Alabaster Tree (Planar Flora) (CC)\\
7 & Black Dragon, Young\\
7 & Blood Mage (CC)\\
7 & \textbf{Deathwisp} (TOB1)\\
7 & Dhampir Commander (CC)\\
7 & Fire Dancers Swarm (TOB1)\\
7 & Graveyard Dragon (TOB2)\\
7 & \textbf{Herald of Darkness} (TOB1)\\
7 & Lambent Witchfyre (TOB2)\\
7 & Mangrove Treant (TOB2)\\
7 & \textbf{Red Hag} (TOB1)\\
7 & \textbf{Risen Reaver} (TOB1)\\
7 & Scrag, Greater (CC)\\
7 & Shadow Fey Enchantress (Elves) (TOB1)\\
7 & \textbf{Spark} (TOB1)\\
7 & Swamp Lily (TOB2)\\
7 & \textbf{Umbral Vampire} (TOB1)\\
7 & Void Speaker (CC)\\
8 & Blood Giant (CC)\\
8 & \textbf{Hydra}\\
8 & \textbf{Liosalfar} (TOB1)\\
8 & \textbf{Lunar Devil} (TOB1)\\
8 & Marsh Dire (TOB2)\\
8 & Nachzehrer (CC)\\
8 & \textbf{Qwyllion} (TOB1)\\
8 & Voidwracked Mage (CC)\\
8 & War Priest (CC)\\
8 & \textbf{Xhkarsh} (TOB1)\\
9 & \textbf{Bukavac} (TOB1)\\
9 & Hierophant Lich (CC)\\
9 & Monarch Skeleton (CC)\\
9 & Murgrik (TOB2)\\
9 & Plague Spirit (TOB2)\\
9 & \textbf{Vine Troll Skeleton} (TOB1)\\
10 & \textbf{Bone Swarm} (TOB1)\\
10 & Dragon Zombie (TOB2)\\
10 & \textbf{Fear Smith} (TOB1)\\
10 & Foxfire Ooze (CC)\\
10 & Nightgaunt (CC)\\
10 & Peat Mammoth (TOB2)\\
11 & \textbf{Abominable Beauty} (TOB1)\\
11 & Ghost Dragon (CC)\\
11 & Herald of Undeath (CC)\\
11 & Vampiric Knight (CC)\\
11 & Void Giant (CC)\\
11 & \textbf{Voidling} (TOB1)\\
12 & Culicoid Demon (TOB2)\\
12 & \textbf{Flutterflesh} (TOB1)\\
12 & Venom Maw Hydra (CC)\\
13 & \textbf{Mask Wight} (TOB1)\\
13 & Pustulent Shambler (TOB2)\\
13 & Vampire Warlock (Variant) (TOB)\\
14 & Black Dragon, Adult\\
16 & Infernal Knight (Devil) (CC)\\
17 & Moonlit King (Fey Lord) (TOB1)\\
18 & Lord of the Hunt (Fey Lord) (TOB1)\\
\end{DndTable}

\section{Rivers and Other Waterways}

\begin{DndTable}{cX}

CR & Creature\\
0 & \textbf{Commoner}\\
0 & \textbf{Crab}\\
0 & \textbf{Eagle}\\
0 & \textbf{Octopus}\\
0 & \textbf{Quipper}\\
0 & \textbf{Sea Horse}\\
1/8 & \textbf{Bandit}\\
1/8 & \textbf{Blood Hawk}\\
1/8 & \textbf{Giant Crab}\\
1/8 & \textbf{Guard}\\
1/8 & \textbf{Kobold}\\
1/8 & Leonino (CC)\\
1/8 & \textbf{Merfolk}\\
1/8 & \textbf{Poisonous Snake}\\
1/8 & \textbf{Stirge}\\
1/8 & \textbf{Wharfling} (TOB1)\\
1/4 & \textbf{Garroter Crab} (TOB1)\\
1/4 & \textbf{Giant Lizard}\\
1/4 & \textbf{Giant Wolf Spider}\\
1/4 & \textbf{Pseudodragon}\\
1/4 & \textbf{Steam Mephit}\\
1/2 & \textbf{Bone Crab} (TOB1)\\
1/2 & \textbf{Giant Sea Horse}\\
1/2 & \textbf{Morphoi} (TOB1)\\
1/2 & \textbf{Reef Shark}\\
1/2 & \textbf{Rum Gremlin} (TOB1)\\
1/2 & \textbf{Sahuagin}\\
1 & \textbf{Gerridae} (TOB1)\\
1 & \textbf{Giant Eagle}\\
1 & \textbf{Giant Octopus}\\
1 & \textbf{Giant Toad}\\
1 & \textbf{Glass Gator} (TOB1)\\
1 & Grindylow (CC)\\
1 & \textbf{Harpy}\\
1 & \textbf{Swarm of Quippers}\\
2 & \textbf{Bandit Captain}\\
2 & Bearmit Crab (CC)\\
2 & \textbf{Berserker}\\
2 & \textbf{Bronze Dragon Wyrmling}\\
2 & \textbf{Deep One} (TOB1)\\
2 & \textbf{Eel Hound} (TOB1)\\
2 & \textbf{Giant Constrictor Snake}\\
2 & \textbf{Griffon}\\
2 & \textbf{Hunter Shark}\\
2 & \textbf{Merrow}\\
2 & \textbf{Ogre}\\
2 & Scrag, Lesser (CC)\\
2 & \textbf{Sea Dragon Wyrmling} (TOB1)\\
2 & \textbf{Sea Hag}\\
2 & Shadow Skeleton (CC)\\
2 & \textbf{Shellycoat} (TOB1)\\
2 & Yann-An-Oed (CC)\\
3 & \textbf{Blue Dragon Wyrmling}\\
3 & Kelp Drake (TOB2)\\
3 & \textbf{Killer Whale}\\
3 & \textbf{Mahoru} (TOB1)\\
3 & \textbf{Manticore}\\
3 & Purple Slime (CC)\\
3 & \textbf{Veteran}\\
4 & Deep One Hybrid Priest (TOB1)\\
4 & Giant Water Scorpion (TOB2)\\
4 & Laestrigonian Giant (CC)\\
4 & \textbf{River Spirit} (see \textbf{Chapter 9})\\
4 & Shoreline Scrapper (CC)\\
4 & Water Horse (CC)\\
4 & \textbf{Water Leaper} (TOB1)\\
4 & Wharfling Swarm (TOB1)\\
5 & \textbf{Drowned Maiden} (TOB1)\\
5 & Gargoctopus (CC)\\
5 & Garroter Crab, Giant (TOB)\\
5 & \textbf{Giant Shark}\\
5 & \textbf{Lorelei} (TOB1)\\
5 & \textbf{Subek} (TOB1)\\
5 & Undine (CC)\\
5 & \textbf{Water Elemental}\\
5 & Wereshark (Lycanthrope) (TOB2)\\
5 & Windy Wailer (TOB2)\\
6 & \textbf{River Giant} (see \textbf{Chapter 9})\\
6 & Rusalka (TOB1)\\
6 & Tidehunter (TOB2)\\
7 & \textbf{Coral Drake} (TOB1)\\
7 & Haleshi (TOB2)\\
7 & \textbf{Lake Troll} (TOB1)\\
7 & Light Dragon, Young (CC)\\
7 & Scrag, Greater (CC)\\
7 & Two-Headed Eagle (CC)\\
8 & Bronze Dragon, Young\\
8 & \textbf{Deep One Archimandrite} (TOB1)\\
8 & Ikuchi (TOB2)\\
8 & Ziphius (CC)\\
9 & Blue Dragon, Young\\
9 & Keelbreaker Crab (TOB2)\\
9 & Sea Dragon, Young (TOB1)\\
9 & Shadow River Lord (CC)\\
10 & \textbf{Aboleth}\\
11 & \textbf{Roc}\\
12 & \textbf{Dragon Eel} (TOB1)\\
13 & Devil Shark (CC)\\
13 & \textbf{Storm Giant}\\
13 & Storm Lord (Elemental) (CC)\\
14 & \textbf{Isonade} (TOB1)\\
14 & \textbf{Krake Spawn} (TOB1)\\
15 & Bronze Dragon, Adult\\
16 & Blue Dragon, Adult\\
16 & Light Dragon, Adult (CC)\\
16 & River King (Fey Lord) (TOB1)\\
16 & Sea Dragon, Adult (TOB1)\\
17 & \textbf{Dragon Turtle}\\
23 & \textbf{Kraken}\\
\end{DndTable}

\section{Shadow Roads}

\begin{DndTable}{cX}

CR & Creature\\
1/8 & Lone Sluagh (Sluagh Swarm) (TOB1)\\
1/8 & \textbf{Shroud} (TOB1)\\
1/8 & \textbf{Stryx} (TOB1)\\
1/4 & Alliumite (CC)\\
1/4 & \textbf{Clurichaun} (TOB1)\\
1/4 & Erina Scrounger (TOB1)\\
1/4 & Goreling (CC)\\
1/4 & Living Shade (CC)\\
1/4 & \textbf{Living Wick} (TOB1)\\
1/4 & \textbf{Molefolk} (see \textbf{Chapter 9})\\
1/4 & \textbf{Ramag} (TOB1)\\
1/4 & \textbf{Ratfolk} (TOB1)\\
1/4 & \textbf{Roachling Skirmisher} (TOB1)\\
1/4 & Shadow Fey (Elves) (TOB1)\\
1/4 & \textbf{Treacle} (TOB1)\\
1/4 & \textbf{Witchlight} (TOB1)\\
1/2 & \textbf{Alehouse Drake} (TOB1)\\
1/2 & \textbf{Beggar Ghoul} (TOB1)\\
1/2 & \textbf{Bucca} (TOB1)\\
1/2 & \textbf{Burrowling} (TOB1)\\
1/2 & Chupacabra (CC)\\
1/2 & \textbf{Clockwork Beetle} (TOB1)\\
1/2 & Corpse Thief (CC)\\
1/2 & \textbf{Empty Cloak} (TOB1)\\
1/2 & Execrable Shrub (Planar Flora) (CC)\\
1/2 & Fire Imp (Devil) (CC)\\
1/2 & Flesh Reaver (CC)\\
1/2 & Foxin (CC)\\
1/2 & \textbf{Fraughashar} (TOB1)\\
1/2 & Keyhole Dragonette (TOB2)\\
1/2 & Kobold Spellclerk (TOB2)\\
1/2 & Nalusa Falaya (CC)\\
1/2 & \textbf{Nkosi} (TOB1)\\
1/2 & \textbf{Ravenfolk Scout} (TOB1)\\
1/2 & Rubble Swarm (Shard Swarm) (CC)\\
1/2 & Shard Swarm (CC)\\
1/2 & \textbf{Skin bat} (TOB1)\\
1/2 & Sooze (Shoth) (CC)\\
1/2 & \textbf{Tosculi Drone} (TOB1)\\
1/2 & Void Cultist (CC)\\
1/2 & Weirding Scroll (CC)\\
1 & Alchemical Apprentice Ooze (CC)\\
1 & \textbf{Chernomoi} (TOB1)\\
1 & Dark Servant (Dark Folk) (CC)\\
1 & Dhampir (CC)\\
1 & Dragonborn Light Cavalry (CC)\\
1 & Elite Kobold (CC)\\
1 & \textbf{Emerald Eye} (TOB1)\\
1 & \textbf{Eonic Drifter} (TOB1)\\
1 & \textbf{Erina Defender} (TOB1)\\
1 & Flithidir (TOB2)\\
1 & \textbf{Kobold Trapsmith} (TOB1)\\
1 & Nodosaurus (Dinosaur) (CC)\\
1 & Pact Drake (CC)\\
1 & Phantom (CC)\\
1 & \textbf{Ratfolk Rogue} (TOB1)\\
1 & Ratfolk Warlock (CC)\\
1 & Serpentfolk of Yig (CC)\\
1 & Suppurating Ooze (CC)\\
1 & \textbf{Wind Dragon Wyrmling} (TOB1)\\
2 & \textbf{Beli} (TOB1)\\
2 & \textbf{Clockwork Hound} (TOB1)\\
2 & \textbf{Cobbleswarm} (TOB1)\\
2 & \textbf{Doppelrat} (TOB1)\\
2 & Dragonborn Elementalist (CC)\\
2 & Dragonborn Heavy Cavalry (CC)\\
2 & Dream Squire (CC)\\
2 & \textbf{Eala} (TOB1)\\
2 & \textbf{Firegeist} (TOB1)\\
2 & \textbf{Folk of Leng} (TOB1)\\
2 & \textbf{Ink Devil} (TOB1)\\
2 & Kitsune (CC)\\
2 & \textbf{Kobold Alchemist} (TOB1)\\
2 & Kobold Junk Shaman (CC)\\
2 & Lady in White (CC)\\
2 & Light Dragon Wyrmling (CC)\\
2 & Lystrosaurus (Dinosaur) (CC)\\
2 & \textbf{Myling} (TOB1)\\
2 & \textbf{Noctiny} (TOB1)\\
2 & \textbf{Paper Drake} (TOB1)\\
2 & Ratfolk Mercenary (CC)\\
2 & \textbf{Roachling Lord} (TOB1)\\
2 & Roachling Scout (CC)\\
2 & Scitalis (CC)\\
2 & Seeping Death Skeleton (Suppurating Ooze) (CC)\\
2 & \textbf{Shadhavar} (TOB1)\\
2 & Shadow Fey Pattern Dancer (Elves) (CC)\\
2 & Sigilian (CC)\\
2 & \textbf{Spider Thief} (TOB1)\\
2 & \textbf{Tosculi Warrior} (TOB1)\\
2 & Trollkin Grunt (CC)\\
2 & \textbf{Uraeus} (TOB1)\\
2 & Vampire Thrall (CC)\\
2 & \textbf{Vile Barber} (TOB1)\\
2 & Wind Demon (CC)\\
2 & Wind Eater (CC)\\
3 & Ahu-Nixta (CC)\\
3 & Bleakheart (TOB2)\\
3 & Bloody Bones (CC)\\
3 & Bronze Golem (CC)\\
3 & Catscratch (TOB2)\\
3 & \textbf{Clockwork Huntsman} (TOB1)\\
3 & Dark Eye Dark Folk (CC)\\
3 & Dragonborn Edjet (CC)\\
3 & \textbf{Elvish Veteran Archer} (TOB1)\\
3 & Fang of the Great Wolf (Worg) (CC)\\
3 & \textbf{Far Darrig} (TOB1)\\
3 & Far Wanderer (CC)\\
3 & Freezing Shadow Ooze (CC)\\
3 & Ghast of Leng (CC)\\
3 & Ghost Boar (CC)\\
3 & Lesser Wood Golem (CC)\\
3 & \textbf{Mbielu Dinosaur} (TOB1)\\
3 & \textbf{Monolith Footman} (TOB1)\\
3 & Orthrus (CC)\\
3 & Plaresh Demon (CC)\\
3 & Preta (Hungry Ghost) (CC)\\
3 & Shadow Boxer (TOB2)\\
3 & Shadow Ooze (CC)\\
3 & Skull Drake (CC)\\
3 & Snake with a Hundred Mage Hands (TOB2)\\
3 & Thursir Giant (TOB1)\\
3 & War Chaplain (CC)\\
3 & Werebat (Lycanthrope) (CC)\\
3 & Werehyena (Lycanthrope) (CC)\\
3 & \textbf{Wolf Reaver Dwarf} (TOB1)\\
4 & \textbf{Ash Drake} (TOB1)\\
4 & \textbf{Bandit Lord} (TOB1)\\
4 & Battle Mage (CC)\\
4 & Cats of Ulthar (CC)\\
4 & \textbf{City Watch Captain} (TOB1)\\
4 & \textbf{Death Butterfly Swarm} (TOB1)\\
4 & Dwarf Graveslayer (CC)\\
4 & \textbf{Firebird} (TOB1)\\
4 & Gear Mage (CC)\\
4 & Gremlin Rum Lord (TOB2)\\
4 & Ink Guardian Ooze (CC)\\
4 & \textbf{Kobold Chieftain} (TOB1)\\
4 & \textbf{Lich Hound} (TOB1)\\
4 & Necrophage Ghast (Ghoul) (CC)\\
4 & \textbf{Nkosi Pridelord} (TOB1)\\
4 & Ortifex (TOB2)\\
4 & Philosopher's Ghost (CC)\\
4 & \textbf{Ratatosk} (TOB1)\\
4 & Rattok Demon (CC)\\
4 & Servant of Yig (CC)\\
4 & Shadow Fey Guardian (Elves) (TOB1)\\
4 & Tar Ghoul (CC)\\
4 & Tosculi Jeweled Drone (CC)\\
4 & \textbf{Trollkin Reaver} (TOB1)\\
4 & Trollkin Shaman (CC)\\
4 & \textbf{Tusked Skyfish} (TOB1)\\
4 & Vaettir (TOB1)\\
4 & Yek Demon (CC)\\
5 & \textbf{Aridni} (TOB1)\\
5 & \textbf{Black Knight Commander} (TOB1)\\
5 & \textbf{Bouda} (TOB1)\\
5 & Centaur Chieftan (CC)\\
5 & Cipactli Demon (CC)\\
5 & \textbf{Clockwork Abomination} (TOB1)\\
5 & \textbf{Corrupting Ooze} (TOB1)\\
5 & Dark Voice (Dark Folk) (CC)\\
5 & \textbf{Dream Eater} (TOB1)\\
5 & Dream Wraith (CC)\\
5 & Dvarapala (CC)\\
5 & Dwarf Cleric of the Brew (CC)\\
5 & \textbf{Fellforged} (TOB1)\\
5 & Fierstjerren (CC)\\
5 & \textbf{Gnomish Distiller} (see \textbf{Chapter 9})\\
5 & Gnomish Knife Cultist (CC)\\
5 & \textbf{Hound of the Night} (TOB1)\\
5 & \textbf{Hulking Whelp} (TOB1)\\
5 & Karakura (CC)\\
5 & \textbf{Kikimora} (TOB1)\\
5 & Lamia Serpentine Matriarch (CC)\\
5 & Lord Zombie (CC)\\
5 & \textbf{Mi-Go} (TOB1)\\
5 & Oth (Shoth) (CC)\\
5 & \textbf{Owl Harpy} (TOB1)\\
5 & Quickstep (CC)\\
5 & Ramag Portal Master (CC)\\
5 & Ratatosk Warlord (CC)\\
5 & \textbf{Ravenfolk Doom Croaker} (TOB1)\\
5 & \textbf{Sandman} (TOB1)\\
5 & Shadow Fey Forest Hunter (Elves) (TOB1)\\
5 & Shadow Fey Knight of the Road (Elves) (CC)\\
5 & Sleipnir (CC)\\
5 & Spawn of Akyishigal (Demon) (TOB1)\\
5 & Spawn of Arbeyach (Devil) (TOB1)\\
5 & Spawn of Parzelon (Devil) (CC)\\
5 & Spree (Demon) (CC)\\
5 & \textbf{Temple Dog} (TOB1)\\
5 & Three-Headed Cobra (CC)\\
5 & Tusked Crimson Ogre (CC)\\
5 & \textbf{Vapor Lynx} (TOB1)\\
5 & \textbf{Wormhearted Suffragan} (TOB1)\\
5 & Xiphus (CC)\\
6 & Bearfolk Chieftain (CC)\\
6 & Blood Ooze (CC)\\
6 & Clockwork Assassin (CC)\\
6 & Crimson Mist (CC)\\
6 & \textbf{Crystalline Devil} (TOB1)\\
6 & Death Butterfly Swarm, Greater (TOB1)\\
6 & Doomspeaker (CC)\\
6 & Dracotaur (CC)\\
6 & Elf Deathsworn (CC)\\
6 & Elf Servant of the Vine (CC)\\
6 & \textbf{Fate Eater} (TOB1)\\
6 & \textbf{Fext} (TOB1)\\
6 & Fey Drake (CC)\\
6 & \textbf{Gearforged Templar} (TOB1)\\
6 & Ghost Boar, Elder (CC)\\
6 & Ghost Dwarf (CC)\\
6 & \textbf{Ghost Knight} (TOB1)\\
6 & Kobold King (CC)\\
6 & \textbf{Loxoda} (TOB1)\\
6 & \textbf{Malphas} (TOB1)\\
6 & \textbf{Mirror Hag} (TOB1)\\
6 & Moonchild Naga (CC)\\
6 & Necromancer (CC)\\
6 & \textbf{Nichny} (TOB1)\\
6 & \textbf{Nightgarm} (TOB1)\\
6 & Ouroboros (CC)\\
6 & Shadow Fey Duelist (Elves) (TOB1)\\
6 & \textbf{Spectral Guardian} (TOB1)\\
6 & Wind Dragon, Young (TOB1)\\
6 & Wolf Spirit Swarm (TOB1)\\
7 & Alabaster Tree (Planar Flora) (CC)\\
7 & Blood Mage (CC)\\
7 & Dhampir Commander (CC)\\
7 & \textbf{Dissimortuum} (TOB1)\\
7 & \textbf{Dwarven ringmage} (TOB1)\\
7 & \textbf{Einherjar} (TOB1)\\
7 & Elf Enchanter (CC)\\
7 & Fire Dancer Swarm (TOB1)\\
7 & Giant Sloth (CC)\\
7 & \textbf{Gilded Devil} (TOB1)\\
7 & \textbf{Herald of Darkness} (TOB1)\\
7 & \textbf{Psoglav Demon} (TOB1)\\
7 & \textbf{Red Hag} (TOB1)\\
7 & \textbf{Risen Reaver} (TOB1)\\
7 & Sentinel in Darkness (CC)\\
7 & Shadow Fey Enchantress (Elves) (TOB1)\\
7 & Silenal (TOB2)\\
7 & \textbf{Soul Eater} (TOB1)\\
7 & \textbf{Spark} (TOB1)\\
7 & \textbf{Umbral Vampire} (TOB1)\\
7 & Vellso Demon (CC)\\
7 & Void Speaker (CC)\\
8 & Ankou Soul Seeker (CC)\\
8 & Arborcyte (CC)\\
8 & Arcane Spectral Guardian (TOB1)\\
8 & Blood Giant (CC)\\
8 & \textbf{Bone Collective} (TOB1)\\
8 & \textbf{Chronalmental} (TOB1)\\
8 & Clockwork Dragon (CC)\\
8 & \textbf{Degenerate Titan} (TOB1)\\
8 & \textbf{Dragonleaf Tree} (TOB1)\\
8 & \textbf{Emerald Order Cult Leader} (TOB1)\\
8 & \textbf{Feyward Tree} (TOB1)\\
8 & Fractal Golem (CC)\\
8 & Hound of Tindalos (CC)\\
8 & \textbf{Kishi Demon} (TOB1)\\
8 & \textbf{Liosalfar} (TOB1)\\
8 & \textbf{Lunar Devil} (TOB1)\\
8 & Manastorm Golem (CC)\\
8 & \textbf{Monolith Champion} (TOB1)\\
8 & Moon Nymph (CC)\\
8 & Nachzehrer (CC)\\
8 & Neophron Demon (CC)\\
8 & Oliphaunt (CC)\\
8 & \textbf{Qwyllion} (TOB1)\\
8 & Ring Servant (CC)\\
8 & \textbf{Rust Drake} (TOB1)\\
8 & Tetomatli (TOB2)\\
8 & Thief Lord (CC)\\
8 & Vampire Priestess (CC)\\
8 & Vines of Nemthyr (CC)\\
8 & Voidwracked Mage (CC)\\
8 & War Priest (CC)\\
8 & \textbf{Xhkarsh} (TOB1)\\
9 & Alchemical Golem (CC)\\
9 & Alnaar Demon (CC)\\
9 & Devilbound Gnomish Prince (TOB1)\\
9 & \textbf{Eater of Dust} (TOB1)\\
9 & Fulad-Zereh Demon (CC)\\
9 & Gaki-Jikininki (Hungry Ghost) (CC)\\
9 & Hierophant Lich (CC)\\
9 & Lotus Golem (CC)\\
9 & \textbf{Malakbel Demon} (TOB1)\\
9 & Monarch Skeleton (CC)\\
9 & Righteous Sentinel (TOB2)\\
9 & Shadow Fey Ambassador (Elves) (CC)\\
9 & \textbf{Vine Troll Skeleton} (TOB1)\\
9 & Yek Demon, Alpha (CC)\\
10 & \textbf{Algorith} (TOB1)\\
10 & Altar Flame Golem (CC)\\
10 & \textbf{Automata Devil} (TOB1)\\
10 & \textbf{Bone Swarm} (TOB1)\\
10 & Doom Golem (CC)\\
10 & Elf Alchemist Archer (CC)\\
10 & \textbf{Fear Smith} (TOB1)\\
10 & Grave Behemoth (CC)\\
10 & \textbf{Hundun} (TOB1)\\
10 & Lamassu (CC)\\
10 & Nightgaunt (CC)\\
10 & \textbf{Rubezahl Demon} (TOB1)\\
10 & \textbf{Thuellai} (TOB1)\\
10 & Vampire Patrician (CC)\\
11 & \textbf{Abominable Beauty} (TOB1)\\
11 & \textbf{Berstuc Demon} (TOB1)\\
11 & \textbf{Blood Hag} (TOB1)\\
11 & \textbf{Buraq} (TOB1)\\
11 & \textbf{Corpse Mound} (TOB1)\\
11 & Dragonborn Ouroban (CC)\\
11 & \textbf{Dullahan} (TOB1)\\
11 & \textbf{Eye Golem} (TOB1)\\
11 & Ghost Dragon (CC)\\
11 & \textbf{Grim Jester} (TOB1)\\
11 & Herald of Slaughter (TOB2)\\
11 & Herald of Undeath (CC)\\
11 & \textbf{Koralk Devil} (TOB1)\\
11 & Mouse King (Animal Lord) (CC)\\
11 & \textbf{Naina} (TOB1)\\
11 & Shadow Fey Poisoner (Elves) (CC)\\
11 & \textbf{Valkyrie} (TOB1)\\
11 & Vampiric Knight (CC)\\
11 & Void Giant (CC)\\
11 & \textbf{Voidling} (TOB1)\\
12 & Akyishigal of Cockroaches (Demon Lord) (TOB1)\\
12 & \textbf{Ancient Titan} (TOB1)\\
12 & \textbf{Chort Devil} (TOB1)\\
12 & Chronomatic Enhancer (TOB2)\\
12 & Droth (Shoth) (CC)\\
12 & \textbf{Herald of Blood} (TOB1)\\
12 & Herald of Fire (CC)\\
12 & \textbf{Knight of Shadows} (see \textbf{Chapter 9})\\
12 & \textbf{Mavka} (TOB1)\\
12 & \textbf{Skein Witch} (TOB1)\\
12 & \textbf{Son of Fenris} (TOB1)\\
12 & \textbf{Tosculi Hive Queen} (TOB1)\\
12 & Virtuoso Lich (TOB2)\\
13 & Bat King (Animal Lord) (CC)\\
13 & Chittr'k'k of Rats (Demon Lord) (CC)\\
13 & Clockwork Archon (TOB2)\\
13 & \textbf{Haugbui} (TOB1)\\
13 & \textbf{Mask Wight} (TOB1)\\
13 & Psychic Vampire (TOB2)\\
13 & \textbf{Steam Golem} (TOB1)\\
13 & Vampire Warlock (TOB)\\
14 & \textbf{Cambium} (TOB1)\\
14 & \textbf{Orobas Devil} (TOB1)\\
14 & Three-Headed Clockwork Dragon (CC)\\
15 & Herald of the Great Wyrm (CC)\\
15 & Ia'Affrat the Insatiable (Arch-Devil) (TOB1)\\
15 & Pact Lich (CC)\\
15 & Queen of Cats (Animal Lord) (CC)\\
15 & Rakshasa, Greater (CC)\\
15 & \textbf{Star-Spawn of Cthulhu} (TOB1)\\
15 & Whisperer in Darkness (CC)\\
16 & Berchta (Fey Lady) (CC)\\
16 & Infernal Knight (Devil) (CC)\\
16 & Snow Queen (Fey Lady) (TOB1)\\
17 & Barong (CC)\\
17 & Koschei (TOB1)\\
17 & Moonlit King (Fey Lord) (TOB1)\\
17 & Queen of Witches (Fey Lady) (TOB1)\\
17 & Rangda of Witches (Demon Lord) (CC)\\
17 & Wind Dragon, Adult (TOB1)\\
18 & Lord of the Hunt (Fey Lord) (TOB1)\\
19 & Chemosh (Demon Lord) (CC)\\
19 & Hraesvelgr the Corpse Swallower (Giant) (TOB1)\\
19 & Living Star (CC)\\
19 & \textbf{Shoggoth} (TOB1)\\
20+ & Alquam (Demon Lord) (TOB1)\\
20+ & Avatar of Shoth (Shoth) (CC)\\
20+ & Camazotz (Demon Lord) (TOB1)\\
20+ & Emperor of the Ghouls (Ghouls) (TOB1)\\
20+ & Jotun Giant (TOB1)\\
20+ & Mammon (Arch-Devil) (TOB1)\\
20+ & Monkey King (Simian) (CC)\\
20+ & Parzalon (Arch-Devil) (CC)\\
20+ & Qorgeth (Demon Lord) (TOB1)\\
20+ & Queen of Night and Magic (Fey Lady) (TOB1)\\
20+ & Totivillus (Arch-Devil) (TOB1)\\
20+ & Typhon (Demon Lord) (CC)\\
20+ & Wind Dragon, Ancient (TOB1)\\
\end{DndTable}

\begin{DndSidebar}[float=!b]{On the Road}
An extremely wide range of creatures from lowly stryx to Demon lords may be found on the Shadow Roads. Surprisingly, the more dangerous they are, the more often they seem to be in a hurry to get somewhere. Some encounters on a Shadow Road are nothing more than a passing nod to Camazotz, Demon Lord of Bats, as he flutters past.



\end{DndSidebar}

\section{Underground}

\begin{DndTable}{cX}

CR & Creature\\
0 & \textbf{Giant Fire Beetle}\\
0 & Lantern Beetle (TOB2)\\
0 & \textbf{Shrieker}\\
0 & Zoog (CC)\\
1/8 & \textbf{Giant Rat}\\
1/8 & \textbf{Kobold}\\
1/8 & Pixie's Umbrella Fungi (TOB2)\\
1/8 & Sniffer Beetle (TOB2)\\
1/8 & \textbf{Stirge}\\
1/8 & Wicked Skull (TOB2)\\
1/4 & Cave Goat (TOB2)\\
1/4 & Ghoul Bat (TOB2)\\
1/4 & \textbf{Giant Bat}\\
1/4 & \textbf{Giant Centipede}\\
1/4 & \textbf{Giant Lizard}\\
1/4 & \textbf{Giant Poisonous Snake}\\
1/4 & \textbf{Goblin}\\
1/4 & \textbf{Grimlock}\\
1/4 & \textbf{Living Wick} (TOB1)\\
1/4 & \textbf{Ratfolk} (TOB1)\\
1/4 & Shiftshroom Fungi (TOB2)\\
1/4 & \textbf{Skeleton}\\
1/4 & \textbf{Swarm of Bats}\\
1/4 & \textbf{Violet Fungus}\\
1/2 & \textbf{Beggar Ghoul} (TOB1)\\
1/2 & Boomer Fungi (TOB2)\\
1/2 & Boot Grabber (CC)\\
1/2 & Crimson Shambler (TOB2)\\
1/2 & \textbf{Darkmantle}\\
1/2 & Derro Guard (TOB2)\\
1/2 & \textbf{Gray Ooze}\\
1/2 & \textbf{Hobgoblin}\\
1/2 & \textbf{Kobold Trapsmith} (TOB1)\\
1/2 & Light Drake (TOB2)\\
1/2 & \textbf{Magma Mephit}\\
1/2 & Mydnari (TOB2)\\
1/2 & \textbf{Orc}\\
1/2 & \textbf{Rust Monster}\\
1/2 & \textbf{Shadow}\\
1/2 & Shadow Goblin (CC)\\
1 & \textbf{Bugbear}\\
1 & \textbf{Dogmole} (TOB1)\\
1 & \textbf{Duergar}\\
1 & \textbf{Ghoul}\\
1 & \textbf{Giant Spider}\\
1 & \textbf{Giant Toad}\\
1 & Kobold Elite (CC)\\
1 & \textbf{Kobold Trapsmith} (TOB1)\\
1 & Putrescent Slime (TOB2)\\
1 & \textbf{Ratfolk Rogue} (TOB1)\\
1 & Ratfolk Warlock (CC)\\
1 & \textbf{Specter}\\
1 & Stone Creeper (TOB2)\\
1 & Suppurating Ooze (CC)\\
1 & Wirbeln Fungi (CC)\\
2 & \textbf{Ankheg}\\
2 & Black Sun Orc (CC)\\
2 & \textbf{Cave Dragon Wyrmling} (TOB1)\\
2 & Clacking Skeleton (CC)\\
2 & \textbf{Cobbleswarm} (TOB1)\\
2 & Crypt Spider (CC)\\
2 & Derro Explorer (TOB2)\\
2 & \textbf{Doppelrat} (TOB1)\\
2 & Dream Squire (CC)\\
2 & Flame Eater Swarm (CC)\\
2 & \textbf{Gargoyle}\\
2 & \textbf{Gelatinous Cube}\\
2 & \textbf{Ghast}\\
2 & Giant Vampire Bat (CC)\\
2 & \textbf{Gibbering Mouther}\\
2 & \textbf{Grick}\\
2 & \textbf{Ink Devil} (TOB1)\\
2 & \textbf{Kobold Alchemist} (TOB1)\\
2 & Kobold Junk Shaman (CC)\\
2 & \textbf{Mimic}\\
2 & \textbf{Ochre Jelly}\\
2 & \textbf{Ogre}\\
2 & Ophinix (TOB2)\\
2 & Pech (CC)\\
2 & Ratfolk Mercenary (CC)\\
2 & Roachling Scout (CC)\\
2 & Rock Roach (TOB2)\\
2 & Seeping Death Skeleton (Suppurating Ooze) (CC)\\
2 & Shadow Fey Pattern Dancer (Elves) (CC)\\
2 & \textbf{Shellycoat} (TOB1)\\
2 & Swarm of Esteron (TOB2)\\
3 & \textbf{Basilisk}\\
3 & Corpse Worm (TOB2)\\
3 & \textbf{Darakhul} (TOB1)\\
3 & \textbf{Doppelganger}\\
3 & Gaunt One (TOB2)\\
3 & Ghoulsteed (CC)\\
3 & Giant Albino Bat (CC)\\
3 & \textbf{Hell Hound}\\
3 & Keg Golem (CC)\\
3 & Megapede (CC)\\
3 & Mineral Ooze (TOB2)\\
3 & \textbf{Minotaur}\\
3 & \textbf{Night Scorpion} (TOB1)\\
3 & Orc Black Sun Priestess (CC)\\
3 & \textbf{Phase Spider}\\
3 & Swolbold (Kobold) (CC)\\
3 & \textbf{Venomous Mummy} (TOB1)\\
3 & \textbf{Veteran}\\
3 & \textbf{Wight}\\
3 & Woe Siphon (TOB2)\\
4 & \textbf{Angler Worm} (TOB1)\\
4 & \textbf{Black Pudding}\\
4 & \textbf{Cavelight Moss} (TOB1)\\
4 & \textbf{Chuul}\\
4 & Darakhul Shadowmancer (Ghoul) (CC)\\
4 & \textbf{Deathcap Myconid} (TOB1)\\
4 & Deep Troll (TOB2)\\
4 & \textbf{Derro Fetal Savant} (TOB1)\\
4 & Emperor's Hyena (TOB2)\\
4 & \textbf{Ettin}\\
4 & \textbf{Ghost}\\
4 & \textbf{Imperial Ghoul} (TOB1)\\
4 & \textbf{Kobold Chieftain} (TOB1)\\
4 & Necrophage Ghast (Ghoul) (CC)\\
4 & \textbf{Oculo Swarm} (TOB1)\\
4 & Overshadow (TOB2)\\
4 & Pech Stonemaster (CC)\\
4 & \textbf{Skitterhaunt} (TOB1)\\
4 & Strobing Fungi (TOB2)\\
5 & Darakhul Spy (Ghoul) (TOB2)\\
5 & Deathweaver (TOB2)\\
5 & \textbf{Derro Shadow Antipaladin} (TOB1)\\
5 & Derro Speaker to Darkness (CC)\\
5 & Derro Witch Queen (CC)\\
5 & \textbf{Dogmole Juggernaut} (TOB1)\\
5 & \textbf{Dream Eater} (TOB1)\\
5 & \textbf{Earth Elemental}\\
5 & Eldritch Ooze (TOB2)\\
5 & \textbf{Iron Ghoul} (TOB1)\\
5 & Kobold Wizard (CC)\\
5 & Lunarian (TOB2)\\
5 & \textbf{Otyugh}\\
5 & Phase Giant (TOB2)\\
5 & \textbf{Rat King} (TOB1)\\
5 & \textbf{Roper}\\
5 & \textbf{Salamander}\\
5 & \textbf{Sarcophagus Slime} (TOB1)\\
5 & \textbf{Troll}\\
5 & \textbf{Vampire Spawn}\\
5 & \textbf{Wraith}\\
5 & \textbf{Xorn}\\
6 & Bloated Ghoul (TOB2)\\
6 & Blood Ooze (CC)\\
6 & \textbf{Chimera}\\
6 & \textbf{Crystalline Devil} (TOB1)\\
6 & Derro Shadowseeker (TOB2)\\
6 & \textbf{Drider}\\
6 & Echo Demon (CC)\\
6 & \textbf{Gnarljak} (TOB1)\\
6 & King Kobold (CC)\\
6 & Quiet Soul (CC)\\
6 & Ruby Ooze (CC)\\
6 & \textbf{Scheznyki} (TOB1)\\
7 & Crystalline Monolith (TOB2)\\
7 & Fey Revenant (TOB2)\\
7 & Pech Lithlord (CC)\\
7 & Sentinel in Darkness (CC)\\
7 & \textbf{Shadow Beast} (TOB1)\\
7 & \textbf{Stone Giant}\\
7 & Walled Horror (TOB2)\\
7 & Yathon (TOB2)\\
8 & Bloodstone Sentinel (TOB2)\\
8 & Cave Dragon, Young (TOB1)\\
8 & \textbf{Cloaker}\\
8 & Giant Shark Bowl Ooze (CC)\\
8 & \textbf{Idolic Deity} (TOB1)\\
8 & Lost Minotaur (CC)\\
8 & Onyx Magistrate (TOB2)\\
8 & \textbf{Rust Drake} (TOB1)\\
8 & \textbf{Spirit Naga}\\
9 & Darakhul Captain (Ghoul) (TOB2)\\
9 & Darakhul High Priestess (Ghoul) (CC)\\
9 & \textbf{Deep Drake} (TOB1)\\
9 & \textbf{Fire Giant}\\
9 & \textbf{Horakh} (TOB1)\\
9 & Nharyth (TOB2)\\
9 & Shadow Fey Ambassador (Elves) (CC)\\
9 & Shriekbat (TOB2)\\
9 & Spider Drake (CC)\\
9 & Thin Giant (TOB2)\\
10 & \textbf{Aboleth}\\
10 & Cave Giant (CC)\\
10 & Conjoined Queen (TOB2)\\
10 & \textbf{Sathaq Worm} (TOB1)\\
11 & \textbf{Behir}\\
11 & Shadow Fey Poisoner (Elves) (CC)\\
12 & \textbf{Bonepowder Ghoul} (TOB1)\\
12 & \textbf{Gug} (TOB1)\\
12 & Nihileth (Aboleth) (TOB1)\\
13 & Bone Colossus (Necrotech) (TOB2)\\
14 & \textbf{Cambium} (TOB1)\\
14 & Cave Giant Shaman (TOB2)\\
14 & \textbf{Gypsosphinx} (TOB1)\\
15 & \textbf{Purple Worm} (TOB1)\\
16 & Cave Dragon, Adult (TOB1)\\
16 & Flame Dragon, Adult (TOB1)\\
17 & \textbf{Urochar} (TOB1)\\
20 & Emperor of the Ghouls (TOB1)\\
\end{DndTable}

\chapter{Tactical Maps for the Shadow Realm}
The maps below include a fey village, a shadow road, a stretch of the River Styx, and other areas with the twilight shades and native flora emblematic of the Shadow Realm. The maps can be used for any encounter on the plane to enhance its shadowy atmosphere and fey landscapes.

\section{Fey Court Audience Chamber}
\section{Shadow Realm Forest}
\section{Shadow Realm Marsh}
\section{Shadow Realm Road}
\section{Shadow Realm River}
\section{Shadow Realm Village}
\chapter{Credits}

\begin{description}
\begin{multicols}{2}
\item[Design.]
Wolfgang Baur, Celeste Conowitch
\item[Additional Design.]
Richard Green, Sarah Madsen, Kelly Pawlik, Brian Suskind
\item[Development and Editing.]
Scott Gable
\item[Proofreading.]
Jeff Quick
\item[Cover Artist.]
Marcel Mercado
\item[Limited Edition Cover Artist.]
Addison Rankin
\item[Interior Artists.]
George Johnstone, Mike Pape, David Auden Nash, William O'Brien, Roberto Pitturru, Addison Rankin, Kiki Moch Rizki, Florian Stitz, Bryan Syme, Egil Thompson, Alexander Yakovlev
\item[Cartography.]
Jon Pintar, Samantha Senn, Dean Spencer
\item[Graphic Design.]
Marc Radle
\item[Layout.]
Marc Radle
\item[Additional Layout.]
Amber Seger
\item[Fantasy Grounds Digital Edition.]
Linda Buth
\item[Roll20 Digital Edition.]
Nic Bradley
\item[Shard Tabletop Digital Edition.]
Hal Howard
\item[Director of Digital Growth.]
Blaine McNutt
\item[Art Director.]
Marc Radle
\item[Editorial Director.]
Thomas M. Reid
\item[Director of Operations.]
T. Alexander Stangroom
\item[Sales Manager.]
Kym Weiler
\item[Community Manager.]
Victoria Rogers
\item[Twitch Producer.]
Chelsea "Dot" Steverson
\item[Publisher.]
Wolfgang Baur
\item A special thanks to the designers and contributors to earlier printings of some material that appears in this book, including, but not limited to: Wolfgang Baur and Dan Dillon
\item \textbf{An extra special thanks to the 2,706 backers who made this volume possible!}
\end{multicols}

\end{description}

\end{document}
